\documentclass{llncs}

\providecommand{\event}{RuleML 2015} % Name of the event you are submitting to
\usepackage{aeguill}
\usepackage[utf8]{inputenc}
\usepackage[T1]{fontenc}
\usepackage{textcomp}
\usepackage{lmodern}
\usepackage[final]{microtype}
\usepackage{slashbox}
\usepackage{graphicx,xcolor}
\usepackage[draft]{changes}
\usepackage[linktocpage, colorlinks=true, linkcolor=blue, urlcolor=blue, citecolor=blue, %
             pdfkeywords={
                genomic annotation,
                network reconstruction,
                business rules,
                oriented object logic,
                four valued logic}]{hyperref}
\usepackage{filecontents}
\usepackage{caption}
\usepackage{tikz}

\bibliographystyle{splncs03}
\pagestyle{headings}


\title{Reactive Graph Reasoning for Genomic Annotation}

\institute{Direction des Sciences du Vivant, CEA, Institut de Génomique, Genoscope, France \\
\and
CNRS-UMR8030, Evry, France \\
\and
Université d’Evry Val d’Essonne, Evry, France}

\author{Jonathan Mercier\inst{1}  \inst{3} \and David Vallenet\inst{1} \and Claudine Medigue\inst{1}  \inst{2}}

\begin{filecontents}{\jobname.bib}
    @article{mazumder2010community,
        title={Community annotation in biology},
        author={Mazumder, Raja and Natale, Darren A and Julio, Jessica Anne Ecalnir and Yeh, Lai-Su and Wu, Cathy H},
        journal={Biol Direct},
        volume={5},
        pages={12},
        year={2010}
    }
    @Article{vallenet2012microscope,
        author="Vallenet, D.  and Belda, E.  and Calteau, A.  and Cruveiller, S.  and Engelen, S.  and Lajus, A.  and Le Fevre, F.  and Longin, C.  and Mornico, D.  and Roche, D.  and Rouy, Z.  and Salvignol, G.  and Scarpelli, C.  and Thil Smith, A. A.  and Weiman, M.  and Medigue, C. ",
        title="{{M}icro{S}cope--an integrated microbial resource for the curation and comparative analysis of genomic and metabolic data}",
        journal="Nucleic Acids Res.",
        year="2013",
        volume="41",
        number="Database issue",
        pages="D636--647",
    }
    @incollection{belnap1977useful,
        title={A useful four-valued logic},
        author={Belnap Jr, Nuel D},
        booktitle={Modern uses of multiple-valued logic},
        pages={5--37},
        year={1977},
        publisher={Springer}
    }
    @misc{biolog,
        title={Biolog},
        author={Bochner, Barry R. and Barnby, Donald W.},
        journal={},
        howpublished="Website",
        url={http://www.biolog.com/products-static/phenotype\_microbial\_cells\_overview.php},
        year={}
    }
    @article{morgat2011unipathway,
        title={UniPathway: a resource for the exploration and annotation of metabolic pathways},
        author={Morgat, Anne and Coissac, Eric and Coudert, Elisabeth and Axelsen, Kristian B and Keller, Guillaume and Bairoch, Amos and Bridge, Alan and Bougueleret, Lydie and Xenarios, Ioannis and Viari, Alain},
        journal={Nucleic acids research},
        pages={gkr1023},
        year={2011},
        publisher={Oxford Univ Press}
    }
    @article{amir1999object,
        title={Object-oriented first-order logic},
        author={Amir, Eyal},
        journal={Linkiping Electronic Articles in Computer and Information Science (http://www.ida.liu.se/ext/etai)},
        url={http://www.ep.liu.se/ej/etai/1999/008/},
        pages={63--84},
        volume={4},
        year={1999}
    }
    @article{mcwhirter2001drools,
        title={Drools},
        author={McWhirter, Bob and others},
        howpublished="Website",
        pages={\href{http://drools.jboss.org}{drools.jboss.org}},
        year={2001}
    }
    @article{francke2005reconstructing,
        title={Reconstructing the metabolic network of a bacterium from its genome},
        author={Francke, Christof and Siezen, Roland J and Teusink, Bas},
        journal={Trends in microbiology},
        volume={13},
        number={11},
        pages={550--558},
        year={2005},
        publisher={Elsevier}
    }
    @article{proctor2008drools,
        title={Drools documentation},
        author={Proctor, Mark and Neale, Michael and Lin, Peter and Frandsen, Michael},
        journal={JBoss. org, Tech. Rep},
        year={2008}
    }
    @misc{phreak,
        title={PHREAK algorithm logic},
        author={Proctor, Mark},
        journal={},
        howpublished="Website",
        url={ https://www.google.fr/webhp?sourceid=chrome-instant\&ion=1\&espv=2\&ie=UTF-8#q=phreak+algorithm+logic},
        year={2013}
    }
    @article{sorokina2014profiling,
  title={Profiling the orphan enzymes},
  author={Sorokina, Maria and Stam, Mark and M{\'e}digue, Claudine and Lespinet, Olivier and Vallenet, David},
  journal={Biol. Direct},
  volume={9},
  number={10},
  year={2014}
}
    
\end{filecontents}

\begin{document}
    
\definecolor{plantucolor0000}{RGB}{254,254,206}
\definecolor{plantucolor0001}{RGB}{168,0,54}
\definecolor{plantucolor0002}{RGB}{180,167,229}
\definecolor{plantucolor0003}{RGB}{0,0,0}
\definecolor{plantucolor0004}{RGB}{132,190,132}
\definecolor{plantucolor0005}{RGB}{3,128,72}
\definecolor{plantucolor0006}{RGB}{173,209,178}
\definecolor{plantucolor0007}{RGB}{200,41,48}
    
\maketitle

\begin{abstract}
\replaced{Today's functional genomic knowledge is built from an incredible amount of genomes}{The emergence of the next generation sequencing generates an incredible amount of genomes}, whereas curation efforts to annotate them tend to decrease despite some community initiatives \cite{mazumder2010community}. Providing human expert annotation as gold standard for all of these genomes became challenging.\\
To ease this manual process, we develop a Genomic Rule Oriented Object Logic System named GROOLS. \added{Reconstructing metabolite networks with automatic tasks firstly to improves the expert genomic annotation process. Following this bio-annotator has a global view of predicted functions. And could to use in more the network context to annotate missing proteins.}\\
The combined use of modern paradigm as oriented object programming with the logic fields provide an interesting space of research. Complex biological information are pulled from various resources and represented as a structured object. These objects are bivalent, for a computer specialist, is an instance of a defined class, while for a logician is a theory.\\
Business Rule is used to coerce, control and making business decision. It separates clearly
the business logic from the application. Theses rules are applied to facts throw a Business To reason over these theories is need to use a Rule Management System (BRMS). It evaluates facts against a set of rules in a reproducible and efficient way.
This paper show a standardized way to analyze and annotate genomes using a four valued logic and a hierarchical knowledge representation.
\keywords{genomic annotation, network reconstruction, business rules, oriented object logic, four valued logic, hierarchical knowledge representation}
\end{abstract}

\section{General description}
GROOLS is an applied research project mixing  biology, informatics and logic. It aim to standardize
the use of biological results to annotate a genome. Highlighting bio-annotator on expected gene or not and discovery new knowledge. During analysis of a newly sequenced genome, the bio-annotator is faced to the difficult task of checking the annotation consistency. In such cases the bio-annotator should to check the proteome against the biological knowledge about the organisms under study. However this task is too huge, bio-annotator usually check sequence similarity to an already annotated protein to infer the gene annotation. The early stages are to understand the bio-annotator reasoning and the identification of a generic model establishing a vocabularies.\\
A bio-annotator use at least three kinds of facts. The first one is a prediction fact, found by bioinformatics computing or bio-annotator analysis. The second one is an asserted fact, an empirical evidence on organisms. The thirds is a trusted knowledge, a knowledge known to occur on one ore more organisms. A bio-annotator knowning for the studied organisms that is a tryptophan autotroph. This human expert will expect to found tryptophan bio-synthesis pathway. The corresponding knowledge will be marked as required. Reciprocally, if is known to not grow on tryptophan free media, the tryptophan bio-synthesis pathway is avoided. \\
Such reasoning system need a generic representation of the biological knowledge to be used on genomic annotation. And, on the other hand a BRMS able to reason over complex and structured data. \\
In order to have this kind of reasoning we need to select resources, where facts come from. An asserted fact can come from biolog data \cite{biolog}, a trusted knowledge from unipathway \cite{morgat2011unipathway} and prediction from bioinformatics results or bio-annotator expertise from the Prokaryotic Genome DataBase (PkGDB) \cite{vallenet2012microscope} \\
Trusted knowledge are linked together to form an graph. Prediction will influence the leaf of graph while assertion influence the roots of graph. The predictions influences are propagate from the leaf to the roots. Assertions and predictions facts are four valued. An assertion can take one of these values:  present, absent, both and unknown and a prediction: required, avoided, both, unknown. These state are inferred to knowledge giving ability to get a precise conclusion. This aims to help bio-annotator to focus on missing gene annotation using a reconstruction of the metabolic network from a sequenced genome \cite{francke2005reconstructing}.\\
This kind of logic is described by Belnap \cite{belnap1977useful} as a four valued logic. These theories are connected via an interface vocabularies. A set of rules is applied on a huge amount of biological knowledge to notify the completeness and the consistencies for each genome. To achieve this we have applied the Object-Oriented First Order Logic (OOFOL) \cite{amir1999object} and extended to a four valued logic. The four valued logic give a finer analysis than classic First Order Logic. In GROOLS, knowledge objects are theories which takes the form of a directed acyclic graph. Others are simple typed fact describing what exists that means all computed prediction or human prediction. Or the others kind of fact what should to exists, an assertion to a theory.

\section{Genomic annotation}

Genomic annotation task is one of most important, at ends organisms molecular physiology are decoded increasing our understanding of that particular organism. This field aims to understand what an organisms is able to do from his DNA sequence. The unit in genome annotation process is the product of gene (protein) and understanding its function. Some of these proteins are defined as enzymes, that means they are able to catalyze at least one chemical reaction. The enzymatic product of theses chemicals reactions will be used by an enzyme and so on. This enzyme chained process is described by a metabolic pathway (Figure~\ref{fig:pathway}). An enzyme can to be depicted in more than one metabolism pathway to form a metabolic network.


\begin{figure}
    \centering
    \includegraphics[scale=0.6]{img/unipathway.pdf}
    \caption{A metabolic pathway with two path to go from 2-oxoglutarate metabolite to L-Lysine . UER means unipathway enzymatic reaction which are involved by an enzymatic protein. Box with a same color are in a same block of reaction. This pathway has three reactions block: red green and blue. An organism could have a path with red and green blocks and/or red and blue block. Source: \href{http://www.grenoble.prabi.fr}{unipathway} .}
    \label{fig:pathway}
\end{figure}

A bio-annotator usually defined a coding region, search an homologous sequence and take a look to the genomic context to infer an annotation. In our opinion this task is more productive to do after an automatic genomic annotation. This allow the reconstruction of metabolic process (see \cite{francke2005reconstructing}), filled by automatics tasks, and help the bio-annotator to focus on missing gene. This automatics gene annotation can to be displayed with some others information
enabling a better understanding on the organism abilities (depicted in Figure~\ref{fig:metabolicNetworkReconstruction}). A first approach lead to conclusion the Pathway X is not present. But if is known that studied organism has the pathway X with some experimental information. Bio-annotator point of view the pathway X is certainly present as a specific reaction occur in the pathway X. Indeed biological knowledge has some hole. Some enzymatic reaction are known to occur into an  organisms but this protein are bound to any gene, usually described ad orphan enzyme \cite{sorokina2014profiling}. This alternative reasoning are optimistic and fit better with our empirical knowledge.

\begin{figure}
    \centering
    \def\svgwidth{\columnwidth}
    \input{pathway2.pdf_tex}
    \caption{Guided metabolic network reconstruction. Octagon shape is used for ``or node '' and a rounded box for ``and node ''. A red backgrounds means a knowledge is not present as opposite a cyan color means a present knowledge. }
    \label{fig:metabolicNetworkReconstruction}
\end{figure}

This kind of reasoning should to be done over more than four thousand genes by organisms. Moreover genome to annotate come quickly than bio-annotator community can do. On the other hand human expertise can lead to some inconsistencies which will be reused by others bio-annotator. To avoid these issues we develop an expert system to automatize the bio-annotator reasoning. An automated reasoner can manage billions of facts in a standard and reproducible way.

\section{Benefits of applying an oriented-object expert system}
Biological facts are stored into many bio-data warehouse, which are cross-referenced to reason over. For a computer scientist such information fit an oriented object approach. An oriented object programming ease in handling the data. Each kind of data corresponds to a class. With the use of an object relational mapping, objects are automatically pull to a typed information from a data warehouse. An object is an instance of the corresponding class. A computer scientist will apply a formula, an algorithm, a reasoning, depending of the current information type used. \\
We do not know in advance which type object the end used will put. To communicate between these different type of objects, the vocabulary is described by the use of interface. Interface provides a low coupling to have a well-structured computer system. Interface impose communication protocol, a partial set of attributes from classes are shared through any type of object. In other hand interface avoid multi-inheritance issues and provide a standardized vocabulary through any type of object. The use of inheritance Interface define common accessor and mutator following the lowerCamelCase convention syntax (i.e getX, setX \dots), attribute name are extratable in a standard way from various objects. We describe three generics type of facts: Knowledge, Prediction and Assertion (Figure~\ref{fig:fact}) .

\begin{figure}
    \centering
    \resizebox{0.8\textwidth}{!}{%
        %LaTeX with PSTricks extensions
%%Creator: inkscape 0.47
%%Please note this file requires PSTricks extensions
\psset{xunit=.5pt,yunit=.5pt,runit=.5pt}
\begin{pspicture}(5493.1640625,5841.796875)
{
\newrgbcolor{curcolor}{0.99607843 0.99607843 0.80784315}
\pscustom[linestyle=none,fillstyle=solid,fillcolor=curcolor]
{
\newpath
\moveto(8.74478125,477.5390625)
\lineto(1122.02603125,477.5390625)
\lineto(1122.02603125,8.7890625)
\lineto(8.74478125,8.7890625)
\lineto(8.74478125,477.5390625)
\closepath
}
}
{
\newrgbcolor{curcolor}{0.65882355 0 0.21176471}
\pscustom[linewidth=17.578125,linecolor=curcolor]
{
\newpath
\moveto(8.74478125,477.5390625)
\lineto(1122.02603125,477.5390625)
\lineto(1122.02603125,8.7890625)
\lineto(8.74478125,8.7890625)
\lineto(8.74478125,477.5390625)
\closepath
}
}
{
\newrgbcolor{curcolor}{0.70588237 0.65490198 0.89803922}
\pscustom[linestyle=none,fillstyle=solid,fillcolor=curcolor]
{
\newpath
\moveto(313.43221875,290.039125)
\curveto(313.43221875,218.84616897)(255.71892478,161.132875)(184.52596875,161.132875)
\curveto(113.33301272,161.132875)(55.61971875,218.84616897)(55.61971875,290.039125)
\curveto(55.61971875,361.23208103)(113.33301272,418.945375)(184.52596875,418.945375)
\curveto(255.71892478,418.945375)(313.43221875,361.23208103)(313.43221875,290.039125)
\closepath
}
}
{
\newrgbcolor{curcolor}{0.65882355 0 0.21176471}
\pscustom[linewidth=11.71875,linecolor=curcolor]
{
\newpath
\moveto(313.43221875,290.039125)
\curveto(313.43221875,218.84616897)(255.71892478,161.132875)(184.52596875,161.132875)
\curveto(113.33301272,161.132875)(55.61971875,218.84616897)(55.61971875,290.039125)
\curveto(55.61971875,361.23208103)(113.33301272,418.945375)(184.52596875,418.945375)
\curveto(255.71892478,418.945375)(313.43221875,361.23208103)(313.43221875,290.039125)
\closepath
}
}
{
\newrgbcolor{curcolor}{0 0 0}
\pscustom[linestyle=none,fillstyle=solid,fillcolor=curcolor]
{
\newpath
\moveto(137.83415625,219.726625)
\lineto(137.83415625,238.95275)
\lineto(171.52553125,238.95275)
\lineto(171.52553125,344.6045)
\lineto(137.83415625,344.6045)
\lineto(137.83415625,363.6475)
\lineto(233.96453125,363.6475)
\lineto(233.96453125,344.6045)
\lineto(200.27315625,344.6045)
\lineto(200.27315625,238.95275)
\lineto(233.96453125,238.95275)
\lineto(233.96453125,219.726625)
\lineto(137.83415625,219.726625)
\closepath
}
}
{
\newrgbcolor{curcolor}{0 0 0}
\pscustom[linestyle=none,fillstyle=solid,fillcolor=curcolor]
{
\newpath
\moveto(412.37789893,252.6171875)
\lineto(407.98336768,278.8470459)
\lineto(369.32522559,278.8470459)
\lineto(354.83700537,252.6171875)
\lineto(341.37875342,252.6171875)
\lineto(394.38778662,345.2456665)
\lineto(409.28799414,345.2456665)
\lineto(425.69882178,252.6171875)
\lineto(412.37789893,252.6171875)
\moveto(402.21554541,318.12316895)
\curveto(401.94083385,319.86260365)(401.66617593,321.62499203)(401.3915708,323.41033936)
\curveto(401.11686007,325.1955451)(400.8650903,326.88926899)(400.63626074,328.49151611)
\curveto(400.45310341,330.09361149)(400.26999812,331.48978929)(400.08694434,332.68005371)
\curveto(399.94956387,333.91593433)(399.83512307,334.80857259)(399.74362158,335.35797119)
\curveto(399.4689125,334.80857259)(399.01114928,333.91593433)(398.37033057,332.68005371)
\curveto(397.77518861,331.44401297)(397.06565562,330.02494701)(396.24172949,328.42285156)
\curveto(395.41770805,326.82060451)(394.52506979,325.12688062)(393.56381201,323.3416748)
\curveto(392.64824061,321.55632755)(391.73271418,319.81682733)(390.81722998,318.12316895)
\lineto(374.81838965,289.07806396)
\lineto(406.61007666,289.07806396)
\lineto(402.21554541,318.12316895)
}
}
{
\newrgbcolor{curcolor}{0 0 0}
\pscustom[linestyle=none,fillstyle=solid,fillcolor=curcolor]
{
\newpath
\moveto(502.60311865,273.83453369)
\curveto(501.91640414,270.35551205)(500.63466714,267.21983403)(498.75790381,264.42749023)
\curveto(496.88100878,261.63512282)(494.45486375,259.2547541)(491.47946143,257.28637695)
\curveto(488.54971828,255.36376678)(485.11649418,253.87603634)(481.17977881,252.82318115)
\curveto(477.2429669,251.77032555)(472.8713282,251.24389786)(468.06484961,251.24389648)
\curveto(463.76184024,251.24389786)(459.9395174,251.56433211)(456.59786963,252.2052002)
\curveto(453.25617448,252.80029279)(450.41804255,253.80737186)(448.08346533,255.22644043)
\curveto(445.74885777,256.64550378)(443.94069307,258.52233296)(442.65896582,260.85693359)
\curveto(441.37721907,263.23729407)(440.64479793,266.16697864)(440.4617002,269.64599609)
\lineto(451.79135107,271.77459717)
\curveto(451.88288483,269.25688033)(452.38642437,267.24272219)(453.30197119,265.7321167)
\curveto(454.21747723,264.2672613)(455.4534379,263.14574142)(457.00985693,262.36755371)
\curveto(458.61200408,261.5893465)(460.48883326,261.0629188)(462.6403501,260.78826904)
\curveto(464.83758379,260.55937926)(467.24084067,260.44493846)(469.85012793,260.44494629)
\curveto(472.50511763,260.44493846)(474.99992715,260.62804375)(477.33456396,260.9942627)
\curveto(479.66911193,261.40624121)(481.75193456,262.06999787)(483.58303809,262.98553467)
\curveto(485.41404027,263.94682705)(486.94754704,265.15989957)(488.18356299,266.62475586)
\curveto(489.41946839,268.13536046)(490.24344218,269.98930148)(490.65548682,272.18658447)
\curveto(490.79275803,272.96476237)(490.86142252,273.65140719)(490.86148047,274.246521)
\curveto(490.86142252,276.0317759)(490.33499482,277.49661819)(489.2821958,278.64105225)
\curveto(488.22928403,279.83121058)(486.78732991,280.83828965)(484.9563291,281.66229248)
\curveto(483.17100052,282.48623722)(481.04240157,283.21865836)(478.57052588,283.85955811)
\curveto(476.14433518,284.54617168)(473.5121967,285.27859283)(470.67410254,286.05682373)
\curveto(468.1563671,286.74343511)(465.66155758,287.47585625)(463.1896665,288.25408936)
\curveto(460.76349119,289.0780275)(458.58911592,290.08510658)(456.66653418,291.27532959)
\curveto(454.74390492,292.51125161)(453.18750999,294.02187022)(451.99734473,295.80718994)
\curveto(450.80714128,297.63819961)(450.2120491,299.90412752)(450.21206641,302.60498047)
\curveto(450.2120491,303.84089116)(450.34937807,305.16840448)(450.62405371,306.58752441)
\curveto(451.2649045,309.97491823)(452.45508885,312.99615544)(454.19461035,315.65124512)
\curveto(455.9798656,318.30620872)(458.24579351,320.54924847)(460.99240088,322.38037109)
\curveto(463.73895208,324.21135419)(466.96618274,325.60753199)(470.67410254,326.56890869)
\curveto(474.42772313,327.57591381)(478.59336838,328.07945335)(483.17105078,328.07952881)
\curveto(486.78732991,328.07945335)(490.10611321,327.73613094)(493.12741064,327.04956055)
\curveto(496.14858764,326.40861761)(498.75783796,325.35576222)(500.95516943,323.89099121)
\curveto(503.15236481,322.47185397)(504.86897687,320.61791296)(506.10501074,318.3291626)
\curveto(507.34089822,316.08605714)(507.95887856,313.38525417)(507.95895361,310.22674561)
\lineto(496.56063818,308.85345459)
\curveto(496.60635085,310.91333282)(496.24014028,312.58416855)(495.46200537,313.8659668)
\curveto(494.68374535,315.14764255)(493.6080018,316.15472162)(492.23477148,316.88720703)
\curveto(490.86142252,317.66534023)(489.23636311,318.19176792)(487.35958838,318.4664917)
\curveto(485.48270475,318.74108378)(483.49143477,318.87841274)(481.38577246,318.878479)
\curveto(470.17052524,318.87841274)(463.87628104,315.42230048)(462.503021,308.51013184)
\curveto(462.36566244,307.7776548)(462.29699796,307.09100998)(462.29702734,306.45019531)
\curveto(462.29699796,304.80219391)(462.73187301,303.42890427)(463.60165381,302.33032227)
\curveto(464.51714954,301.27741716)(465.77599838,300.38477889)(467.3782041,299.65240479)
\curveto(468.98034088,298.91993661)(470.88005822,298.25617995)(473.07736182,297.66113281)
\curveto(475.27458508,297.11177191)(477.65495379,296.49379157)(480.2184751,295.80718994)
\curveto(482.78190179,295.12050193)(485.41404027,294.34230447)(488.11489844,293.47259521)
\curveto(490.81564619,292.64858057)(493.26467939,291.57283702)(495.46200537,290.24536133)
\curveto(497.65920625,288.91781038)(499.46737094,287.24697465)(500.88650488,285.23284912)
\curveto(502.30550287,283.26443468)(503.01503585,280.76962517)(503.01510596,277.74841309)
\curveto(503.01503585,276.51242728)(502.87770689,275.20780212)(502.60311865,273.83453369)
}
}
{
\newrgbcolor{curcolor}{0 0 0}
\pscustom[linestyle=none,fillstyle=solid,fillcolor=curcolor]
{
\newpath
\moveto(586.92318701,273.83453369)
\curveto(586.2364725,270.35551205)(584.9547355,267.21983403)(583.07797217,264.42749023)
\curveto(581.20107714,261.63512282)(578.77493211,259.2547541)(575.79952979,257.28637695)
\curveto(572.86978664,255.36376678)(569.43656254,253.87603634)(565.49984717,252.82318115)
\curveto(561.56303526,251.77032555)(557.19139656,251.24389786)(552.38491797,251.24389648)
\curveto(548.0819086,251.24389786)(544.25958576,251.56433211)(540.91793799,252.2052002)
\curveto(537.57624284,252.80029279)(534.73811091,253.80737186)(532.40353369,255.22644043)
\curveto(530.06892612,256.64550378)(528.26076143,258.52233296)(526.97903418,260.85693359)
\curveto(525.69728743,263.23729407)(524.96486629,266.16697864)(524.78176855,269.64599609)
\lineto(536.11141943,271.77459717)
\curveto(536.20295319,269.25688033)(536.70649273,267.24272219)(537.62203955,265.7321167)
\curveto(538.53754559,264.2672613)(539.77350626,263.14574142)(541.32992529,262.36755371)
\curveto(542.93207244,261.5893465)(544.80890162,261.0629188)(546.96041846,260.78826904)
\curveto(549.15765215,260.55937926)(551.56090903,260.44493846)(554.17019629,260.44494629)
\curveto(556.82518599,260.44493846)(559.31999551,260.62804375)(561.65463232,260.9942627)
\curveto(563.98918029,261.40624121)(566.07200291,262.06999787)(567.90310645,262.98553467)
\curveto(569.73410863,263.94682705)(571.26761539,265.15989957)(572.50363135,266.62475586)
\curveto(573.73953675,268.13536046)(574.56351054,269.98930148)(574.97555518,272.18658447)
\curveto(575.11282639,272.96476237)(575.18149088,273.65140719)(575.18154883,274.246521)
\curveto(575.18149088,276.0317759)(574.65506318,277.49661819)(573.60226416,278.64105225)
\curveto(572.54935239,279.83121058)(571.10739827,280.83828965)(569.27639746,281.66229248)
\curveto(567.49106888,282.48623722)(565.36246993,283.21865836)(562.89059424,283.85955811)
\curveto(560.46440354,284.54617168)(557.83226506,285.27859283)(554.9941709,286.05682373)
\curveto(552.47643546,286.74343511)(549.98162594,287.47585625)(547.50973486,288.25408936)
\curveto(545.08355955,289.0780275)(542.90918428,290.08510658)(540.98660254,291.27532959)
\curveto(539.06397328,292.51125161)(537.50757835,294.02187022)(536.31741309,295.80718994)
\curveto(535.12720964,297.63819961)(534.53211746,299.90412752)(534.53213477,302.60498047)
\curveto(534.53211746,303.84089116)(534.66944643,305.16840448)(534.94412207,306.58752441)
\curveto(535.58497285,309.97491823)(536.77515721,312.99615544)(538.51467871,315.65124512)
\curveto(540.29993396,318.30620872)(542.56586187,320.54924847)(545.31246924,322.38037109)
\curveto(548.05902044,324.21135419)(551.2862511,325.60753199)(554.9941709,326.56890869)
\curveto(558.74779149,327.57591381)(562.91343674,328.07945335)(567.49111914,328.07952881)
\curveto(571.10739827,328.07945335)(574.42618157,327.73613094)(577.447479,327.04956055)
\curveto(580.468656,326.40861761)(583.07790632,325.35576222)(585.27523779,323.89099121)
\curveto(587.47243317,322.47185397)(589.18904523,320.61791296)(590.4250791,318.3291626)
\curveto(591.66096658,316.08605714)(592.27894692,313.38525417)(592.27902197,310.22674561)
\lineto(580.88070654,308.85345459)
\curveto(580.92641921,310.91333282)(580.56020864,312.58416855)(579.78207373,313.8659668)
\curveto(579.00381371,315.14764255)(577.92807016,316.15472162)(576.55483984,316.88720703)
\curveto(575.18149088,317.66534023)(573.55643147,318.19176792)(571.67965674,318.4664917)
\curveto(569.80277311,318.74108378)(567.81150313,318.87841274)(565.70584082,318.878479)
\curveto(554.4905936,318.87841274)(548.1963494,315.42230048)(546.82308936,308.51013184)
\curveto(546.6857308,307.7776548)(546.61706631,307.09100998)(546.6170957,306.45019531)
\curveto(546.61706631,304.80219391)(547.05194137,303.42890427)(547.92172217,302.33032227)
\curveto(548.8372179,301.27741716)(550.09606674,300.38477889)(551.69827246,299.65240479)
\curveto(553.30040924,298.91993661)(555.20012658,298.25617995)(557.39743018,297.66113281)
\curveto(559.59465343,297.11177191)(561.97502215,296.49379157)(564.53854346,295.80718994)
\curveto(567.10197015,295.12050193)(569.73410863,294.34230447)(572.4349668,293.47259521)
\curveto(575.13571455,292.64858057)(577.58474775,291.57283702)(579.78207373,290.24536133)
\curveto(581.9792746,288.91781038)(583.7874393,287.24697465)(585.20657324,285.23284912)
\curveto(586.62557123,283.26443468)(587.33510421,280.76962517)(587.33517432,277.74841309)
\curveto(587.33510421,276.51242728)(587.19777525,275.20780212)(586.92318701,273.83453369)
}
}
{
\newrgbcolor{curcolor}{0 0 0}
\pscustom[linestyle=none,fillstyle=solid,fillcolor=curcolor]
{
\newpath
\moveto(623.17806982,287.15545654)
\curveto(622.90339027,285.69057972)(622.67450866,284.27151376)(622.49142432,282.89825439)
\curveto(622.35407441,281.52493447)(622.28540993,280.17453299)(622.28543066,278.8470459)
\curveto(622.28540993,276.10044038)(622.60584418,273.58274271)(623.24673437,271.29394531)
\curveto(623.9333575,269.05088689)(624.96332473,267.12828139)(626.33663916,265.52612305)
\curveto(627.70990401,263.92393889)(629.40362791,262.68797821)(631.41781592,261.8182373)
\curveto(633.47772051,260.948478)(635.88097739,260.51360294)(638.62759375,260.51361084)
\curveto(641.14525435,260.51360294)(643.52562306,260.78826087)(645.76870703,261.33758545)
\curveto(648.01170256,261.93266891)(650.07163702,262.73375453)(651.9485166,263.74084473)
\curveto(653.8710717,264.74791267)(655.54190743,265.93809703)(656.96102881,267.31140137)
\curveto(658.38003936,268.73045264)(659.5244474,270.2639594)(660.39425635,271.91192627)
\lineto(670.62527441,268.82202148)
\curveto(669.4807973,266.76207082)(667.9701787,264.67924819)(666.09341406,262.57354736)
\curveto(664.21652034,260.51360294)(661.90481611,258.63677377)(659.15829443,256.9430542)
\curveto(656.41165754,255.2951023)(653.18442688,253.92181266)(649.47659277,252.82318115)
\curveto(645.81443913,251.77032555)(641.60301756,251.24389786)(636.84231543,251.24389648)
\curveto(632.4935296,251.24389786)(628.62543044,251.79321371)(625.23800635,252.8918457)
\curveto(621.85053487,254.03625346)(618.98951479,255.72997736)(616.6549375,257.97302246)
\curveto(614.32033,260.21605685)(612.53505347,263.00841246)(611.29910254,266.35009766)
\curveto(610.10890843,269.69175539)(609.51381625,273.58274271)(609.51382422,278.02307129)
\curveto(609.51381625,281.82248056)(609.94869131,285.89657317)(610.81845068,290.24536133)
\curveto(612.19173106,297.38642984)(614.25166552,303.38312794)(616.99826025,308.23547363)
\curveto(619.74482409,313.08770808)(622.92627843,317.00158356)(626.54263281,319.97711182)
\curveto(630.20471353,322.95250535)(634.16436533,325.08110429)(638.4216001,326.36291504)
\curveto(642.67876112,327.64457829)(647.00462349,328.28544679)(651.3992002,328.28552246)
\curveto(656.15988777,328.28544679)(660.21109222,327.66746645)(663.55282568,326.43157959)
\curveto(666.89443514,325.24132142)(669.61812627,323.50182121)(671.72390723,321.21307373)
\curveto(673.87532416,318.96996539)(675.43171909,316.24627426)(676.39309668,313.04199219)
\curveto(677.40010091,309.83758926)(677.90364044,306.24414803)(677.9037168,302.26165771)
\curveto(677.90364044,298.09596282)(677.42298907,293.60988332)(676.46176123,288.80340576)
\lineto(676.18710303,287.15545654)
\lineto(623.17806982,287.15545654)
\moveto(649.75125098,319.15313721)
\curveto(647.64549199,319.15307067)(645.37956408,318.80974826)(642.95346045,318.12316895)
\curveto(640.52727401,317.48223494)(638.16979346,316.31493874)(635.88101172,314.62127686)
\curveto(633.59216132,312.92749096)(631.48645053,310.61578673)(629.56387305,307.68615723)
\curveto(627.68701585,304.80219391)(626.22217357,301.11720004)(625.1693418,296.63116455)
\lineto(665.26943945,296.63116455)
\curveto(665.36092838,297.54664697)(665.42959286,298.39350891)(665.47543311,299.17175293)
\curveto(665.5211455,299.99568016)(665.54403366,300.79676579)(665.54409766,301.57501221)
\curveto(665.54403366,304.91663471)(665.13204677,307.70899032)(664.30813574,309.9520874)
\curveto(663.4840992,312.19506982)(662.36257932,314.00323451)(660.94357275,315.37658691)
\curveto(659.5244474,316.7498138)(657.85361166,317.71111655)(655.93106055,318.26049805)
\curveto(654.00840067,318.85552458)(651.9484662,319.15307067)(649.75125098,319.15313721)
}
}
{
\newrgbcolor{curcolor}{0 0 0}
\pscustom[linestyle=none,fillstyle=solid,fillcolor=curcolor]
{
\newpath
\moveto(762.63577246,315.65124512)
\curveto(760.57576123,316.01739265)(758.3784978,316.31493874)(756.04397559,316.54388428)
\curveto(753.75508934,316.81847828)(751.19161534,316.95580724)(748.3535459,316.95587158)
\curveto(744.50827242,316.95580724)(740.89194302,316.17760978)(737.50454687,314.62127686)
\curveto(734.16282378,313.06481992)(731.16447472,310.93622098)(728.50949072,308.23547363)
\curveto(725.90019776,305.58039137)(723.68004617,302.46760152)(721.8490293,298.89709473)
\curveto(720.06371678,295.32649538)(718.78197978,291.5270607)(718.00381445,287.4987793)
\lineto(711.20602393,252.6171875)
\lineto(698.84640479,252.6171875)
\lineto(708.18478369,300.7510376)
\curveto(708.82562988,304.22998989)(709.26050493,307.70899032)(709.48941016,311.18804932)
\curveto(709.76404446,314.71276749)(709.90137343,317.84844551)(709.90139746,320.59509277)
\curveto(709.90137343,321.92253812)(709.87848527,323.11272247)(709.83273291,324.16564941)
\curveto(709.78693262,325.26420958)(709.7411563,326.17973601)(709.69540381,326.91223145)
\lineto(721.43704199,326.91223145)
\curveto(721.57433539,325.76774911)(721.66588803,324.41734763)(721.7117002,322.86102295)
\curveto(721.803217,321.3503341)(721.84899332,319.77105101)(721.8490293,318.12316895)
\curveto(721.84899332,316.65826115)(721.803217,315.21630703)(721.7117002,313.79730225)
\curveto(721.66588803,312.3781751)(721.57433539,311.02777362)(721.43704199,309.74609375)
\lineto(721.78036475,309.74609375)
\curveto(723.61138169,312.58416855)(725.44243455,315.14764255)(727.27352881,317.43652344)
\curveto(729.15031658,319.72527469)(731.21025105,321.67076835)(733.45333838,323.27301025)
\curveto(735.74210686,324.87511085)(738.2826927,326.11107153)(741.07510352,326.980896)
\curveto(743.91318024,327.85057174)(747.20907538,328.28544679)(750.96279883,328.28552246)
\curveto(753.52620773,328.28544679)(756.02101725,328.17100599)(758.44723486,327.94219971)
\curveto(760.91908364,327.71324277)(763.07057075,327.43858485)(764.90170264,327.1182251)
\lineto(762.63577246,315.65124512)
}
}
{
\newrgbcolor{curcolor}{0 0 0}
\pscustom[linestyle=none,fillstyle=solid,fillcolor=curcolor]
{
\newpath
\moveto(788.59097266,317.16186523)
\lineto(790.44491553,326.91223145)
\lineto(802.11788916,326.91223145)
\lineto(809.8769834,346.27563477)
\lineto(818.11672949,346.27563477)
\lineto(814.3401792,326.91223145)
\lineto(844.00326514,326.91223145)
\lineto(842.14932227,317.16186523)
\lineto(812.48623633,317.16186523)
\lineto(803.76583838,272.39257812)
\curveto(803.67425216,271.84324249)(803.58269952,271.3168148)(803.49118018,270.81329346)
\curveto(803.44537055,270.30973573)(803.42248239,269.82908435)(803.42251562,269.37133789)
\curveto(803.42248239,266.76207082)(804.2922325,264.83946532)(806.03176855,263.60351562)
\curveto(807.77123293,262.36754396)(810.49492405,261.74956362)(814.2028501,261.74957275)
\curveto(816.03385894,261.74956362)(817.93357628,261.8182281)(819.90200781,261.95556641)
\curveto(821.87033992,262.13866235)(823.77005726,262.3446558)(825.60116553,262.57354736)
\curveto(827.47793929,262.80241901)(829.21743951,263.05418878)(830.81967139,263.32885742)
\curveto(832.46755833,263.60350464)(833.86373613,263.85527441)(835.00820898,264.08416748)
\lineto(833.22293066,254.67712402)
\curveto(832.17001224,254.40246404)(830.7967226,254.08202979)(829.10305762,253.71582031)
\curveto(827.45505113,253.34960864)(825.60111012,253.00628623)(823.541229,252.68585205)
\curveto(821.48124119,252.36541773)(819.2610896,252.0907598)(816.88076758,251.86187744)
\curveto(814.50035217,251.63299659)(812.07420714,251.51855579)(809.6023252,251.51855469)
\curveto(803.23937711,251.51855579)(798.47863968,252.64007566)(795.32009863,254.88311768)
\curveto(792.16150732,257.17193148)(790.58222424,260.74248455)(790.58224463,265.5947876)
\curveto(790.58222424,266.41874841)(790.6051124,267.28849851)(790.65090918,268.20404053)
\curveto(790.74244136,269.11955137)(790.90265849,270.08085412)(791.13156104,271.08795166)
\lineto(800.12661719,317.16186523)
\lineto(788.59097266,317.16186523)
}
}
{
\newrgbcolor{curcolor}{0 0 0}
\pscustom[linestyle=none,fillstyle=solid,fillcolor=curcolor]
{
\newpath
\moveto(900.37686133,262.36755371)
\lineto(926.46939062,262.36755371)
\lineto(924.5467832,252.6171875)
\lineto(857.11819434,252.6171875)
\lineto(859.04080176,262.36755371)
\lineto(888.01724219,262.36755371)
\lineto(898.66024756,317.16186523)
\lineto(876.75625586,317.16186523)
\lineto(878.61019873,326.91223145)
\lineto(912.87380957,326.91223145)
\lineto(900.37686133,262.36755371)
\moveto(901.956146,341.33178711)
\lineto(904.49673437,354.51538086)
\lineto(918.22964453,354.51538086)
\lineto(915.68905615,341.33178711)
\lineto(901.956146,341.33178711)
}
}
{
\newrgbcolor{curcolor}{0 0 0}
\pscustom[linestyle=none,fillstyle=solid,fillcolor=curcolor]
{
\newpath
\moveto(1014.15402197,289.83337402)
\curveto(1012.91798597,283.33309917)(1011.06404495,277.65683531)(1008.59219336,272.80456543)
\curveto(1006.12020224,267.99803149)(1003.12185319,263.99260337)(999.59713721,260.78826904)
\curveto(996.11807601,257.58391837)(992.13553605,255.1806615)(987.64950537,253.57849121)
\curveto(983.20915337,252.02209532)(978.37975146,251.24389786)(973.16128516,251.24389648)
\curveto(969.0413819,251.24389786)(965.33349986,251.77032555)(962.03762793,252.82318115)
\curveto(958.74170958,253.92181266)(955.92646581,255.56976023)(953.59188818,257.76702881)
\curveto(951.30305735,259.96428709)(949.54066898,262.73375453)(948.30471777,266.07543945)
\curveto(947.11452394,269.41709746)(946.51943176,273.33097294)(946.51943945,277.81707764)
\curveto(946.51943176,281.43338183)(946.93141866,285.43880995)(947.75540137,289.83337402)
\curveto(950.27309012,302.60493048)(955.03382755,312.19506982)(962.03762793,318.6038208)
\curveto(969.08715822,325.05821613)(977.92198825,328.28544679)(988.54214453,328.28552246)
\curveto(997.65158278,328.28544679)(1004.42647835,326.17973601)(1008.86685156,321.96838379)
\curveto(1013.3070847,317.75689287)(1015.52723629,311.39398419)(1015.52731299,302.87963867)
\curveto(1015.52723629,300.91120659)(1015.41279549,298.85127213)(1015.18399023,296.6998291)
\curveto(1015.00080859,294.54829791)(1014.65748618,292.25948184)(1014.15402197,289.83337402)
\moveto(1001.17642187,289.83337402)
\curveto(1001.63412274,292.21370552)(1001.97744515,294.41096895)(1002.20639014,296.4251709)
\curveto(1002.43520837,298.43928523)(1002.54964917,300.31611441)(1002.54971289,302.05566406)
\curveto(1002.54964917,305.35150977)(1002.16055044,308.07520089)(1001.38241553,310.22674561)
\curveto(1000.64993183,312.42395142)(999.59707644,314.1863398)(998.22384619,315.51391602)
\curveto(996.85049715,316.84136644)(995.20254958,317.77978103)(993.27999854,318.3291626)
\curveto(991.40311491,318.87841274)(989.29740412,319.15307067)(986.96285986,319.15313721)
\curveto(983.98735084,319.15307067)(981.14921891,318.6724193)(978.44845557,317.71118164)
\curveto(975.74761298,316.79559012)(973.25280347,315.19341887)(970.96401953,312.90466309)
\curveto(968.72094765,310.66156305)(966.72967767,307.663214)(964.99020361,303.90960693)
\curveto(963.25067724,300.15589729)(961.83161127,295.46382434)(960.73300146,289.83337402)
\curveto(960.27521635,287.59029706)(959.9547821,285.50747443)(959.77169775,283.5848999)
\curveto(959.58857153,281.66226344)(959.49701888,279.8769869)(959.49703955,278.22906494)
\curveto(959.49701888,274.79581522)(959.88611761,271.93479514)(960.66433691,269.64599609)
\curveto(961.48828886,267.35716299)(962.58692058,265.52611014)(963.96023535,264.15283203)
\curveto(965.37927618,262.77953085)(967.00433559,261.79533994)(968.83541846,261.20025635)
\curveto(970.71221763,260.65093191)(972.70348761,260.37627398)(974.80923437,260.37628174)
\curveto(977.78465929,260.37627398)(980.62279121,260.81114903)(983.32363867,261.6809082)
\curveto(986.02439714,262.59642557)(988.4963185,264.17570866)(990.73941016,266.41876221)
\curveto(993.02817432,268.66178815)(995.06522062,271.68302537)(996.85055518,275.48248291)
\curveto(998.63577369,279.28189472)(1000.07772781,284.06552031)(1001.17642187,289.83337402)
}
}
{
\newrgbcolor{curcolor}{0 0 0}
\pscustom[linestyle=none,fillstyle=solid,fillcolor=curcolor]
{
\newpath
\moveto(1075.54013037,252.6171875)
\lineto(1084.80984473,300.33905029)
\curveto(1085.3133226,302.99402921)(1085.56509237,305.23706896)(1085.56515479,307.06817627)
\curveto(1085.56509237,309.35693789)(1085.26754628,311.23376707)(1084.67251562,312.69866943)
\curveto(1084.07736192,314.20922796)(1083.23049998,315.39941232)(1082.13192725,316.26922607)
\curveto(1081.03323655,317.18468885)(1079.72861139,317.82555735)(1078.21804785,318.1918335)
\curveto(1076.70737418,318.55797849)(1075.05942661,318.74108378)(1073.2742002,318.7411499)
\curveto(1070.57334711,318.74108378)(1067.94120863,318.21465608)(1065.37777686,317.16186523)
\curveto(1062.81426063,316.15472162)(1060.45678008,314.66699117)(1058.30532812,312.69866943)
\curveto(1056.19958218,310.77600385)(1054.36852933,308.37274698)(1052.81216406,305.4888916)
\curveto(1051.25573947,302.6507068)(1050.11133144,299.37769982)(1049.37893652,295.66986084)
\lineto(1041.00186133,252.6171875)
\lineto(1028.64224219,252.6171875)
\lineto(1039.97189307,311.05072021)
\curveto(1040.29231049,312.60705671)(1040.58985658,314.23211612)(1040.86453223,315.92590332)
\curveto(1041.18494876,317.66534023)(1041.45960669,319.29039964)(1041.68850684,320.80108643)
\curveto(1041.96314622,322.35741317)(1042.19202783,323.68492649)(1042.37515234,324.78363037)
\curveto(1042.5582384,325.88218992)(1042.64979105,326.5917229)(1042.64981055,326.91223145)
\lineto(1054.32278418,326.91223145)
\curveto(1054.32275301,326.68327554)(1054.25408852,326.0652952)(1054.11679053,325.05828857)
\curveto(1053.9794306,324.05113706)(1053.81921347,322.88384087)(1053.63613867,321.55639648)
\curveto(1053.4530029,320.27459055)(1053.24700945,318.94707722)(1053.01815771,317.57385254)
\curveto(1052.78924624,316.24627426)(1052.58325279,315.12475439)(1052.40017676,314.20928955)
\lineto(1052.60617041,314.20928955)
\curveto(1055.71893081,318.55797849)(1059.40392468,321.9912026)(1063.66116309,324.50897217)
\curveto(1067.91832047,327.02659795)(1073.06815662,328.28544679)(1079.11068701,328.28552246)
\curveto(1085.38198708,328.28544679)(1090.18850083,326.88926899)(1093.53024268,324.09698486)
\curveto(1096.87184376,321.3503341)(1098.54267949,316.93291908)(1098.54275488,310.84472656)
\curveto(1098.54267949,309.51715501)(1098.45112685,308.12097721)(1098.26809668,306.65618896)
\curveto(1098.1306926,305.23706896)(1097.90181099,303.72645036)(1097.58145117,302.12432861)
\lineto(1087.96841406,252.6171875)
\lineto(1075.54013037,252.6171875)
}
}
{
\newrgbcolor{curcolor}{0 0 0}
\pscustom[linestyle=none,fillstyle=solid,fillcolor=curcolor]
{
\newpath
\moveto(20.46353125,102.5390625)
\lineto(1110.30728125,102.5390625)
}
}
{
\newrgbcolor{curcolor}{0.65882355 0 0.21176471}
\pscustom[linewidth=17.578125,linecolor=curcolor]
{
\newpath
\moveto(20.46353125,102.5390625)
\lineto(1110.30728125,102.5390625)
}
}
{
\newrgbcolor{curcolor}{0.99607843 0.99607843 0.80784315}
\pscustom[linestyle=none,fillstyle=solid,fillcolor=curcolor]
{
\newpath
\moveto(243.11978125,2446.2890625)
\lineto(2457.96353125,2446.2890625)
\lineto(2457.96353125,1677.55508423)
\lineto(243.11978125,1677.55508423)
\lineto(243.11978125,2446.2890625)
\closepath
}
}
{
\newrgbcolor{curcolor}{0.65882355 0 0.21176471}
\pscustom[linewidth=17.578125,linecolor=curcolor]
{
\newpath
\moveto(243.11978125,2446.2890625)
\lineto(2457.96353125,2446.2890625)
\lineto(2457.96353125,1677.55508423)
\lineto(243.11978125,1677.55508423)
\lineto(243.11978125,2446.2890625)
\closepath
}
}
{
\newrgbcolor{curcolor}{0.70588237 0.65490198 0.89803922}
\pscustom[linestyle=none,fillstyle=solid,fillcolor=curcolor]
{
\newpath
\moveto(978.47128125,2258.789125)
\curveto(978.47128125,2187.59616897)(920.75798728,2129.882875)(849.56503125,2129.882875)
\curveto(778.37207522,2129.882875)(720.65878125,2187.59616897)(720.65878125,2258.789125)
\curveto(720.65878125,2329.98208103)(778.37207522,2387.695375)(849.56503125,2387.695375)
\curveto(920.75798728,2387.695375)(978.47128125,2329.98208103)(978.47128125,2258.789125)
\closepath
}
}
{
\newrgbcolor{curcolor}{0.65882355 0 0.21176471}
\pscustom[linewidth=11.71875,linecolor=curcolor]
{
\newpath
\moveto(978.47128125,2258.789125)
\curveto(978.47128125,2187.59616897)(920.75798728,2129.882875)(849.56503125,2129.882875)
\curveto(778.37207522,2129.882875)(720.65878125,2187.59616897)(720.65878125,2258.789125)
\curveto(720.65878125,2329.98208103)(778.37207522,2387.695375)(849.56503125,2387.695375)
\curveto(920.75798728,2387.695375)(978.47128125,2329.98208103)(978.47128125,2258.789125)
\closepath
}
}
{
\newrgbcolor{curcolor}{0 0 0}
\pscustom[linestyle=none,fillstyle=solid,fillcolor=curcolor]
{
\newpath
\moveto(802.87315625,2188.4765)
\lineto(802.87315625,2207.70275)
\lineto(836.56465625,2207.70275)
\lineto(836.56465625,2313.354375)
\lineto(802.87315625,2313.354375)
\lineto(802.87315625,2332.397375)
\lineto(899.00353125,2332.397375)
\lineto(899.00353125,2313.354375)
\lineto(865.31215625,2313.354375)
\lineto(865.31215625,2207.70275)
\lineto(899.00353125,2207.70275)
\lineto(899.00353125,2188.4765)
\lineto(802.87315625,2188.4765)
\closepath
}
}
{
\newrgbcolor{curcolor}{0 0 0}
\pscustom[linestyle=none,fillstyle=solid,fillcolor=curcolor]
{
\newpath
\moveto(1168.9696958,2268.12774658)
\curveto(1167.45899802,2260.20839622)(1165.30751092,2253.25039536)(1162.51522803,2247.25372314)
\curveto(1159.76857603,2241.25699915)(1156.47268089,2236.2216038)(1152.62753271,2232.14752197)
\curveto(1148.82803521,2228.11919491)(1144.54794916,2225.07506953)(1139.78726172,2223.01513672)
\curveto(1135.0264743,2221.00097693)(1129.89952631,2219.99389786)(1124.40640234,2219.99389648)
\curveto(1115.15955081,2219.99389786)(1108.24732628,2222.6489245)(1103.66970801,2227.95898438)
\curveto(1099.092062,2233.31480739)(1096.80324593,2240.95945306)(1096.80325293,2250.89294434)
\curveto(1096.80324593,2256.06563913)(1097.42122626,2261.81056747)(1098.6571958,2268.12774658)
\curveto(1100.21358187,2276.04700343)(1102.34218082,2282.95922796)(1105.04299902,2288.86444092)
\curveto(1107.74378674,2294.76951888)(1110.94812924,2299.69047343)(1114.65603613,2303.62731934)
\curveto(1118.40966963,2307.56400072)(1122.66686752,2310.51657345)(1127.42764258,2312.48504639)
\curveto(1132.18834238,2314.45333709)(1137.40684302,2315.437528)(1143.08316016,2315.43762207)
\curveto(1152.14681851,2315.437528)(1159.03615488,2312.89694216)(1163.75118994,2307.81585693)
\curveto(1168.4660771,2302.78037513)(1170.82355765,2295.20439394)(1170.82363867,2285.08789062)
\curveto(1170.82355765,2280.05243155)(1170.20557731,2274.39905585)(1168.9696958,2268.12774658)
\moveto(1155.16812109,2268.12774658)
\curveto(1155.7631479,2271.19471336)(1156.19802296,2274.03284528)(1156.47274756,2276.64215088)
\curveto(1156.79311514,2279.29712225)(1156.95333226,2281.7690436)(1156.95339941,2284.05792236)
\curveto(1156.95333226,2291.01586053)(1155.60293078,2296.18858485)(1152.90219092,2299.57611084)
\curveto(1150.24710117,2303.00925674)(1146.28744937,2304.72586879)(1141.02322363,2304.72595215)
\curveto(1133.88206627,2304.72586879)(1127.93114449,2301.70463158)(1123.17044043,2295.66223145)
\curveto(1118.45544595,2289.66545905)(1114.88489288,2280.4873066)(1112.45877051,2268.12774658)
\curveto(1111.86365567,2265.06068629)(1111.40589246,2262.17677804)(1111.08547949,2259.47601318)
\curveto(1110.81080028,2256.82094843)(1110.67347131,2254.32613892)(1110.67349219,2251.99157715)
\curveto(1110.67347131,2244.89621671)(1112.00098463,2239.56327526)(1114.65603613,2235.99273682)
\curveto(1117.31103792,2232.42216912)(1121.27068972,2230.63689259)(1126.53500342,2230.63690186)
\curveto(1133.99650707,2230.63689259)(1140.0389815,2233.84123508)(1144.66244482,2240.24993896)
\curveto(1149.28579842,2246.65860508)(1152.78768701,2255.95119833)(1155.16812109,2268.12774658)
}
}
{
\newrgbcolor{curcolor}{0 0 0}
\pscustom[linestyle=none,fillstyle=solid,fillcolor=curcolor]
{
\newpath
\moveto(1249.03256201,2258.85803223)
\curveto(1246.51478942,2245.85751945)(1242.30336785,2236.13005116)(1236.39828467,2229.67559814)
\curveto(1230.49307693,2223.22112852)(1223.21464182,2219.99389786)(1214.56295752,2219.99389648)
\curveto(1209.25286379,2219.99389786)(1204.83544878,2220.97808877)(1201.31069922,2222.94647217)
\curveto(1197.8316716,2224.91485241)(1195.42841473,2228.14208307)(1194.10092139,2232.62817383)
\lineto(1193.96359229,2232.62817383)
\curveto(1193.73469083,2231.43797821)(1193.46003291,2230.20201753)(1193.13961768,2228.92028809)
\curveto(1192.86494073,2227.68431985)(1192.5902828,2226.53991182)(1192.31564307,2225.48706055)
\curveto(1192.04096694,2224.43420103)(1191.78919717,2223.51867461)(1191.56033301,2222.74047852)
\curveto(1191.37721028,2222.008056)(1191.23988132,2221.55029279)(1191.1483457,2221.3671875)
\lineto(1179.20071387,2221.3671875)
\curveto(1179.33803775,2221.77917439)(1179.52114304,2222.46581921)(1179.75003027,2223.42712402)
\curveto(1179.97890625,2224.43420103)(1180.23067602,2225.62438539)(1180.50534033,2226.99768066)
\curveto(1180.8257682,2228.37096468)(1181.16909061,2229.88158328)(1181.53530859,2231.52954102)
\curveto(1181.90151175,2233.17747842)(1182.24483416,2234.89409048)(1182.56527686,2236.67938232)
\lineto(1199.3880918,2323.26538086)
\lineto(1211.74771094,2323.26538086)
\lineto(1206.11721777,2294.22027588)
\curveto(1205.88830417,2292.8926897)(1205.61364624,2291.56517638)(1205.29324316,2290.23773193)
\curveto(1205.01855406,2288.95592606)(1204.76678429,2287.81151803)(1204.53793311,2286.80450439)
\curveto(1204.26324476,2285.6142546)(1204.01147499,2284.49273473)(1203.78262305,2283.43994141)
\lineto(1204.05728125,2283.43994141)
\curveto(1207.17004117,2288.24639308)(1210.8779232,2291.72539351)(1215.18093848,2293.87695312)
\curveto(1219.48387163,2296.07414404)(1224.31327354,2297.17277576)(1229.66915869,2297.17285156)
\curveto(1236.627104,2297.17277576)(1241.8684928,2295.2272821)(1245.39334082,2291.33636475)
\curveto(1248.9180463,2287.44530746)(1250.68043467,2281.54016199)(1250.68051123,2273.62091064)
\curveto(1250.68043467,2269.3636605)(1250.13111881,2264.44270595)(1249.03256201,2258.85803223)
\moveto(1236.19229102,2258.44604492)
\curveto(1236.69576848,2261.05525816)(1237.06197905,2263.45851504)(1237.29092383,2265.65582275)
\curveto(1237.56551859,2267.89881821)(1237.70284755,2269.98164084)(1237.70291113,2271.90429688)
\curveto(1237.70284755,2277.39740491)(1236.60421584,2281.37994487)(1234.4070127,2283.85192871)
\curveto(1232.2554653,2286.3695639)(1228.936682,2287.62841274)(1224.45065283,2287.628479)
\curveto(1221.42936529,2287.62841274)(1218.61412152,2287.10198505)(1216.00491309,2286.04919434)
\curveto(1213.39562088,2285.04205058)(1210.99236401,2283.34832669)(1208.79513525,2280.96801758)
\curveto(1206.64361347,2278.63336559)(1204.74389613,2275.52057573)(1203.09597754,2271.62963867)
\curveto(1201.44800099,2267.78437741)(1200.07471135,2263.02363998)(1198.97610449,2257.34741211)
\curveto(1198.51831642,2255.01278374)(1198.17499401,2252.86129663)(1197.94613623,2250.89294434)
\curveto(1197.76300712,2248.92453299)(1197.67145448,2247.11636829)(1197.67147803,2245.46844482)
\curveto(1197.67145448,2242.40140719)(1198.01477689,2239.81504503)(1198.70144629,2237.70935059)
\curveto(1199.43384285,2235.60362346)(1200.46381008,2233.90989957)(1201.79135107,2232.62817383)
\curveto(1203.11883672,2231.39220189)(1204.72100797,2230.49956362)(1206.59786963,2229.95025635)
\curveto(1208.47466633,2229.40093191)(1210.58037711,2229.12627398)(1212.9150083,2229.12628174)
\curveto(1215.79887775,2229.12627398)(1218.43101624,2229.65270168)(1220.81143164,2230.70556641)
\curveto(1223.23752998,2231.75841246)(1225.41190525,2233.45213635)(1227.33456396,2235.78674316)
\curveto(1229.25711625,2238.12132114)(1230.95084014,2241.14255835)(1232.41574072,2244.85046387)
\curveto(1233.92630103,2248.60409874)(1235.18514987,2253.13595456)(1236.19229102,2258.44604492)
}
}
{
\newrgbcolor{curcolor}{0 0 0}
\pscustom[linestyle=none,fillstyle=solid,fillcolor=curcolor]
{
\newpath
\moveto(1328.13412451,2242.58453369)
\curveto(1327.44741,2239.10551205)(1326.165673,2235.96983403)(1324.28890967,2233.17749023)
\curveto(1322.41201464,2230.38512282)(1319.98586961,2228.0047541)(1317.01046729,2226.03637695)
\curveto(1314.08072414,2224.11376678)(1310.64750004,2222.62603634)(1306.71078467,2221.57318115)
\curveto(1302.77397276,2220.52032555)(1298.40233406,2219.99389786)(1293.59585547,2219.99389648)
\curveto(1289.2928461,2219.99389786)(1285.47052326,2220.31433211)(1282.12887549,2220.9552002)
\curveto(1278.78718034,2221.55029279)(1275.94904841,2222.55737186)(1273.61447119,2223.97644043)
\curveto(1271.27986362,2225.39550378)(1269.47169893,2227.27233296)(1268.18997168,2229.60693359)
\curveto(1266.90822493,2231.98729407)(1266.17580379,2234.91697864)(1265.99270605,2238.39599609)
\lineto(1277.32235693,2240.52459717)
\curveto(1277.41389069,2238.00688033)(1277.91743023,2235.99272219)(1278.83297705,2234.4821167)
\curveto(1279.74848309,2233.0172613)(1280.98444376,2231.89574142)(1282.54086279,2231.11755371)
\curveto(1284.14300994,2230.3393465)(1286.01983912,2229.8129188)(1288.17135596,2229.53826904)
\curveto(1290.36858965,2229.30937926)(1292.77184653,2229.19493846)(1295.38113379,2229.19494629)
\curveto(1298.03612349,2229.19493846)(1300.53093301,2229.37804375)(1302.86556982,2229.7442627)
\curveto(1305.20011779,2230.15624121)(1307.28294041,2230.81999787)(1309.11404395,2231.73553467)
\curveto(1310.94504613,2232.69682705)(1312.47855289,2233.90989957)(1313.71456885,2235.37475586)
\curveto(1314.95047425,2236.88536046)(1315.77444804,2238.73930148)(1316.18649268,2240.93658447)
\curveto(1316.32376389,2241.71476237)(1316.39242838,2242.40140719)(1316.39248633,2242.996521)
\curveto(1316.39242838,2244.7817759)(1315.86600068,2246.24661819)(1314.81320166,2247.39105225)
\curveto(1313.76028989,2248.58121058)(1312.31833577,2249.58828965)(1310.48733496,2250.41229248)
\curveto(1308.70200638,2251.23623722)(1306.57340743,2251.96865836)(1304.10153174,2252.60955811)
\curveto(1301.67534104,2253.29617168)(1299.04320256,2254.02859283)(1296.2051084,2254.80682373)
\curveto(1293.68737296,2255.49343511)(1291.19256344,2256.22585625)(1288.72067236,2257.00408936)
\curveto(1286.29449705,2257.8280275)(1284.12012178,2258.83510658)(1282.19754004,2260.02532959)
\curveto(1280.27491078,2261.26125161)(1278.71851585,2262.77187022)(1277.52835059,2264.55718994)
\curveto(1276.33814714,2266.38819961)(1275.74305496,2268.65412752)(1275.74307227,2271.35498047)
\curveto(1275.74305496,2272.59089116)(1275.88038393,2273.91840448)(1276.15505957,2275.33752441)
\curveto(1276.79591035,2278.72491823)(1277.98609471,2281.74615544)(1279.72561621,2284.40124512)
\curveto(1281.51087146,2287.05620872)(1283.77679937,2289.29924847)(1286.52340674,2291.13037109)
\curveto(1289.26995794,2292.96135419)(1292.4971886,2294.35753199)(1296.2051084,2295.31890869)
\curveto(1299.95872899,2296.32591381)(1304.12437424,2296.82945335)(1308.70205664,2296.82952881)
\curveto(1312.31833577,2296.82945335)(1315.63711907,2296.48613094)(1318.6584165,2295.79956055)
\curveto(1321.6795935,2295.15861761)(1324.28884382,2294.10576222)(1326.48617529,2292.64099121)
\curveto(1328.68337067,2291.22185397)(1330.39998273,2289.36791296)(1331.6360166,2287.0791626)
\curveto(1332.87190408,2284.83605714)(1333.48988442,2282.13525417)(1333.48995947,2278.97674561)
\lineto(1322.09164404,2277.60345459)
\curveto(1322.13735671,2279.66333282)(1321.77114614,2281.33416855)(1320.99301123,2282.6159668)
\curveto(1320.21475121,2283.89764255)(1319.13900766,2284.90472162)(1317.76577734,2285.63720703)
\curveto(1316.39242838,2286.41534023)(1314.76736897,2286.94176792)(1312.89059424,2287.2164917)
\curveto(1311.01371061,2287.49108378)(1309.02244063,2287.62841274)(1306.91677832,2287.628479)
\curveto(1295.7015311,2287.62841274)(1289.4072869,2284.17230048)(1288.03402686,2277.26013184)
\curveto(1287.8966683,2276.5276548)(1287.82800381,2275.84100998)(1287.8280332,2275.20019531)
\curveto(1287.82800381,2273.55219391)(1288.26287887,2272.17890427)(1289.13265967,2271.08032227)
\curveto(1290.0481554,2270.02741716)(1291.30700424,2269.13477889)(1292.90920996,2268.40240479)
\curveto(1294.51134674,2267.66993661)(1296.41106408,2267.00617995)(1298.60836768,2266.41113281)
\curveto(1300.80559093,2265.86177191)(1303.18595965,2265.24379157)(1305.74948096,2264.55718994)
\curveto(1308.31290765,2263.87050193)(1310.94504613,2263.09230447)(1313.6459043,2262.22259521)
\curveto(1316.34665205,2261.39858057)(1318.79568525,2260.32283702)(1320.99301123,2258.99536133)
\curveto(1323.1902121,2257.66781038)(1324.9983768,2255.99697465)(1326.41751074,2253.98284912)
\curveto(1327.83650873,2252.01443468)(1328.54604171,2249.51962517)(1328.54611182,2246.49841309)
\curveto(1328.54604171,2245.26242728)(1328.40871275,2243.95780212)(1328.13412451,2242.58453369)
}
}
{
\newrgbcolor{curcolor}{0 0 0}
\pscustom[linestyle=none,fillstyle=solid,fillcolor=curcolor]
{
\newpath
\moveto(1364.38900732,2255.90545654)
\curveto(1364.11432777,2254.44057972)(1363.88544616,2253.02151376)(1363.70236182,2251.64825439)
\curveto(1363.56501191,2250.27493447)(1363.49634743,2248.92453299)(1363.49636816,2247.5970459)
\curveto(1363.49634743,2244.85044038)(1363.81678168,2242.33274271)(1364.45767187,2240.04394531)
\curveto(1365.144295,2237.80088689)(1366.17426223,2235.87828139)(1367.54757666,2234.27612305)
\curveto(1368.92084151,2232.67393889)(1370.61456541,2231.43797821)(1372.62875342,2230.5682373)
\curveto(1374.68865801,2229.698478)(1377.09191489,2229.26360294)(1379.83853125,2229.26361084)
\curveto(1382.35619185,2229.26360294)(1384.73656056,2229.53826087)(1386.97964453,2230.08758545)
\curveto(1389.22264006,2230.68266891)(1391.28257452,2231.48375453)(1393.1594541,2232.49084473)
\curveto(1395.0820092,2233.49791267)(1396.75284493,2234.68809703)(1398.17196631,2236.06140137)
\curveto(1399.59097686,2237.48045264)(1400.7353849,2239.0139594)(1401.60519385,2240.66192627)
\lineto(1411.83621191,2237.57202148)
\curveto(1410.6917348,2235.51207082)(1409.1811162,2233.42924819)(1407.30435156,2231.32354736)
\curveto(1405.42745784,2229.26360294)(1403.11575361,2227.38677377)(1400.36923193,2225.6930542)
\curveto(1397.62259504,2224.0451023)(1394.39536438,2222.67181266)(1390.68753027,2221.57318115)
\curveto(1387.02537663,2220.52032555)(1382.81395506,2219.99389786)(1378.05325293,2219.99389648)
\curveto(1373.7044671,2219.99389786)(1369.83636794,2220.54321371)(1366.44894385,2221.6418457)
\curveto(1363.06147237,2222.78625346)(1360.20045229,2224.47997736)(1357.865875,2226.72302246)
\curveto(1355.5312675,2228.96605685)(1353.74599097,2231.75841246)(1352.51004004,2235.10009766)
\curveto(1351.31984593,2238.44175539)(1350.72475375,2242.33274271)(1350.72476172,2246.77307129)
\curveto(1350.72475375,2250.57248056)(1351.15962881,2254.64657317)(1352.02938818,2258.99536133)
\curveto(1353.40266856,2266.13642984)(1355.46260302,2272.13312794)(1358.20919775,2276.98547363)
\curveto(1360.95576159,2281.83770808)(1364.13721593,2285.75158356)(1367.75357031,2288.72711182)
\curveto(1371.41565103,2291.70250535)(1375.37530283,2293.83110429)(1379.6325376,2295.11291504)
\curveto(1383.88969862,2296.39457829)(1388.21556099,2297.03544679)(1392.6101377,2297.03552246)
\curveto(1397.37082527,2297.03544679)(1401.42202972,2296.41746645)(1404.76376318,2295.18157959)
\curveto(1408.10537264,2293.99132142)(1410.82906377,2292.25182121)(1412.93484473,2289.96307373)
\curveto(1415.08626166,2287.71996539)(1416.64265659,2284.99627426)(1417.60403418,2281.79199219)
\curveto(1418.61103841,2278.58758926)(1419.11457794,2274.99414803)(1419.1146543,2271.01165771)
\curveto(1419.11457794,2266.84596282)(1418.63392657,2262.35988332)(1417.67269873,2257.55340576)
\lineto(1417.39804053,2255.90545654)
\lineto(1364.38900732,2255.90545654)
\moveto(1390.96218848,2287.90313721)
\curveto(1388.85642949,2287.90307067)(1386.59050158,2287.55974826)(1384.16439795,2286.87316895)
\curveto(1381.73821151,2286.23223494)(1379.38073096,2285.06493874)(1377.09194922,2283.37127686)
\curveto(1374.80309882,2281.67749096)(1372.69738803,2279.36578673)(1370.77481055,2276.43615723)
\curveto(1368.89795335,2273.55219391)(1367.43311107,2269.86720004)(1366.3802793,2265.38116455)
\lineto(1406.48037695,2265.38116455)
\curveto(1406.57186588,2266.29664697)(1406.64053036,2267.14350891)(1406.68637061,2267.92175293)
\curveto(1406.732083,2268.74568016)(1406.75497116,2269.54676579)(1406.75503516,2270.32501221)
\curveto(1406.75497116,2273.66663471)(1406.34298427,2276.45899032)(1405.51907324,2278.7020874)
\curveto(1404.6950367,2280.94506982)(1403.57351682,2282.75323451)(1402.15451025,2284.12658691)
\curveto(1400.7353849,2285.4998138)(1399.06454916,2286.46111655)(1397.14199805,2287.01049805)
\curveto(1395.21933817,2287.60552458)(1393.1594037,2287.90307067)(1390.96218848,2287.90313721)
}
}
{
\newrgbcolor{curcolor}{0 0 0}
\pscustom[linestyle=none,fillstyle=solid,fillcolor=curcolor]
{
\newpath
\moveto(1503.84670996,2284.40124512)
\curveto(1501.78669873,2284.76739265)(1499.5894353,2285.06493874)(1497.25491309,2285.29388428)
\curveto(1494.96602684,2285.56847828)(1492.40255284,2285.70580724)(1489.5644834,2285.70587158)
\curveto(1485.71920992,2285.70580724)(1482.10288052,2284.92760978)(1478.71548437,2283.37127686)
\curveto(1475.37376128,2281.81481992)(1472.37541222,2279.68622098)(1469.72042822,2276.98547363)
\curveto(1467.11113526,2274.33039137)(1464.89098367,2271.21760152)(1463.0599668,2267.64709473)
\curveto(1461.27465428,2264.07649538)(1459.99291728,2260.2770607)(1459.21475195,2256.2487793)
\lineto(1452.41696143,2221.3671875)
\lineto(1440.05734229,2221.3671875)
\lineto(1449.39572119,2269.5010376)
\curveto(1450.03656738,2272.97998989)(1450.47144243,2276.45899032)(1450.70034766,2279.93804932)
\curveto(1450.97498196,2283.46276749)(1451.11231093,2286.59844551)(1451.11233496,2289.34509277)
\curveto(1451.11231093,2290.67253812)(1451.08942277,2291.86272247)(1451.04367041,2292.91564941)
\curveto(1450.99787012,2294.01420958)(1450.9520938,2294.92973601)(1450.90634131,2295.66223145)
\lineto(1462.64797949,2295.66223145)
\curveto(1462.78527289,2294.51774911)(1462.87682553,2293.16734763)(1462.9226377,2291.61102295)
\curveto(1463.0141545,2290.1003341)(1463.05993082,2288.52105101)(1463.0599668,2286.87316895)
\curveto(1463.05993082,2285.40826115)(1463.0141545,2283.96630703)(1462.9226377,2282.54730225)
\curveto(1462.87682553,2281.1281751)(1462.78527289,2279.77777362)(1462.64797949,2278.49609375)
\lineto(1462.99130225,2278.49609375)
\curveto(1464.82231919,2281.33416855)(1466.65337205,2283.89764255)(1468.48446631,2286.18652344)
\curveto(1470.36125408,2288.47527469)(1472.42118855,2290.42076835)(1474.66427588,2292.02301025)
\curveto(1476.95304436,2293.62511085)(1479.4936302,2294.86107153)(1482.28604102,2295.730896)
\curveto(1485.12411774,2296.60057174)(1488.42001288,2297.03544679)(1492.17373633,2297.03552246)
\curveto(1494.73714523,2297.03544679)(1497.23195475,2296.92100599)(1499.65817236,2296.69219971)
\curveto(1502.13002114,2296.46324277)(1504.28150825,2296.18858485)(1506.11264014,2295.8682251)
\lineto(1503.84670996,2284.40124512)
}
}
{
\newrgbcolor{curcolor}{0 0 0}
\pscustom[linestyle=none,fillstyle=solid,fillcolor=curcolor]
{
\newpath
\moveto(1553.28518652,2221.3671875)
\lineto(1538.65963721,2221.3671875)
\lineto(1523.34744238,2295.66223145)
\lineto(1537.21768164,2295.66223145)
\lineto(1546.21273779,2247.3223877)
\curveto(1546.35003194,2246.31528267)(1546.48736091,2245.07932199)(1546.6247251,2243.61450195)
\curveto(1546.80779516,2242.19541374)(1546.96801228,2240.7076833)(1547.10537695,2239.15130615)
\curveto(1547.24267021,2237.59489344)(1547.35711102,2236.08427483)(1547.44869971,2234.6194458)
\curveto(1547.58599262,2233.20036658)(1547.70043343,2232.01018223)(1547.79202246,2231.04888916)
\curveto(1548.2955256,2232.01018223)(1548.89061778,2233.20036658)(1549.57730078,2234.6194458)
\curveto(1550.26390742,2236.08427483)(1550.99632857,2237.57200528)(1551.77456641,2239.0826416)
\curveto(1552.5527235,2240.59324249)(1553.3080328,2242.05808478)(1554.04049658,2243.47717285)
\curveto(1554.81865141,2244.94199303)(1555.48240807,2246.1779537)(1556.03176855,2247.18505859)
\lineto(1584.39022803,2295.66223145)
\lineto(1598.19180273,2295.66223145)
\lineto(1553.28518652,2221.3671875)
}
}
{
\newrgbcolor{curcolor}{0 0 0}
\pscustom[linestyle=none,fillstyle=solid,fillcolor=curcolor]
{
\newpath
\moveto(1665.6203916,2228.98895264)
\curveto(1666.30696652,2228.98894501)(1666.99361134,2229.03472134)(1667.68032812,2229.12628174)
\curveto(1668.41267731,2229.21782662)(1669.09932213,2229.33226743)(1669.74026465,2229.46960449)
\lineto(1668.22964453,2221.7791748)
\curveto(1666.62740077,2221.41296382)(1665.02522952,2221.13830589)(1663.42312598,2220.9552002)
\curveto(1661.86666335,2220.77209532)(1660.21871577,2220.68054268)(1658.47927832,2220.68054199)
\curveto(1654.4966756,2220.68054268)(1651.65854367,2221.55029279)(1649.96487402,2223.28979492)
\curveto(1648.31687221,2225.02929321)(1647.49289842,2227.63854353)(1647.4929502,2231.11755371)
\curveto(1647.49289842,2231.80418878)(1647.51578658,2232.51372176)(1647.56161475,2233.24615479)
\curveto(1647.65311555,2234.02434037)(1647.74466819,2234.80253783)(1647.83627295,2235.58074951)
\lineto(1647.42428564,2235.58074951)
\curveto(1645.68473373,2233.20036658)(1643.85368087,2231.04887948)(1641.93112158,2229.12628174)
\curveto(1640.00846987,2227.20366848)(1637.90275909,2225.55572091)(1635.61398291,2224.18243408)
\curveto(1633.32512695,2222.85491795)(1630.80742927,2221.82495071)(1628.06088232,2221.0925293)
\curveto(1625.3142707,2220.36010843)(1622.24725716,2219.99389786)(1618.85983252,2219.99389648)
\curveto(1612.77155863,2219.99389786)(1608.21681465,2221.41296382)(1605.19558691,2224.25109863)
\curveto(1602.22011655,2227.08922768)(1600.7323861,2231.00310316)(1600.73239111,2235.99273682)
\curveto(1600.7323861,2237.86955137)(1600.93837955,2239.90659767)(1601.35037207,2242.10388184)
\curveto(1602.31166919,2247.00192749)(1603.98250492,2251.00735561)(1606.36288428,2254.12017822)
\curveto(1608.78901867,2257.27871165)(1611.6271506,2259.750633)(1614.87728857,2261.53594971)
\curveto(1618.17316456,2263.32118607)(1621.74371763,2264.55714675)(1625.5889585,2265.24383545)
\curveto(1629.47991595,2265.97621272)(1633.37090327,2266.36531145)(1637.26193213,2266.41113281)
\lineto(1653.46676611,2266.68579102)
\lineto(1654.29074072,2270.73699951)
\curveto(1654.51956376,2271.88135818)(1654.67978088,2272.95710173)(1654.77139258,2273.9642334)
\curveto(1654.90866249,2274.97125987)(1654.97732697,2275.93256262)(1654.97738623,2276.84814453)
\curveto(1654.97732697,2280.73907637)(1653.92447158,2283.50854382)(1651.81881689,2285.15655518)
\curveto(1649.71305001,2286.80443896)(1646.55448383,2287.62841274)(1642.34310889,2287.628479)
\curveto(1640.05424619,2287.62841274)(1637.94853541,2287.44530746)(1636.02597021,2287.0791626)
\curveto(1634.10332441,2286.75866264)(1632.38671236,2286.1406823)(1630.87612891,2285.22521973)
\curveto(1629.36547514,2284.35540576)(1628.03796182,2283.14233324)(1626.89358496,2281.58599854)
\curveto(1625.79492207,2280.07531971)(1624.85650748,2278.15271421)(1624.07833838,2275.81817627)
\lineto(1611.37539648,2276.98547363)
\curveto(1612.42823622,2279.91510258)(1613.80152586,2282.59301739)(1615.49526953,2285.01922607)
\curveto(1617.23474997,2287.49108378)(1619.40912524,2289.61968272)(1622.01840186,2291.4050293)
\curveto(1624.6734022,2293.19023579)(1627.83196838,2294.56352544)(1631.49410986,2295.52490234)
\curveto(1635.1561798,2296.53190726)(1639.43626585,2297.03544679)(1644.33438086,2297.03552246)
\curveto(1652.16208321,2297.03544679)(1657.99856419,2295.47905186)(1661.84384131,2292.36633301)
\curveto(1665.7347625,2289.25347215)(1667.68025616,2284.69872817)(1667.68032812,2278.7020874)
\curveto(1667.68025616,2277.64917467)(1667.61159168,2276.57343112)(1667.47433447,2275.47485352)
\curveto(1667.38271008,2274.37616769)(1667.22249295,2273.23175966)(1666.99368262,2272.04162598)
\lineto(1660.81387305,2240.04394531)
\curveto(1660.40182106,2238.07554482)(1660.19582761,2236.3360446)(1660.19589209,2234.82543945)
\curveto(1660.19582761,2232.9028205)(1660.60781451,2231.43797821)(1661.431854,2230.4309082)
\curveto(1662.25576208,2229.46959639)(1663.65193988,2228.98894501)(1665.6203916,2228.98895264)
\moveto(1623.87234473,2229.40093994)
\curveto(1627.67175125,2229.40093191)(1631.10497536,2230.11046489)(1634.17202734,2231.52954102)
\curveto(1637.28477875,2232.99437314)(1639.96269355,2234.82542599)(1642.20577979,2237.02270508)
\curveto(1644.49454937,2239.21995285)(1646.34849039,2241.64609789)(1647.7676084,2244.30114746)
\curveto(1649.18662231,2246.95615117)(1650.1250369,2249.49673701)(1650.58285498,2251.9229126)
\lineto(1651.81881689,2258.03405762)
\lineto(1638.7038877,2257.75939941)
\curveto(1636.18614703,2257.7135867)(1633.62267303,2257.50759325)(1631.01345801,2257.14141846)
\curveto(1628.40417239,2256.82094843)(1625.9551392,2256.11141545)(1623.66635107,2255.01281738)
\curveto(1621.42328338,2253.91415202)(1619.45490156,2252.31198077)(1617.76119971,2250.20629883)
\curveto(1616.11323009,2248.14633553)(1614.9459339,2245.37686808)(1614.25930762,2241.89788818)
\curveto(1614.03040747,2240.61613065)(1613.91596667,2239.51749894)(1613.91598486,2238.60198975)
\curveto(1613.91596667,2235.76384058)(1614.73994045,2233.52080083)(1616.38790869,2231.87286377)
\curveto(1618.08161192,2230.22490569)(1620.57642143,2229.40093191)(1623.87234473,2229.40093994)
}
}
{
\newrgbcolor{curcolor}{0 0 0}
\pscustom[linestyle=none,fillstyle=solid,fillcolor=curcolor]
{
\newpath
\moveto(1698.44204687,2285.91186523)
\lineto(1700.29598975,2295.66223145)
\lineto(1711.96896338,2295.66223145)
\lineto(1719.72805762,2315.02563477)
\lineto(1727.96780371,2315.02563477)
\lineto(1724.19125342,2295.66223145)
\lineto(1753.85433936,2295.66223145)
\lineto(1752.00039648,2285.91186523)
\lineto(1722.33731055,2285.91186523)
\lineto(1713.6169126,2241.14257812)
\curveto(1713.52532638,2240.59324249)(1713.43377374,2240.0668148)(1713.34225439,2239.56329346)
\curveto(1713.29644477,2239.05973573)(1713.27355661,2238.57908435)(1713.27358984,2238.12133789)
\curveto(1713.27355661,2235.51207082)(1714.14330672,2233.58946532)(1715.88284277,2232.35351562)
\curveto(1717.62230714,2231.11754396)(1720.34599827,2230.49956362)(1724.05392432,2230.49957275)
\curveto(1725.88493316,2230.49956362)(1727.7846505,2230.5682281)(1729.75308203,2230.70556641)
\curveto(1731.72141414,2230.88866235)(1733.62113148,2231.0946558)(1735.45223975,2231.32354736)
\curveto(1737.32901351,2231.55241901)(1739.06851373,2231.80418878)(1740.67074561,2232.07885742)
\curveto(1742.31863255,2232.35350464)(1743.71481035,2232.60527441)(1744.8592832,2232.83416748)
\lineto(1743.07400488,2223.42712402)
\curveto(1742.02108646,2223.15246404)(1740.64779681,2222.83202979)(1738.95413184,2222.46582031)
\curveto(1737.30612535,2222.09960864)(1735.45218433,2221.75628623)(1733.39230322,2221.43585205)
\curveto(1731.33231541,2221.11541773)(1729.11216382,2220.8407598)(1726.7318418,2220.61187744)
\curveto(1724.35142639,2220.38299659)(1721.92528136,2220.26855579)(1719.45339941,2220.26855469)
\curveto(1713.09045132,2220.26855579)(1708.3297139,2221.39007566)(1705.17117285,2223.63311768)
\curveto(1702.01258154,2225.92193148)(1700.43329845,2229.49248455)(1700.43331885,2234.3447876)
\curveto(1700.43329845,2235.16874841)(1700.45618661,2236.03849851)(1700.5019834,2236.95404053)
\curveto(1700.59351558,2237.86955137)(1700.7537327,2238.83085412)(1700.98263525,2239.83795166)
\lineto(1709.97769141,2285.91186523)
\lineto(1698.44204687,2285.91186523)
}
}
{
\newrgbcolor{curcolor}{0 0 0}
\pscustom[linestyle=none,fillstyle=solid,fillcolor=curcolor]
{
\newpath
\moveto(1810.22793555,2231.11755371)
\lineto(1836.32046484,2231.11755371)
\lineto(1834.39785742,2221.3671875)
\lineto(1766.96926855,2221.3671875)
\lineto(1768.89187598,2231.11755371)
\lineto(1797.86831641,2231.11755371)
\lineto(1808.51132178,2285.91186523)
\lineto(1786.60733008,2285.91186523)
\lineto(1788.46127295,2295.66223145)
\lineto(1822.72488379,2295.66223145)
\lineto(1810.22793555,2231.11755371)
\moveto(1811.80722021,2310.08178711)
\lineto(1814.34780859,2323.26538086)
\lineto(1828.08071875,2323.26538086)
\lineto(1825.54013037,2310.08178711)
\lineto(1811.80722021,2310.08178711)
}
}
{
\newrgbcolor{curcolor}{0 0 0}
\pscustom[linestyle=none,fillstyle=solid,fillcolor=curcolor]
{
\newpath
\moveto(1924.00509619,2258.58337402)
\curveto(1922.76906019,2252.08309917)(1920.91511917,2246.40683531)(1918.44326758,2241.55456543)
\curveto(1915.97127646,2236.74803149)(1912.97292741,2232.74260337)(1909.44821143,2229.53826904)
\curveto(1905.96915023,2226.33391837)(1901.98661027,2223.9306615)(1897.50057959,2222.32849121)
\curveto(1893.06022759,2220.77209532)(1888.23082568,2219.99389786)(1883.01235937,2219.99389648)
\curveto(1878.89245612,2219.99389786)(1875.18457408,2220.52032555)(1871.88870215,2221.57318115)
\curveto(1868.5927838,2222.67181266)(1865.77754003,2224.31976023)(1863.4429624,2226.51702881)
\curveto(1861.15413157,2228.71428709)(1859.39174319,2231.48375453)(1858.15579199,2234.82543945)
\curveto(1856.96559816,2238.16709746)(1856.37050598,2242.08097294)(1856.37051367,2246.56707764)
\curveto(1856.37050598,2250.18338183)(1856.78249287,2254.18880995)(1857.60647559,2258.58337402)
\curveto(1860.12416434,2271.35493048)(1864.88490176,2280.94506982)(1871.88870215,2287.3538208)
\curveto(1878.93823244,2293.80821613)(1887.77306247,2297.03544679)(1898.39321875,2297.03552246)
\curveto(1907.502657,2297.03544679)(1914.27755257,2294.92973601)(1918.71792578,2290.71838379)
\curveto(1923.15815892,2286.50689287)(1925.37831051,2280.14398419)(1925.37838721,2271.62963867)
\curveto(1925.37831051,2269.66120659)(1925.26386971,2267.60127213)(1925.03506445,2265.4498291)
\curveto(1924.85188281,2263.29829791)(1924.5085604,2261.00948184)(1924.00509619,2258.58337402)
\moveto(1911.02749609,2258.58337402)
\curveto(1911.48519696,2260.96370552)(1911.82851937,2263.16096895)(1912.05746436,2265.1751709)
\curveto(1912.28628259,2267.18928523)(1912.40072339,2269.06611441)(1912.40078711,2270.80566406)
\curveto(1912.40072339,2274.10150977)(1912.01162466,2276.82520089)(1911.23348975,2278.97674561)
\curveto(1910.50100605,2281.17395142)(1909.44815066,2282.9363398)(1908.07492041,2284.26391602)
\curveto(1906.70157137,2285.59136644)(1905.0536238,2286.52978103)(1903.13107275,2287.0791626)
\curveto(1901.25418913,2287.62841274)(1899.14847834,2287.90307067)(1896.81393408,2287.90313721)
\curveto(1893.83842506,2287.90307067)(1891.00029313,2287.4224193)(1888.29952979,2286.46118164)
\curveto(1885.5986872,2285.54559012)(1883.10387769,2283.94341887)(1880.81509375,2281.65466309)
\curveto(1878.57202187,2279.41156305)(1876.58075188,2276.413214)(1874.84127783,2272.65960693)
\curveto(1873.10175146,2268.90589729)(1871.68268549,2264.21382434)(1870.58407568,2258.58337402)
\curveto(1870.12629057,2256.34029706)(1869.80585632,2254.25747443)(1869.62277197,2252.3348999)
\curveto(1869.43964574,2250.41226344)(1869.3480931,2248.6269869)(1869.34811377,2246.97906494)
\curveto(1869.3480931,2243.54581522)(1869.73719183,2240.68479514)(1870.51541113,2238.39599609)
\curveto(1871.33936308,2236.10716299)(1872.4379948,2234.27611014)(1873.81130957,2232.90283203)
\curveto(1875.2303504,2231.52953085)(1876.85540981,2230.54533994)(1878.68649268,2229.95025635)
\curveto(1880.56329185,2229.40093191)(1882.55456183,2229.12627398)(1884.66030859,2229.12628174)
\curveto(1887.63573351,2229.12627398)(1890.47386543,2229.56114903)(1893.17471289,2230.4309082)
\curveto(1895.87547136,2231.34642557)(1898.34739272,2232.92570866)(1900.59048437,2235.16876221)
\curveto(1902.87924854,2237.41178815)(1904.91629484,2240.43302537)(1906.70162939,2244.23248291)
\curveto(1908.48684791,2248.03189472)(1909.92880203,2252.81552031)(1911.02749609,2258.58337402)
}
}
{
\newrgbcolor{curcolor}{0 0 0}
\pscustom[linestyle=none,fillstyle=solid,fillcolor=curcolor]
{
\newpath
\moveto(1985.39120459,2221.3671875)
\lineto(1994.66091895,2269.08905029)
\curveto(1995.16439682,2271.74402921)(1995.41616659,2273.98706896)(1995.416229,2275.81817627)
\curveto(1995.41616659,2278.10693789)(1995.1186205,2279.98376707)(1994.52358984,2281.44866943)
\curveto(1993.92843614,2282.95922796)(1993.0815742,2284.14941232)(1991.98300146,2285.01922607)
\curveto(1990.88431077,2285.93468885)(1989.57968561,2286.57555735)(1988.06912207,2286.9418335)
\curveto(1986.55844839,2287.30797849)(1984.91050082,2287.49108378)(1983.12527441,2287.4911499)
\curveto(1980.42442133,2287.49108378)(1977.79228284,2286.96465608)(1975.22885107,2285.91186523)
\curveto(1972.66533485,2284.90472162)(1970.30785429,2283.41699117)(1968.15640234,2281.44866943)
\curveto(1966.0506564,2279.52600385)(1964.21960355,2277.12274698)(1962.66323828,2274.2388916)
\curveto(1961.10681369,2271.4007068)(1959.96240565,2268.12769982)(1959.23001074,2264.41986084)
\lineto(1950.85293555,2221.3671875)
\lineto(1938.49331641,2221.3671875)
\lineto(1949.82296729,2279.80072021)
\curveto(1950.14338471,2281.35705671)(1950.4409308,2282.98211612)(1950.71560645,2284.67590332)
\curveto(1951.03602298,2286.41534023)(1951.31068091,2288.04039964)(1951.53958105,2289.55108643)
\curveto(1951.81422044,2291.10741317)(1952.04310205,2292.43492649)(1952.22622656,2293.53363037)
\curveto(1952.40931262,2294.63218992)(1952.50086526,2295.3417229)(1952.50088477,2295.66223145)
\lineto(1964.1738584,2295.66223145)
\curveto(1964.17382722,2295.43327554)(1964.10516274,2294.8152952)(1963.96786475,2293.80828857)
\curveto(1963.83050481,2292.80113706)(1963.67028769,2291.63384087)(1963.48721289,2290.30639648)
\curveto(1963.30407712,2289.02459055)(1963.09808367,2287.69707722)(1962.86923193,2286.32385254)
\curveto(1962.64032046,2284.99627426)(1962.43432701,2283.87475439)(1962.25125098,2282.95928955)
\lineto(1962.45724463,2282.95928955)
\curveto(1965.57000503,2287.30797849)(1969.2549989,2290.7412026)(1973.5122373,2293.25897217)
\curveto(1977.76939468,2295.77659795)(1982.91923084,2297.03544679)(1988.96176123,2297.03552246)
\curveto(1995.2330613,2297.03544679)(2000.03957505,2295.63926899)(2003.38131689,2292.84698486)
\curveto(2006.72291798,2290.1003341)(2008.39375371,2285.68291908)(2008.3938291,2279.59472656)
\curveto(2008.39375371,2278.26715501)(2008.30220107,2276.87097721)(2008.1191709,2275.40618896)
\curveto(2007.98176682,2273.98706896)(2007.75288521,2272.47645036)(2007.43252539,2270.87432861)
\lineto(1997.81948828,2221.3671875)
\lineto(1985.39120459,2221.3671875)
}
}
{
\newrgbcolor{curcolor}{0 0 0}
\pscustom[linestyle=none,fillstyle=solid,fillcolor=curcolor]
{
\newpath
\moveto(254.83853125,2071.2890625)
\lineto(2446.24478125,2071.2890625)
}
}
{
\newrgbcolor{curcolor}{0.65882355 0 0.21176471}
\pscustom[linewidth=17.578125,linecolor=curcolor]
{
\newpath
\moveto(254.83853125,2071.2890625)
\lineto(2446.24478125,2071.2890625)
}
}
{
\newrgbcolor{curcolor}{0.51764709 0.74509805 0.51764709}
\pscustom[linestyle=none,fillstyle=solid,fillcolor=curcolor]
{
\newpath
\moveto(407.18221875,1931.84027246)
\curveto(407.18221875,1912.42401172)(391.44222949,1896.68402246)(372.02596875,1896.68402246)
\curveto(352.60970801,1896.68402246)(336.86971875,1912.42401172)(336.86971875,1931.84027246)
\curveto(336.86971875,1951.2565332)(352.60970801,1966.99652246)(372.02596875,1966.99652246)
\curveto(391.44222949,1966.99652246)(407.18221875,1951.2565332)(407.18221875,1931.84027246)
\closepath
}
}
{
\newrgbcolor{curcolor}{0.01176471 0.50196081 0.28235295}
\pscustom[linewidth=11.71875,linecolor=curcolor]
{
\newpath
\moveto(407.18221875,1931.84027246)
\curveto(407.18221875,1912.42401172)(391.44222949,1896.68402246)(372.02596875,1896.68402246)
\curveto(352.60970801,1896.68402246)(336.86971875,1912.42401172)(336.86971875,1931.84027246)
\curveto(336.86971875,1951.2565332)(352.60970801,1966.99652246)(372.02596875,1966.99652246)
\curveto(391.44222949,1966.99652246)(407.18221875,1951.2565332)(407.18221875,1931.84027246)
\closepath
}
}
{
\newrgbcolor{curcolor}{0 0 0}
\pscustom[linestyle=none,fillstyle=solid,fillcolor=curcolor]
{
\newpath
\moveto(536.03131079,1937.99415588)
\curveto(536.03125225,1946.38644541)(534.28984469,1952.89049774)(530.80708289,1957.5063324)
\curveto(527.36617608,1962.12205589)(522.51960805,1964.42994543)(516.26736426,1964.43000793)
\curveto(510.05700455,1964.42994543)(505.21043652,1962.12205589)(501.72764563,1957.5063324)
\curveto(498.2867679,1952.89049774)(496.56634116,1946.38644541)(496.56636023,1937.99415588)
\curveto(496.56634116,1929.64375585)(498.2867679,1923.16068433)(501.72764563,1918.54492188)
\curveto(505.21043652,1913.92912618)(510.05700455,1911.62123664)(516.26736426,1911.62124634)
\curveto(522.51960805,1911.62123664)(527.36617608,1913.92912618)(530.80708289,1918.54492188)
\curveto(534.28984469,1923.16068433)(536.03125225,1929.64375585)(536.03131079,1937.99415588)
\moveto(547.61273169,1910.67710876)
\curveto(547.61266157,1898.67607442)(544.9480982,1889.75922848)(539.61903357,1883.92654419)
\curveto(534.28984469,1878.05193428)(526.12830806,1875.11462032)(515.13439917,1875.11459351)
\curveto(511.06408362,1875.11462032)(507.22459466,1875.42933253)(503.61592078,1876.05873108)
\curveto(500.00719465,1876.64621974)(496.50339872,1877.56937556)(493.10452246,1878.82820129)
\lineto(493.10452246,1890.09490967)
\curveto(496.50339872,1888.24860987)(499.86032895,1886.88485696)(503.17532324,1886.00364685)
\curveto(506.49026617,1885.12246859)(509.86817722,1884.68187149)(513.30906653,1884.68185425)
\curveto(520.90408537,1884.68187149)(526.58988596,1886.67504882)(530.36648535,1890.66139221)
\curveto(534.142979,1894.60579651)(536.03125225,1900.58532849)(536.03131079,1908.6000061)
\lineto(536.03131079,1914.32777405)
\curveto(533.63943946,1910.17356048)(530.57624062,1907.06840001)(526.84170508,1905.01228333)
\curveto(523.10707084,1902.95616047)(518.63815746,1901.92810059)(513.43495154,1901.92810059)
\curveto(504.79082024,1901.92810059)(497.82519,1905.22208838)(492.53803992,1911.81007385)
\curveto(487.25085975,1918.39803956)(484.60727719,1927.12605818)(484.6072843,1937.99415588)
\curveto(484.60727719,1948.90414309)(487.25085975,1957.65314252)(492.53803992,1964.24118042)
\curveto(497.82519,1970.8290937)(504.79082024,1974.12308149)(513.43495154,1974.12315369)
\curveto(518.63815746,1974.12308149)(523.10707084,1973.09502161)(526.84170508,1971.03897095)
\curveto(530.57624062,1968.98278207)(533.63943946,1965.8776216)(536.03131079,1961.72348022)
\lineto(536.03131079,1972.42370605)
\lineto(547.61273169,1972.42370605)
\lineto(547.61273169,1910.67710876)
}
}
{
\newrgbcolor{curcolor}{0 0 0}
\pscustom[linestyle=none,fillstyle=solid,fillcolor=curcolor]
{
\newpath
\moveto(631.76686072,1940.07125854)
\lineto(631.76686072,1934.40643311)
\lineto(578.51750159,1934.40643311)
\curveto(579.02102192,1926.43369131)(581.41283472,1920.34925526)(585.69294714,1916.15310669)
\curveto(590.01496845,1911.9988913)(596.01548125,1909.92179071)(603.69450354,1909.92179871)
\curveto(608.14239173,1909.92179071)(612.4434586,1910.46729188)(616.59771704,1911.55830383)
\curveto(620.79382256,1912.64929653)(624.94802373,1914.28580002)(629.06033301,1916.46781921)
\lineto(629.06033301,1905.51582336)
\curveto(624.9060621,1903.7534314)(620.64695686,1902.41065931)(616.28300452,1901.48750305)
\curveto(611.91893825,1900.56434768)(607.4919865,1900.10276977)(603.00213599,1900.10276794)
\curveto(591.75637601,1900.10276977)(582.83953007,1903.37577675)(576.25157141,1909.92179871)
\curveto(569.70554052,1916.46780467)(566.43253354,1925.32170817)(566.43254065,1936.48353577)
\curveto(566.43253354,1948.0229489)(569.53769401,1957.1705838)(575.74803137,1963.9264679)
\curveto(582.00029751,1970.72418963)(590.41360392,1974.12308149)(600.98797583,1974.12315369)
\curveto(610.47126208,1974.12308149)(617.96141267,1971.05988265)(623.45845007,1964.93354797)
\curveto(628.99732083,1958.84904891)(631.76678827,1950.56162739)(631.76686072,1940.07125854)
\moveto(620.18543982,1943.47015381)
\curveto(620.1014557,1949.80631809)(618.31808651,1954.86269426)(614.8353269,1958.63929749)
\curveto(611.3944179,1962.41578729)(606.82060045,1964.30406055)(601.11386084,1964.30412292)
\curveto(594.65172834,1964.30406055)(589.46946729,1962.47872973)(585.56706213,1958.828125)
\curveto(581.70656611,1955.17740647)(579.48259983,1950.03710704)(578.89515662,1943.4072113)
\lineto(620.18543982,1943.47015381)
}
}
{
\newrgbcolor{curcolor}{0 0 0}
\pscustom[linestyle=none,fillstyle=solid,fillcolor=curcolor]
{
\newpath
\moveto(661.97926306,1992.43942261)
\lineto(661.97926306,1972.42370605)
\lineto(685.83447241,1972.42370605)
\lineto(685.83447241,1963.42292786)
\lineto(661.97926306,1963.42292786)
\lineto(661.97926306,1925.15388489)
\curveto(661.97923946,1919.40511863)(662.75552958,1915.71249537)(664.30813574,1914.07600403)
\curveto(665.90265167,1912.43948839)(669.11271621,1911.62123664)(673.93833899,1911.62124634)
\lineto(685.83447241,1911.62124634)
\lineto(685.83447241,1901.92810059)
\lineto(673.93833899,1901.92810059)
\curveto(665.00047667,1901.92810059)(658.83211736,1903.58558489)(655.43324255,1906.90055847)
\curveto(652.03433363,1910.25748374)(650.3348877,1916.34191979)(650.33489966,1925.15388489)
\lineto(650.33489966,1963.42292786)
\lineto(641.8376615,1963.42292786)
\lineto(641.8376615,1972.42370605)
\lineto(650.33489966,1972.42370605)
\lineto(650.33489966,1992.43942261)
\lineto(661.97926306,1992.43942261)
}
}
{
\newrgbcolor{curcolor}{0 0 0}
\pscustom[linestyle=none,fillstyle=solid,fillcolor=curcolor]
{
\newpath
\moveto(701.63304114,1995.90126038)
\lineto(714.34742712,1995.90126038)
\lineto(714.34742712,1956.18453979)
\lineto(756.5189054,1995.90126038)
\lineto(772.88395667,1995.90126038)
\lineto(726.24356055,1952.09327698)
\lineto(776.21990942,1901.92810059)
\lineto(759.47720312,1901.92810059)
\lineto(714.34742712,1947.1837616)
\lineto(714.34742712,1901.92810059)
\lineto(701.63304114,1901.92810059)
\lineto(701.63304114,1995.90126038)
}
}
{
\newrgbcolor{curcolor}{0 0 0}
\pscustom[linestyle=none,fillstyle=solid,fillcolor=curcolor]
{
\newpath
\moveto(844.32369971,1944.47723389)
\lineto(844.32369971,1901.92810059)
\lineto(832.74227881,1901.92810059)
\lineto(832.74227881,1944.09957886)
\curveto(832.74221964,1950.77143553)(831.44140918,1955.76486926)(828.83984351,1959.07989502)
\curveto(826.23816731,1962.39480648)(822.33573591,1964.05229078)(817.1325376,1964.05235291)
\curveto(810.88021148,1964.05229078)(805.94972019,1962.05911345)(802.34104895,1958.07281494)
\curveto(798.73232018,1954.08640414)(796.92797018,1948.65237332)(796.92799353,1941.77070618)
\lineto(796.92799353,1901.92810059)
\lineto(785.28363013,1901.92810059)
\lineto(785.28363013,1972.42370605)
\lineto(796.92799353,1972.42370605)
\lineto(796.92799353,1961.47171021)
\curveto(799.69743762,1965.70977509)(802.94946379,1968.877878)(806.68408179,1970.97602844)
\curveto(810.4605952,1973.07404079)(814.80362369,1974.12308149)(819.7131803,1974.12315369)
\curveto(827.81172836,1974.12308149)(833.93812604,1971.60538381)(838.09239172,1966.5700531)
\curveto(842.24652838,1961.57655473)(844.32362896,1954.21228902)(844.32369971,1944.47723389)
}
}
{
\newrgbcolor{curcolor}{0 0 0}
\pscustom[linestyle=none,fillstyle=solid,fillcolor=curcolor]
{
\newpath
\moveto(894.86653113,1964.30412292)
\curveto(888.65617072,1964.30406055)(883.74666025,1961.87028613)(880.13798499,1957.00279236)
\curveto(876.52926024,1952.17715007)(874.72491024,1945.54721285)(874.72492957,1937.11296082)
\curveto(874.72491024,1928.67863841)(876.50827943,1922.02772038)(880.07504248,1917.16018677)
\curveto(883.68371781,1912.33458432)(888.6142091,1909.92179071)(894.86653113,1909.92179871)
\curveto(901.03485097,1909.92179071)(905.92338063,1912.35556513)(909.53213477,1917.22312927)
\curveto(913.14078064,1922.09066282)(914.94513064,1928.72060004)(914.94519019,1937.11296082)
\curveto(914.94513064,1945.46328959)(913.14078064,1952.072246)(909.53213477,1956.93984985)
\curveto(905.92338063,1961.84930531)(901.03485097,1964.30406055)(894.86653113,1964.30412292)
\moveto(894.86653113,1974.12315369)
\curveto(904.93728237,1974.12308149)(912.84704924,1970.85007451)(918.59585547,1964.30412292)
\curveto(924.34453531,1957.75804659)(927.21890682,1948.69433495)(927.21897864,1937.11296082)
\curveto(927.21890682,1925.57347794)(924.34453531,1916.5097663)(918.59585547,1909.92179871)
\curveto(912.84704924,1903.37577675)(904.93728237,1900.10276977)(894.86653113,1900.10276794)
\curveto(884.75373932,1900.10276977)(876.82299164,1903.37577675)(871.07426428,1909.92179871)
\curveto(865.36746721,1916.5097663)(862.51407651,1925.57347794)(862.51408362,1937.11296082)
\curveto(862.51407651,1948.69433495)(865.36746721,1957.75804659)(871.07426428,1964.30412292)
\curveto(876.82299164,1970.85007451)(884.75373932,1974.12308149)(894.86653113,1974.12315369)
}
}
{
\newrgbcolor{curcolor}{0 0 0}
\pscustom[linestyle=none,fillstyle=solid,fillcolor=curcolor]
{
\newpath
\moveto(939.6186521,1972.42370605)
\lineto(951.200073,1972.42370605)
\lineto(965.67684912,1917.41195679)
\lineto(980.09068274,1972.42370605)
\lineto(993.7492063,1972.42370605)
\lineto(1008.22598242,1917.41195679)
\lineto(1022.63981604,1972.42370605)
\lineto(1034.22123694,1972.42370605)
\lineto(1015.77908301,1901.92810059)
\lineto(1002.12055945,1901.92810059)
\lineto(986.95141577,1959.70932007)
\lineto(971.71932959,1901.92810059)
\lineto(958.06080603,1901.92810059)
\lineto(939.6186521,1972.42370605)
}
}
{
\newrgbcolor{curcolor}{0 0 0}
\pscustom[linestyle=none,fillstyle=solid,fillcolor=curcolor]
{
\newpath
\moveto(1051.84513831,1999.86663818)
\lineto(1063.4265592,1999.86663818)
\lineto(1063.4265592,1901.92810059)
\lineto(1051.84513831,1901.92810059)
\lineto(1051.84513831,1999.86663818)
}
}
{
\newrgbcolor{curcolor}{0 0 0}
\pscustom[linestyle=none,fillstyle=solid,fillcolor=curcolor]
{
\newpath
\moveto(1147.89540076,1940.07125854)
\lineto(1147.89540076,1934.40643311)
\lineto(1094.64604163,1934.40643311)
\curveto(1095.14956196,1926.43369131)(1097.54137476,1920.34925526)(1101.82148718,1916.15310669)
\curveto(1106.14350849,1911.9988913)(1112.14402129,1909.92179071)(1119.82304358,1909.92179871)
\curveto(1124.27093177,1909.92179071)(1128.57199863,1910.46729188)(1132.72625708,1911.55830383)
\curveto(1136.9223626,1912.64929653)(1141.07656377,1914.28580002)(1145.18887305,1916.46781921)
\lineto(1145.18887305,1905.51582336)
\curveto(1141.03460214,1903.7534314)(1136.7754969,1902.41065931)(1132.41154456,1901.48750305)
\curveto(1128.04747828,1900.56434768)(1123.62052654,1900.10276977)(1119.13067603,1900.10276794)
\curveto(1107.88491605,1900.10276977)(1098.96807011,1903.37577675)(1092.38011145,1909.92179871)
\curveto(1085.83408056,1916.46780467)(1082.56107358,1925.32170817)(1082.56108069,1936.48353577)
\curveto(1082.56107358,1948.0229489)(1085.66623405,1957.1705838)(1091.87657141,1963.9264679)
\curveto(1098.12883755,1970.72418963)(1106.54214396,1974.12308149)(1117.11651587,1974.12315369)
\curveto(1126.59980212,1974.12308149)(1134.08995271,1971.05988265)(1139.58699011,1964.93354797)
\curveto(1145.12586086,1958.84904891)(1147.89532831,1950.56162739)(1147.89540076,1940.07125854)
\moveto(1136.31397986,1943.47015381)
\curveto(1136.22999574,1949.80631809)(1134.44662655,1954.86269426)(1130.96386694,1958.63929749)
\curveto(1127.52295794,1962.41578729)(1122.94914049,1964.30406055)(1117.24240088,1964.30412292)
\curveto(1110.78026838,1964.30406055)(1105.59800733,1962.47872973)(1101.69560217,1958.828125)
\curveto(1097.83510615,1955.17740647)(1095.61113987,1950.03710704)(1095.02369666,1943.4072113)
\lineto(1136.31397986,1943.47015381)
}
}
{
\newrgbcolor{curcolor}{0 0 0}
\pscustom[linestyle=none,fillstyle=solid,fillcolor=curcolor]
{
\newpath
\moveto(1213.04089331,1961.72348022)
\lineto(1213.04089331,1999.86663818)
\lineto(1224.62231421,1999.86663818)
\lineto(1224.62231421,1901.92810059)
\lineto(1213.04089331,1901.92810059)
\lineto(1213.04089331,1912.50244141)
\curveto(1210.60706035,1908.30626804)(1207.5228807,1905.18012675)(1203.78834509,1903.12400818)
\curveto(1200.09567255,1901.10984884)(1195.64773998,1900.10276977)(1190.44453406,1900.10276794)
\curveto(1181.92628764,1900.10276977)(1174.98163821,1903.50166163)(1169.61056494,1910.29945374)
\curveto(1164.28142308,1917.09722909)(1161.61685971,1926.03505585)(1161.61686682,1937.11296082)
\curveto(1161.61685971,1948.19079541)(1164.28142308,1957.12862217)(1169.61056494,1963.9264679)
\curveto(1174.98163821,1970.72418963)(1181.92628764,1974.12308149)(1190.44453406,1974.12315369)
\curveto(1195.64773998,1974.12308149)(1200.09567255,1973.09502161)(1203.78834509,1971.03897095)
\curveto(1207.5228807,1969.02474369)(1210.60706035,1965.91958323)(1213.04089331,1961.72348022)
\moveto(1173.57594275,1937.11296082)
\curveto(1173.57592368,1928.59471515)(1175.31733124,1921.90183549)(1178.80017065,1917.03430176)
\curveto(1182.32492311,1912.20869944)(1187.15051032,1909.79590583)(1193.27694678,1909.7959137)
\curveto(1199.40330569,1909.79590583)(1204.2288929,1912.20869944)(1207.7537229,1917.03430176)
\curveto(1211.2784464,1921.90183549)(1213.04083477,1928.59471515)(1213.04089331,1937.11296082)
\curveto(1213.04083477,1945.63113611)(1211.2784464,1952.30303495)(1207.7537229,1957.12867737)
\curveto(1204.2288929,1961.99617101)(1199.40330569,1964.42994543)(1193.27694678,1964.43000793)
\curveto(1187.15051032,1964.42994543)(1182.32492311,1961.99617101)(1178.80017065,1957.12867737)
\curveto(1175.31733124,1952.30303495)(1173.57592368,1945.63113611)(1173.57594275,1937.11296082)
}
}
{
\newrgbcolor{curcolor}{0 0 0}
\pscustom[linestyle=none,fillstyle=solid,fillcolor=curcolor]
{
\newpath
\moveto(1294.86614966,1937.99415588)
\curveto(1294.86609112,1946.38644541)(1293.12468356,1952.89049774)(1289.64192175,1957.5063324)
\curveto(1286.20101495,1962.12205589)(1281.35444692,1964.42994543)(1275.10220312,1964.43000793)
\curveto(1268.89184341,1964.42994543)(1264.04527538,1962.12205589)(1260.5624845,1957.5063324)
\curveto(1257.12160677,1952.89049774)(1255.40118003,1946.38644541)(1255.4011991,1937.99415588)
\curveto(1255.40118003,1929.64375585)(1257.12160677,1923.16068433)(1260.5624845,1918.54492188)
\curveto(1264.04527538,1913.92912618)(1268.89184341,1911.62123664)(1275.10220312,1911.62124634)
\curveto(1281.35444692,1911.62123664)(1286.20101495,1913.92912618)(1289.64192175,1918.54492188)
\curveto(1293.12468356,1923.16068433)(1294.86609112,1929.64375585)(1294.86614966,1937.99415588)
\moveto(1306.44757056,1910.67710876)
\curveto(1306.44750044,1898.67607442)(1303.78293706,1889.75922848)(1298.45387244,1883.92654419)
\curveto(1293.12468356,1878.05193428)(1284.96314692,1875.11462032)(1273.96923804,1875.11459351)
\curveto(1269.89892249,1875.11462032)(1266.05943353,1875.42933253)(1262.45075964,1876.05873108)
\curveto(1258.84203352,1876.64621974)(1255.33823758,1877.56937556)(1251.93936133,1878.82820129)
\lineto(1251.93936133,1890.09490967)
\curveto(1255.33823758,1888.24860987)(1258.69516782,1886.88485696)(1262.01016211,1886.00364685)
\curveto(1265.32510504,1885.12246859)(1268.70301609,1884.68187149)(1272.1439054,1884.68185425)
\curveto(1279.73892424,1884.68187149)(1285.42472483,1886.67504882)(1289.20132422,1890.66139221)
\curveto(1292.97781786,1894.60579651)(1294.86609112,1900.58532849)(1294.86614966,1908.6000061)
\lineto(1294.86614966,1914.32777405)
\curveto(1292.47427833,1910.17356048)(1289.41107949,1907.06840001)(1285.67654395,1905.01228333)
\curveto(1281.94190971,1902.95616047)(1277.47299633,1901.92810059)(1272.26979041,1901.92810059)
\curveto(1263.62565911,1901.92810059)(1256.66002886,1905.22208838)(1251.37287878,1911.81007385)
\curveto(1246.08569862,1918.39803956)(1243.44211606,1927.12605818)(1243.44212317,1937.99415588)
\curveto(1243.44211606,1948.90414309)(1246.08569862,1957.65314252)(1251.37287878,1964.24118042)
\curveto(1256.66002886,1970.8290937)(1263.62565911,1974.12308149)(1272.26979041,1974.12315369)
\curveto(1277.47299633,1974.12308149)(1281.94190971,1973.09502161)(1285.67654395,1971.03897095)
\curveto(1289.41107949,1968.98278207)(1292.47427833,1965.8776216)(1294.86614966,1961.72348022)
\lineto(1294.86614966,1972.42370605)
\lineto(1306.44757056,1972.42370605)
\lineto(1306.44757056,1910.67710876)
}
}
{
\newrgbcolor{curcolor}{0 0 0}
\pscustom[linestyle=none,fillstyle=solid,fillcolor=curcolor]
{
\newpath
\moveto(1390.60169958,1940.07125854)
\lineto(1390.60169958,1934.40643311)
\lineto(1337.35234045,1934.40643311)
\curveto(1337.85586079,1926.43369131)(1340.24767359,1920.34925526)(1344.52778601,1916.15310669)
\curveto(1348.84980732,1911.9988913)(1354.85032012,1909.92179071)(1362.52934241,1909.92179871)
\curveto(1366.9772306,1909.92179071)(1371.27829746,1910.46729188)(1375.43255591,1911.55830383)
\curveto(1379.62866143,1912.64929653)(1383.78286259,1914.28580002)(1387.89517187,1916.46781921)
\lineto(1387.89517187,1905.51582336)
\curveto(1383.74090097,1903.7534314)(1379.48179573,1902.41065931)(1375.11784338,1901.48750305)
\curveto(1370.75377711,1900.56434768)(1366.32682536,1900.10276977)(1361.83697485,1900.10276794)
\curveto(1350.59121488,1900.10276977)(1341.67436894,1903.37577675)(1335.08641028,1909.92179871)
\curveto(1328.54037938,1916.46780467)(1325.2673724,1925.32170817)(1325.26737952,1936.48353577)
\curveto(1325.2673724,1948.0229489)(1328.37253287,1957.1705838)(1334.58287024,1963.9264679)
\curveto(1340.83513638,1970.72418963)(1349.24844278,1974.12308149)(1359.8228147,1974.12315369)
\curveto(1369.30610095,1974.12308149)(1376.79625154,1971.05988265)(1382.29328894,1964.93354797)
\curveto(1387.83215969,1958.84904891)(1390.60162714,1950.56162739)(1390.60169958,1940.07125854)
\moveto(1379.02027869,1943.47015381)
\curveto(1378.93629457,1949.80631809)(1377.15292538,1954.86269426)(1373.67016577,1958.63929749)
\curveto(1370.22925676,1962.41578729)(1365.65543932,1964.30406055)(1359.94869971,1964.30412292)
\curveto(1353.48656721,1964.30406055)(1348.30430615,1962.47872973)(1344.401901,1958.828125)
\curveto(1340.54140498,1955.17740647)(1338.3174387,1950.03710704)(1337.72999548,1943.4072113)
\lineto(1379.02027869,1943.47015381)
}
}
{
\newrgbcolor{curcolor}{0 0 0}
\pscustom[linestyle=none,fillstyle=solid,fillcolor=curcolor]
{
\newpath
\moveto(1409.86210608,1995.90126038)
\lineto(1426.98246741,1995.90126038)
\lineto(1468.65040564,1917.28607178)
\lineto(1468.65040564,1995.90126038)
\lineto(1480.9871366,1995.90126038)
\lineto(1480.9871366,1901.92810059)
\lineto(1463.86677527,1901.92810059)
\lineto(1422.19883704,1980.54328918)
\lineto(1422.19883704,1901.92810059)
\lineto(1409.86210608,1901.92810059)
\lineto(1409.86210608,1995.90126038)
}
}
{
\newrgbcolor{curcolor}{0 0 0}
\pscustom[linestyle=none,fillstyle=solid,fillcolor=curcolor]
{
\newpath
\moveto(1537.82421851,1937.36473083)
\curveto(1528.46673129,1937.3646954)(1521.98365977,1936.29467389)(1518.3749845,1934.15466309)
\curveto(1514.76625976,1932.01458783)(1512.96190975,1928.3639262)(1512.96192908,1923.20266724)
\curveto(1512.96190975,1919.09040642)(1514.30468185,1915.81739944)(1516.99024939,1913.38363647)
\curveto(1519.71773186,1910.99181223)(1523.41035512,1909.79590583)(1528.06813025,1909.7959137)
\curveto(1534.4882249,1909.79590583)(1539.62852432,1912.06183374)(1543.48904395,1916.59370422)
\curveto(1547.3914255,1921.16750701)(1549.3426412,1927.23096225)(1549.3426969,1934.78408813)
\lineto(1549.3426969,1937.36473083)
\lineto(1537.82421851,1937.36473083)
\moveto(1560.9241178,1942.14836121)
\lineto(1560.9241178,1901.92810059)
\lineto(1549.3426969,1901.92810059)
\lineto(1549.3426969,1912.62832642)
\curveto(1546.69905863,1908.34822966)(1543.40507084,1905.18012675)(1539.46072363,1903.12400818)
\curveto(1535.51628478,1901.10984884)(1530.69069757,1900.10276977)(1524.98394751,1900.10276794)
\curveto(1517.76651616,1900.10276977)(1512.01777312,1902.11692791)(1507.73770117,1906.14524841)
\curveto(1503.49956265,1910.21552211)(1501.38050044,1915.64955293)(1501.38050818,1922.44735718)
\curveto(1501.38050044,1930.37808434)(1504.024083,1936.35761633)(1509.31126379,1940.38597107)
\curveto(1514.64037487,1944.4142489)(1522.57112256,1946.42840704)(1533.10353064,1946.42845154)
\lineto(1549.3426969,1946.42845154)
\lineto(1549.3426969,1947.56141663)
\curveto(1549.3426412,1952.89049774)(1547.58025282,1957.00273728)(1544.05552649,1959.89814758)
\curveto(1540.57266095,1962.83540357)(1535.66315048,1964.30406055)(1529.32698035,1964.30412292)
\curveto(1525.29862837,1964.30406055)(1521.37521616,1963.82150183)(1517.55673193,1962.85644531)
\curveto(1513.73819987,1961.89126694)(1510.06655742,1960.44359078)(1506.54179358,1958.51341248)
\lineto(1506.54179358,1969.21363831)
\curveto(1510.7799051,1970.85007451)(1514.89214464,1972.06696172)(1518.87852454,1972.86430359)
\curveto(1522.86485395,1973.70346521)(1526.74630454,1974.12308149)(1530.52288794,1974.12315369)
\curveto(1540.71952665,1974.12308149)(1548.33556212,1971.47949893)(1553.37101721,1966.19239807)
\curveto(1558.40635283,1960.90516868)(1560.92405051,1952.89049774)(1560.9241178,1942.14836121)
}
}
{
\newrgbcolor{curcolor}{0 0 0}
\pscustom[linestyle=none,fillstyle=solid,fillcolor=curcolor]
{
\newpath
\moveto(1639.97990393,1958.8910675)
\curveto(1642.87518923,1964.09425241)(1646.33702353,1967.93374137)(1650.36541724,1970.4095459)
\curveto(1654.3936561,1972.88521347)(1659.13532006,1974.12308149)(1664.59042334,1974.12315369)
\curveto(1671.93361659,1974.12308149)(1677.59843636,1971.54244137)(1681.58489966,1966.38122559)
\curveto(1685.57114568,1961.26184252)(1687.564323,1953.96051926)(1687.56443762,1944.47723389)
\lineto(1687.56443762,1901.92810059)
\lineto(1675.92007422,1901.92810059)
\lineto(1675.92007422,1944.09957886)
\curveto(1675.91997124,1950.85535879)(1674.72406485,1955.86977333)(1672.33235144,1959.14283752)
\curveto(1669.94043926,1962.41578729)(1666.28977763,1964.05229078)(1661.38035559,1964.05235291)
\curveto(1655.37975436,1964.05229078)(1650.6380904,1962.05911345)(1647.15534949,1958.07281494)
\curveto(1643.67246016,1954.08640414)(1641.9310526,1948.65237332)(1641.93112158,1941.77070618)
\lineto(1641.93112158,1901.92810059)
\lineto(1630.28675818,1901.92810059)
\lineto(1630.28675818,1944.09957886)
\curveto(1630.28670084,1950.89732042)(1629.09079444,1955.91173496)(1626.6990354,1959.14283752)
\curveto(1624.30716885,1962.41578729)(1620.61454559,1964.05229078)(1615.62115454,1964.05235291)
\curveto(1609.70452232,1964.05229078)(1605.00481999,1962.03813264)(1601.52203345,1958.00987244)
\curveto(1598.03918975,1954.0234617)(1596.29778219,1948.61041169)(1596.29780554,1941.77070618)
\lineto(1596.29780554,1901.92810059)
\lineto(1584.65344214,1901.92810059)
\lineto(1584.65344214,1972.42370605)
\lineto(1596.29780554,1972.42370605)
\lineto(1596.29780554,1961.47171021)
\curveto(1598.94136475,1965.79369834)(1602.10946766,1968.98278207)(1605.80212378,1971.03897095)
\curveto(1609.49471418,1973.09502161)(1613.87970431,1974.12308149)(1618.9571073,1974.12315369)
\curveto(1624.0763799,1974.12308149)(1628.41940839,1972.82227102)(1631.98620581,1970.22071838)
\curveto(1635.59484678,1967.61902916)(1638.25941015,1963.84248264)(1639.97990393,1958.8910675)
}
}
{
\newrgbcolor{curcolor}{0 0 0}
\pscustom[linestyle=none,fillstyle=solid,fillcolor=curcolor]
{
\newpath
\moveto(1770.77442908,1940.07125854)
\lineto(1770.77442908,1934.40643311)
\lineto(1717.52506995,1934.40643311)
\curveto(1718.02859028,1926.43369131)(1720.42040308,1920.34925526)(1724.7005155,1916.15310669)
\curveto(1729.02253681,1911.9988913)(1735.02304961,1909.92179071)(1742.7020719,1909.92179871)
\curveto(1747.14996009,1909.92179071)(1751.45102695,1910.46729188)(1755.6052854,1911.55830383)
\curveto(1759.80139092,1912.64929653)(1763.95559209,1914.28580002)(1768.06790137,1916.46781921)
\lineto(1768.06790137,1905.51582336)
\curveto(1763.91363046,1903.7534314)(1759.65452522,1902.41065931)(1755.29057288,1901.48750305)
\curveto(1750.92650661,1900.56434768)(1746.49955486,1900.10276977)(1742.00970435,1900.10276794)
\curveto(1730.76394437,1900.10276977)(1721.84709843,1903.37577675)(1715.25913977,1909.92179871)
\curveto(1708.71310888,1916.46780467)(1705.4401019,1925.32170817)(1705.44010901,1936.48353577)
\curveto(1705.4401019,1948.0229489)(1708.54526237,1957.1705838)(1714.75559973,1963.9264679)
\curveto(1721.00786587,1970.72418963)(1729.42117228,1974.12308149)(1739.99554419,1974.12315369)
\curveto(1749.47883044,1974.12308149)(1756.96898103,1971.05988265)(1762.46601843,1964.93354797)
\curveto(1768.00488918,1958.84904891)(1770.77435663,1950.56162739)(1770.77442908,1940.07125854)
\moveto(1759.19300818,1943.47015381)
\curveto(1759.10902406,1949.80631809)(1757.32565487,1954.86269426)(1753.84289526,1958.63929749)
\curveto(1750.40198626,1962.41578729)(1745.82816881,1964.30406055)(1740.1214292,1964.30412292)
\curveto(1733.6592967,1964.30406055)(1728.47703565,1962.47872973)(1724.57463049,1958.828125)
\curveto(1720.71413447,1955.17740647)(1718.49016819,1950.03710704)(1717.90272498,1943.4072113)
\lineto(1759.19300818,1943.47015381)
}
}
{
\newrgbcolor{curcolor}{0 0 0}
\pscustom[linestyle=none,fillstyle=solid,fillcolor=curcolor]
{
\newpath
\moveto(1817.35188269,1999.74075317)
\curveto(1811.72898458,1990.08948093)(1807.55380259,1980.54321057)(1804.82632422,1971.10191345)
\curveto(1802.09879096,1961.66047799)(1800.73503805,1952.09322681)(1800.7350614,1942.40013123)
\curveto(1800.73503805,1932.70695469)(1802.09879096,1923.09774189)(1804.82632422,1913.57246399)
\curveto(1807.59576422,1904.08912443)(1811.7709462,1894.54285406)(1817.35188269,1884.93362427)
\lineto(1807.28108191,1884.93362427)
\curveto(1800.98680782,1894.79462383)(1796.26612467,1904.48775989)(1793.11901831,1914.01306152)
\curveto(1790.01384211,1923.53833899)(1788.46126187,1933.00068609)(1788.46127295,1942.40013123)
\curveto(1788.46126187,1951.75753379)(1790.01384211,1961.17791927)(1793.11901831,1970.66131592)
\curveto(1796.22416304,1980.1445751)(1800.94484619,1989.83771116)(1807.28108191,1999.74075317)
\lineto(1817.35188269,1999.74075317)
}
}
{
\newrgbcolor{curcolor}{0 0 0}
\pscustom[linestyle=none,fillstyle=solid,fillcolor=curcolor]
{
\newpath
\moveto(1838.0599668,1999.74075317)
\lineto(1848.13076758,1999.74075317)
\curveto(1854.42499138,1989.83771116)(1859.12469371,1980.1445751)(1862.22988867,1970.66131592)
\curveto(1865.37697628,1961.17791927)(1866.95053732,1951.75753379)(1866.95057654,1942.40013123)
\curveto(1866.95053732,1933.00068609)(1865.37697628,1923.53833899)(1862.22988867,1914.01306152)
\curveto(1859.12469371,1904.48775989)(1854.42499138,1894.79462383)(1848.13076758,1884.93362427)
\lineto(1838.0599668,1884.93362427)
\curveto(1843.64085299,1894.54285406)(1847.79505416,1904.08912443)(1850.52258276,1913.57246399)
\curveto(1853.29202742,1923.09774189)(1854.67676115,1932.70695469)(1854.67678809,1942.40013123)
\curveto(1854.67676115,1952.09322681)(1853.29202742,1961.66047799)(1850.52258276,1971.10191345)
\curveto(1847.79505416,1980.54321057)(1843.64085299,1990.08948093)(1838.0599668,1999.74075317)
}
}
{
\newrgbcolor{curcolor}{0 0 0}
\pscustom[linestyle=none,fillstyle=solid,fillcolor=curcolor]
{
}
}
{
\newrgbcolor{curcolor}{0 0 0}
\pscustom[linestyle=none,fillstyle=solid,fillcolor=curcolor]
{
\newpath
\moveto(1934.23611426,1917.91549683)
\lineto(1947.51698279,1917.91549683)
\lineto(1947.51698279,1901.92810059)
\lineto(1934.23611426,1901.92810059)
\lineto(1934.23611426,1917.91549683)
\moveto(1934.23611426,1968.58421326)
\lineto(1947.51698279,1968.58421326)
\lineto(1947.51698279,1952.59681702)
\lineto(1934.23611426,1952.59681702)
\lineto(1934.23611426,1968.58421326)
}
}
{
\newrgbcolor{curcolor}{0 0 0}
\pscustom[linestyle=none,fillstyle=solid,fillcolor=curcolor]
{
}
}
{
\newrgbcolor{curcolor}{0 0 0}
\pscustom[linestyle=none,fillstyle=solid,fillcolor=curcolor]
{
\newpath
\moveto(2072.709625,1992.81707764)
\lineto(2072.709625,1980.41740417)
\curveto(2067.8839688,1982.72521522)(2063.33113217,1984.44564197)(2059.05110144,1985.57868958)
\curveto(2054.77096006,1986.71156988)(2050.63773971,1987.27805186)(2046.65142798,1987.27813721)
\curveto(2039.72771644,1987.27805186)(2034.37760887,1985.93527976)(2030.60108923,1983.24981689)
\curveto(2026.86647747,1980.56419138)(2024.99918502,1976.74568324)(2024.9992063,1971.79428101)
\curveto(2024.99918502,1967.64000997)(2026.23705305,1964.49288787)(2028.71281409,1962.35290527)
\curveto(2031.23048678,1960.25476345)(2035.97215074,1958.55531752)(2042.93782019,1957.25456238)
\lineto(2050.61680579,1955.68099976)
\curveto(2060.10008681,1953.876596)(2067.08669787,1950.68751228)(2071.57665991,1946.11373901)
\curveto(2076.10844788,1941.58183901)(2078.37437579,1935.49740295)(2078.37445044,1927.8604126)
\curveto(2078.37437579,1918.7547134)(2075.31117695,1911.8520256)(2069.18484473,1907.15232849)
\curveto(2063.10034321,1902.45262094)(2054.16251646,1900.10276977)(2042.37133765,1900.10276794)
\curveto(2037.92336644,1900.10276977)(2033.18170248,1900.60630931)(2028.14633154,1901.61338806)
\curveto(2023.15287339,1902.62046745)(2017.97061234,1904.11010524)(2012.59953284,1906.08230591)
\lineto(2012.59953284,1919.17434692)
\curveto(2017.7608042,1916.27897735)(2022.81718037,1914.09697269)(2027.76867651,1912.62832642)
\curveto(2032.72012457,1911.15965874)(2037.58767341,1910.42533025)(2042.37133765,1910.42533875)
\curveto(2049.63066064,1910.42533025)(2055.23253797,1911.8520256)(2059.17698645,1914.70542908)
\curveto(2063.12132403,1917.558807)(2065.09352054,1921.62908491)(2065.09358191,1926.91627502)
\curveto(2065.09352054,1931.53202911)(2063.66682519,1935.14072912)(2060.81349158,1937.74238586)
\curveto(2058.00200541,1940.34397098)(2053.36524553,1942.29518668)(2046.903198,1943.59603882)
\lineto(2039.1612699,1945.10665894)
\curveto(2029.67790654,1946.99488902)(2022.81718037,1949.95318379)(2018.5790708,1953.98155212)
\curveto(2014.34093152,1958.00981635)(2012.22186931,1963.61169369)(2012.22187781,1970.78720093)
\curveto(2012.22186931,1979.0955344)(2015.13820245,1985.64154837)(2020.97088599,1990.42526245)
\curveto(2026.84549665,1995.20879954)(2034.92311004,1997.60061234)(2045.20375037,1997.60070801)
\curveto(2049.60967982,1997.60061234)(2054.09957401,1997.20197687)(2058.67344641,1996.40480042)
\curveto(2063.24720891,1995.60743501)(2067.92593043,1994.41152861)(2072.709625,1992.81707764)
}
}
{
\newrgbcolor{curcolor}{0 0 0}
\pscustom[linestyle=none,fillstyle=solid,fillcolor=curcolor]
{
\newpath
\moveto(2109.40510535,1992.43942261)
\lineto(2109.40510535,1972.42370605)
\lineto(2133.2603147,1972.42370605)
\lineto(2133.2603147,1963.42292786)
\lineto(2109.40510535,1963.42292786)
\lineto(2109.40510535,1925.15388489)
\curveto(2109.40508174,1919.40511863)(2110.18137186,1915.71249537)(2111.73397803,1914.07600403)
\curveto(2113.32849396,1912.43948839)(2116.5385585,1911.62123664)(2121.36418127,1911.62124634)
\lineto(2133.2603147,1911.62124634)
\lineto(2133.2603147,1901.92810059)
\lineto(2121.36418127,1901.92810059)
\curveto(2112.42631896,1901.92810059)(2106.25795965,1903.58558489)(2102.85908484,1906.90055847)
\curveto(2099.46017592,1910.25748374)(2097.76072998,1916.34191979)(2097.76074194,1925.15388489)
\lineto(2097.76074194,1963.42292786)
\lineto(2089.26350378,1963.42292786)
\lineto(2089.26350378,1972.42370605)
\lineto(2097.76074194,1972.42370605)
\lineto(2097.76074194,1992.43942261)
\lineto(2109.40510535,1992.43942261)
}
}
{
\newrgbcolor{curcolor}{0 0 0}
\pscustom[linestyle=none,fillstyle=solid,fillcolor=curcolor]
{
\newpath
\moveto(2189.40502905,1961.59759521)
\curveto(2188.10416559,1962.35284485)(2186.67747024,1962.89834601)(2185.12493872,1963.23410034)
\curveto(2183.6142714,1963.61169369)(2181.93580628,1963.80052101)(2180.08953833,1963.80058289)
\curveto(2173.54348069,1963.80052101)(2168.50808533,1961.66047799)(2164.98333716,1957.38044739)
\curveto(2161.50049346,1953.14226751)(2159.7590859,1947.03685064)(2159.75910925,1939.06417847)
\lineto(2159.75910925,1901.92810059)
\lineto(2148.11474585,1901.92810059)
\lineto(2148.11474585,1972.42370605)
\lineto(2159.75910925,1972.42370605)
\lineto(2159.75910925,1961.47171021)
\curveto(2162.19286032,1965.75173671)(2165.36096323,1968.91983962)(2169.26342749,1970.97602844)
\curveto(2173.16582603,1973.07404079)(2177.90748999,1974.12308149)(2183.48843359,1974.12315369)
\curveto(2184.28565744,1974.12308149)(2185.16685163,1974.06013905)(2186.1320188,1973.93432617)
\curveto(2187.09708652,1973.85033091)(2188.16710803,1973.70346521)(2189.34208655,1973.49372864)
\lineto(2189.40502905,1961.59759521)
}
}
{
\newrgbcolor{curcolor}{0 0 0}
\pscustom[linestyle=none,fillstyle=solid,fillcolor=curcolor]
{
\newpath
\moveto(2201.6788175,1972.42370605)
\lineto(2213.2602384,1972.42370605)
\lineto(2213.2602384,1901.92810059)
\lineto(2201.6788175,1901.92810059)
\lineto(2201.6788175,1972.42370605)
\moveto(2201.6788175,1999.86663818)
\lineto(2213.2602384,1999.86663818)
\lineto(2213.2602384,1985.20103455)
\lineto(2201.6788175,1985.20103455)
\lineto(2201.6788175,1999.86663818)
}
}
{
\newrgbcolor{curcolor}{0 0 0}
\pscustom[linestyle=none,fillstyle=solid,fillcolor=curcolor]
{
\newpath
\moveto(2296.02963232,1944.47723389)
\lineto(2296.02963232,1901.92810059)
\lineto(2284.44821143,1901.92810059)
\lineto(2284.44821143,1944.09957886)
\curveto(2284.44815226,1950.77143553)(2283.14734179,1955.76486926)(2280.54577612,1959.07989502)
\curveto(2277.94409993,1962.39480648)(2274.04166853,1964.05229078)(2268.83847021,1964.05235291)
\curveto(2262.58614409,1964.05229078)(2257.65565281,1962.05911345)(2254.04698157,1958.07281494)
\curveto(2250.4382528,1954.08640414)(2248.6339028,1948.65237332)(2248.63392615,1941.77070618)
\lineto(2248.63392615,1901.92810059)
\lineto(2236.98956274,1901.92810059)
\lineto(2236.98956274,1972.42370605)
\lineto(2248.63392615,1972.42370605)
\lineto(2248.63392615,1961.47171021)
\curveto(2251.40337024,1965.70977509)(2254.65539641,1968.877878)(2258.3900144,1970.97602844)
\curveto(2262.16652781,1973.07404079)(2266.50955631,1974.12308149)(2271.41911292,1974.12315369)
\curveto(2279.51766097,1974.12308149)(2285.64405866,1971.60538381)(2289.79832434,1966.5700531)
\curveto(2293.95246099,1961.57655473)(2296.02956158,1954.21228902)(2296.02963232,1944.47723389)
}
}
{
\newrgbcolor{curcolor}{0 0 0}
\pscustom[linestyle=none,fillstyle=solid,fillcolor=curcolor]
{
\newpath
\moveto(2365.64404272,1937.99415588)
\curveto(2365.64398419,1946.38644541)(2363.90257663,1952.89049774)(2360.41981482,1957.5063324)
\curveto(2356.97890801,1962.12205589)(2352.13233998,1964.42994543)(2345.88009619,1964.43000793)
\curveto(2339.66973648,1964.42994543)(2334.82316845,1962.12205589)(2331.34037756,1957.5063324)
\curveto(2327.89949984,1952.89049774)(2326.17907309,1946.38644541)(2326.17909216,1937.99415588)
\curveto(2326.17907309,1929.64375585)(2327.89949984,1923.16068433)(2331.34037756,1918.54492188)
\curveto(2334.82316845,1913.92912618)(2339.66973648,1911.62123664)(2345.88009619,1911.62124634)
\curveto(2352.13233998,1911.62123664)(2356.97890801,1913.92912618)(2360.41981482,1918.54492188)
\curveto(2363.90257663,1923.16068433)(2365.64398419,1929.64375585)(2365.64404272,1937.99415588)
\moveto(2377.22546362,1910.67710876)
\curveto(2377.22539351,1898.67607442)(2374.56083013,1889.75922848)(2369.2317655,1883.92654419)
\curveto(2363.90257663,1878.05193428)(2355.74103999,1875.11462032)(2344.7471311,1875.11459351)
\curveto(2340.67681555,1875.11462032)(2336.83732659,1875.42933253)(2333.22865271,1876.05873108)
\curveto(2329.61992658,1876.64621974)(2326.11613065,1877.56937556)(2322.71725439,1878.82820129)
\lineto(2322.71725439,1890.09490967)
\curveto(2326.11613065,1888.24860987)(2329.47306089,1886.88485696)(2332.78805518,1886.00364685)
\curveto(2336.1029981,1885.12246859)(2339.48090915,1884.68187149)(2342.92179846,1884.68185425)
\curveto(2350.51681731,1884.68187149)(2356.2026179,1886.67504882)(2359.97921729,1890.66139221)
\curveto(2363.75571093,1894.60579651)(2365.64398419,1900.58532849)(2365.64404272,1908.6000061)
\lineto(2365.64404272,1914.32777405)
\curveto(2363.25217139,1910.17356048)(2360.18897255,1907.06840001)(2356.45443701,1905.01228333)
\curveto(2352.71980278,1902.95616047)(2348.2508894,1901.92810059)(2343.04768347,1901.92810059)
\curveto(2334.40355217,1901.92810059)(2327.43792193,1905.22208838)(2322.15077185,1911.81007385)
\curveto(2316.86359168,1918.39803956)(2314.22000912,1927.12605818)(2314.22001624,1937.99415588)
\curveto(2314.22000912,1948.90414309)(2316.86359168,1957.65314252)(2322.15077185,1964.24118042)
\curveto(2327.43792193,1970.8290937)(2334.40355217,1974.12308149)(2343.04768347,1974.12315369)
\curveto(2348.2508894,1974.12308149)(2352.71980278,1973.09502161)(2356.45443701,1971.03897095)
\curveto(2360.18897255,1968.98278207)(2363.25217139,1965.8776216)(2365.64404272,1961.72348022)
\lineto(2365.64404272,1972.42370605)
\lineto(2377.22546362,1972.42370605)
\lineto(2377.22546362,1910.67710876)
}
}
{
\newrgbcolor{curcolor}{0.51764709 0.74509805 0.51764709}
\pscustom[linestyle=none,fillstyle=solid,fillcolor=curcolor]
{
\newpath
\moveto(407.18221875,1781.84820703)
\curveto(407.18221875,1762.4319463)(391.44222949,1746.69195703)(372.02596875,1746.69195703)
\curveto(352.60970801,1746.69195703)(336.86971875,1762.4319463)(336.86971875,1781.84820703)
\curveto(336.86971875,1801.26446777)(352.60970801,1817.00445703)(372.02596875,1817.00445703)
\curveto(391.44222949,1817.00445703)(407.18221875,1801.26446777)(407.18221875,1781.84820703)
\closepath
}
}
{
\newrgbcolor{curcolor}{0.01176471 0.50196081 0.28235295}
\pscustom[linewidth=11.71875,linecolor=curcolor]
{
\newpath
\moveto(407.18221875,1781.84820703)
\curveto(407.18221875,1762.4319463)(391.44222949,1746.69195703)(372.02596875,1746.69195703)
\curveto(352.60970801,1746.69195703)(336.86971875,1762.4319463)(336.86971875,1781.84820703)
\curveto(336.86971875,1801.26446777)(352.60970801,1817.00445703)(372.02596875,1817.00445703)
\curveto(391.44222949,1817.00445703)(407.18221875,1801.26446777)(407.18221875,1781.84820703)
\closepath
}
}
{
\newrgbcolor{curcolor}{0 0 0}
\pscustom[linestyle=none,fillstyle=solid,fillcolor=curcolor]
{
\newpath
\moveto(536.03131079,1788.00209045)
\curveto(536.03125225,1796.39437998)(534.28984469,1802.89843231)(530.80708289,1807.51426697)
\curveto(527.36617608,1812.12999047)(522.51960805,1814.43788)(516.26736426,1814.4379425)
\curveto(510.05700455,1814.43788)(505.21043652,1812.12999047)(501.72764563,1807.51426697)
\curveto(498.2867679,1802.89843231)(496.56634116,1796.39437998)(496.56636023,1788.00209045)
\curveto(496.56634116,1779.65169042)(498.2867679,1773.1686189)(501.72764563,1768.55285645)
\curveto(505.21043652,1763.93706075)(510.05700455,1761.62917122)(516.26736426,1761.62918091)
\curveto(522.51960805,1761.62917122)(527.36617608,1763.93706075)(530.80708289,1768.55285645)
\curveto(534.28984469,1773.1686189)(536.03125225,1779.65169042)(536.03131079,1788.00209045)
\moveto(547.61273169,1760.68504333)
\curveto(547.61266157,1748.68400899)(544.9480982,1739.76716305)(539.61903357,1733.93447876)
\curveto(534.28984469,1728.05986885)(526.12830806,1725.12255489)(515.13439917,1725.12252808)
\curveto(511.06408362,1725.12255489)(507.22459466,1725.4372671)(503.61592078,1726.06666565)
\curveto(500.00719465,1726.65415431)(496.50339872,1727.57731013)(493.10452246,1728.83613586)
\lineto(493.10452246,1740.10284424)
\curveto(496.50339872,1738.25654444)(499.86032895,1736.89279153)(503.17532324,1736.01158142)
\curveto(506.49026617,1735.13040316)(509.86817722,1734.68980606)(513.30906653,1734.68978882)
\curveto(520.90408537,1734.68980606)(526.58988596,1736.68298339)(530.36648535,1740.66932678)
\curveto(534.142979,1744.61373108)(536.03125225,1750.59326306)(536.03131079,1758.60794067)
\lineto(536.03131079,1764.33570862)
\curveto(533.63943946,1760.18149505)(530.57624062,1757.07633458)(526.84170508,1755.0202179)
\curveto(523.10707084,1752.96409504)(518.63815746,1751.93603516)(513.43495154,1751.93603516)
\curveto(504.79082024,1751.93603516)(497.82519,1755.23002295)(492.53803992,1761.81800842)
\curveto(487.25085975,1768.40597413)(484.60727719,1777.13399275)(484.6072843,1788.00209045)
\curveto(484.60727719,1798.91207766)(487.25085975,1807.66107709)(492.53803992,1814.24911499)
\curveto(497.82519,1820.83702827)(504.79082024,1824.13101606)(513.43495154,1824.13108826)
\curveto(518.63815746,1824.13101606)(523.10707084,1823.10295618)(526.84170508,1821.04690552)
\curveto(530.57624062,1818.99071664)(533.63943946,1815.88555617)(536.03131079,1811.73141479)
\lineto(536.03131079,1822.43164062)
\lineto(547.61273169,1822.43164062)
\lineto(547.61273169,1760.68504333)
}
}
{
\newrgbcolor{curcolor}{0 0 0}
\pscustom[linestyle=none,fillstyle=solid,fillcolor=curcolor]
{
\newpath
\moveto(631.76686072,1790.07919312)
\lineto(631.76686072,1784.41436768)
\lineto(578.51750159,1784.41436768)
\curveto(579.02102192,1776.44162589)(581.41283472,1770.35718983)(585.69294714,1766.16104126)
\curveto(590.01496845,1762.00682587)(596.01548125,1759.92972528)(603.69450354,1759.92973328)
\curveto(608.14239173,1759.92972528)(612.4434586,1760.47522645)(616.59771704,1761.5662384)
\curveto(620.79382256,1762.6572311)(624.94802373,1764.29373459)(629.06033301,1766.47575378)
\lineto(629.06033301,1755.52375793)
\curveto(624.9060621,1753.76136597)(620.64695686,1752.41859388)(616.28300452,1751.49543762)
\curveto(611.91893825,1750.57228225)(607.4919865,1750.11070434)(603.00213599,1750.11070251)
\curveto(591.75637601,1750.11070434)(582.83953007,1753.38371132)(576.25157141,1759.92973328)
\curveto(569.70554052,1766.47573924)(566.43253354,1775.32964274)(566.43254065,1786.49147034)
\curveto(566.43253354,1798.03088347)(569.53769401,1807.17851837)(575.74803137,1813.93440247)
\curveto(582.00029751,1820.7321242)(590.41360392,1824.13101606)(600.98797583,1824.13108826)
\curveto(610.47126208,1824.13101606)(617.96141267,1821.06781722)(623.45845007,1814.94148254)
\curveto(628.99732083,1808.85698348)(631.76678827,1800.56956196)(631.76686072,1790.07919312)
\moveto(620.18543982,1793.47808838)
\curveto(620.1014557,1799.81425266)(618.31808651,1804.87062883)(614.8353269,1808.64723206)
\curveto(611.3944179,1812.42372186)(606.82060045,1814.31199512)(601.11386084,1814.3120575)
\curveto(594.65172834,1814.31199512)(589.46946729,1812.4866643)(585.56706213,1808.83605957)
\curveto(581.70656611,1805.18534104)(579.48259983,1800.04504161)(578.89515662,1793.41514587)
\lineto(620.18543982,1793.47808838)
}
}
{
\newrgbcolor{curcolor}{0 0 0}
\pscustom[linestyle=none,fillstyle=solid,fillcolor=curcolor]
{
\newpath
\moveto(661.97926306,1842.44735718)
\lineto(661.97926306,1822.43164062)
\lineto(685.83447241,1822.43164062)
\lineto(685.83447241,1813.43086243)
\lineto(661.97926306,1813.43086243)
\lineto(661.97926306,1775.16181946)
\curveto(661.97923946,1769.4130532)(662.75552958,1765.72042994)(664.30813574,1764.0839386)
\curveto(665.90265167,1762.44742296)(669.11271621,1761.62917122)(673.93833899,1761.62918091)
\lineto(685.83447241,1761.62918091)
\lineto(685.83447241,1751.93603516)
\lineto(673.93833899,1751.93603516)
\curveto(665.00047667,1751.93603516)(658.83211736,1753.59351946)(655.43324255,1756.90849304)
\curveto(652.03433363,1760.26541831)(650.3348877,1766.34985436)(650.33489966,1775.16181946)
\lineto(650.33489966,1813.43086243)
\lineto(641.8376615,1813.43086243)
\lineto(641.8376615,1822.43164062)
\lineto(650.33489966,1822.43164062)
\lineto(650.33489966,1842.44735718)
\lineto(661.97926306,1842.44735718)
}
}
{
\newrgbcolor{curcolor}{0 0 0}
\pscustom[linestyle=none,fillstyle=solid,fillcolor=curcolor]
{
\newpath
\moveto(701.63304114,1845.90919495)
\lineto(761.05076575,1845.90919495)
\lineto(761.05076575,1835.20896912)
\lineto(714.34742712,1835.20896912)
\lineto(714.34742712,1807.38838196)
\lineto(759.0995481,1807.38838196)
\lineto(759.0995481,1796.68815613)
\lineto(714.34742712,1796.68815613)
\lineto(714.34742712,1762.63626099)
\lineto(762.18373083,1762.63626099)
\lineto(762.18373083,1751.93603516)
\lineto(701.63304114,1751.93603516)
\lineto(701.63304114,1845.90919495)
}
}
{
\newrgbcolor{curcolor}{0 0 0}
\pscustom[linestyle=none,fillstyle=solid,fillcolor=curcolor]
{
\newpath
\moveto(774.14280676,1822.43164062)
\lineto(786.41659521,1822.43164062)
\lineto(808.44647192,1763.26568604)
\lineto(830.47634863,1822.43164062)
\lineto(842.75013708,1822.43164062)
\lineto(816.31428503,1751.93603516)
\lineto(800.57865881,1751.93603516)
\lineto(774.14280676,1822.43164062)
}
}
{
\newrgbcolor{curcolor}{0 0 0}
\pscustom[linestyle=none,fillstyle=solid,fillcolor=curcolor]
{
\newpath
\moveto(858.73753333,1822.43164062)
\lineto(870.31895422,1822.43164062)
\lineto(870.31895422,1751.93603516)
\lineto(858.73753333,1751.93603516)
\lineto(858.73753333,1822.43164062)
\moveto(858.73753333,1849.87457275)
\lineto(870.31895422,1849.87457275)
\lineto(870.31895422,1835.20896912)
\lineto(858.73753333,1835.20896912)
\lineto(858.73753333,1849.87457275)
}
}
{
\newrgbcolor{curcolor}{0 0 0}
\pscustom[linestyle=none,fillstyle=solid,fillcolor=curcolor]
{
\newpath
\moveto(940.8775022,1811.73141479)
\lineto(940.8775022,1849.87457275)
\lineto(952.4589231,1849.87457275)
\lineto(952.4589231,1751.93603516)
\lineto(940.8775022,1751.93603516)
\lineto(940.8775022,1762.51037598)
\curveto(938.44366924,1758.31420261)(935.35948958,1755.18806132)(931.62495398,1753.13194275)
\curveto(927.93228144,1751.11778341)(923.48434887,1750.11070434)(918.28114294,1750.11070251)
\curveto(909.76289653,1750.11070434)(902.8182471,1753.5095962)(897.44717383,1760.30738831)
\curveto(892.11803197,1767.10516366)(889.4534686,1776.04299042)(889.45347571,1787.12089539)
\curveto(889.4534686,1798.19872998)(892.11803197,1807.13655674)(897.44717383,1813.93440247)
\curveto(902.8182471,1820.7321242)(909.76289653,1824.13101606)(918.28114294,1824.13108826)
\curveto(923.48434887,1824.13101606)(927.93228144,1823.10295618)(931.62495398,1821.04690552)
\curveto(935.35948958,1819.03267826)(938.44366924,1815.9275178)(940.8775022,1811.73141479)
\moveto(901.41255164,1787.12089539)
\curveto(901.41253256,1778.60264972)(903.15394012,1771.90977007)(906.63677954,1767.04223633)
\curveto(910.16153199,1762.21663401)(914.98711921,1759.8038404)(921.11355566,1759.80384827)
\curveto(927.23991457,1759.8038404)(932.06550179,1762.21663401)(935.59033179,1767.04223633)
\curveto(939.11505529,1771.90977007)(940.87744366,1778.60264972)(940.8775022,1787.12089539)
\curveto(940.87744366,1795.63907068)(939.11505529,1802.31096952)(935.59033179,1807.13661194)
\curveto(932.06550179,1812.00410558)(927.23991457,1814.43788)(921.11355566,1814.4379425)
\curveto(914.98711921,1814.43788)(910.16153199,1812.00410558)(906.63677954,1807.13661194)
\curveto(903.15394012,1802.31096952)(901.41253256,1795.63907068)(901.41255164,1787.12089539)
}
}
{
\newrgbcolor{curcolor}{0 0 0}
\pscustom[linestyle=none,fillstyle=solid,fillcolor=curcolor]
{
\newpath
\moveto(1036.61305212,1790.07919312)
\lineto(1036.61305212,1784.41436768)
\lineto(983.36369299,1784.41436768)
\curveto(983.86721333,1776.44162589)(986.25902612,1770.35718983)(990.53913855,1766.16104126)
\curveto(994.86115986,1762.00682587)(1000.86167266,1759.92972528)(1008.54069495,1759.92973328)
\curveto(1012.98858314,1759.92972528)(1017.28965,1760.47522645)(1021.44390845,1761.5662384)
\curveto(1025.64001397,1762.6572311)(1029.79421513,1764.29373459)(1033.90652441,1766.47575378)
\lineto(1033.90652441,1755.52375793)
\curveto(1029.75225351,1753.76136597)(1025.49314827,1752.41859388)(1021.12919592,1751.49543762)
\curveto(1016.76512965,1750.57228225)(1012.3381779,1750.11070434)(1007.84832739,1750.11070251)
\curveto(996.60256742,1750.11070434)(987.68572148,1753.38371132)(981.09776282,1759.92973328)
\curveto(974.55173192,1766.47573924)(971.27872494,1775.32964274)(971.27873206,1786.49147034)
\curveto(971.27872494,1798.03088347)(974.38388541,1807.17851837)(980.59422278,1813.93440247)
\curveto(986.84648892,1820.7321242)(995.25979532,1824.13101606)(1005.83416724,1824.13108826)
\curveto(1015.31745349,1824.13101606)(1022.80760408,1821.06781722)(1028.30464148,1814.94148254)
\curveto(1033.84351223,1808.85698348)(1036.61297968,1800.56956196)(1036.61305212,1790.07919312)
\moveto(1025.03163123,1793.47808838)
\curveto(1024.9476471,1799.81425266)(1023.16427792,1804.87062883)(1019.68151831,1808.64723206)
\curveto(1016.2406093,1812.42372186)(1011.66679185,1814.31199512)(1005.96005225,1814.3120575)
\curveto(999.49791975,1814.31199512)(994.31565869,1812.4866643)(990.41325354,1808.83605957)
\curveto(986.55275752,1805.18534104)(984.32879124,1800.04504161)(983.74134802,1793.41514587)
\lineto(1025.03163123,1793.47808838)
}
}
{
\newrgbcolor{curcolor}{0 0 0}
\pscustom[linestyle=none,fillstyle=solid,fillcolor=curcolor]
{
\newpath
\moveto(1113.96939062,1794.48516846)
\lineto(1113.96939062,1751.93603516)
\lineto(1102.38796973,1751.93603516)
\lineto(1102.38796973,1794.10751343)
\curveto(1102.38791056,1800.7793701)(1101.08710009,1805.77280383)(1098.48553442,1809.08782959)
\curveto(1095.88385823,1812.40274105)(1091.98142683,1814.06022535)(1086.77822852,1814.06028748)
\curveto(1080.52590239,1814.06022535)(1075.59541111,1812.06704802)(1071.98673987,1808.08074951)
\curveto(1068.3780111,1804.09433871)(1066.5736611,1798.66030789)(1066.57368445,1791.77864075)
\lineto(1066.57368445,1751.93603516)
\lineto(1054.92932104,1751.93603516)
\lineto(1054.92932104,1822.43164062)
\lineto(1066.57368445,1822.43164062)
\lineto(1066.57368445,1811.47964478)
\curveto(1069.34312854,1815.71770966)(1072.59515471,1818.88581257)(1076.32977271,1820.98396301)
\curveto(1080.10628611,1823.08197536)(1084.44931461,1824.13101606)(1089.35887122,1824.13108826)
\curveto(1097.45741928,1824.13101606)(1103.58381696,1821.61331838)(1107.73808264,1816.57798767)
\curveto(1111.89221929,1811.5844893)(1113.96931988,1804.22022359)(1113.96939062,1794.48516846)
}
}
{
\newrgbcolor{curcolor}{0 0 0}
\pscustom[linestyle=none,fillstyle=solid,fillcolor=curcolor]
{
\newpath
\moveto(1187.92683386,1819.72511292)
\lineto(1187.92683386,1808.89900208)
\curveto(1184.653764,1810.70329511)(1181.35977621,1812.04606721)(1178.0448606,1812.92732239)
\curveto(1174.77180062,1813.85041721)(1171.45683201,1814.31199512)(1168.09994482,1814.3120575)
\curveto(1160.58877037,1814.31199512)(1154.75610408,1811.92018233)(1150.60192847,1807.13661194)
\curveto(1146.44770174,1802.39489278)(1144.37060116,1795.72299393)(1144.37062048,1787.12089539)
\curveto(1144.37060116,1778.51872647)(1146.44770174,1771.82584681)(1150.60192847,1767.04223633)
\curveto(1154.75610408,1762.30055726)(1160.58877037,1759.92972528)(1168.09994482,1759.92973328)
\curveto(1171.45683201,1759.92972528)(1174.77180062,1760.37032238)(1178.0448606,1761.25152588)
\curveto(1181.35977621,1762.17467238)(1184.653764,1763.53842529)(1187.92683386,1765.3427887)
\lineto(1187.92683386,1754.64256287)
\curveto(1184.69572563,1753.13194155)(1181.33879539,1751.9989776)(1177.85603308,1751.2436676)
\curveto(1174.41512678,1750.48835899)(1170.74348433,1750.11070434)(1166.84109473,1750.11070251)
\curveto(1156.22476106,1750.11070434)(1147.79047384,1753.44665376)(1141.53820776,1760.11856079)
\curveto(1135.28590871,1766.79045145)(1132.15976742,1775.79122065)(1132.15977454,1787.12089539)
\curveto(1132.15976742,1798.61834626)(1135.30688952,1807.66107709)(1141.60115027,1814.24911499)
\curveto(1147.93733954,1820.83702827)(1156.60241571,1824.13101606)(1167.59640479,1824.13108826)
\curveto(1171.16310061,1824.13101606)(1174.64591573,1823.75336141)(1178.0448606,1822.99812317)
\curveto(1181.44369946,1822.28470443)(1184.73768726,1821.1937021)(1187.92683386,1819.72511292)
}
}
{
\newrgbcolor{curcolor}{0 0 0}
\pscustom[linestyle=none,fillstyle=solid,fillcolor=curcolor]
{
\newpath
\moveto(1268.24147009,1790.07919312)
\lineto(1268.24147009,1784.41436768)
\lineto(1214.99211096,1784.41436768)
\curveto(1215.4956313,1776.44162589)(1217.88744409,1770.35718983)(1222.16755652,1766.16104126)
\curveto(1226.48957783,1762.00682587)(1232.49009062,1759.92972528)(1240.16911292,1759.92973328)
\curveto(1244.6170011,1759.92972528)(1248.91806797,1760.47522645)(1253.07232642,1761.5662384)
\curveto(1257.26843193,1762.6572311)(1261.4226331,1764.29373459)(1265.53494238,1766.47575378)
\lineto(1265.53494238,1755.52375793)
\curveto(1261.38067147,1753.76136597)(1257.12156624,1752.41859388)(1252.75761389,1751.49543762)
\curveto(1248.39354762,1750.57228225)(1243.96659587,1750.11070434)(1239.47674536,1750.11070251)
\curveto(1228.23098539,1750.11070434)(1219.31413944,1753.38371132)(1212.72618079,1759.92973328)
\curveto(1206.18014989,1766.47573924)(1202.90714291,1775.32964274)(1202.90715002,1786.49147034)
\curveto(1202.90714291,1798.03088347)(1206.01230338,1807.17851837)(1212.22264075,1813.93440247)
\curveto(1218.47490689,1820.7321242)(1226.88821329,1824.13101606)(1237.46258521,1824.13108826)
\curveto(1246.94587146,1824.13101606)(1254.43602205,1821.06781722)(1259.93305945,1814.94148254)
\curveto(1265.4719302,1808.85698348)(1268.24139765,1800.56956196)(1268.24147009,1790.07919312)
\moveto(1256.66004919,1793.47808838)
\curveto(1256.57606507,1799.81425266)(1254.79269588,1804.87062883)(1251.30993628,1808.64723206)
\curveto(1247.86902727,1812.42372186)(1243.29520982,1814.31199512)(1237.58847021,1814.3120575)
\curveto(1231.12633772,1814.31199512)(1225.94407666,1812.4866643)(1222.04167151,1808.83605957)
\curveto(1218.18117549,1805.18534104)(1215.95720921,1800.04504161)(1215.36976599,1793.41514587)
\lineto(1256.66004919,1793.47808838)
}
}
{
\newrgbcolor{curcolor}{0 0 0}
\pscustom[linestyle=none,fillstyle=solid,fillcolor=curcolor]
{
\newpath
\moveto(1314.81892371,1849.74868774)
\curveto(1309.19602559,1840.0974155)(1305.02084361,1830.55114514)(1302.29336523,1821.10984802)
\curveto(1299.56583197,1811.66841256)(1298.20207907,1802.10116138)(1298.20210242,1792.4080658)
\curveto(1298.20207907,1782.71488927)(1299.56583197,1773.10567646)(1302.29336523,1763.58039856)
\curveto(1305.06280524,1754.097059)(1309.23798722,1744.55078864)(1314.81892371,1734.94155884)
\lineto(1304.74812292,1734.94155884)
\curveto(1298.45384883,1744.8025584)(1293.73316569,1754.49569446)(1290.58605933,1764.02099609)
\curveto(1287.48088312,1773.54627356)(1285.92830289,1783.00862066)(1285.92831396,1792.4080658)
\curveto(1285.92830289,1801.76546836)(1287.48088312,1811.18585384)(1290.58605933,1820.66925049)
\curveto(1293.69120406,1830.15250967)(1298.41188721,1839.84564573)(1304.74812292,1849.74868774)
\lineto(1314.81892371,1849.74868774)
}
}
{
\newrgbcolor{curcolor}{0 0 0}
\pscustom[linestyle=none,fillstyle=solid,fillcolor=curcolor]
{
\newpath
\moveto(1335.52700781,1849.74868774)
\lineto(1345.59780859,1849.74868774)
\curveto(1351.89203239,1839.84564573)(1356.59173473,1830.15250967)(1359.69692969,1820.66925049)
\curveto(1362.84401729,1811.18585384)(1364.41757834,1801.76546836)(1364.41761755,1792.4080658)
\curveto(1364.41757834,1783.00862066)(1362.84401729,1773.54627356)(1359.69692969,1764.02099609)
\curveto(1356.59173473,1754.49569446)(1351.89203239,1744.8025584)(1345.59780859,1734.94155884)
\lineto(1335.52700781,1734.94155884)
\curveto(1341.10789401,1744.55078864)(1345.26209518,1754.097059)(1347.98962378,1763.58039856)
\curveto(1350.75906844,1773.10567646)(1352.14380216,1782.71488927)(1352.1438291,1792.4080658)
\curveto(1352.14380216,1802.10116138)(1350.75906844,1811.66841256)(1347.98962378,1821.10984802)
\curveto(1345.26209518,1830.55114514)(1341.10789401,1840.0974155)(1335.52700781,1849.74868774)
}
}
{
\newrgbcolor{curcolor}{0 0 0}
\pscustom[linestyle=none,fillstyle=solid,fillcolor=curcolor]
{
}
}
{
\newrgbcolor{curcolor}{0 0 0}
\pscustom[linestyle=none,fillstyle=solid,fillcolor=curcolor]
{
\newpath
\moveto(1431.70315527,1767.9234314)
\lineto(1444.9840238,1767.9234314)
\lineto(1444.9840238,1751.93603516)
\lineto(1431.70315527,1751.93603516)
\lineto(1431.70315527,1767.9234314)
\moveto(1431.70315527,1818.59214783)
\lineto(1444.9840238,1818.59214783)
\lineto(1444.9840238,1802.60475159)
\lineto(1431.70315527,1802.60475159)
\lineto(1431.70315527,1818.59214783)
}
}
{
\newrgbcolor{curcolor}{0 0 0}
\pscustom[linestyle=none,fillstyle=solid,fillcolor=curcolor]
{
}
}
{
\newrgbcolor{curcolor}{0 0 0}
\pscustom[linestyle=none,fillstyle=solid,fillcolor=curcolor]
{
\newpath
\moveto(1513.84312415,1845.90919495)
\lineto(1573.26084875,1845.90919495)
\lineto(1573.26084875,1835.20896912)
\lineto(1526.55751013,1835.20896912)
\lineto(1526.55751013,1807.38838196)
\lineto(1571.3096311,1807.38838196)
\lineto(1571.3096311,1796.68815613)
\lineto(1526.55751013,1796.68815613)
\lineto(1526.55751013,1762.63626099)
\lineto(1574.39381384,1762.63626099)
\lineto(1574.39381384,1751.93603516)
\lineto(1513.84312415,1751.93603516)
\lineto(1513.84312415,1845.90919495)
}
}
{
\newrgbcolor{curcolor}{0 0 0}
\pscustom[linestyle=none,fillstyle=solid,fillcolor=curcolor]
{
\newpath
\moveto(1586.35288977,1822.43164062)
\lineto(1598.62667822,1822.43164062)
\lineto(1620.65655493,1763.26568604)
\lineto(1642.68643164,1822.43164062)
\lineto(1654.96022009,1822.43164062)
\lineto(1628.52436804,1751.93603516)
\lineto(1612.78874182,1751.93603516)
\lineto(1586.35288977,1822.43164062)
}
}
{
\newrgbcolor{curcolor}{0 0 0}
\pscustom[linestyle=none,fillstyle=solid,fillcolor=curcolor]
{
\newpath
\moveto(1670.94761633,1822.43164062)
\lineto(1682.52903723,1822.43164062)
\lineto(1682.52903723,1751.93603516)
\lineto(1670.94761633,1751.93603516)
\lineto(1670.94761633,1822.43164062)
\moveto(1670.94761633,1849.87457275)
\lineto(1682.52903723,1849.87457275)
\lineto(1682.52903723,1835.20896912)
\lineto(1670.94761633,1835.20896912)
\lineto(1670.94761633,1849.87457275)
}
}
{
\newrgbcolor{curcolor}{0 0 0}
\pscustom[linestyle=none,fillstyle=solid,fillcolor=curcolor]
{
\newpath
\moveto(1753.08758521,1811.73141479)
\lineto(1753.08758521,1849.87457275)
\lineto(1764.6690061,1849.87457275)
\lineto(1764.6690061,1751.93603516)
\lineto(1753.08758521,1751.93603516)
\lineto(1753.08758521,1762.51037598)
\curveto(1750.65375225,1758.31420261)(1747.56957259,1755.18806132)(1743.83503699,1753.13194275)
\curveto(1740.14236444,1751.11778341)(1735.69443188,1750.11070434)(1730.49122595,1750.11070251)
\curveto(1721.97297954,1750.11070434)(1715.02833011,1753.5095962)(1709.65725684,1760.30738831)
\curveto(1704.32811498,1767.10516366)(1701.6635516,1776.04299042)(1701.66355872,1787.12089539)
\curveto(1701.6635516,1798.19872998)(1704.32811498,1807.13655674)(1709.65725684,1813.93440247)
\curveto(1715.02833011,1820.7321242)(1721.97297954,1824.13101606)(1730.49122595,1824.13108826)
\curveto(1735.69443188,1824.13101606)(1740.14236444,1823.10295618)(1743.83503699,1821.04690552)
\curveto(1747.56957259,1819.03267826)(1750.65375225,1815.9275178)(1753.08758521,1811.73141479)
\moveto(1713.62263464,1787.12089539)
\curveto(1713.62261557,1778.60264972)(1715.36402313,1771.90977007)(1718.84686255,1767.04223633)
\curveto(1722.371615,1762.21663401)(1727.19720222,1759.8038404)(1733.32363867,1759.80384827)
\curveto(1739.44999758,1759.8038404)(1744.2755848,1762.21663401)(1747.80041479,1767.04223633)
\curveto(1751.32513829,1771.90977007)(1753.08752667,1778.60264972)(1753.08758521,1787.12089539)
\curveto(1753.08752667,1795.63907068)(1751.32513829,1802.31096952)(1747.80041479,1807.13661194)
\curveto(1744.2755848,1812.00410558)(1739.44999758,1814.43788)(1733.32363867,1814.4379425)
\curveto(1727.19720222,1814.43788)(1722.371615,1812.00410558)(1718.84686255,1807.13661194)
\curveto(1715.36402313,1802.31096952)(1713.62261557,1795.63907068)(1713.62263464,1787.12089539)
}
}
{
\newrgbcolor{curcolor}{0 0 0}
\pscustom[linestyle=none,fillstyle=solid,fillcolor=curcolor]
{
\newpath
\moveto(1848.82313513,1790.07919312)
\lineto(1848.82313513,1784.41436768)
\lineto(1795.573776,1784.41436768)
\curveto(1796.07729634,1776.44162589)(1798.46910913,1770.35718983)(1802.74922156,1766.16104126)
\curveto(1807.07124286,1762.00682587)(1813.07175566,1759.92972528)(1820.75077795,1759.92973328)
\curveto(1825.19866614,1759.92972528)(1829.49973301,1760.47522645)(1833.65399146,1761.5662384)
\curveto(1837.85009697,1762.6572311)(1842.00429814,1764.29373459)(1846.11660742,1766.47575378)
\lineto(1846.11660742,1755.52375793)
\curveto(1841.96233651,1753.76136597)(1837.70323128,1752.41859388)(1833.33927893,1751.49543762)
\curveto(1828.97521266,1750.57228225)(1824.54826091,1750.11070434)(1820.0584104,1750.11070251)
\curveto(1808.81265042,1750.11070434)(1799.89580448,1753.38371132)(1793.30784583,1759.92973328)
\curveto(1786.76181493,1766.47573924)(1783.48880795,1775.32964274)(1783.48881506,1786.49147034)
\curveto(1783.48880795,1798.03088347)(1786.59396842,1807.17851837)(1792.80430579,1813.93440247)
\curveto(1799.05657192,1820.7321242)(1807.46987833,1824.13101606)(1818.04425024,1824.13108826)
\curveto(1827.5275365,1824.13101606)(1835.01768709,1821.06781722)(1840.51472449,1814.94148254)
\curveto(1846.05359524,1808.85698348)(1848.82306269,1800.56956196)(1848.82313513,1790.07919312)
\moveto(1837.24171423,1793.47808838)
\curveto(1837.15773011,1799.81425266)(1835.37436092,1804.87062883)(1831.89160132,1808.64723206)
\curveto(1828.45069231,1812.42372186)(1823.87687486,1814.31199512)(1818.17013525,1814.3120575)
\curveto(1811.70800275,1814.31199512)(1806.5257417,1812.4866643)(1802.62333655,1808.83605957)
\curveto(1798.76284053,1805.18534104)(1796.53887425,1800.04504161)(1795.95143103,1793.41514587)
\lineto(1837.24171423,1793.47808838)
}
}
{
\newrgbcolor{curcolor}{0 0 0}
\pscustom[linestyle=none,fillstyle=solid,fillcolor=curcolor]
{
\newpath
\moveto(1926.17947363,1794.48516846)
\lineto(1926.17947363,1751.93603516)
\lineto(1914.59805273,1751.93603516)
\lineto(1914.59805273,1794.10751343)
\curveto(1914.59799357,1800.7793701)(1913.2971831,1805.77280383)(1910.69561743,1809.08782959)
\curveto(1908.09394123,1812.40274105)(1904.19150983,1814.06022535)(1898.98831152,1814.06028748)
\curveto(1892.7359854,1814.06022535)(1887.80549412,1812.06704802)(1884.19682288,1808.08074951)
\curveto(1880.58809411,1804.09433871)(1878.7837441,1798.66030789)(1878.78376746,1791.77864075)
\lineto(1878.78376746,1751.93603516)
\lineto(1867.13940405,1751.93603516)
\lineto(1867.13940405,1822.43164062)
\lineto(1878.78376746,1822.43164062)
\lineto(1878.78376746,1811.47964478)
\curveto(1881.55321155,1815.71770966)(1884.80523772,1818.88581257)(1888.53985571,1820.98396301)
\curveto(1892.31636912,1823.08197536)(1896.65939762,1824.13101606)(1901.56895422,1824.13108826)
\curveto(1909.66750228,1824.13101606)(1915.79389997,1821.61331838)(1919.94816565,1816.57798767)
\curveto(1924.1023023,1811.5844893)(1926.17940289,1804.22022359)(1926.17947363,1794.48516846)
}
}
{
\newrgbcolor{curcolor}{0 0 0}
\pscustom[linestyle=none,fillstyle=solid,fillcolor=curcolor]
{
\newpath
\moveto(2000.13691687,1819.72511292)
\lineto(2000.13691687,1808.89900208)
\curveto(1996.86384701,1810.70329511)(1993.56985921,1812.04606721)(1990.2549436,1812.92732239)
\curveto(1986.98188363,1813.85041721)(1983.66691502,1814.31199512)(1980.31002783,1814.3120575)
\curveto(1972.79885337,1814.31199512)(1966.96618709,1811.92018233)(1962.81201147,1807.13661194)
\curveto(1958.65778475,1802.39489278)(1956.58068417,1795.72299393)(1956.58070349,1787.12089539)
\curveto(1956.58068417,1778.51872647)(1958.65778475,1771.82584681)(1962.81201147,1767.04223633)
\curveto(1966.96618709,1762.30055726)(1972.79885337,1759.92972528)(1980.31002783,1759.92973328)
\curveto(1983.66691502,1759.92972528)(1986.98188363,1760.37032238)(1990.2549436,1761.25152588)
\curveto(1993.56985921,1762.17467238)(1996.86384701,1763.53842529)(2000.13691687,1765.3427887)
\lineto(2000.13691687,1754.64256287)
\curveto(1996.90580864,1753.13194155)(1993.5488784,1751.9989776)(1990.06611609,1751.2436676)
\curveto(1986.62520979,1750.48835899)(1982.95356734,1750.11070434)(1979.05117773,1750.11070251)
\curveto(1968.43484407,1750.11070434)(1960.00055685,1753.44665376)(1953.74829077,1760.11856079)
\curveto(1947.49599171,1766.79045145)(1944.36985043,1775.79122065)(1944.36985754,1787.12089539)
\curveto(1944.36985043,1798.61834626)(1947.51697253,1807.66107709)(1953.81123328,1814.24911499)
\curveto(1960.14742254,1820.83702827)(1968.81249872,1824.13101606)(1979.80648779,1824.13108826)
\curveto(1983.37318362,1824.13101606)(1986.85599874,1823.75336141)(1990.2549436,1822.99812317)
\curveto(1993.65378247,1822.28470443)(1996.94777027,1821.1937021)(2000.13691687,1819.72511292)
}
}
{
\newrgbcolor{curcolor}{0 0 0}
\pscustom[linestyle=none,fillstyle=solid,fillcolor=curcolor]
{
\newpath
\moveto(2080.4515531,1790.07919312)
\lineto(2080.4515531,1784.41436768)
\lineto(2027.20219397,1784.41436768)
\curveto(2027.70571431,1776.44162589)(2030.0975271,1770.35718983)(2034.37763953,1766.16104126)
\curveto(2038.69966083,1762.00682587)(2044.70017363,1759.92972528)(2052.37919592,1759.92973328)
\curveto(2056.82708411,1759.92972528)(2061.12815098,1760.47522645)(2065.28240942,1761.5662384)
\curveto(2069.47851494,1762.6572311)(2073.63271611,1764.29373459)(2077.74502539,1766.47575378)
\lineto(2077.74502539,1755.52375793)
\curveto(2073.59075448,1753.76136597)(2069.33164924,1752.41859388)(2064.9676969,1751.49543762)
\curveto(2060.60363063,1750.57228225)(2056.17667888,1750.11070434)(2051.68682837,1750.11070251)
\curveto(2040.44106839,1750.11070434)(2031.52422245,1753.38371132)(2024.93626379,1759.92973328)
\curveto(2018.3902329,1766.47573924)(2015.11722592,1775.32964274)(2015.11723303,1786.49147034)
\curveto(2015.11722592,1798.03088347)(2018.22238639,1807.17851837)(2024.43272375,1813.93440247)
\curveto(2030.68498989,1820.7321242)(2039.0982963,1824.13101606)(2049.67266821,1824.13108826)
\curveto(2059.15595446,1824.13101606)(2066.64610505,1821.06781722)(2072.14314246,1814.94148254)
\curveto(2077.68201321,1808.85698348)(2080.45148065,1800.56956196)(2080.4515531,1790.07919312)
\moveto(2068.8701322,1793.47808838)
\curveto(2068.78614808,1799.81425266)(2067.00277889,1804.87062883)(2063.52001929,1808.64723206)
\curveto(2060.07911028,1812.42372186)(2055.50529283,1814.31199512)(2049.79855322,1814.3120575)
\curveto(2043.33642072,1814.31199512)(2038.15415967,1812.4866643)(2034.25175452,1808.83605957)
\curveto(2030.3912585,1805.18534104)(2028.16729222,1800.04504161)(2027.579849,1793.41514587)
\lineto(2068.8701322,1793.47808838)
}
}
{
\newrgbcolor{curcolor}{0.99607843 0.99607843 0.80784315}
\pscustom[linestyle=none,fillstyle=solid,fillcolor=curcolor]
{
\newpath
\moveto(1092.72915625,5833.0078125)
\lineto(2944.29165625,5833.0078125)
\lineto(2944.29165625,4614.29779053)
\lineto(1092.72915625,4614.29779053)
\lineto(1092.72915625,5833.0078125)
\closepath
}
}
{
\newrgbcolor{curcolor}{0.65882355 0 0.21176471}
\pscustom[linewidth=17.578125,linecolor=curcolor]
{
\newpath
\moveto(1092.72915625,5833.0078125)
\lineto(2944.29165625,5833.0078125)
\lineto(2944.29165625,4614.29779053)
\lineto(1092.72915625,4614.29779053)
\lineto(1092.72915625,5833.0078125)
\closepath
}
}
{
\newrgbcolor{curcolor}{0.70588237 0.65490198 0.89803922}
\pscustom[linestyle=none,fillstyle=solid,fillcolor=curcolor]
{
\newpath
\moveto(1933.54940625,5645.507875)
\curveto(1933.54940625,5574.31491897)(1875.83611228,5516.601625)(1804.64315625,5516.601625)
\curveto(1733.45020022,5516.601625)(1675.73690625,5574.31491897)(1675.73690625,5645.507875)
\curveto(1675.73690625,5716.70083103)(1733.45020022,5774.414125)(1804.64315625,5774.414125)
\curveto(1875.83611228,5774.414125)(1933.54940625,5716.70083103)(1933.54940625,5645.507875)
\closepath
}
}
{
\newrgbcolor{curcolor}{0.65882355 0 0.21176471}
\pscustom[linewidth=11.71875,linecolor=curcolor]
{
\newpath
\moveto(1933.54940625,5645.507875)
\curveto(1933.54940625,5574.31491897)(1875.83611228,5516.601625)(1804.64315625,5516.601625)
\curveto(1733.45020022,5516.601625)(1675.73690625,5574.31491897)(1675.73690625,5645.507875)
\curveto(1675.73690625,5716.70083103)(1733.45020022,5774.414125)(1804.64315625,5774.414125)
\curveto(1875.83611228,5774.414125)(1933.54940625,5716.70083103)(1933.54940625,5645.507875)
\closepath
}
}
{
\newrgbcolor{curcolor}{0 0 0}
\pscustom[linestyle=none,fillstyle=solid,fillcolor=curcolor]
{
\newpath
\moveto(1757.95134375,5575.1953125)
\lineto(1757.95134375,5594.4214375)
\lineto(1791.64284375,5594.4214375)
\lineto(1791.64284375,5700.0731875)
\lineto(1757.95134375,5700.0731875)
\lineto(1757.95134375,5719.1161875)
\lineto(1854.08171875,5719.1161875)
\lineto(1854.08171875,5700.0731875)
\lineto(1820.39034375,5700.0731875)
\lineto(1820.39034375,5594.4214375)
\lineto(1854.08171875,5594.4214375)
\lineto(1854.08171875,5575.1953125)
\lineto(1757.95134375,5575.1953125)
\closepath
}
}
{
\newrgbcolor{curcolor}{0 0 0}
\pscustom[linestyle=none,fillstyle=solid,fillcolor=curcolor]
{
\newpath
\moveto(2080.03384375,5690.00244141)
\lineto(2073.44204687,5656.08215332)
\lineto(2119.8592832,5656.08215332)
\lineto(2117.73068213,5645.2331543)
\lineto(2071.3134458,5645.2331543)
\lineto(2064.10366797,5608.08563232)
\lineto(2050.98873877,5608.08563232)
\lineto(2068.97885107,5700.71411133)
\lineto(2130.15896582,5700.71411133)
\lineto(2128.0990293,5690.00244141)
\lineto(2080.03384375,5690.00244141)
}
}
{
\newrgbcolor{curcolor}{0 0 0}
\pscustom[linestyle=none,fillstyle=solid,fillcolor=curcolor]
{
\newpath
\moveto(2199.0981748,5615.70739746)
\curveto(2199.78474973,5615.70738984)(2200.47139455,5615.75316616)(2201.15811133,5615.84472656)
\curveto(2201.89046051,5615.93627145)(2202.57710533,5616.05071225)(2203.21804785,5616.18804932)
\lineto(2201.70742773,5608.49761963)
\curveto(2200.10518398,5608.13140865)(2198.50301273,5607.85675072)(2196.90090918,5607.67364502)
\curveto(2195.34444655,5607.49054015)(2193.69649898,5607.3989875)(2191.95706152,5607.39898682)
\curveto(2187.9744588,5607.3989875)(2185.13632687,5608.26873761)(2183.44265723,5610.00823975)
\curveto(2181.79465541,5611.74773804)(2180.97068163,5614.35698836)(2180.9707334,5617.83599854)
\curveto(2180.97068163,5618.52263361)(2180.99356979,5619.23216659)(2181.03939795,5619.96459961)
\curveto(2181.13089875,5620.74278519)(2181.22245139,5621.52098266)(2181.31405615,5622.29919434)
\lineto(2180.90206885,5622.29919434)
\curveto(2179.16251693,5619.91881141)(2177.33146407,5617.7673243)(2175.40890479,5615.84472656)
\curveto(2173.48625307,5613.9221133)(2171.38054229,5612.27416573)(2169.09176611,5610.90087891)
\curveto(2166.80291015,5609.57336277)(2164.28521247,5608.54339554)(2161.53866553,5607.81097412)
\curveto(2158.7920539,5607.07855325)(2155.72504037,5606.71234268)(2152.33761572,5606.71234131)
\curveto(2146.24934184,5606.71234268)(2141.69459785,5608.13140865)(2138.67337012,5610.96954346)
\curveto(2135.69789975,5613.8076725)(2134.2101693,5617.72154798)(2134.21017432,5622.71118164)
\curveto(2134.2101693,5624.58799619)(2134.41616275,5626.6250425)(2134.82815527,5628.82232666)
\curveto(2135.78945239,5633.72037231)(2137.46028812,5637.72580044)(2139.84066748,5640.83862305)
\curveto(2142.26680187,5643.99715647)(2145.1049338,5646.46907783)(2148.35507178,5648.25439453)
\curveto(2151.65094776,5650.0396309)(2155.22150083,5651.27559158)(2159.0667417,5651.96228027)
\curveto(2162.95769915,5652.69465754)(2166.84868647,5653.08375627)(2170.73971533,5653.12957764)
\lineto(2186.94454932,5653.40423584)
\lineto(2187.76852393,5657.45544434)
\curveto(2187.99734696,5658.599803)(2188.15756409,5659.67554655)(2188.24917578,5660.68267822)
\curveto(2188.38644569,5661.6897047)(2188.45511018,5662.65100745)(2188.45516943,5663.56658936)
\curveto(2188.45511018,5667.45752119)(2187.40225478,5670.22698864)(2185.2966001,5671.875)
\curveto(2183.19083321,5673.52288378)(2180.03226704,5674.34685757)(2175.82089209,5674.34692383)
\curveto(2173.5320294,5674.34685757)(2171.42631861,5674.16375228)(2169.50375342,5673.79760742)
\curveto(2167.58110761,5673.47710746)(2165.86449556,5672.85912712)(2164.35391211,5671.94366455)
\curveto(2162.84325835,5671.07385059)(2161.51574503,5669.86077807)(2160.37136816,5668.30444336)
\curveto(2159.27270528,5666.79376453)(2158.33429069,5664.87115903)(2157.55612158,5662.53662109)
\lineto(2144.85317969,5663.70391846)
\curveto(2145.90601942,5666.63354741)(2147.27930907,5669.31146221)(2148.97305273,5671.7376709)
\curveto(2150.71253317,5674.2095286)(2152.88690844,5676.33812755)(2155.49618506,5678.12347412)
\curveto(2158.1511854,5679.90868062)(2161.30975158,5681.28197026)(2164.97189307,5682.24334717)
\curveto(2168.63396301,5683.25035208)(2172.91404906,5683.75389162)(2177.81216406,5683.75396729)
\curveto(2185.63986641,5683.75389162)(2191.47634739,5682.19749669)(2195.32162451,5679.08477783)
\curveto(2199.21254571,5675.97191698)(2201.15803937,5671.417173)(2201.15811133,5665.42053223)
\curveto(2201.15803937,5664.3676195)(2201.08937489,5663.29187595)(2200.95211768,5662.19329834)
\curveto(2200.86049328,5661.09461252)(2200.70027615,5659.95020448)(2200.47146582,5658.7600708)
\lineto(2194.29165625,5626.76239014)
\curveto(2193.87960426,5624.79398964)(2193.67361082,5623.05448943)(2193.67367529,5621.54388428)
\curveto(2193.67361082,5619.62126532)(2194.08559771,5618.15642303)(2194.90963721,5617.14935303)
\curveto(2195.73354528,5616.18804121)(2197.12972308,5615.70738984)(2199.0981748,5615.70739746)
\moveto(2157.35012793,5616.11938477)
\curveto(2161.14953445,5616.11937673)(2164.58275856,5616.82890971)(2167.64981055,5618.24798584)
\curveto(2170.76256195,5619.71281796)(2173.44047675,5621.54387082)(2175.68356299,5623.7411499)
\curveto(2177.97233257,5625.93839767)(2179.82627359,5628.36454271)(2181.2453916,5631.01959229)
\curveto(2182.66440552,5633.67459599)(2183.60282011,5636.21518183)(2184.06063818,5638.64135742)
\lineto(2185.2966001,5644.75250244)
\lineto(2172.1816709,5644.47784424)
\curveto(2169.66393024,5644.43203152)(2167.10045624,5644.22603808)(2164.49124121,5643.85986328)
\curveto(2161.8819556,5643.53939326)(2159.4329224,5642.82986028)(2157.14413428,5641.73126221)
\curveto(2154.90106658,5640.63259685)(2152.93268476,5639.0304256)(2151.23898291,5636.92474365)
\curveto(2149.5910133,5634.86478035)(2148.4237171,5632.0953129)(2147.73709082,5628.61633301)
\curveto(2147.50819067,5627.33457548)(2147.39374987,5626.23594376)(2147.39376807,5625.32043457)
\curveto(2147.39374987,5622.48228541)(2148.21772366,5620.23924566)(2149.86569189,5618.59130859)
\curveto(2151.55939512,5616.94335052)(2154.05420464,5616.11937673)(2157.35012793,5616.11938477)
}
}
{
\newrgbcolor{curcolor}{0 0 0}
\pscustom[linestyle=none,fillstyle=solid,fillcolor=curcolor]
{
\newpath
\moveto(2222.44412207,5645.30181885)
\curveto(2223.90895543,5652.76332202)(2226.10621886,5658.96601357)(2229.03591895,5663.90991211)
\curveto(2231.965588,5668.853699)(2235.30725946,5672.79046264)(2239.06094336,5675.72021484)
\curveto(2242.86035249,5678.64983178)(2246.9344451,5680.70976624)(2251.2832334,5681.90002441)
\curveto(2255.67772249,5683.13591128)(2260.04936118,5683.75389162)(2264.3981626,5683.75396729)
\curveto(2268.88419122,5683.75389162)(2272.77517854,5683.1816876)(2276.07113623,5682.03735352)
\curveto(2279.41274514,5680.89287153)(2282.15932442,5679.31358844)(2284.31088232,5677.29949951)
\curveto(2286.46229864,5675.33104848)(2288.06446989,5672.99645609)(2289.11740088,5670.29571533)
\curveto(2290.17018067,5667.59485016)(2290.69660837,5664.68805375)(2290.69668555,5661.57531738)
\lineto(2277.3757627,5660.61401367)
\curveto(2277.46725148,5664.73383007)(2276.27706712,5668.00683705)(2273.80520605,5670.43304443)
\curveto(2271.33322441,5672.85912712)(2267.39646077,5674.07219964)(2261.99490332,5674.07226562)
\curveto(2258.0580912,5674.07219964)(2254.6248671,5673.43133114)(2251.6952207,5672.1496582)
\curveto(2248.81127428,5670.91363346)(2246.31646476,5669.0825806)(2244.21078467,5666.65649414)
\curveto(2242.10504319,5664.27606686)(2240.34265482,5661.30060596)(2238.92361426,5657.73010254)
\curveto(2237.55029921,5654.20527615)(2236.40589117,5650.1540717)(2235.49038672,5645.57647705)
\curveto(2234.62061464,5641.22768903)(2234.18573959,5637.24514906)(2234.18576025,5633.62884521)
\curveto(2234.18573959,5630.88224039)(2234.48328568,5628.41031903)(2235.07839941,5626.21307373)
\curveto(2235.67347003,5624.0615685)(2236.63477278,5622.2534038)(2237.96231055,5620.78857422)
\curveto(2239.28979942,5619.32371923)(2240.98352332,5618.20219936)(2243.0434873,5617.42401123)
\curveto(2245.14916856,5616.64580443)(2247.6897544,5616.2567057)(2250.66525244,5616.25671387)
\curveto(2255.60905801,5616.2567057)(2259.93492038,5617.46977821)(2263.64285254,5619.89593506)
\curveto(2267.39646077,5622.32206828)(2270.14304005,5626.02995032)(2271.88259863,5631.01959229)
\lineto(2284.79153418,5630.19561768)
\curveto(2283.78438383,5627.12858203)(2282.22798891,5624.1760093)(2280.12234473,5621.33789062)
\curveto(2278.01656734,5618.49974545)(2275.45309334,5616.00493593)(2272.43191504,5613.85345459)
\curveto(2269.41061891,5611.70196172)(2265.95450665,5609.9624615)(2262.06356787,5608.63494873)
\curveto(2258.17253201,5607.35321118)(2253.93822227,5606.71234268)(2249.36062598,5606.71234131)
\curveto(2244.27941846,5606.71234268)(2239.95355608,5607.35321118)(2236.38302588,5608.63494873)
\curveto(2232.81244994,5609.9624615)(2229.90565353,5611.79351436)(2227.66262793,5614.12811279)
\curveto(2225.41957404,5616.50847546)(2223.77162646,5619.32371923)(2222.71878027,5622.57385254)
\curveto(2221.711692,5625.86973319)(2221.20815247,5629.48606258)(2221.20816016,5633.42285156)
\curveto(2221.20815247,5637.17648458)(2221.62013936,5641.13613638)(2222.44412207,5645.30181885)
}
}
{
\newrgbcolor{curcolor}{0 0 0}
\pscustom[linestyle=none,fillstyle=solid,fillcolor=curcolor]
{
\newpath
\moveto(2316.23989844,5672.63031006)
\lineto(2318.09384131,5682.38067627)
\lineto(2329.76681494,5682.38067627)
\lineto(2337.52590918,5701.74407959)
\lineto(2345.76565527,5701.74407959)
\lineto(2341.98910498,5682.38067627)
\lineto(2371.65219092,5682.38067627)
\lineto(2369.79824805,5672.63031006)
\lineto(2340.13516211,5672.63031006)
\lineto(2331.41476416,5627.86102295)
\curveto(2331.32317794,5627.31168732)(2331.2316253,5626.78525962)(2331.14010596,5626.28173828)
\curveto(2331.09429633,5625.77818055)(2331.07140817,5625.29752917)(2331.07144141,5624.83978271)
\curveto(2331.07140817,5622.23051564)(2331.94115828,5620.30791014)(2333.68069434,5619.07196045)
\curveto(2335.42015871,5617.83598878)(2338.14384983,5617.21800845)(2341.85177588,5617.21801758)
\curveto(2343.68278472,5617.21800845)(2345.58250206,5617.28667293)(2347.55093359,5617.42401123)
\curveto(2349.5192657,5617.60710718)(2351.41898304,5617.81310062)(2353.25009131,5618.04199219)
\curveto(2355.12686507,5618.27086384)(2356.86636529,5618.52263361)(2358.46859717,5618.79730225)
\curveto(2360.11648411,5619.07194946)(2361.51266191,5619.32371923)(2362.65713477,5619.5526123)
\lineto(2360.87185645,5610.14556885)
\curveto(2359.81893802,5609.87090886)(2358.44564838,5609.55047461)(2356.7519834,5609.18426514)
\curveto(2355.10397691,5608.81805347)(2353.2500359,5608.47473106)(2351.19015479,5608.15429688)
\curveto(2349.13016697,5607.83386256)(2346.91001538,5607.55920463)(2344.52969336,5607.33032227)
\curveto(2342.14927795,5607.10144141)(2339.72313292,5606.98700061)(2337.25125098,5606.98699951)
\curveto(2330.88830289,5606.98700061)(2326.12756546,5608.10852048)(2322.96902441,5610.3515625)
\curveto(2319.81043311,5612.6403763)(2318.23115002,5616.21092937)(2318.23117041,5621.06323242)
\curveto(2318.23115002,5621.88719323)(2318.25403818,5622.75694334)(2318.29983496,5623.67248535)
\curveto(2318.39136714,5624.58799619)(2318.55158427,5625.54929894)(2318.78048682,5626.55639648)
\lineto(2327.77554297,5672.63031006)
\lineto(2316.23989844,5672.63031006)
}
}
{
\newrgbcolor{curcolor}{0 0 0}
\pscustom[linestyle=none,fillstyle=solid,fillcolor=curcolor]
{
\newpath
\moveto(1104.44790625,5458.0078125)
\lineto(2932.57290625,5458.0078125)
}
}
{
\newrgbcolor{curcolor}{0.65882355 0 0.21176471}
\pscustom[linewidth=17.578125,linecolor=curcolor]
{
\newpath
\moveto(1104.44790625,5458.0078125)
\lineto(2932.57290625,5458.0078125)
}
}
{
\newrgbcolor{curcolor}{0.51764709 0.74509805 0.51764709}
\pscustom[linestyle=none,fillstyle=solid,fillcolor=curcolor]
{
\newpath
\moveto(1256.79159375,5318.55871729)
\curveto(1256.79159375,5299.14245655)(1241.05160449,5283.40246729)(1221.63534375,5283.40246729)
\curveto(1202.21908301,5283.40246729)(1186.47909375,5299.14245655)(1186.47909375,5318.55871729)
\curveto(1186.47909375,5337.97497802)(1202.21908301,5353.71496729)(1221.63534375,5353.71496729)
\curveto(1241.05160449,5353.71496729)(1256.79159375,5337.97497802)(1256.79159375,5318.55871729)
\closepath
}
}
{
\newrgbcolor{curcolor}{0.01176471 0.50196081 0.28235295}
\pscustom[linewidth=11.71875,linecolor=curcolor]
{
\newpath
\moveto(1256.79159375,5318.55871729)
\curveto(1256.79159375,5299.14245655)(1241.05160449,5283.40246729)(1221.63534375,5283.40246729)
\curveto(1202.21908301,5283.40246729)(1186.47909375,5299.14245655)(1186.47909375,5318.55871729)
\curveto(1186.47909375,5337.97497802)(1202.21908301,5353.71496729)(1221.63534375,5353.71496729)
\curveto(1241.05160449,5353.71496729)(1256.79159375,5337.97497802)(1256.79159375,5318.55871729)
\closepath
}
}
{
\newrgbcolor{curcolor}{0 0 0}
\pscustom[linestyle=none,fillstyle=solid,fillcolor=curcolor]
{
\newpath
\moveto(1385.64068579,5324.7127533)
\curveto(1385.64062725,5333.10504282)(1383.89921969,5339.60909516)(1380.41645789,5344.22492981)
\curveto(1376.97555108,5348.84065331)(1372.12898305,5351.14854284)(1365.87673926,5351.14860535)
\curveto(1359.66637955,5351.14854284)(1354.81981152,5348.84065331)(1351.33702063,5344.22492981)
\curveto(1347.8961429,5339.60909516)(1346.17571616,5333.10504282)(1346.17573523,5324.7127533)
\curveto(1346.17571616,5316.36235327)(1347.8961429,5309.87928175)(1351.33702063,5305.26351929)
\curveto(1354.81981152,5300.64772359)(1359.66637955,5298.33983406)(1365.87673926,5298.33984375)
\curveto(1372.12898305,5298.33983406)(1376.97555108,5300.64772359)(1380.41645789,5305.26351929)
\curveto(1383.89921969,5309.87928175)(1385.64062725,5316.36235327)(1385.64068579,5324.7127533)
\moveto(1397.22210669,5297.39570618)
\curveto(1397.22203657,5285.39467183)(1394.5574732,5276.47782589)(1389.22840857,5270.6451416)
\curveto(1383.89921969,5264.77053169)(1375.73768306,5261.83321773)(1364.74377417,5261.83319092)
\curveto(1360.67345862,5261.83321773)(1356.83396966,5262.14792994)(1353.22529578,5262.77732849)
\curveto(1349.61656965,5263.36481715)(1346.11277372,5264.28797297)(1342.71389746,5265.54679871)
\lineto(1342.71389746,5276.81350708)
\curveto(1346.11277372,5274.96720728)(1349.46970395,5273.60345437)(1352.78469824,5272.72224426)
\curveto(1356.09964117,5271.841066)(1359.47755222,5271.40046891)(1362.91844153,5271.40045166)
\curveto(1370.51346037,5271.40046891)(1376.19926096,5273.39364623)(1379.97586035,5277.37998962)
\curveto(1383.752354,5281.32439392)(1385.64062725,5287.3039259)(1385.64068579,5295.31860352)
\lineto(1385.64068579,5301.04637146)
\curveto(1383.24881446,5296.89215789)(1380.18561562,5293.78699742)(1376.45108008,5291.73088074)
\curveto(1372.71644584,5289.67475788)(1368.24753246,5288.646698)(1363.04432654,5288.646698)
\curveto(1354.40019524,5288.646698)(1347.434565,5291.94068579)(1342.14741492,5298.52867126)
\curveto(1336.86023475,5305.11663697)(1334.21665219,5313.84465559)(1334.2166593,5324.7127533)
\curveto(1334.21665219,5335.6227405)(1336.86023475,5344.37173993)(1342.14741492,5350.95977783)
\curveto(1347.434565,5357.54769111)(1354.40019524,5360.8416789)(1363.04432654,5360.8417511)
\curveto(1368.24753246,5360.8416789)(1372.71644584,5359.81361902)(1376.45108008,5357.75756836)
\curveto(1380.18561562,5355.70137948)(1383.24881446,5352.59621901)(1385.64068579,5348.44207764)
\lineto(1385.64068579,5359.14230347)
\lineto(1397.22210669,5359.14230347)
\lineto(1397.22210669,5297.39570618)
}
}
{
\newrgbcolor{curcolor}{0 0 0}
\pscustom[linestyle=none,fillstyle=solid,fillcolor=curcolor]
{
\newpath
\moveto(1481.37623572,5326.78985596)
\lineto(1481.37623572,5321.12503052)
\lineto(1428.12687659,5321.12503052)
\curveto(1428.63039692,5313.15228873)(1431.02220972,5307.06785267)(1435.30232214,5302.8717041)
\curveto(1439.62434345,5298.71748871)(1445.62485625,5296.64038812)(1453.30387854,5296.64039612)
\curveto(1457.75176673,5296.64038812)(1462.0528336,5297.18588929)(1466.20709204,5298.27690125)
\curveto(1470.40319756,5299.36789394)(1474.55739873,5301.00439743)(1478.66970801,5303.18641663)
\lineto(1478.66970801,5292.23442078)
\curveto(1474.5154371,5290.47202881)(1470.25633186,5289.12925672)(1465.89237952,5288.20610046)
\curveto(1461.52831325,5287.28294509)(1457.1013615,5286.82136718)(1452.61151099,5286.82136536)
\curveto(1441.36575101,5286.82136718)(1432.44890507,5290.09437416)(1425.86094641,5296.64039612)
\curveto(1419.31491552,5303.18640209)(1416.04190854,5312.04030559)(1416.04191565,5323.20213318)
\curveto(1416.04190854,5334.74154631)(1419.14706901,5343.88918121)(1425.35740637,5350.64506531)
\curveto(1431.60967251,5357.44278704)(1440.02297892,5360.8416789)(1450.59735083,5360.8417511)
\curveto(1460.08063708,5360.8416789)(1467.57078767,5357.77848006)(1473.06782507,5351.65214539)
\curveto(1478.60669583,5345.56764633)(1481.37616327,5337.2802248)(1481.37623572,5326.78985596)
\moveto(1469.79481482,5330.18875122)
\curveto(1469.7108307,5336.5249155)(1467.92746151,5341.58129167)(1464.4447019,5345.3578949)
\curveto(1461.0037929,5349.1343847)(1456.42997545,5351.02265796)(1450.72323584,5351.02272034)
\curveto(1444.26110334,5351.02265796)(1439.07884229,5349.19732714)(1435.17643713,5345.54672241)
\curveto(1431.31594111,5341.89600388)(1429.09197483,5336.75570445)(1428.50453162,5330.12580872)
\lineto(1469.79481482,5330.18875122)
}
}
{
\newrgbcolor{curcolor}{0 0 0}
\pscustom[linestyle=none,fillstyle=solid,fillcolor=curcolor]
{
\newpath
\moveto(1511.58863806,5379.15802002)
\lineto(1511.58863806,5359.14230347)
\lineto(1535.44384741,5359.14230347)
\lineto(1535.44384741,5350.14152527)
\lineto(1511.58863806,5350.14152527)
\lineto(1511.58863806,5311.8724823)
\curveto(1511.58861446,5306.12371604)(1512.36490458,5302.43109278)(1513.91751074,5300.79460144)
\curveto(1515.51202667,5299.1580858)(1518.72209121,5298.33983406)(1523.54771399,5298.33984375)
\lineto(1535.44384741,5298.33984375)
\lineto(1535.44384741,5288.646698)
\lineto(1523.54771399,5288.646698)
\curveto(1514.60985167,5288.646698)(1508.44149236,5290.3041823)(1505.04261755,5293.61915588)
\curveto(1501.64370863,5296.97608115)(1499.9442627,5303.0605172)(1499.94427466,5311.8724823)
\lineto(1499.94427466,5350.14152527)
\lineto(1491.4470365,5350.14152527)
\lineto(1491.4470365,5359.14230347)
\lineto(1499.94427466,5359.14230347)
\lineto(1499.94427466,5379.15802002)
\lineto(1511.58863806,5379.15802002)
}
}
{
\newrgbcolor{curcolor}{0 0 0}
\pscustom[linestyle=none,fillstyle=solid,fillcolor=curcolor]
{
\newpath
\moveto(1551.24241614,5382.61985779)
\lineto(1563.95680212,5382.61985779)
\lineto(1563.95680212,5288.646698)
\lineto(1551.24241614,5288.646698)
\lineto(1551.24241614,5382.61985779)
}
}
{
\newrgbcolor{curcolor}{0 0 0}
\pscustom[linestyle=none,fillstyle=solid,fillcolor=curcolor]
{
\newpath
\moveto(1635.14477515,5348.44207764)
\lineto(1635.14477515,5386.5852356)
\lineto(1646.72619604,5386.5852356)
\lineto(1646.72619604,5288.646698)
\lineto(1635.14477515,5288.646698)
\lineto(1635.14477515,5299.22103882)
\curveto(1632.71094219,5295.02486545)(1629.62676253,5291.89872416)(1625.89222693,5289.84260559)
\curveto(1622.19955438,5287.82844625)(1617.75162182,5286.82136718)(1612.54841589,5286.82136536)
\curveto(1604.03016948,5286.82136718)(1597.08552005,5290.22025905)(1591.71444678,5297.01805115)
\curveto(1586.38530492,5303.81582651)(1583.72074154,5312.75365326)(1583.72074866,5323.83155823)
\curveto(1583.72074154,5334.90939282)(1586.38530492,5343.84721958)(1591.71444678,5350.64506531)
\curveto(1597.08552005,5357.44278704)(1604.03016948,5360.8416789)(1612.54841589,5360.8417511)
\curveto(1617.75162182,5360.8416789)(1622.19955438,5359.81361902)(1625.89222693,5357.75756836)
\curveto(1629.62676253,5355.74334111)(1632.71094219,5352.63818064)(1635.14477515,5348.44207764)
\moveto(1595.67982458,5323.83155823)
\curveto(1595.67980551,5315.31331257)(1597.42121307,5308.62043291)(1600.90405249,5303.75289917)
\curveto(1604.42880494,5298.92729685)(1609.25439216,5296.51450324)(1615.38082861,5296.51451111)
\curveto(1621.50718752,5296.51450324)(1626.33277474,5298.92729685)(1629.85760474,5303.75289917)
\curveto(1633.38232824,5308.62043291)(1635.14471661,5315.31331257)(1635.14477515,5323.83155823)
\curveto(1635.14471661,5332.34973352)(1633.38232824,5339.02163236)(1629.85760474,5343.84727478)
\curveto(1626.33277474,5348.71476842)(1621.50718752,5351.14854284)(1615.38082861,5351.14860535)
\curveto(1609.25439216,5351.14854284)(1604.42880494,5348.71476842)(1600.90405249,5343.84727478)
\curveto(1597.42121307,5339.02163236)(1595.67980551,5332.34973352)(1595.67982458,5323.83155823)
}
}
{
\newrgbcolor{curcolor}{0 0 0}
\pscustom[linestyle=none,fillstyle=solid,fillcolor=curcolor]
{
\newpath
\moveto(1698.40199255,5386.45935059)
\curveto(1692.77909444,5376.80807834)(1688.60391246,5367.26180798)(1685.87643408,5357.82051086)
\curveto(1683.14890082,5348.3790754)(1681.78514791,5338.81182422)(1681.78517126,5329.11872864)
\curveto(1681.78514791,5319.42555211)(1683.14890082,5309.8163393)(1685.87643408,5300.2910614)
\curveto(1688.64587408,5290.80772184)(1692.82105607,5281.26145148)(1698.40199255,5271.65222168)
\lineto(1688.33119177,5271.65222168)
\curveto(1682.03691768,5281.51322124)(1677.31623454,5291.2063573)(1674.16912817,5300.73165894)
\curveto(1671.06395197,5310.2569364)(1669.51137173,5319.7192835)(1669.51138281,5329.11872864)
\curveto(1669.51137173,5338.4761312)(1671.06395197,5347.89651668)(1674.16912817,5357.37991333)
\curveto(1677.27427291,5366.86317252)(1681.99495605,5376.55630857)(1688.33119177,5386.45935059)
\lineto(1698.40199255,5386.45935059)
}
}
{
\newrgbcolor{curcolor}{0 0 0}
\pscustom[linestyle=none,fillstyle=solid,fillcolor=curcolor]
{
\newpath
\moveto(1719.11007666,5386.45935059)
\lineto(1729.18087744,5386.45935059)
\curveto(1735.47510124,5376.55630857)(1740.17480357,5366.86317252)(1743.27999854,5357.37991333)
\curveto(1746.42708614,5347.89651668)(1748.00064719,5338.4761312)(1748.0006864,5329.11872864)
\curveto(1748.00064719,5319.7192835)(1746.42708614,5310.2569364)(1743.27999854,5300.73165894)
\curveto(1740.17480357,5291.2063573)(1735.47510124,5281.51322124)(1729.18087744,5271.65222168)
\lineto(1719.11007666,5271.65222168)
\curveto(1724.69096286,5281.26145148)(1728.84516402,5290.80772184)(1731.57269263,5300.2910614)
\curveto(1734.34213729,5309.8163393)(1735.72687101,5319.42555211)(1735.72689795,5329.11872864)
\curveto(1735.72687101,5338.81182422)(1734.34213729,5348.3790754)(1731.57269263,5357.82051086)
\curveto(1728.84516402,5367.26180798)(1724.69096286,5376.80807834)(1719.11007666,5386.45935059)
}
}
{
\newrgbcolor{curcolor}{0 0 0}
\pscustom[linestyle=none,fillstyle=solid,fillcolor=curcolor]
{
}
}
{
\newrgbcolor{curcolor}{0 0 0}
\pscustom[linestyle=none,fillstyle=solid,fillcolor=curcolor]
{
\newpath
\moveto(1815.28622412,5304.63409424)
\lineto(1828.56709265,5304.63409424)
\lineto(1828.56709265,5288.646698)
\lineto(1815.28622412,5288.646698)
\lineto(1815.28622412,5304.63409424)
\moveto(1815.28622412,5355.30281067)
\lineto(1828.56709265,5355.30281067)
\lineto(1828.56709265,5339.31541443)
\lineto(1815.28622412,5339.31541443)
\lineto(1815.28622412,5355.30281067)
}
}
{
\newrgbcolor{curcolor}{0 0 0}
\pscustom[linestyle=none,fillstyle=solid,fillcolor=curcolor]
{
}
}
{
\newrgbcolor{curcolor}{0 0 0}
\pscustom[linestyle=none,fillstyle=solid,fillcolor=curcolor]
{
\newpath
\moveto(1953.75973486,5379.53567505)
\lineto(1953.75973486,5367.13600159)
\curveto(1948.93407866,5369.44381264)(1944.38124203,5371.16423938)(1940.1012113,5372.29728699)
\curveto(1935.82106993,5373.43016729)(1931.68784957,5373.99664927)(1927.70153784,5373.99673462)
\curveto(1920.7778263,5373.99664927)(1915.42771874,5372.65387717)(1911.6511991,5369.96841431)
\curveto(1907.91658733,5367.2827888)(1906.04929489,5363.46428065)(1906.04931616,5358.51287842)
\curveto(1906.04929489,5354.35860738)(1907.28716291,5351.21148529)(1909.76292395,5349.07150269)
\curveto(1912.28059664,5346.97336086)(1917.0222606,5345.27391493)(1923.98793005,5343.97315979)
\lineto(1931.66691565,5342.39959717)
\curveto(1941.15019668,5340.59519341)(1948.13680773,5337.40610969)(1952.62676978,5332.83233643)
\curveto(1957.15855774,5328.30043642)(1959.42448565,5322.21600037)(1959.4245603,5314.57901001)
\curveto(1959.42448565,5305.47331081)(1956.36128681,5298.57062301)(1950.23495459,5293.8709259)
\curveto(1944.15045308,5289.17121835)(1935.21262632,5286.82136718)(1923.42144751,5286.82136536)
\curveto(1918.9734763,5286.82136718)(1914.23181234,5287.32490672)(1909.19644141,5288.33198547)
\curveto(1904.20298326,5289.33906486)(1899.0207222,5290.82870265)(1893.6496427,5292.80090332)
\lineto(1893.6496427,5305.89294434)
\curveto(1898.81091406,5302.99757476)(1903.86729023,5300.81557011)(1908.81878638,5299.34692383)
\curveto(1913.77023443,5297.87825615)(1918.63778328,5297.14392766)(1923.42144751,5297.14393616)
\curveto(1930.6807705,5297.14392766)(1936.28264783,5298.57062301)(1940.22709631,5301.42402649)
\curveto(1944.17143389,5304.27740441)(1946.1436304,5308.34768233)(1946.14369177,5313.63487244)
\curveto(1946.1436304,5318.25062652)(1944.71693505,5321.85932653)(1941.86360144,5324.46098328)
\curveto(1939.05211528,5327.0625684)(1934.41535539,5329.0137841)(1927.95330786,5330.31463623)
\lineto(1920.21137976,5331.82525635)
\curveto(1910.72801641,5333.71348643)(1903.86729023,5336.6717812)(1899.62918066,5340.70014954)
\curveto(1895.39104139,5344.72841377)(1893.27197917,5350.3302911)(1893.27198767,5357.50579834)
\curveto(1893.27197917,5365.81413182)(1896.18831232,5372.36014578)(1902.02099585,5377.14385986)
\curveto(1907.89560652,5381.92739695)(1915.9732199,5384.31920975)(1926.25386023,5384.31930542)
\curveto(1930.65978969,5384.31920975)(1935.14968388,5383.92057428)(1939.72355627,5383.12339783)
\curveto(1944.29731877,5382.32603242)(1948.97604029,5381.13012602)(1953.75973486,5379.53567505)
}
}
{
\newrgbcolor{curcolor}{0 0 0}
\pscustom[linestyle=none,fillstyle=solid,fillcolor=curcolor]
{
\newpath
\moveto(1990.45521521,5379.15802002)
\lineto(1990.45521521,5359.14230347)
\lineto(2014.31042456,5359.14230347)
\lineto(2014.31042456,5350.14152527)
\lineto(1990.45521521,5350.14152527)
\lineto(1990.45521521,5311.8724823)
\curveto(1990.45519161,5306.12371604)(1991.23148172,5302.43109278)(1992.78408789,5300.79460144)
\curveto(1994.37860382,5299.1580858)(1997.58866836,5298.33983406)(2002.41429114,5298.33984375)
\lineto(2014.31042456,5298.33984375)
\lineto(2014.31042456,5288.646698)
\lineto(2002.41429114,5288.646698)
\curveto(1993.47642882,5288.646698)(1987.30806951,5290.3041823)(1983.9091947,5293.61915588)
\curveto(1980.51028578,5296.97608115)(1978.81083985,5303.0605172)(1978.81085181,5311.8724823)
\lineto(1978.81085181,5350.14152527)
\lineto(1970.31361365,5350.14152527)
\lineto(1970.31361365,5359.14230347)
\lineto(1978.81085181,5359.14230347)
\lineto(1978.81085181,5379.15802002)
\lineto(1990.45521521,5379.15802002)
}
}
{
\newrgbcolor{curcolor}{0 0 0}
\pscustom[linestyle=none,fillstyle=solid,fillcolor=curcolor]
{
\newpath
\moveto(2070.45513892,5348.31619263)
\curveto(2069.15427545,5349.07144226)(2067.7275801,5349.61694342)(2066.17504858,5349.95269775)
\curveto(2064.66438126,5350.3302911)(2062.98591614,5350.51911843)(2061.13964819,5350.5191803)
\curveto(2054.59359055,5350.51911843)(2049.55819519,5348.3790754)(2046.03344702,5344.0990448)
\curveto(2042.55060332,5339.86086492)(2040.80919576,5333.75544806)(2040.80921912,5325.78277588)
\lineto(2040.80921912,5288.646698)
\lineto(2029.16485571,5288.646698)
\lineto(2029.16485571,5359.14230347)
\lineto(2040.80921912,5359.14230347)
\lineto(2040.80921912,5348.19030762)
\curveto(2043.24297019,5352.47033413)(2046.4110731,5355.63843704)(2050.31353735,5357.69462585)
\curveto(2054.2159359,5359.7926382)(2058.95759986,5360.8416789)(2064.53854346,5360.8417511)
\curveto(2065.33576731,5360.8416789)(2066.21696149,5360.77873646)(2067.18212866,5360.65292358)
\curveto(2068.14719638,5360.56892832)(2069.21721789,5360.42206262)(2070.39219641,5360.21232605)
\lineto(2070.45513892,5348.31619263)
}
}
{
\newrgbcolor{curcolor}{0 0 0}
\pscustom[linestyle=none,fillstyle=solid,fillcolor=curcolor]
{
\newpath
\moveto(2082.72892737,5359.14230347)
\lineto(2094.31034827,5359.14230347)
\lineto(2094.31034827,5288.646698)
\lineto(2082.72892737,5288.646698)
\lineto(2082.72892737,5359.14230347)
\moveto(2082.72892737,5386.5852356)
\lineto(2094.31034827,5386.5852356)
\lineto(2094.31034827,5371.91963196)
\lineto(2082.72892737,5371.91963196)
\lineto(2082.72892737,5386.5852356)
}
}
{
\newrgbcolor{curcolor}{0 0 0}
\pscustom[linestyle=none,fillstyle=solid,fillcolor=curcolor]
{
\newpath
\moveto(2177.07974219,5331.1958313)
\lineto(2177.07974219,5288.646698)
\lineto(2165.49832129,5288.646698)
\lineto(2165.49832129,5330.81817627)
\curveto(2165.49826212,5337.49003294)(2164.19745166,5342.48346667)(2161.59588599,5345.79849243)
\curveto(2158.99420979,5349.11340389)(2155.09177839,5350.77088819)(2149.88858008,5350.77095032)
\curveto(2143.63625396,5350.77088819)(2138.70576267,5348.77771087)(2135.09709143,5344.79141235)
\curveto(2131.48836266,5340.80500155)(2129.68401266,5335.37097073)(2129.68403601,5328.48930359)
\lineto(2129.68403601,5288.646698)
\lineto(2118.03967261,5288.646698)
\lineto(2118.03967261,5359.14230347)
\lineto(2129.68403601,5359.14230347)
\lineto(2129.68403601,5348.19030762)
\curveto(2132.4534801,5352.4283725)(2135.70550627,5355.59647541)(2139.44012427,5357.69462585)
\curveto(2143.21663768,5359.7926382)(2147.55966617,5360.8416789)(2152.46922278,5360.8417511)
\curveto(2160.56777084,5360.8416789)(2166.69416852,5358.32398123)(2170.8484342,5353.28865051)
\curveto(2175.00257086,5348.29515214)(2177.07967144,5340.93088644)(2177.07974219,5331.1958313)
}
}
{
\newrgbcolor{curcolor}{0 0 0}
\pscustom[linestyle=none,fillstyle=solid,fillcolor=curcolor]
{
\newpath
\moveto(2246.69415259,5324.7127533)
\curveto(2246.69409405,5333.10504282)(2244.95268649,5339.60909516)(2241.46992468,5344.22492981)
\curveto(2238.02901788,5348.84065331)(2233.18244985,5351.14854284)(2226.93020605,5351.14860535)
\curveto(2220.71984634,5351.14854284)(2215.87327831,5348.84065331)(2212.39048743,5344.22492981)
\curveto(2208.9496097,5339.60909516)(2207.22918295,5333.10504282)(2207.22920203,5324.7127533)
\curveto(2207.22918295,5316.36235327)(2208.9496097,5309.87928175)(2212.39048743,5305.26351929)
\curveto(2215.87327831,5300.64772359)(2220.71984634,5298.33983406)(2226.93020605,5298.33984375)
\curveto(2233.18244985,5298.33983406)(2238.02901788,5300.64772359)(2241.46992468,5305.26351929)
\curveto(2244.95268649,5309.87928175)(2246.69409405,5316.36235327)(2246.69415259,5324.7127533)
\moveto(2258.27557349,5297.39570618)
\curveto(2258.27550337,5285.39467183)(2255.61093999,5276.47782589)(2250.28187537,5270.6451416)
\curveto(2244.95268649,5264.77053169)(2236.79114985,5261.83321773)(2225.79724097,5261.83319092)
\curveto(2221.72692542,5261.83321773)(2217.88743646,5262.14792994)(2214.27876257,5262.77732849)
\curveto(2210.67003645,5263.36481715)(2207.16624051,5264.28797297)(2203.76736426,5265.54679871)
\lineto(2203.76736426,5276.81350708)
\curveto(2207.16624051,5274.96720728)(2210.52317075,5273.60345437)(2213.83816504,5272.72224426)
\curveto(2217.15310797,5271.841066)(2220.53101902,5271.40046891)(2223.97190833,5271.40045166)
\curveto(2231.56692717,5271.40046891)(2237.25272776,5273.39364623)(2241.02932715,5277.37998962)
\curveto(2244.80582079,5281.32439392)(2246.69409405,5287.3039259)(2246.69415259,5295.31860352)
\lineto(2246.69415259,5301.04637146)
\curveto(2244.30228126,5296.89215789)(2241.23908242,5293.78699742)(2237.50454687,5291.73088074)
\curveto(2233.76991264,5289.67475788)(2229.30099926,5288.646698)(2224.09779333,5288.646698)
\curveto(2215.45366203,5288.646698)(2208.48803179,5291.94068579)(2203.20088171,5298.52867126)
\curveto(2197.91370155,5305.11663697)(2195.27011899,5313.84465559)(2195.2701261,5324.7127533)
\curveto(2195.27011899,5335.6227405)(2197.91370155,5344.37173993)(2203.20088171,5350.95977783)
\curveto(2208.48803179,5357.54769111)(2215.45366203,5360.8416789)(2224.09779333,5360.8417511)
\curveto(2229.30099926,5360.8416789)(2233.76991264,5359.81361902)(2237.50454687,5357.75756836)
\curveto(2241.23908242,5355.70137948)(2244.30228126,5352.59621901)(2246.69415259,5348.44207764)
\lineto(2246.69415259,5359.14230347)
\lineto(2258.27557349,5359.14230347)
\lineto(2258.27557349,5297.39570618)
}
}
{
\newrgbcolor{curcolor}{0.51764709 0.74509805 0.51764709}
\pscustom[linestyle=none,fillstyle=solid,fillcolor=curcolor]
{
\newpath
\moveto(1256.79159375,5168.56680444)
\curveto(1256.79159375,5149.15054371)(1241.05160449,5133.41055444)(1221.63534375,5133.41055444)
\curveto(1202.21908301,5133.41055444)(1186.47909375,5149.15054371)(1186.47909375,5168.56680444)
\curveto(1186.47909375,5187.98306518)(1202.21908301,5203.72305444)(1221.63534375,5203.72305444)
\curveto(1241.05160449,5203.72305444)(1256.79159375,5187.98306518)(1256.79159375,5168.56680444)
\closepath
}
}
{
\newrgbcolor{curcolor}{0.01176471 0.50196081 0.28235295}
\pscustom[linewidth=11.71875,linecolor=curcolor]
{
\newpath
\moveto(1256.79159375,5168.56680444)
\curveto(1256.79159375,5149.15054371)(1241.05160449,5133.41055444)(1221.63534375,5133.41055444)
\curveto(1202.21908301,5133.41055444)(1186.47909375,5149.15054371)(1186.47909375,5168.56680444)
\curveto(1186.47909375,5187.98306518)(1202.21908301,5203.72305444)(1221.63534375,5203.72305444)
\curveto(1241.05160449,5203.72305444)(1256.79159375,5187.98306518)(1256.79159375,5168.56680444)
\closepath
}
}
{
\newrgbcolor{curcolor}{0 0 0}
\pscustom[linestyle=none,fillstyle=solid,fillcolor=curcolor]
{
\newpath
\moveto(1385.64068579,5174.72068787)
\curveto(1385.64062725,5183.11297739)(1383.89921969,5189.61702973)(1380.41645789,5194.23286438)
\curveto(1376.97555108,5198.84858788)(1372.12898305,5201.15647742)(1365.87673926,5201.15653992)
\curveto(1359.66637955,5201.15647742)(1354.81981152,5198.84858788)(1351.33702063,5194.23286438)
\curveto(1347.8961429,5189.61702973)(1346.17571616,5183.11297739)(1346.17573523,5174.72068787)
\curveto(1346.17571616,5166.37028784)(1347.8961429,5159.88721632)(1351.33702063,5155.27145386)
\curveto(1354.81981152,5150.65565816)(1359.66637955,5148.34776863)(1365.87673926,5148.34777832)
\curveto(1372.12898305,5148.34776863)(1376.97555108,5150.65565816)(1380.41645789,5155.27145386)
\curveto(1383.89921969,5159.88721632)(1385.64062725,5166.37028784)(1385.64068579,5174.72068787)
\moveto(1397.22210669,5147.40364075)
\curveto(1397.22203657,5135.4026064)(1394.5574732,5126.48576046)(1389.22840857,5120.65307617)
\curveto(1383.89921969,5114.77846626)(1375.73768306,5111.8411523)(1364.74377417,5111.84112549)
\curveto(1360.67345862,5111.8411523)(1356.83396966,5112.15586451)(1353.22529578,5112.78526306)
\curveto(1349.61656965,5113.37275172)(1346.11277372,5114.29590754)(1342.71389746,5115.55473328)
\lineto(1342.71389746,5126.82144165)
\curveto(1346.11277372,5124.97514185)(1349.46970395,5123.61138894)(1352.78469824,5122.73017883)
\curveto(1356.09964117,5121.84900057)(1359.47755222,5121.40840348)(1362.91844153,5121.40838623)
\curveto(1370.51346037,5121.40840348)(1376.19926096,5123.4015808)(1379.97586035,5127.38792419)
\curveto(1383.752354,5131.33232849)(1385.64062725,5137.31186047)(1385.64068579,5145.32653809)
\lineto(1385.64068579,5151.05430603)
\curveto(1383.24881446,5146.90009246)(1380.18561562,5143.79493199)(1376.45108008,5141.73881531)
\curveto(1372.71644584,5139.68269245)(1368.24753246,5138.65463257)(1363.04432654,5138.65463257)
\curveto(1354.40019524,5138.65463257)(1347.434565,5141.94862036)(1342.14741492,5148.53660583)
\curveto(1336.86023475,5155.12457154)(1334.21665219,5163.85259016)(1334.2166593,5174.72068787)
\curveto(1334.21665219,5185.63067507)(1336.86023475,5194.3796745)(1342.14741492,5200.9677124)
\curveto(1347.434565,5207.55562568)(1354.40019524,5210.84961347)(1363.04432654,5210.84968567)
\curveto(1368.24753246,5210.84961347)(1372.71644584,5209.82155359)(1376.45108008,5207.76550293)
\curveto(1380.18561562,5205.70931405)(1383.24881446,5202.60415358)(1385.64068579,5198.45001221)
\lineto(1385.64068579,5209.15023804)
\lineto(1397.22210669,5209.15023804)
\lineto(1397.22210669,5147.40364075)
}
}
{
\newrgbcolor{curcolor}{0 0 0}
\pscustom[linestyle=none,fillstyle=solid,fillcolor=curcolor]
{
\newpath
\moveto(1481.37623572,5176.79779053)
\lineto(1481.37623572,5171.13296509)
\lineto(1428.12687659,5171.13296509)
\curveto(1428.63039692,5163.1602233)(1431.02220972,5157.07578724)(1435.30232214,5152.87963867)
\curveto(1439.62434345,5148.72542328)(1445.62485625,5146.64832269)(1453.30387854,5146.64833069)
\curveto(1457.75176673,5146.64832269)(1462.0528336,5147.19382386)(1466.20709204,5148.28483582)
\curveto(1470.40319756,5149.37582851)(1474.55739873,5151.012332)(1478.66970801,5153.1943512)
\lineto(1478.66970801,5142.24235535)
\curveto(1474.5154371,5140.47996338)(1470.25633186,5139.13719129)(1465.89237952,5138.21403503)
\curveto(1461.52831325,5137.29087966)(1457.1013615,5136.82930175)(1452.61151099,5136.82929993)
\curveto(1441.36575101,5136.82930175)(1432.44890507,5140.10230873)(1425.86094641,5146.64833069)
\curveto(1419.31491552,5153.19433666)(1416.04190854,5162.04824016)(1416.04191565,5173.21006775)
\curveto(1416.04190854,5184.74948088)(1419.14706901,5193.89711578)(1425.35740637,5200.65299988)
\curveto(1431.60967251,5207.45072161)(1440.02297892,5210.84961347)(1450.59735083,5210.84968567)
\curveto(1460.08063708,5210.84961347)(1467.57078767,5207.78641463)(1473.06782507,5201.66007996)
\curveto(1478.60669583,5195.5755809)(1481.37616327,5187.28815937)(1481.37623572,5176.79779053)
\moveto(1469.79481482,5180.19668579)
\curveto(1469.7108307,5186.53285007)(1467.92746151,5191.58922624)(1464.4447019,5195.36582947)
\curveto(1461.0037929,5199.14231927)(1456.42997545,5201.03059253)(1450.72323584,5201.03065491)
\curveto(1444.26110334,5201.03059253)(1439.07884229,5199.20526171)(1435.17643713,5195.55465698)
\curveto(1431.31594111,5191.90393845)(1429.09197483,5186.76363902)(1428.50453162,5180.13374329)
\lineto(1469.79481482,5180.19668579)
}
}
{
\newrgbcolor{curcolor}{0 0 0}
\pscustom[linestyle=none,fillstyle=solid,fillcolor=curcolor]
{
\newpath
\moveto(1511.58863806,5229.16595459)
\lineto(1511.58863806,5209.15023804)
\lineto(1535.44384741,5209.15023804)
\lineto(1535.44384741,5200.14945984)
\lineto(1511.58863806,5200.14945984)
\lineto(1511.58863806,5161.88041687)
\curveto(1511.58861446,5156.13165061)(1512.36490458,5152.43902735)(1513.91751074,5150.80253601)
\curveto(1515.51202667,5149.16602037)(1518.72209121,5148.34776863)(1523.54771399,5148.34777832)
\lineto(1535.44384741,5148.34777832)
\lineto(1535.44384741,5138.65463257)
\lineto(1523.54771399,5138.65463257)
\curveto(1514.60985167,5138.65463257)(1508.44149236,5140.31211687)(1505.04261755,5143.62709045)
\curveto(1501.64370863,5146.98401572)(1499.9442627,5153.06845177)(1499.94427466,5161.88041687)
\lineto(1499.94427466,5200.14945984)
\lineto(1491.4470365,5200.14945984)
\lineto(1491.4470365,5209.15023804)
\lineto(1499.94427466,5209.15023804)
\lineto(1499.94427466,5229.16595459)
\lineto(1511.58863806,5229.16595459)
}
}
{
\newrgbcolor{curcolor}{0 0 0}
\pscustom[linestyle=none,fillstyle=solid,fillcolor=curcolor]
{
\newpath
\moveto(1551.24241614,5232.62779236)
\lineto(1568.36277747,5232.62779236)
\lineto(1610.0307157,5154.01260376)
\lineto(1610.0307157,5232.62779236)
\lineto(1622.36744666,5232.62779236)
\lineto(1622.36744666,5138.65463257)
\lineto(1605.24708533,5138.65463257)
\lineto(1563.57914709,5217.26982117)
\lineto(1563.57914709,5138.65463257)
\lineto(1551.24241614,5138.65463257)
\lineto(1551.24241614,5232.62779236)
}
}
{
\newrgbcolor{curcolor}{0 0 0}
\pscustom[linestyle=none,fillstyle=solid,fillcolor=curcolor]
{
\newpath
\moveto(1679.20452856,5174.09126282)
\curveto(1669.84704134,5174.09122738)(1663.36396982,5173.02120587)(1659.75529456,5170.88119507)
\curveto(1656.14656981,5168.74111982)(1654.34221981,5165.09045818)(1654.34223914,5159.92919922)
\curveto(1654.34221981,5155.8169384)(1655.68499191,5152.54393142)(1658.37055945,5150.11016846)
\curveto(1661.09804191,5147.71834421)(1664.79066517,5146.52243781)(1669.44844031,5146.52244568)
\curveto(1675.86853496,5146.52243781)(1681.00883438,5148.78836572)(1684.869354,5153.32023621)
\curveto(1688.77173555,5157.89403899)(1690.72295125,5163.95749423)(1690.72300696,5171.51062012)
\lineto(1690.72300696,5174.09126282)
\lineto(1679.20452856,5174.09126282)
\moveto(1702.30442786,5178.87489319)
\lineto(1702.30442786,5138.65463257)
\lineto(1690.72300696,5138.65463257)
\lineto(1690.72300696,5149.3548584)
\curveto(1688.07936869,5145.07476165)(1684.7853809,5141.90665874)(1680.84103369,5139.85054016)
\curveto(1676.89659484,5137.83638082)(1672.07100763,5136.82930175)(1666.36425757,5136.82929993)
\curveto(1659.14682621,5136.82930175)(1653.39808318,5138.84345989)(1649.11801123,5142.8717804)
\curveto(1644.87987271,5146.94205409)(1642.7608105,5152.37608491)(1642.76081824,5159.17388916)
\curveto(1642.7608105,5167.10461633)(1645.40439306,5173.08414831)(1650.69157385,5177.11250305)
\curveto(1656.02068493,5181.14078088)(1663.95143262,5183.15493902)(1674.4838407,5183.15498352)
\lineto(1690.72300696,5183.15498352)
\lineto(1690.72300696,5184.28794861)
\curveto(1690.72295125,5189.61702973)(1688.96056288,5193.72926927)(1685.43583655,5196.62467957)
\curveto(1681.95297101,5199.56193555)(1677.04346054,5201.03059253)(1670.70729041,5201.03065491)
\curveto(1666.67893843,5201.03059253)(1662.75552622,5200.54803381)(1658.93704199,5199.58297729)
\curveto(1655.11850993,5198.61779892)(1651.44686748,5197.17012276)(1647.92210364,5195.23994446)
\lineto(1647.92210364,5205.94017029)
\curveto(1652.16021516,5207.57660649)(1656.2724547,5208.7934937)(1660.25883459,5209.59083557)
\curveto(1664.24516401,5210.42999719)(1668.1266146,5210.84961347)(1671.903198,5210.84968567)
\curveto(1682.09983671,5210.84961347)(1689.71587218,5208.20603091)(1694.75132727,5202.91893005)
\curveto(1699.78666289,5197.63170067)(1702.30436057,5189.61702973)(1702.30442786,5178.87489319)
}
}
{
\newrgbcolor{curcolor}{0 0 0}
\pscustom[linestyle=none,fillstyle=solid,fillcolor=curcolor]
{
\newpath
\moveto(1781.36021399,5195.61759949)
\curveto(1784.25549928,5200.82078439)(1787.71733359,5204.66027335)(1791.74572729,5207.13607788)
\curveto(1795.77396616,5209.61174545)(1800.51563012,5210.84961347)(1805.9707334,5210.84968567)
\curveto(1813.31392665,5210.84961347)(1818.97874642,5208.26897335)(1822.96520972,5203.10775757)
\curveto(1826.95145573,5197.9883745)(1828.94463306,5190.68705124)(1828.94474768,5181.20376587)
\lineto(1828.94474768,5138.65463257)
\lineto(1817.30038428,5138.65463257)
\lineto(1817.30038428,5180.82611084)
\curveto(1817.3002813,5187.58189077)(1816.10437491,5192.59630531)(1813.7126615,5195.86936951)
\curveto(1811.32074932,5199.14231927)(1807.67008769,5200.77882276)(1802.76066565,5200.77888489)
\curveto(1796.76006442,5200.77882276)(1792.01840046,5198.78564544)(1788.53565955,5194.79934692)
\curveto(1785.05277022,5190.81293612)(1783.31136266,5185.3789053)(1783.31143164,5178.49723816)
\lineto(1783.31143164,5138.65463257)
\lineto(1771.66706824,5138.65463257)
\lineto(1771.66706824,5180.82611084)
\curveto(1771.6670109,5187.6238524)(1770.4711045,5192.63826694)(1768.07934546,5195.86936951)
\curveto(1765.68747891,5199.14231927)(1761.99485565,5200.77882276)(1757.0014646,5200.77888489)
\curveto(1751.08483238,5200.77882276)(1746.38513005,5198.76466462)(1742.90234351,5194.73640442)
\curveto(1739.41949981,5190.74999368)(1737.67809225,5185.33694367)(1737.6781156,5178.49723816)
\lineto(1737.6781156,5138.65463257)
\lineto(1726.0337522,5138.65463257)
\lineto(1726.0337522,5209.15023804)
\lineto(1737.6781156,5209.15023804)
\lineto(1737.6781156,5198.19824219)
\curveto(1740.32167481,5202.52023032)(1743.48977772,5205.70931405)(1747.18243384,5207.76550293)
\curveto(1750.87502424,5209.82155359)(1755.26001436,5210.84961347)(1760.33741736,5210.84968567)
\curveto(1765.45668996,5210.84961347)(1769.79971845,5209.54880301)(1773.36651587,5206.94725037)
\curveto(1776.97515683,5204.34556114)(1779.63972021,5200.56901462)(1781.36021399,5195.61759949)
}
}
{
\newrgbcolor{curcolor}{0 0 0}
\pscustom[linestyle=none,fillstyle=solid,fillcolor=curcolor]
{
\newpath
\moveto(1912.15473914,5176.79779053)
\lineto(1912.15473914,5171.13296509)
\lineto(1858.90538,5171.13296509)
\curveto(1859.40890034,5163.1602233)(1861.80071314,5157.07578724)(1866.08082556,5152.87963867)
\curveto(1870.40284687,5148.72542328)(1876.40335967,5146.64832269)(1884.08238196,5146.64833069)
\curveto(1888.53027015,5146.64832269)(1892.83133701,5147.19382386)(1896.98559546,5148.28483582)
\curveto(1901.18170098,5149.37582851)(1905.33590215,5151.012332)(1909.44821143,5153.1943512)
\lineto(1909.44821143,5142.24235535)
\curveto(1905.29394052,5140.47996338)(1901.03483528,5139.13719129)(1896.67088293,5138.21403503)
\curveto(1892.30681666,5137.29087966)(1887.87986491,5136.82930175)(1883.3900144,5136.82929993)
\curveto(1872.14425443,5136.82930175)(1863.22740849,5140.10230873)(1856.63944983,5146.64833069)
\curveto(1850.09341894,5153.19433666)(1846.82041195,5162.04824016)(1846.82041907,5173.21006775)
\curveto(1846.82041195,5184.74948088)(1849.92557242,5193.89711578)(1856.13590979,5200.65299988)
\curveto(1862.38817593,5207.45072161)(1870.80148233,5210.84961347)(1881.37585425,5210.84968567)
\curveto(1890.8591405,5210.84961347)(1898.34929109,5207.78641463)(1903.84632849,5201.66007996)
\curveto(1909.38519924,5195.5755809)(1912.15466669,5187.28815937)(1912.15473914,5176.79779053)
\moveto(1900.57331824,5180.19668579)
\curveto(1900.48933412,5186.53285007)(1898.70596493,5191.58922624)(1895.22320532,5195.36582947)
\curveto(1891.78229631,5199.14231927)(1887.20847887,5201.03059253)(1881.50173926,5201.03065491)
\curveto(1875.03960676,5201.03059253)(1869.85734571,5199.20526171)(1865.95494055,5195.55465698)
\curveto(1862.09444453,5191.90393845)(1859.87047825,5186.76363902)(1859.28303503,5180.13374329)
\lineto(1900.57331824,5180.19668579)
}
}
{
\newrgbcolor{curcolor}{0 0 0}
\pscustom[linestyle=none,fillstyle=solid,fillcolor=curcolor]
{
\newpath
\moveto(1958.73219275,5236.46728516)
\curveto(1953.10929463,5226.81601291)(1948.93411265,5217.26974255)(1946.20663428,5207.82844543)
\curveto(1943.47910102,5198.38700997)(1942.11534811,5188.81975879)(1942.11537146,5179.12666321)
\curveto(1942.11534811,5169.43348668)(1943.47910102,5159.82427387)(1946.20663428,5150.29899597)
\curveto(1948.97607428,5140.81565641)(1953.15125626,5131.26938605)(1958.73219275,5121.66015625)
\lineto(1948.66139197,5121.66015625)
\curveto(1942.36711788,5131.52115582)(1937.64643473,5141.21429187)(1934.49932837,5150.73959351)
\curveto(1931.39415216,5160.26487097)(1929.84157193,5169.72721807)(1929.84158301,5179.12666321)
\curveto(1929.84157193,5188.48406577)(1931.39415216,5197.90445125)(1934.49932837,5207.3878479)
\curveto(1937.6044731,5216.87110709)(1942.32515625,5226.56424314)(1948.66139197,5236.46728516)
\lineto(1958.73219275,5236.46728516)
}
}
{
\newrgbcolor{curcolor}{0 0 0}
\pscustom[linestyle=none,fillstyle=solid,fillcolor=curcolor]
{
\newpath
\moveto(1979.44027686,5236.46728516)
\lineto(1989.51107764,5236.46728516)
\curveto(1995.80530144,5226.56424314)(2000.50500377,5216.87110709)(2003.61019873,5207.3878479)
\curveto(2006.75728633,5197.90445125)(2008.33084738,5188.48406577)(2008.3308866,5179.12666321)
\curveto(2008.33084738,5169.72721807)(2006.75728633,5160.26487097)(2003.61019873,5150.73959351)
\curveto(2000.50500377,5141.21429187)(1995.80530144,5131.52115582)(1989.51107764,5121.66015625)
\lineto(1979.44027686,5121.66015625)
\curveto(1985.02116305,5131.26938605)(1989.17536422,5140.81565641)(1991.90289282,5150.29899597)
\curveto(1994.67233748,5159.82427387)(1996.05707121,5169.43348668)(1996.05709814,5179.12666321)
\curveto(1996.05707121,5188.81975879)(1994.67233748,5198.38700997)(1991.90289282,5207.82844543)
\curveto(1989.17536422,5217.26974255)(1985.02116305,5226.81601291)(1979.44027686,5236.46728516)
}
}
{
\newrgbcolor{curcolor}{0 0 0}
\pscustom[linestyle=none,fillstyle=solid,fillcolor=curcolor]
{
}
}
{
\newrgbcolor{curcolor}{0 0 0}
\pscustom[linestyle=none,fillstyle=solid,fillcolor=curcolor]
{
\newpath
\moveto(2075.61642432,5154.64202881)
\lineto(2088.89729285,5154.64202881)
\lineto(2088.89729285,5138.65463257)
\lineto(2075.61642432,5138.65463257)
\lineto(2075.61642432,5154.64202881)
\moveto(2075.61642432,5205.31074524)
\lineto(2088.89729285,5205.31074524)
\lineto(2088.89729285,5189.323349)
\lineto(2075.61642432,5189.323349)
\lineto(2075.61642432,5205.31074524)
}
}
{
\newrgbcolor{curcolor}{0 0 0}
\pscustom[linestyle=none,fillstyle=solid,fillcolor=curcolor]
{
}
}
{
\newrgbcolor{curcolor}{0 0 0}
\pscustom[linestyle=none,fillstyle=solid,fillcolor=curcolor]
{
\newpath
\moveto(2214.08993506,5229.54360962)
\lineto(2214.08993506,5217.14393616)
\curveto(2209.26427886,5219.45174721)(2204.71144222,5221.17217395)(2200.4314115,5222.30522156)
\curveto(2196.15127012,5223.43810186)(2192.01804977,5224.00458384)(2188.03173804,5224.00466919)
\curveto(2181.1080265,5224.00458384)(2175.75791893,5222.66181174)(2171.98139929,5219.97634888)
\curveto(2168.24678753,5217.29072337)(2166.37949508,5213.47221522)(2166.37951636,5208.52081299)
\curveto(2166.37949508,5204.36654195)(2167.61736311,5201.21941986)(2170.09312415,5199.07943726)
\curveto(2172.61079683,5196.98129543)(2177.35246079,5195.2818495)(2184.31813025,5193.98109436)
\lineto(2191.99711584,5192.40753174)
\curveto(2201.48039687,5190.60312798)(2208.46700793,5187.41404426)(2212.95696997,5182.840271)
\curveto(2217.48875794,5178.30837099)(2219.75468585,5172.22393494)(2219.7547605,5164.58694458)
\curveto(2219.75468585,5155.48124538)(2216.69148701,5148.57855758)(2210.56515479,5143.87886047)
\curveto(2204.48065327,5139.17915292)(2195.54282652,5136.82930175)(2183.75164771,5136.82929993)
\curveto(2179.30367649,5136.82930175)(2174.56201254,5137.33284129)(2169.5266416,5138.33992004)
\curveto(2164.53318345,5139.34699943)(2159.3509224,5140.83663722)(2153.9798429,5142.80883789)
\lineto(2153.9798429,5155.90087891)
\curveto(2159.14111426,5153.00550933)(2164.19749043,5150.82350468)(2169.14898657,5149.3548584)
\curveto(2174.10043463,5147.88619072)(2178.96798347,5147.15186223)(2183.75164771,5147.15187073)
\curveto(2191.0109707,5147.15186223)(2196.61284803,5148.57855758)(2200.55729651,5151.43196106)
\curveto(2204.50163408,5154.28533898)(2206.4738306,5158.3556169)(2206.47389197,5163.64280701)
\curveto(2206.4738306,5168.25856109)(2205.04713525,5171.8672611)(2202.19380164,5174.46891785)
\curveto(2199.38231547,5177.07050297)(2194.74555558,5179.02171867)(2188.28350806,5180.3225708)
\lineto(2180.54157996,5181.83319092)
\curveto(2171.0582166,5183.721421)(2164.19749043,5186.67971577)(2159.95938086,5190.70808411)
\curveto(2155.72124158,5194.73634834)(2153.60217937,5200.33822567)(2153.60218787,5207.51373291)
\curveto(2153.60217937,5215.82206639)(2156.51851251,5222.36808035)(2162.35119604,5227.15179443)
\curveto(2168.22580671,5231.93533152)(2176.3034201,5234.32714432)(2186.58406042,5234.32723999)
\curveto(2190.98998988,5234.32714432)(2195.47988407,5233.92850885)(2200.05375647,5233.1313324)
\curveto(2204.62751897,5232.33396699)(2209.30624049,5231.13806059)(2214.08993506,5229.54360962)
}
}
{
\newrgbcolor{curcolor}{0 0 0}
\pscustom[linestyle=none,fillstyle=solid,fillcolor=curcolor]
{
\newpath
\moveto(2250.78541541,5229.16595459)
\lineto(2250.78541541,5209.15023804)
\lineto(2274.64062476,5209.15023804)
\lineto(2274.64062476,5200.14945984)
\lineto(2250.78541541,5200.14945984)
\lineto(2250.78541541,5161.88041687)
\curveto(2250.7853918,5156.13165061)(2251.56168192,5152.43902735)(2253.11428809,5150.80253601)
\curveto(2254.70880402,5149.16602037)(2257.91886856,5148.34776863)(2262.74449133,5148.34777832)
\lineto(2274.64062476,5148.34777832)
\lineto(2274.64062476,5138.65463257)
\lineto(2262.74449133,5138.65463257)
\curveto(2253.80662901,5138.65463257)(2247.6382697,5140.31211687)(2244.2393949,5143.62709045)
\curveto(2240.84048598,5146.98401572)(2239.14104004,5153.06845177)(2239.141052,5161.88041687)
\lineto(2239.141052,5200.14945984)
\lineto(2230.64381384,5200.14945984)
\lineto(2230.64381384,5209.15023804)
\lineto(2239.141052,5209.15023804)
\lineto(2239.141052,5229.16595459)
\lineto(2250.78541541,5229.16595459)
}
}
{
\newrgbcolor{curcolor}{0 0 0}
\pscustom[linestyle=none,fillstyle=solid,fillcolor=curcolor]
{
\newpath
\moveto(2330.78533911,5198.3241272)
\curveto(2329.48447565,5199.07937683)(2328.0577803,5199.62487799)(2326.50524878,5199.96063232)
\curveto(2324.99458146,5200.33822567)(2323.31611634,5200.527053)(2321.46984839,5200.52711487)
\curveto(2314.92379074,5200.527053)(2309.88839539,5198.38700997)(2306.36364722,5194.10697937)
\curveto(2302.88080352,5189.86879949)(2301.13939596,5183.76338263)(2301.13941931,5175.79071045)
\lineto(2301.13941931,5138.65463257)
\lineto(2289.49505591,5138.65463257)
\lineto(2289.49505591,5209.15023804)
\lineto(2301.13941931,5209.15023804)
\lineto(2301.13941931,5198.19824219)
\curveto(2303.57317038,5202.4782687)(2306.74127329,5205.64637161)(2310.64373755,5207.70256042)
\curveto(2314.54613609,5209.80057277)(2319.28780005,5210.84961347)(2324.86874365,5210.84968567)
\curveto(2325.6659675,5210.84961347)(2326.54716169,5210.78667103)(2327.51232886,5210.66085815)
\curveto(2328.47739658,5210.57686289)(2329.54741809,5210.42999719)(2330.72239661,5210.22026062)
\lineto(2330.78533911,5198.3241272)
}
}
{
\newrgbcolor{curcolor}{0 0 0}
\pscustom[linestyle=none,fillstyle=solid,fillcolor=curcolor]
{
\newpath
\moveto(2343.05912756,5209.15023804)
\lineto(2354.64054846,5209.15023804)
\lineto(2354.64054846,5138.65463257)
\lineto(2343.05912756,5138.65463257)
\lineto(2343.05912756,5209.15023804)
\moveto(2343.05912756,5236.59317017)
\lineto(2354.64054846,5236.59317017)
\lineto(2354.64054846,5221.92756653)
\lineto(2343.05912756,5221.92756653)
\lineto(2343.05912756,5236.59317017)
}
}
{
\newrgbcolor{curcolor}{0 0 0}
\pscustom[linestyle=none,fillstyle=solid,fillcolor=curcolor]
{
\newpath
\moveto(2437.40994238,5181.20376587)
\lineto(2437.40994238,5138.65463257)
\lineto(2425.82852148,5138.65463257)
\lineto(2425.82852148,5180.82611084)
\curveto(2425.82846232,5187.49796751)(2424.52765185,5192.49140124)(2421.92608618,5195.806427)
\curveto(2419.32440998,5199.12133846)(2415.42197858,5200.77882276)(2410.21878027,5200.77888489)
\curveto(2403.96645415,5200.77882276)(2399.03596287,5198.78564544)(2395.42729163,5194.79934692)
\curveto(2391.81856286,5190.81293612)(2390.01421285,5185.3789053)(2390.01423621,5178.49723816)
\lineto(2390.01423621,5138.65463257)
\lineto(2378.3698728,5138.65463257)
\lineto(2378.3698728,5209.15023804)
\lineto(2390.01423621,5209.15023804)
\lineto(2390.01423621,5198.19824219)
\curveto(2392.7836803,5202.43630707)(2396.03570647,5205.60440998)(2399.77032446,5207.70256042)
\curveto(2403.54683787,5209.80057277)(2407.88986637,5210.84961347)(2412.79942297,5210.84968567)
\curveto(2420.89797103,5210.84961347)(2427.02436872,5208.3319158)(2431.1786344,5203.29658508)
\curveto(2435.33277105,5198.30308671)(2437.40987164,5190.93882101)(2437.40994238,5181.20376587)
}
}
{
\newrgbcolor{curcolor}{0 0 0}
\pscustom[linestyle=none,fillstyle=solid,fillcolor=curcolor]
{
\newpath
\moveto(2507.02435278,5174.72068787)
\curveto(2507.02429425,5183.11297739)(2505.28288669,5189.61702973)(2501.80012488,5194.23286438)
\curveto(2498.35921807,5198.84858788)(2493.51265004,5201.15647742)(2487.26040625,5201.15653992)
\curveto(2481.05004654,5201.15647742)(2476.20347851,5198.84858788)(2472.72068762,5194.23286438)
\curveto(2469.2798099,5189.61702973)(2467.55938315,5183.11297739)(2467.55940222,5174.72068787)
\curveto(2467.55938315,5166.37028784)(2469.2798099,5159.88721632)(2472.72068762,5155.27145386)
\curveto(2476.20347851,5150.65565816)(2481.05004654,5148.34776863)(2487.26040625,5148.34777832)
\curveto(2493.51265004,5148.34776863)(2498.35921807,5150.65565816)(2501.80012488,5155.27145386)
\curveto(2505.28288669,5159.88721632)(2507.02429425,5166.37028784)(2507.02435278,5174.72068787)
\moveto(2518.60577368,5147.40364075)
\curveto(2518.60570356,5135.4026064)(2515.94114019,5126.48576046)(2510.61207556,5120.65307617)
\curveto(2505.28288669,5114.77846626)(2497.12135005,5111.8411523)(2486.12744116,5111.84112549)
\curveto(2482.05712561,5111.8411523)(2478.21763665,5112.15586451)(2474.60896277,5112.78526306)
\curveto(2471.00023664,5113.37275172)(2467.49644071,5114.29590754)(2464.09756445,5115.55473328)
\lineto(2464.09756445,5126.82144165)
\curveto(2467.49644071,5124.97514185)(2470.85337094,5123.61138894)(2474.16836523,5122.73017883)
\curveto(2477.48330816,5121.84900057)(2480.86121921,5121.40840348)(2484.30210852,5121.40838623)
\curveto(2491.89712737,5121.40840348)(2497.58292796,5123.4015808)(2501.35952734,5127.38792419)
\curveto(2505.13602099,5131.33232849)(2507.02429425,5137.31186047)(2507.02435278,5145.32653809)
\lineto(2507.02435278,5151.05430603)
\curveto(2504.63248145,5146.90009246)(2501.56928261,5143.79493199)(2497.83474707,5141.73881531)
\curveto(2494.10011283,5139.68269245)(2489.63119946,5138.65463257)(2484.42799353,5138.65463257)
\curveto(2475.78386223,5138.65463257)(2468.81823199,5141.94862036)(2463.53108191,5148.53660583)
\curveto(2458.24390174,5155.12457154)(2455.60031918,5163.85259016)(2455.60032629,5174.72068787)
\curveto(2455.60031918,5185.63067507)(2458.24390174,5194.3796745)(2463.53108191,5200.9677124)
\curveto(2468.81823199,5207.55562568)(2475.78386223,5210.84961347)(2484.42799353,5210.84968567)
\curveto(2489.63119946,5210.84961347)(2494.10011283,5209.82155359)(2497.83474707,5207.76550293)
\curveto(2501.56928261,5205.70931405)(2504.63248145,5202.60415358)(2507.02435278,5198.45001221)
\lineto(2507.02435278,5209.15023804)
\lineto(2518.60577368,5209.15023804)
\lineto(2518.60577368,5147.40364075)
}
}
{
\newrgbcolor{curcolor}{0.51764709 0.74509805 0.51764709}
\pscustom[linestyle=none,fillstyle=solid,fillcolor=curcolor]
{
\newpath
\moveto(1256.79159375,5018.57473901)
\curveto(1256.79159375,4999.15847828)(1241.05160449,4983.41848901)(1221.63534375,4983.41848901)
\curveto(1202.21908301,4983.41848901)(1186.47909375,4999.15847828)(1186.47909375,5018.57473901)
\curveto(1186.47909375,5037.99099975)(1202.21908301,5053.73098901)(1221.63534375,5053.73098901)
\curveto(1241.05160449,5053.73098901)(1256.79159375,5037.99099975)(1256.79159375,5018.57473901)
\closepath
}
}
{
\newrgbcolor{curcolor}{0.01176471 0.50196081 0.28235295}
\pscustom[linewidth=11.71875,linecolor=curcolor]
{
\newpath
\moveto(1256.79159375,5018.57473901)
\curveto(1256.79159375,4999.15847828)(1241.05160449,4983.41848901)(1221.63534375,4983.41848901)
\curveto(1202.21908301,4983.41848901)(1186.47909375,4999.15847828)(1186.47909375,5018.57473901)
\curveto(1186.47909375,5037.99099975)(1202.21908301,5053.73098901)(1221.63534375,5053.73098901)
\curveto(1241.05160449,5053.73098901)(1256.79159375,5037.99099975)(1256.79159375,5018.57473901)
\closepath
}
}
{
\newrgbcolor{curcolor}{0 0 0}
\pscustom[linestyle=none,fillstyle=solid,fillcolor=curcolor]
{
\newpath
\moveto(1385.64068579,5024.72877502)
\curveto(1385.64062725,5033.12106455)(1383.89921969,5039.62511688)(1380.41645789,5044.24095154)
\curveto(1376.97555108,5048.85667504)(1372.12898305,5051.16456457)(1365.87673926,5051.16462708)
\curveto(1359.66637955,5051.16456457)(1354.81981152,5048.85667504)(1351.33702063,5044.24095154)
\curveto(1347.8961429,5039.62511688)(1346.17571616,5033.12106455)(1346.17573523,5024.72877502)
\curveto(1346.17571616,5016.37837499)(1347.8961429,5009.89530347)(1351.33702063,5005.27954102)
\curveto(1354.81981152,5000.66374532)(1359.66637955,4998.35585579)(1365.87673926,4998.35586548)
\curveto(1372.12898305,4998.35585579)(1376.97555108,5000.66374532)(1380.41645789,5005.27954102)
\curveto(1383.89921969,5009.89530347)(1385.64062725,5016.37837499)(1385.64068579,5024.72877502)
\moveto(1397.22210669,4997.41172791)
\curveto(1397.22203657,4985.41069356)(1394.5574732,4976.49384762)(1389.22840857,4970.66116333)
\curveto(1383.89921969,4964.78655342)(1375.73768306,4961.84923946)(1364.74377417,4961.84921265)
\curveto(1360.67345862,4961.84923946)(1356.83396966,4962.16395167)(1353.22529578,4962.79335022)
\curveto(1349.61656965,4963.38083888)(1346.11277372,4964.3039947)(1342.71389746,4965.56282043)
\lineto(1342.71389746,4976.82952881)
\curveto(1346.11277372,4974.98322901)(1349.46970395,4973.6194761)(1352.78469824,4972.73826599)
\curveto(1356.09964117,4971.85708773)(1359.47755222,4971.41649063)(1362.91844153,4971.41647339)
\curveto(1370.51346037,4971.41649063)(1376.19926096,4973.40966796)(1379.97586035,4977.39601135)
\curveto(1383.752354,4981.34041565)(1385.64062725,4987.31994763)(1385.64068579,4995.33462524)
\lineto(1385.64068579,5001.06239319)
\curveto(1383.24881446,4996.90817962)(1380.18561562,4993.80301915)(1376.45108008,4991.74690247)
\curveto(1372.71644584,4989.69077961)(1368.24753246,4988.66271973)(1363.04432654,4988.66271973)
\curveto(1354.40019524,4988.66271973)(1347.434565,4991.95670752)(1342.14741492,4998.54469299)
\curveto(1336.86023475,5005.1326587)(1334.21665219,5013.86067732)(1334.2166593,5024.72877502)
\curveto(1334.21665219,5035.63876223)(1336.86023475,5044.38776166)(1342.14741492,5050.97579956)
\curveto(1347.434565,5057.56371284)(1354.40019524,5060.85770063)(1363.04432654,5060.85777283)
\curveto(1368.24753246,5060.85770063)(1372.71644584,5059.82964075)(1376.45108008,5057.77359009)
\curveto(1380.18561562,5055.71740121)(1383.24881446,5052.61224074)(1385.64068579,5048.45809937)
\lineto(1385.64068579,5059.1583252)
\lineto(1397.22210669,5059.1583252)
\lineto(1397.22210669,4997.41172791)
}
}
{
\newrgbcolor{curcolor}{0 0 0}
\pscustom[linestyle=none,fillstyle=solid,fillcolor=curcolor]
{
\newpath
\moveto(1481.37623572,5026.80587769)
\lineto(1481.37623572,5021.14105225)
\lineto(1428.12687659,5021.14105225)
\curveto(1428.63039692,5013.16831046)(1431.02220972,5007.0838744)(1435.30232214,5002.88772583)
\curveto(1439.62434345,4998.73351044)(1445.62485625,4996.65640985)(1453.30387854,4996.65641785)
\curveto(1457.75176673,4996.65640985)(1462.0528336,4997.20191102)(1466.20709204,4998.29292297)
\curveto(1470.40319756,4999.38391567)(1474.55739873,5001.02041916)(1478.66970801,5003.20243835)
\lineto(1478.66970801,4992.2504425)
\curveto(1474.5154371,4990.48805054)(1470.25633186,4989.14527845)(1465.89237952,4988.22212219)
\curveto(1461.52831325,4987.29896682)(1457.1013615,4986.83738891)(1452.61151099,4986.83738708)
\curveto(1441.36575101,4986.83738891)(1432.44890507,4990.11039589)(1425.86094641,4996.65641785)
\curveto(1419.31491552,5003.20242381)(1416.04190854,5012.05632731)(1416.04191565,5023.21815491)
\curveto(1416.04190854,5034.75756804)(1419.14706901,5043.90520294)(1425.35740637,5050.66108704)
\curveto(1431.60967251,5057.45880877)(1440.02297892,5060.85770063)(1450.59735083,5060.85777283)
\curveto(1460.08063708,5060.85770063)(1467.57078767,5057.79450179)(1473.06782507,5051.66816711)
\curveto(1478.60669583,5045.58366805)(1481.37616327,5037.29624653)(1481.37623572,5026.80587769)
\moveto(1469.79481482,5030.20477295)
\curveto(1469.7108307,5036.54093723)(1467.92746151,5041.5973134)(1464.4447019,5045.37391663)
\curveto(1461.0037929,5049.15040643)(1456.42997545,5051.03867969)(1450.72323584,5051.03874207)
\curveto(1444.26110334,5051.03867969)(1439.07884229,5049.21334887)(1435.17643713,5045.56274414)
\curveto(1431.31594111,5041.91202561)(1429.09197483,5036.77172618)(1428.50453162,5030.14183044)
\lineto(1469.79481482,5030.20477295)
}
}
{
\newrgbcolor{curcolor}{0 0 0}
\pscustom[linestyle=none,fillstyle=solid,fillcolor=curcolor]
{
\newpath
\moveto(1511.58863806,5079.17404175)
\lineto(1511.58863806,5059.1583252)
\lineto(1535.44384741,5059.1583252)
\lineto(1535.44384741,5050.157547)
\lineto(1511.58863806,5050.157547)
\lineto(1511.58863806,5011.88850403)
\curveto(1511.58861446,5006.13973777)(1512.36490458,5002.44711451)(1513.91751074,5000.81062317)
\curveto(1515.51202667,4999.17410753)(1518.72209121,4998.35585579)(1523.54771399,4998.35586548)
\lineto(1535.44384741,4998.35586548)
\lineto(1535.44384741,4988.66271973)
\lineto(1523.54771399,4988.66271973)
\curveto(1514.60985167,4988.66271973)(1508.44149236,4990.32020403)(1505.04261755,4993.63517761)
\curveto(1501.64370863,4996.99210288)(1499.9442627,5003.07653893)(1499.94427466,5011.88850403)
\lineto(1499.94427466,5050.157547)
\lineto(1491.4470365,5050.157547)
\lineto(1491.4470365,5059.1583252)
\lineto(1499.94427466,5059.1583252)
\lineto(1499.94427466,5079.17404175)
\lineto(1511.58863806,5079.17404175)
}
}
{
\newrgbcolor{curcolor}{0 0 0}
\pscustom[linestyle=none,fillstyle=solid,fillcolor=curcolor]
{
\newpath
\moveto(1607.57595801,5079.55169678)
\lineto(1607.57595801,5067.15202332)
\curveto(1602.75030181,5069.45983436)(1598.19746517,5071.18026111)(1593.91743445,5072.31330872)
\curveto(1589.63729307,5073.44618902)(1585.50407272,5074.012671)(1581.51776099,5074.01275635)
\curveto(1574.59404945,5074.012671)(1569.24394188,5072.6698989)(1565.46742224,5069.98443604)
\curveto(1561.73281048,5067.29881052)(1559.86551803,5063.48030238)(1559.86553931,5058.52890015)
\curveto(1559.86551803,5054.37462911)(1561.10338606,5051.22750702)(1563.57914709,5049.08752441)
\curveto(1566.09681978,5046.98938259)(1570.83848374,5045.28993666)(1577.8041532,5043.98918152)
\lineto(1585.48313879,5042.4156189)
\curveto(1594.96641982,5040.61121514)(1601.95303088,5037.42213142)(1606.44299292,5032.84835815)
\curveto(1610.97478089,5028.31645815)(1613.2407088,5022.23202209)(1613.24078345,5014.59503174)
\curveto(1613.2407088,5005.48933254)(1610.17750996,4998.58664474)(1604.05117773,4993.88694763)
\curveto(1597.96667622,4989.18724008)(1589.02884946,4986.83738891)(1577.23767065,4986.83738708)
\curveto(1572.78969944,4986.83738891)(1568.04803548,4987.34092845)(1563.01266455,4988.3480072)
\curveto(1558.0192064,4989.35508659)(1552.83694535,4990.84472438)(1547.46586584,4992.81692505)
\lineto(1547.46586584,5005.90896606)
\curveto(1552.62713721,5003.01359649)(1557.68351338,5000.83159184)(1562.63500952,4999.36294556)
\curveto(1567.58645758,4997.89427788)(1572.45400642,4997.15994939)(1577.23767065,4997.15995789)
\curveto(1584.49699364,4997.15994939)(1590.09887098,4998.58664474)(1594.04331946,5001.44004822)
\curveto(1597.98765703,5004.29342614)(1599.95985355,5008.36370405)(1599.95991492,5013.65089417)
\curveto(1599.95985355,5018.26664825)(1598.5331582,5021.87534826)(1595.67982458,5024.477005)
\curveto(1592.86833842,5027.07859012)(1588.23157853,5029.02980582)(1581.76953101,5030.33065796)
\lineto(1574.02760291,5031.84127808)
\curveto(1564.54423955,5033.72950816)(1557.68351338,5036.68780293)(1553.44540381,5040.71617126)
\curveto(1549.20726453,5044.7444355)(1547.08820232,5050.34631283)(1547.08821082,5057.52182007)
\curveto(1547.08820232,5065.83015355)(1550.00453546,5072.37616751)(1555.83721899,5077.15988159)
\curveto(1561.71182966,5081.94341868)(1569.78944304,5084.33523148)(1580.07008337,5084.33532715)
\curveto(1584.47601283,5084.33523148)(1588.96590702,5083.93659601)(1593.53977942,5083.13941956)
\curveto(1598.11354192,5082.34205415)(1602.79226344,5081.14614775)(1607.57595801,5079.55169678)
}
}
{
\newrgbcolor{curcolor}{0 0 0}
\pscustom[linestyle=none,fillstyle=solid,fillcolor=curcolor]
{
\newpath
\moveto(1660.13294958,5051.03874207)
\curveto(1653.92258918,5051.03867969)(1649.01307871,5048.60490527)(1645.40440344,5043.7374115)
\curveto(1641.7956787,5038.91176921)(1639.9913287,5032.28183199)(1639.99134802,5023.84757996)
\curveto(1639.9913287,5015.41325755)(1641.77469789,5008.76233952)(1645.34146094,5003.89480591)
\curveto(1648.95013627,4999.06920346)(1653.88062755,4996.65640985)(1660.13294958,4996.65641785)
\curveto(1666.30126943,4996.65640985)(1671.18979909,4999.09018427)(1674.79855322,5003.95774841)
\curveto(1678.4071991,5008.82528196)(1680.2115491,5015.45521918)(1680.21160864,5023.84757996)
\curveto(1680.2115491,5032.19790874)(1678.4071991,5038.80686514)(1674.79855322,5043.67446899)
\curveto(1671.18979909,5048.58392445)(1666.30126943,5051.03867969)(1660.13294958,5051.03874207)
\moveto(1660.13294958,5060.85777283)
\curveto(1670.20370083,5060.85770063)(1678.1134677,5057.58469365)(1683.86227393,5051.03874207)
\curveto(1689.61095376,5044.49266573)(1692.48532528,5035.42895409)(1692.48539709,5023.84757996)
\curveto(1692.48532528,5012.30809708)(1689.61095376,5003.24438544)(1683.86227393,4996.65641785)
\curveto(1678.1134677,4990.11039589)(1670.20370083,4986.83738891)(1660.13294958,4986.83738708)
\curveto(1650.02015778,4986.83738891)(1642.0894101,4990.11039589)(1636.34068274,4996.65641785)
\curveto(1630.63388566,5003.24438544)(1627.78049496,5012.30809708)(1627.78050208,5023.84757996)
\curveto(1627.78049496,5035.42895409)(1630.63388566,5044.49266573)(1636.34068274,5051.03874207)
\curveto(1642.0894101,5057.58469365)(1650.02015778,5060.85770063)(1660.13294958,5060.85777283)
}
}
{
\newrgbcolor{curcolor}{0 0 0}
\pscustom[linestyle=none,fillstyle=solid,fillcolor=curcolor]
{
\newpath
\moveto(1710.42401099,5016.48330688)
\lineto(1710.42401099,5059.1583252)
\lineto(1722.00543188,5059.1583252)
\lineto(1722.00543188,5016.92390442)
\curveto(1722.00540935,5010.25197731)(1723.30621982,5005.23756277)(1725.90786719,5001.88064575)
\curveto(1728.50946169,4998.56566393)(1732.41189309,4996.90817962)(1737.6151731,4996.90818787)
\curveto(1743.86741752,4996.90817962)(1748.7979088,4998.90135695)(1752.40666174,5002.88772583)
\curveto(1756.05727044,5006.87406626)(1757.88260126,5012.30809708)(1757.88265967,5019.18983459)
\lineto(1757.88265967,5059.1583252)
\lineto(1769.46408057,5059.1583252)
\lineto(1769.46408057,4988.66271973)
\lineto(1757.88265967,4988.66271973)
\lineto(1757.88265967,4999.48883057)
\curveto(1755.07117218,4995.20873369)(1751.7981652,4992.01964996)(1748.06362891,4989.92156982)
\curveto(1744.37095705,4987.8654488)(1740.06989019,4986.83738891)(1735.16041541,4986.83738708)
\curveto(1727.06178552,4986.83738891)(1720.91440702,4989.35508659)(1716.71826147,4994.39048767)
\curveto(1712.52208143,4999.4258773)(1710.42400003,5006.79014301)(1710.42401099,5016.48330688)
\moveto(1739.56639075,5060.85777283)
\lineto(1739.56639075,5060.85777283)
}
}
{
\newrgbcolor{curcolor}{0 0 0}
\pscustom[linestyle=none,fillstyle=solid,fillcolor=curcolor]
{
\newpath
\moveto(1834.2948606,5048.33221436)
\curveto(1832.99399713,5049.08746399)(1831.56730178,5049.63296515)(1830.01477026,5049.96871948)
\curveto(1828.50410294,5050.34631283)(1826.82563782,5050.53514015)(1824.97936987,5050.53520203)
\curveto(1818.43331223,5050.53514015)(1813.39791687,5048.39509713)(1809.8731687,5044.11506653)
\curveto(1806.390325,5039.87688665)(1804.64891744,5033.77146978)(1804.6489408,5025.79879761)
\lineto(1804.6489408,4988.66271973)
\lineto(1793.00457739,4988.66271973)
\lineto(1793.00457739,5059.1583252)
\lineto(1804.6489408,5059.1583252)
\lineto(1804.6489408,5048.20632935)
\curveto(1807.08269187,5052.48635585)(1810.25079478,5055.65445877)(1814.15325903,5057.71064758)
\curveto(1818.05565758,5059.80865993)(1822.79732154,5060.85770063)(1828.37826514,5060.85777283)
\curveto(1829.17548899,5060.85770063)(1830.05668317,5060.79475819)(1831.02185034,5060.66894531)
\curveto(1831.98691806,5060.58495005)(1833.05693957,5060.43808435)(1834.23191809,5060.22834778)
\lineto(1834.2948606,5048.33221436)
}
}
{
\newrgbcolor{curcolor}{0 0 0}
\pscustom[linestyle=none,fillstyle=solid,fillcolor=curcolor]
{
\newpath
\moveto(1894.53083777,5056.45179749)
\lineto(1894.53083777,5045.62568665)
\curveto(1891.25776791,5047.42997968)(1887.96378011,5048.77275178)(1884.6488645,5049.65400696)
\curveto(1881.37580452,5050.57710178)(1878.06083591,5051.03867969)(1874.70394873,5051.03874207)
\curveto(1867.19277427,5051.03867969)(1861.36010799,5048.6468669)(1857.20593237,5043.86329651)
\curveto(1853.05170565,5039.12157735)(1850.97460507,5032.4496785)(1850.97462439,5023.84757996)
\curveto(1850.97460507,5015.24541104)(1853.05170565,5008.55253138)(1857.20593237,5003.7689209)
\curveto(1861.36010799,4999.02724183)(1867.19277427,4996.65640985)(1874.70394873,4996.65641785)
\curveto(1878.06083591,4996.65640985)(1881.37580452,4997.09700695)(1884.6488645,4997.97821045)
\curveto(1887.96378011,4998.90135695)(1891.25776791,5000.26510986)(1894.53083777,5002.06947327)
\lineto(1894.53083777,4991.36924744)
\curveto(1891.29972954,4989.85862612)(1887.9427993,4988.72566217)(1884.46003699,4987.97035217)
\curveto(1881.01913069,4987.21504356)(1877.34748824,4986.83738891)(1873.44509863,4986.83738708)
\curveto(1862.82876497,4986.83738891)(1854.39447775,4990.17333833)(1848.14221167,4996.84524536)
\curveto(1841.88991261,5003.51713602)(1838.76377133,5012.51790522)(1838.76377844,5023.84757996)
\curveto(1838.76377133,5035.34503083)(1841.91089343,5044.38776166)(1848.20515417,5050.97579956)
\curveto(1854.54134344,5057.56371284)(1863.20641962,5060.85770063)(1874.20040869,5060.85777283)
\curveto(1877.76710452,5060.85770063)(1881.24991964,5060.48004598)(1884.6488645,5059.72480774)
\curveto(1888.04770337,5059.011389)(1891.34169116,5057.92038667)(1894.53083777,5056.45179749)
}
}
{
\newrgbcolor{curcolor}{0 0 0}
\pscustom[linestyle=none,fillstyle=solid,fillcolor=curcolor]
{
\newpath
\moveto(1974.845474,5026.80587769)
\lineto(1974.845474,5021.14105225)
\lineto(1921.59611487,5021.14105225)
\curveto(1922.09963521,5013.16831046)(1924.491448,5007.0838744)(1928.77156042,5002.88772583)
\curveto(1933.09358173,4998.73351044)(1939.09409453,4996.65640985)(1946.77311682,4996.65641785)
\curveto(1951.22100501,4996.65640985)(1955.52207188,4997.20191102)(1959.67633032,4998.29292297)
\curveto(1963.87243584,4999.38391567)(1968.02663701,5001.02041916)(1972.13894629,5003.20243835)
\lineto(1972.13894629,4992.2504425)
\curveto(1967.98467538,4990.48805054)(1963.72557014,4989.14527845)(1959.3616178,4988.22212219)
\curveto(1954.99755153,4987.29896682)(1950.57059978,4986.83738891)(1946.08074927,4986.83738708)
\curveto(1934.83498929,4986.83738891)(1925.91814335,4990.11039589)(1919.33018469,4996.65641785)
\curveto(1912.7841538,5003.20242381)(1909.51114682,5012.05632731)(1909.51115393,5023.21815491)
\curveto(1909.51114682,5034.75756804)(1912.61630729,5043.90520294)(1918.82664465,5050.66108704)
\curveto(1925.07891079,5057.45880877)(1933.4922172,5060.85770063)(1944.06658911,5060.85777283)
\curveto(1953.54987536,5060.85770063)(1961.04002595,5057.79450179)(1966.53706335,5051.66816711)
\curveto(1972.07593411,5045.58366805)(1974.84540155,5037.29624653)(1974.845474,5026.80587769)
\moveto(1963.2640531,5030.20477295)
\curveto(1963.18006898,5036.54093723)(1961.39669979,5041.5973134)(1957.91394019,5045.37391663)
\curveto(1954.47303118,5049.15040643)(1949.89921373,5051.03867969)(1944.19247412,5051.03874207)
\curveto(1937.73034162,5051.03867969)(1932.54808057,5049.21334887)(1928.64567542,5045.56274414)
\curveto(1924.7851794,5041.91202561)(1922.56121311,5036.77172618)(1921.9737699,5030.14183044)
\lineto(1963.2640531,5030.20477295)
}
}
{
\newrgbcolor{curcolor}{0 0 0}
\pscustom[linestyle=none,fillstyle=solid,fillcolor=curcolor]
{
\newpath
\moveto(2021.42292761,5086.47537231)
\curveto(2015.8000295,5076.82410007)(2011.62484752,5067.27782971)(2008.89736914,5057.83653259)
\curveto(2006.16983588,5048.39509713)(2004.80608297,5038.82784595)(2004.80610632,5029.13475037)
\curveto(2004.80608297,5019.44157384)(2006.16983588,5009.83236103)(2008.89736914,5000.30708313)
\curveto(2011.66680914,4990.82374357)(2015.84199113,4981.27747321)(2021.42292761,4971.66824341)
\lineto(2011.35212683,4971.66824341)
\curveto(2005.05785274,4981.52924297)(2000.33716959,4991.22237903)(1997.19006323,5000.74768066)
\curveto(1994.08488703,5010.27295813)(1992.53230679,5019.73530523)(1992.53231787,5029.13475037)
\curveto(1992.53230679,5038.49215293)(1994.08488703,5047.91253841)(1997.19006323,5057.39593506)
\curveto(2000.29520797,5066.87919424)(2005.01589111,5076.5723303)(2011.35212683,5086.47537231)
\lineto(2021.42292761,5086.47537231)
}
}
{
\newrgbcolor{curcolor}{0 0 0}
\pscustom[linestyle=none,fillstyle=solid,fillcolor=curcolor]
{
\newpath
\moveto(2042.13101172,5086.47537231)
\lineto(2052.2018125,5086.47537231)
\curveto(2058.4960363,5076.5723303)(2063.19573863,5066.87919424)(2066.30093359,5057.39593506)
\curveto(2069.4480212,5047.91253841)(2071.02158225,5038.49215293)(2071.02162146,5029.13475037)
\curveto(2071.02158225,5019.73530523)(2069.4480212,5010.27295813)(2066.30093359,5000.74768066)
\curveto(2063.19573863,4991.22237903)(2058.4960363,4981.52924297)(2052.2018125,4971.66824341)
\lineto(2042.13101172,4971.66824341)
\curveto(2047.71189791,4981.27747321)(2051.86609908,4990.82374357)(2054.59362769,5000.30708313)
\curveto(2057.36307235,5009.83236103)(2058.74780607,5019.44157384)(2058.74783301,5029.13475037)
\curveto(2058.74780607,5038.82784595)(2057.36307235,5048.39509713)(2054.59362769,5057.83653259)
\curveto(2051.86609908,5067.27782971)(2047.71189791,5076.82410007)(2042.13101172,5086.47537231)
}
}
{
\newrgbcolor{curcolor}{0 0 0}
\pscustom[linestyle=none,fillstyle=solid,fillcolor=curcolor]
{
}
}
{
\newrgbcolor{curcolor}{0 0 0}
\pscustom[linestyle=none,fillstyle=solid,fillcolor=curcolor]
{
\newpath
\moveto(2138.30715918,5004.65011597)
\lineto(2151.58802771,5004.65011597)
\lineto(2151.58802771,4988.66271973)
\lineto(2138.30715918,4988.66271973)
\lineto(2138.30715918,5004.65011597)
\moveto(2138.30715918,5055.3188324)
\lineto(2151.58802771,5055.3188324)
\lineto(2151.58802771,5039.33143616)
\lineto(2138.30715918,5039.33143616)
\lineto(2138.30715918,5055.3188324)
}
}
{
\newrgbcolor{curcolor}{0 0 0}
\pscustom[linestyle=none,fillstyle=solid,fillcolor=curcolor]
{
}
}
{
\newrgbcolor{curcolor}{0 0 0}
\pscustom[linestyle=none,fillstyle=solid,fillcolor=curcolor]
{
\newpath
\moveto(2276.78066992,5079.55169678)
\lineto(2276.78066992,5067.15202332)
\curveto(2271.95501372,5069.45983436)(2267.40217709,5071.18026111)(2263.12214636,5072.31330872)
\curveto(2258.84200498,5073.44618902)(2254.70878463,5074.012671)(2250.7224729,5074.01275635)
\curveto(2243.79876136,5074.012671)(2238.4486538,5072.6698989)(2234.67213416,5069.98443604)
\curveto(2230.93752239,5067.29881052)(2229.07022995,5063.48030238)(2229.07025122,5058.52890015)
\curveto(2229.07022995,5054.37462911)(2230.30809797,5051.22750702)(2232.78385901,5049.08752441)
\curveto(2235.3015317,5046.98938259)(2240.04319566,5045.28993666)(2247.00886511,5043.98918152)
\lineto(2254.68785071,5042.4156189)
\curveto(2264.17113173,5040.61121514)(2271.15774279,5037.42213142)(2275.64770483,5032.84835815)
\curveto(2280.1794928,5028.31645815)(2282.44542071,5022.23202209)(2282.44549536,5014.59503174)
\curveto(2282.44542071,5005.48933254)(2279.38222187,4998.58664474)(2273.25588965,4993.88694763)
\curveto(2267.17138813,4989.18724008)(2258.23356138,4986.83738891)(2246.44238257,4986.83738708)
\curveto(2241.99441136,4986.83738891)(2237.2527474,4987.34092845)(2232.21737646,4988.3480072)
\curveto(2227.22391832,4989.35508659)(2222.04165726,4990.84472438)(2216.67057776,4992.81692505)
\lineto(2216.67057776,5005.90896606)
\curveto(2221.83184912,5003.01359649)(2226.88822529,5000.83159184)(2231.83972144,4999.36294556)
\curveto(2236.79116949,4997.89427788)(2241.65871833,4997.15994939)(2246.44238257,4997.15995789)
\curveto(2253.70170556,4997.15994939)(2259.30358289,4998.58664474)(2263.24803137,5001.44004822)
\curveto(2267.19236895,5004.29342614)(2269.16456546,5008.36370405)(2269.16462683,5013.65089417)
\curveto(2269.16456546,5018.26664825)(2267.73787011,5021.87534826)(2264.8845365,5024.477005)
\curveto(2262.07305034,5027.07859012)(2257.43629045,5029.02980582)(2250.97424292,5030.33065796)
\lineto(2243.23231482,5031.84127808)
\curveto(2233.74895146,5033.72950816)(2226.88822529,5036.68780293)(2222.65011572,5040.71617126)
\curveto(2218.41197644,5044.7444355)(2216.29291423,5050.34631283)(2216.29292273,5057.52182007)
\curveto(2216.29291423,5065.83015355)(2219.20924738,5072.37616751)(2225.04193091,5077.15988159)
\curveto(2230.91654158,5081.94341868)(2238.99415496,5084.33523148)(2249.27479529,5084.33532715)
\curveto(2253.68072474,5084.33523148)(2258.17061894,5083.93659601)(2262.74449133,5083.13941956)
\curveto(2267.31825383,5082.34205415)(2271.99697535,5081.14614775)(2276.78066992,5079.55169678)
}
}
{
\newrgbcolor{curcolor}{0 0 0}
\pscustom[linestyle=none,fillstyle=solid,fillcolor=curcolor]
{
\newpath
\moveto(2313.47615027,5079.17404175)
\lineto(2313.47615027,5059.1583252)
\lineto(2337.33135962,5059.1583252)
\lineto(2337.33135962,5050.157547)
\lineto(2313.47615027,5050.157547)
\lineto(2313.47615027,5011.88850403)
\curveto(2313.47612667,5006.13973777)(2314.25241678,5002.44711451)(2315.80502295,5000.81062317)
\curveto(2317.39953888,4999.17410753)(2320.60960342,4998.35585579)(2325.4352262,4998.35586548)
\lineto(2337.33135962,4998.35586548)
\lineto(2337.33135962,4988.66271973)
\lineto(2325.4352262,4988.66271973)
\curveto(2316.49736388,4988.66271973)(2310.32900457,4990.32020403)(2306.93012976,4993.63517761)
\curveto(2303.53122084,4996.99210288)(2301.83177491,5003.07653893)(2301.83178687,5011.88850403)
\lineto(2301.83178687,5050.157547)
\lineto(2293.33454871,5050.157547)
\lineto(2293.33454871,5059.1583252)
\lineto(2301.83178687,5059.1583252)
\lineto(2301.83178687,5079.17404175)
\lineto(2313.47615027,5079.17404175)
}
}
{
\newrgbcolor{curcolor}{0 0 0}
\pscustom[linestyle=none,fillstyle=solid,fillcolor=curcolor]
{
\newpath
\moveto(2393.47607397,5048.33221436)
\curveto(2392.17521051,5049.08746399)(2390.74851516,5049.63296515)(2389.19598364,5049.96871948)
\curveto(2387.68531632,5050.34631283)(2386.0068512,5050.53514015)(2384.16058325,5050.53520203)
\curveto(2377.61452561,5050.53514015)(2372.57913025,5048.39509713)(2369.05438208,5044.11506653)
\curveto(2365.57153838,5039.87688665)(2363.83013082,5033.77146978)(2363.83015417,5025.79879761)
\lineto(2363.83015417,4988.66271973)
\lineto(2352.18579077,4988.66271973)
\lineto(2352.18579077,5059.1583252)
\lineto(2363.83015417,5059.1583252)
\lineto(2363.83015417,5048.20632935)
\curveto(2366.26390524,5052.48635585)(2369.43200816,5055.65445877)(2373.33447241,5057.71064758)
\curveto(2377.23687096,5059.80865993)(2381.97853492,5060.85770063)(2387.55947852,5060.85777283)
\curveto(2388.35670237,5060.85770063)(2389.23789655,5060.79475819)(2390.20306372,5060.66894531)
\curveto(2391.16813144,5060.58495005)(2392.23815295,5060.43808435)(2393.41313147,5060.22834778)
\lineto(2393.47607397,5048.33221436)
}
}
{
\newrgbcolor{curcolor}{0 0 0}
\pscustom[linestyle=none,fillstyle=solid,fillcolor=curcolor]
{
\newpath
\moveto(2405.74986243,5059.1583252)
\lineto(2417.33128333,5059.1583252)
\lineto(2417.33128333,4988.66271973)
\lineto(2405.74986243,4988.66271973)
\lineto(2405.74986243,5059.1583252)
\moveto(2405.74986243,5086.60125732)
\lineto(2417.33128333,5086.60125732)
\lineto(2417.33128333,5071.93565369)
\lineto(2405.74986243,5071.93565369)
\lineto(2405.74986243,5086.60125732)
}
}
{
\newrgbcolor{curcolor}{0 0 0}
\pscustom[linestyle=none,fillstyle=solid,fillcolor=curcolor]
{
\newpath
\moveto(2500.10067725,5031.21185303)
\lineto(2500.10067725,4988.66271973)
\lineto(2488.51925635,4988.66271973)
\lineto(2488.51925635,5030.834198)
\curveto(2488.51919718,5037.50605467)(2487.21838671,5042.4994884)(2484.61682104,5045.81451416)
\curveto(2482.01514485,5049.12942562)(2478.11271345,5050.78690992)(2472.90951514,5050.78697205)
\curveto(2466.65718901,5050.78690992)(2461.72669773,5048.79373259)(2458.11802649,5044.80743408)
\curveto(2454.50929772,5040.82102328)(2452.70494772,5035.38699246)(2452.70497107,5028.50532532)
\lineto(2452.70497107,4988.66271973)
\lineto(2441.06060767,4988.66271973)
\lineto(2441.06060767,5059.1583252)
\lineto(2452.70497107,5059.1583252)
\lineto(2452.70497107,5048.20632935)
\curveto(2455.47441516,5052.44439423)(2458.72644133,5055.61249714)(2462.46105933,5057.71064758)
\curveto(2466.23757273,5059.80865993)(2470.58060123,5060.85770063)(2475.49015784,5060.85777283)
\curveto(2483.5887059,5060.85770063)(2489.71510358,5058.34000295)(2493.86936926,5053.30467224)
\curveto(2498.02350591,5048.31117387)(2500.1006065,5040.94690816)(2500.10067725,5031.21185303)
}
}
{
\newrgbcolor{curcolor}{0 0 0}
\pscustom[linestyle=none,fillstyle=solid,fillcolor=curcolor]
{
\newpath
\moveto(2569.71508765,5024.72877502)
\curveto(2569.71502911,5033.12106455)(2567.97362155,5039.62511688)(2564.49085974,5044.24095154)
\curveto(2561.04995294,5048.85667504)(2556.20338491,5051.16456457)(2549.95114111,5051.16462708)
\curveto(2543.7407814,5051.16456457)(2538.89421337,5048.85667504)(2535.41142249,5044.24095154)
\curveto(2531.97054476,5039.62511688)(2530.25011801,5033.12106455)(2530.25013708,5024.72877502)
\curveto(2530.25011801,5016.37837499)(2531.97054476,5009.89530347)(2535.41142249,5005.27954102)
\curveto(2538.89421337,5000.66374532)(2543.7407814,4998.35585579)(2549.95114111,4998.35586548)
\curveto(2556.20338491,4998.35585579)(2561.04995294,5000.66374532)(2564.49085974,5005.27954102)
\curveto(2567.97362155,5009.89530347)(2569.71502911,5016.37837499)(2569.71508765,5024.72877502)
\moveto(2581.29650854,4997.41172791)
\curveto(2581.29643843,4985.41069356)(2578.63187505,4976.49384762)(2573.30281042,4970.66116333)
\curveto(2567.97362155,4964.78655342)(2559.81208491,4961.84923946)(2548.81817603,4961.84921265)
\curveto(2544.74786047,4961.84923946)(2540.90837152,4962.16395167)(2537.29969763,4962.79335022)
\curveto(2533.69097151,4963.38083888)(2530.18717557,4964.3039947)(2526.78829932,4965.56282043)
\lineto(2526.78829932,4976.82952881)
\curveto(2530.18717557,4974.98322901)(2533.54410581,4973.6194761)(2536.8591001,4972.73826599)
\curveto(2540.17404303,4971.85708773)(2543.55195408,4971.41649063)(2546.99284338,4971.41647339)
\curveto(2554.58786223,4971.41649063)(2560.27366282,4973.40966796)(2564.05026221,4977.39601135)
\curveto(2567.82675585,4981.34041565)(2569.71502911,4987.31994763)(2569.71508765,4995.33462524)
\lineto(2569.71508765,5001.06239319)
\curveto(2567.32321632,4996.90817962)(2564.26001748,4993.80301915)(2560.52548193,4991.74690247)
\curveto(2556.7908477,4989.69077961)(2552.32193432,4988.66271973)(2547.11872839,4988.66271973)
\curveto(2538.47459709,4988.66271973)(2531.50896685,4991.95670752)(2526.22181677,4998.54469299)
\curveto(2520.93463661,5005.1326587)(2518.29105404,5013.86067732)(2518.29106116,5024.72877502)
\curveto(2518.29105404,5035.63876223)(2520.93463661,5044.38776166)(2526.22181677,5050.97579956)
\curveto(2531.50896685,5057.56371284)(2538.47459709,5060.85770063)(2547.11872839,5060.85777283)
\curveto(2552.32193432,5060.85770063)(2556.7908477,5059.82964075)(2560.52548193,5057.77359009)
\curveto(2564.26001748,5055.71740121)(2567.32321632,5052.61224074)(2569.71508765,5048.45809937)
\lineto(2569.71508765,5059.1583252)
\lineto(2581.29650854,5059.1583252)
\lineto(2581.29650854,4997.41172791)
}
}
{
\newrgbcolor{curcolor}{0.51764709 0.74509805 0.51764709}
\pscustom[linestyle=none,fillstyle=solid,fillcolor=curcolor]
{
\newpath
\moveto(1256.79159375,4868.58282617)
\curveto(1256.79159375,4849.16656544)(1241.05160449,4833.42657617)(1221.63534375,4833.42657617)
\curveto(1202.21908301,4833.42657617)(1186.47909375,4849.16656544)(1186.47909375,4868.58282617)
\curveto(1186.47909375,4887.99908691)(1202.21908301,4903.73907617)(1221.63534375,4903.73907617)
\curveto(1241.05160449,4903.73907617)(1256.79159375,4887.99908691)(1256.79159375,4868.58282617)
\closepath
}
}
{
\newrgbcolor{curcolor}{0.01176471 0.50196081 0.28235295}
\pscustom[linewidth=11.71875,linecolor=curcolor]
{
\newpath
\moveto(1256.79159375,4868.58282617)
\curveto(1256.79159375,4849.16656544)(1241.05160449,4833.42657617)(1221.63534375,4833.42657617)
\curveto(1202.21908301,4833.42657617)(1186.47909375,4849.16656544)(1186.47909375,4868.58282617)
\curveto(1186.47909375,4887.99908691)(1202.21908301,4903.73907617)(1221.63534375,4903.73907617)
\curveto(1241.05160449,4903.73907617)(1256.79159375,4887.99908691)(1256.79159375,4868.58282617)
\closepath
}
}
{
\newrgbcolor{curcolor}{0 0 0}
\pscustom[linestyle=none,fillstyle=solid,fillcolor=curcolor]
{
\newpath
\moveto(1385.64068579,4874.73670959)
\curveto(1385.64062725,4883.12899912)(1383.89921969,4889.63305145)(1380.41645789,4894.24888611)
\curveto(1376.97555108,4898.86460961)(1372.12898305,4901.17249914)(1365.87673926,4901.17256165)
\curveto(1359.66637955,4901.17249914)(1354.81981152,4898.86460961)(1351.33702063,4894.24888611)
\curveto(1347.8961429,4889.63305145)(1346.17571616,4883.12899912)(1346.17573523,4874.73670959)
\curveto(1346.17571616,4866.38630956)(1347.8961429,4859.90323804)(1351.33702063,4855.28747559)
\curveto(1354.81981152,4850.67167989)(1359.66637955,4848.36379036)(1365.87673926,4848.36380005)
\curveto(1372.12898305,4848.36379036)(1376.97555108,4850.67167989)(1380.41645789,4855.28747559)
\curveto(1383.89921969,4859.90323804)(1385.64062725,4866.38630956)(1385.64068579,4874.73670959)
\moveto(1397.22210669,4847.41966248)
\curveto(1397.22203657,4835.41862813)(1394.5574732,4826.50178219)(1389.22840857,4820.6690979)
\curveto(1383.89921969,4814.79448799)(1375.73768306,4811.85717403)(1364.74377417,4811.85714722)
\curveto(1360.67345862,4811.85717403)(1356.83396966,4812.17188624)(1353.22529578,4812.80128479)
\curveto(1349.61656965,4813.38877345)(1346.11277372,4814.31192927)(1342.71389746,4815.570755)
\lineto(1342.71389746,4826.83746338)
\curveto(1346.11277372,4824.99116358)(1349.46970395,4823.62741067)(1352.78469824,4822.74620056)
\curveto(1356.09964117,4821.8650223)(1359.47755222,4821.42442521)(1362.91844153,4821.42440796)
\curveto(1370.51346037,4821.42442521)(1376.19926096,4823.41760253)(1379.97586035,4827.40394592)
\curveto(1383.752354,4831.34835022)(1385.64062725,4837.3278822)(1385.64068579,4845.34255981)
\lineto(1385.64068579,4851.07032776)
\curveto(1383.24881446,4846.91611419)(1380.18561562,4843.81095372)(1376.45108008,4841.75483704)
\curveto(1372.71644584,4839.69871418)(1368.24753246,4838.6706543)(1363.04432654,4838.6706543)
\curveto(1354.40019524,4838.6706543)(1347.434565,4841.96464209)(1342.14741492,4848.55262756)
\curveto(1336.86023475,4855.14059327)(1334.21665219,4863.86861189)(1334.2166593,4874.73670959)
\curveto(1334.21665219,4885.6466968)(1336.86023475,4894.39569623)(1342.14741492,4900.98373413)
\curveto(1347.434565,4907.57164741)(1354.40019524,4910.8656352)(1363.04432654,4910.8657074)
\curveto(1368.24753246,4910.8656352)(1372.71644584,4909.83757532)(1376.45108008,4907.78152466)
\curveto(1380.18561562,4905.72533578)(1383.24881446,4902.62017531)(1385.64068579,4898.46603394)
\lineto(1385.64068579,4909.16625977)
\lineto(1397.22210669,4909.16625977)
\lineto(1397.22210669,4847.41966248)
}
}
{
\newrgbcolor{curcolor}{0 0 0}
\pscustom[linestyle=none,fillstyle=solid,fillcolor=curcolor]
{
\newpath
\moveto(1481.37623572,4876.81381226)
\lineto(1481.37623572,4871.14898682)
\lineto(1428.12687659,4871.14898682)
\curveto(1428.63039692,4863.17624503)(1431.02220972,4857.09180897)(1435.30232214,4852.8956604)
\curveto(1439.62434345,4848.74144501)(1445.62485625,4846.66434442)(1453.30387854,4846.66435242)
\curveto(1457.75176673,4846.66434442)(1462.0528336,4847.20984559)(1466.20709204,4848.30085754)
\curveto(1470.40319756,4849.39185024)(1474.55739873,4851.02835373)(1478.66970801,4853.21037292)
\lineto(1478.66970801,4842.25837708)
\curveto(1474.5154371,4840.49598511)(1470.25633186,4839.15321302)(1465.89237952,4838.23005676)
\curveto(1461.52831325,4837.30690139)(1457.1013615,4836.84532348)(1452.61151099,4836.84532166)
\curveto(1441.36575101,4836.84532348)(1432.44890507,4840.11833046)(1425.86094641,4846.66435242)
\curveto(1419.31491552,4853.21035839)(1416.04190854,4862.06426188)(1416.04191565,4873.22608948)
\curveto(1416.04190854,4884.76550261)(1419.14706901,4893.91313751)(1425.35740637,4900.66902161)
\curveto(1431.60967251,4907.46674334)(1440.02297892,4910.8656352)(1450.59735083,4910.8657074)
\curveto(1460.08063708,4910.8656352)(1467.57078767,4907.80243636)(1473.06782507,4901.67610168)
\curveto(1478.60669583,4895.59160262)(1481.37616327,4887.3041811)(1481.37623572,4876.81381226)
\moveto(1469.79481482,4880.21270752)
\curveto(1469.7108307,4886.5488718)(1467.92746151,4891.60524797)(1464.4447019,4895.3818512)
\curveto(1461.0037929,4899.158341)(1456.42997545,4901.04661426)(1450.72323584,4901.04667664)
\curveto(1444.26110334,4901.04661426)(1439.07884229,4899.22128344)(1435.17643713,4895.57067871)
\curveto(1431.31594111,4891.91996018)(1429.09197483,4886.77966075)(1428.50453162,4880.14976501)
\lineto(1469.79481482,4880.21270752)
}
}
{
\newrgbcolor{curcolor}{0 0 0}
\pscustom[linestyle=none,fillstyle=solid,fillcolor=curcolor]
{
\newpath
\moveto(1511.58863806,4929.18197632)
\lineto(1511.58863806,4909.16625977)
\lineto(1535.44384741,4909.16625977)
\lineto(1535.44384741,4900.16548157)
\lineto(1511.58863806,4900.16548157)
\lineto(1511.58863806,4861.8964386)
\curveto(1511.58861446,4856.14767234)(1512.36490458,4852.45504908)(1513.91751074,4850.81855774)
\curveto(1515.51202667,4849.1820421)(1518.72209121,4848.36379036)(1523.54771399,4848.36380005)
\lineto(1535.44384741,4848.36380005)
\lineto(1535.44384741,4838.6706543)
\lineto(1523.54771399,4838.6706543)
\curveto(1514.60985167,4838.6706543)(1508.44149236,4840.3281386)(1505.04261755,4843.64311218)
\curveto(1501.64370863,4847.00003745)(1499.9442627,4853.0844735)(1499.94427466,4861.8964386)
\lineto(1499.94427466,4900.16548157)
\lineto(1491.4470365,4900.16548157)
\lineto(1491.4470365,4909.16625977)
\lineto(1499.94427466,4909.16625977)
\lineto(1499.94427466,4929.18197632)
\lineto(1511.58863806,4929.18197632)
}
}
{
\newrgbcolor{curcolor}{0 0 0}
\pscustom[linestyle=none,fillstyle=solid,fillcolor=curcolor]
{
\newpath
\moveto(1563.95680212,4922.19535828)
\lineto(1563.95680212,4849.11911011)
\lineto(1579.31477332,4849.11911011)
\curveto(1592.28087563,4849.11909966)(1601.76420355,4852.05641362)(1607.76478552,4857.93106079)
\curveto(1613.80719077,4863.80566945)(1616.82842799,4873.07918922)(1616.82850623,4885.75164795)
\curveto(1616.82842799,4898.34008926)(1613.80719077,4907.55066659)(1607.76478552,4913.38340759)
\curveto(1601.76420355,4919.25796079)(1592.28087563,4922.19527475)(1579.31477332,4922.19535828)
\lineto(1563.95680212,4922.19535828)
\moveto(1551.24241614,4932.64381409)
\lineto(1577.36355566,4932.64381409)
\curveto(1595.57486343,4932.64372011)(1608.93964193,4928.84619278)(1617.45793127,4921.2512207)
\curveto(1625.97606288,4913.69804509)(1630.23516812,4901.864866)(1630.23525977,4885.75164795)
\curveto(1630.23516812,4869.55441248)(1625.95508207,4857.65829095)(1617.39498877,4850.06324768)
\curveto(1608.83473786,4842.46818163)(1595.49094017,4838.6706543)(1577.36355566,4838.6706543)
\lineto(1551.24241614,4838.6706543)
\lineto(1551.24241614,4932.64381409)
}
}
{
\newrgbcolor{curcolor}{0 0 0}
\pscustom[linestyle=none,fillstyle=solid,fillcolor=curcolor]
{
\newpath
\moveto(1681.97399878,4874.10728455)
\curveto(1672.61651156,4874.10724911)(1666.13344004,4873.0372276)(1662.52476477,4870.8972168)
\curveto(1658.91604003,4868.75714154)(1657.11169003,4865.10647991)(1657.11170935,4859.94522095)
\curveto(1657.11169003,4855.83296013)(1658.45446212,4852.55995315)(1661.14002966,4850.12619019)
\curveto(1663.86751213,4847.73436594)(1667.56013539,4846.53845954)(1672.21791052,4846.53846741)
\curveto(1678.63800517,4846.53845954)(1683.7783046,4848.80438745)(1687.63882422,4853.33625793)
\curveto(1691.54120577,4857.91006072)(1693.49242147,4863.97351596)(1693.49247717,4871.52664185)
\lineto(1693.49247717,4874.10728455)
\lineto(1681.97399878,4874.10728455)
\moveto(1705.07389807,4878.89091492)
\lineto(1705.07389807,4838.6706543)
\lineto(1693.49247717,4838.6706543)
\lineto(1693.49247717,4849.37088013)
\curveto(1690.84883891,4845.09078337)(1687.55485111,4841.92268046)(1683.61050391,4839.86656189)
\curveto(1679.66606506,4837.85240255)(1674.84047784,4836.84532348)(1669.13372778,4836.84532166)
\curveto(1661.91629643,4836.84532348)(1656.1675534,4838.85948162)(1651.88748145,4842.88780212)
\curveto(1647.64934292,4846.95807582)(1645.53028071,4852.39210664)(1645.53028845,4859.18991089)
\curveto(1645.53028071,4867.12063805)(1648.17386327,4873.10017004)(1653.46104407,4877.12852478)
\curveto(1658.79015515,4881.15680261)(1666.72090283,4883.17096075)(1677.25331091,4883.17100525)
\lineto(1693.49247717,4883.17100525)
\lineto(1693.49247717,4884.30397034)
\curveto(1693.49242147,4889.63305145)(1691.73003309,4893.74529099)(1688.20530676,4896.64070129)
\curveto(1684.72244123,4899.57795728)(1679.81293075,4901.04661426)(1673.47676062,4901.04667664)
\curveto(1669.44840865,4901.04661426)(1665.52499643,4900.56405554)(1661.70651221,4899.59899902)
\curveto(1657.88798014,4898.63382065)(1654.2163377,4897.18614449)(1650.69157385,4895.25596619)
\lineto(1650.69157385,4905.95619202)
\curveto(1654.92968537,4907.59262822)(1659.04192491,4908.80951543)(1663.02830481,4909.6068573)
\curveto(1667.01463423,4910.44601892)(1670.89608481,4910.8656352)(1674.67266821,4910.8657074)
\curveto(1684.86930692,4910.8656352)(1692.4853424,4908.22205264)(1697.52079749,4902.93495178)
\curveto(1702.55613311,4897.64772239)(1705.07383079,4889.63305145)(1705.07389807,4878.89091492)
}
}
{
\newrgbcolor{curcolor}{0 0 0}
\pscustom[linestyle=none,fillstyle=solid,fillcolor=curcolor]
{
\newpath
\moveto(1740.69935583,4929.18197632)
\lineto(1740.69935583,4909.16625977)
\lineto(1764.55456519,4909.16625977)
\lineto(1764.55456519,4900.16548157)
\lineto(1740.69935583,4900.16548157)
\lineto(1740.69935583,4861.8964386)
\curveto(1740.69933223,4856.14767234)(1741.47562235,4852.45504908)(1743.02822852,4850.81855774)
\curveto(1744.62274445,4849.1820421)(1747.83280898,4848.36379036)(1752.65843176,4848.36380005)
\lineto(1764.55456519,4848.36380005)
\lineto(1764.55456519,4838.6706543)
\lineto(1752.65843176,4838.6706543)
\curveto(1743.72056944,4838.6706543)(1737.55221013,4840.3281386)(1734.15333533,4843.64311218)
\curveto(1730.7544264,4847.00003745)(1729.05498047,4853.0844735)(1729.05499243,4861.8964386)
\lineto(1729.05499243,4900.16548157)
\lineto(1720.55775427,4900.16548157)
\lineto(1720.55775427,4909.16625977)
\lineto(1729.05499243,4909.16625977)
\lineto(1729.05499243,4929.18197632)
\lineto(1740.69935583,4929.18197632)
}
}
{
\newrgbcolor{curcolor}{0 0 0}
\pscustom[linestyle=none,fillstyle=solid,fillcolor=curcolor]
{
\newpath
\moveto(1840.14851355,4876.81381226)
\lineto(1840.14851355,4871.14898682)
\lineto(1786.89915442,4871.14898682)
\curveto(1787.40267476,4863.17624503)(1789.79448755,4857.09180897)(1794.07459998,4852.8956604)
\curveto(1798.39662128,4848.74144501)(1804.39713408,4846.66434442)(1812.07615637,4846.66435242)
\curveto(1816.52404456,4846.66434442)(1820.82511143,4847.20984559)(1824.97936987,4848.30085754)
\curveto(1829.17547539,4849.39185024)(1833.32967656,4851.02835373)(1837.44198584,4853.21037292)
\lineto(1837.44198584,4842.25837708)
\curveto(1833.28771493,4840.49598511)(1829.02860969,4839.15321302)(1824.66465735,4838.23005676)
\curveto(1820.30059108,4837.30690139)(1815.87363933,4836.84532348)(1811.38378882,4836.84532166)
\curveto(1800.13802884,4836.84532348)(1791.2211829,4840.11833046)(1784.63322424,4846.66435242)
\curveto(1778.08719335,4853.21035839)(1774.81418637,4862.06426188)(1774.81419348,4873.22608948)
\curveto(1774.81418637,4884.76550261)(1777.91934684,4893.91313751)(1784.1296842,4900.66902161)
\curveto(1790.38195034,4907.46674334)(1798.79525675,4910.8656352)(1809.36962866,4910.8657074)
\curveto(1818.85291491,4910.8656352)(1826.3430655,4907.80243636)(1831.84010291,4901.67610168)
\curveto(1837.37897366,4895.59160262)(1840.1484411,4887.3041811)(1840.14851355,4876.81381226)
\moveto(1828.56709265,4880.21270752)
\curveto(1828.48310853,4886.5488718)(1826.69973934,4891.60524797)(1823.21697974,4895.3818512)
\curveto(1819.77607073,4899.158341)(1815.20225328,4901.04661426)(1809.49551367,4901.04667664)
\curveto(1803.03338117,4901.04661426)(1797.85112012,4899.22128344)(1793.94871497,4895.57067871)
\curveto(1790.08821895,4891.91996018)(1787.86425266,4886.77966075)(1787.27680945,4880.14976501)
\lineto(1828.56709265,4880.21270752)
}
}
{
\newrgbcolor{curcolor}{0 0 0}
\pscustom[linestyle=none,fillstyle=solid,fillcolor=curcolor]
{
\newpath
\moveto(1886.72596716,4936.48330688)
\curveto(1881.10306905,4926.83203464)(1876.92788707,4917.28576428)(1874.20040869,4907.84446716)
\curveto(1871.47287543,4898.4030317)(1870.10912252,4888.83578052)(1870.10914587,4879.14268494)
\curveto(1870.10912252,4869.44950841)(1871.47287543,4859.8402956)(1874.20040869,4850.3150177)
\curveto(1876.96984869,4840.83167814)(1881.14503068,4831.28540778)(1886.72596716,4821.67617798)
\lineto(1876.65516638,4821.67617798)
\curveto(1870.36089229,4831.53717754)(1865.64020914,4841.2303136)(1862.49310278,4850.75561523)
\curveto(1859.38792658,4860.2808927)(1857.83534634,4869.7432398)(1857.83535742,4879.14268494)
\curveto(1857.83534634,4888.5000875)(1859.38792658,4897.92047298)(1862.49310278,4907.40386963)
\curveto(1865.59824752,4916.88712881)(1870.31893066,4926.58026487)(1876.65516638,4936.48330688)
\lineto(1886.72596716,4936.48330688)
}
}
{
\newrgbcolor{curcolor}{0 0 0}
\pscustom[linestyle=none,fillstyle=solid,fillcolor=curcolor]
{
\newpath
\moveto(1907.43405127,4936.48330688)
\lineto(1917.50485205,4936.48330688)
\curveto(1923.79907585,4926.58026487)(1928.49877818,4916.88712881)(1931.60397314,4907.40386963)
\curveto(1934.75106075,4897.92047298)(1936.3246218,4888.5000875)(1936.32466101,4879.14268494)
\curveto(1936.3246218,4869.7432398)(1934.75106075,4860.2808927)(1931.60397314,4850.75561523)
\curveto(1928.49877818,4841.2303136)(1923.79907585,4831.53717754)(1917.50485205,4821.67617798)
\lineto(1907.43405127,4821.67617798)
\curveto(1913.01493747,4831.28540778)(1917.16913863,4840.83167814)(1919.89666724,4850.3150177)
\curveto(1922.6661119,4859.8402956)(1924.05084562,4869.44950841)(1924.05087256,4879.14268494)
\curveto(1924.05084562,4888.83578052)(1922.6661119,4898.4030317)(1919.89666724,4907.84446716)
\curveto(1917.16913863,4917.28576428)(1913.01493747,4926.83203464)(1907.43405127,4936.48330688)
}
}
{
\newrgbcolor{curcolor}{0 0 0}
\pscustom[linestyle=none,fillstyle=solid,fillcolor=curcolor]
{
}
}
{
\newrgbcolor{curcolor}{0 0 0}
\pscustom[linestyle=none,fillstyle=solid,fillcolor=curcolor]
{
\newpath
\moveto(2003.61019873,4854.65805054)
\lineto(2016.89106726,4854.65805054)
\lineto(2016.89106726,4838.6706543)
\lineto(2003.61019873,4838.6706543)
\lineto(2003.61019873,4854.65805054)
\moveto(2003.61019873,4905.32676697)
\lineto(2016.89106726,4905.32676697)
\lineto(2016.89106726,4889.33937073)
\lineto(2003.61019873,4889.33937073)
\lineto(2003.61019873,4905.32676697)
}
}
{
\newrgbcolor{curcolor}{0 0 0}
\pscustom[linestyle=none,fillstyle=solid,fillcolor=curcolor]
{
}
}
{
\newrgbcolor{curcolor}{0 0 0}
\pscustom[linestyle=none,fillstyle=solid,fillcolor=curcolor]
{
\newpath
\moveto(2098.46455359,4922.19535828)
\lineto(2098.46455359,4849.11911011)
\lineto(2113.82252478,4849.11911011)
\curveto(2126.7886271,4849.11909966)(2136.27195502,4852.05641362)(2142.27253699,4857.93106079)
\curveto(2148.31494224,4863.80566945)(2151.33617945,4873.07918922)(2151.33625769,4885.75164795)
\curveto(2151.33617945,4898.34008926)(2148.31494224,4907.55066659)(2142.27253699,4913.38340759)
\curveto(2136.27195502,4919.25796079)(2126.7886271,4922.19527475)(2113.82252478,4922.19535828)
\lineto(2098.46455359,4922.19535828)
\moveto(2085.7501676,4932.64381409)
\lineto(2111.87130713,4932.64381409)
\curveto(2130.08261489,4932.64372011)(2143.4473934,4928.84619278)(2151.96568274,4921.2512207)
\curveto(2160.48381435,4913.69804509)(2164.74291959,4901.864866)(2164.74301123,4885.75164795)
\curveto(2164.74291959,4869.55441248)(2160.46283353,4857.65829095)(2151.90274023,4850.06324768)
\curveto(2143.34248933,4842.46818163)(2129.99869164,4838.6706543)(2111.87130713,4838.6706543)
\lineto(2085.7501676,4838.6706543)
\lineto(2085.7501676,4932.64381409)
}
}
{
\newrgbcolor{curcolor}{0 0 0}
\pscustom[linestyle=none,fillstyle=solid,fillcolor=curcolor]
{
\newpath
\moveto(2216.48175024,4874.10728455)
\curveto(2207.12426302,4874.10724911)(2200.6411915,4873.0372276)(2197.03251624,4870.8972168)
\curveto(2193.42379149,4868.75714154)(2191.61944149,4865.10647991)(2191.61946082,4859.94522095)
\curveto(2191.61944149,4855.83296013)(2192.96221359,4852.55995315)(2195.64778113,4850.12619019)
\curveto(2198.37526359,4847.73436594)(2202.06788685,4846.53845954)(2206.72566199,4846.53846741)
\curveto(2213.14575664,4846.53845954)(2218.28605606,4848.80438745)(2222.14657568,4853.33625793)
\curveto(2226.04895723,4857.91006072)(2228.00017293,4863.97351596)(2228.00022864,4871.52664185)
\lineto(2228.00022864,4874.10728455)
\lineto(2216.48175024,4874.10728455)
\moveto(2239.58164954,4878.89091492)
\lineto(2239.58164954,4838.6706543)
\lineto(2228.00022864,4838.6706543)
\lineto(2228.00022864,4849.37088013)
\curveto(2225.35659037,4845.09078337)(2222.06260258,4841.92268046)(2218.11825537,4839.86656189)
\curveto(2214.17381652,4837.85240255)(2209.34822931,4836.84532348)(2203.64147925,4836.84532166)
\curveto(2196.42404789,4836.84532348)(2190.67530486,4838.85948162)(2186.39523291,4842.88780212)
\curveto(2182.15709439,4846.95807582)(2180.03803218,4852.39210664)(2180.03803992,4859.18991089)
\curveto(2180.03803218,4867.12063805)(2182.68161474,4873.10017004)(2187.96879553,4877.12852478)
\curveto(2193.29790661,4881.15680261)(2201.22865429,4883.17096075)(2211.76106238,4883.17100525)
\lineto(2228.00022864,4883.17100525)
\lineto(2228.00022864,4884.30397034)
\curveto(2228.00017293,4889.63305145)(2226.23778456,4893.74529099)(2222.71305823,4896.64070129)
\curveto(2219.23019269,4899.57795728)(2214.32068222,4901.04661426)(2207.98451208,4901.04667664)
\curveto(2203.95616011,4901.04661426)(2200.0327479,4900.56405554)(2196.21426367,4899.59899902)
\curveto(2192.39573161,4898.63382065)(2188.72408916,4897.18614449)(2185.19932532,4895.25596619)
\lineto(2185.19932532,4905.95619202)
\curveto(2189.43743684,4907.59262822)(2193.54967638,4908.80951543)(2197.53605627,4909.6068573)
\curveto(2201.52238569,4910.44601892)(2205.40383628,4910.8656352)(2209.18041968,4910.8657074)
\curveto(2219.37705839,4910.8656352)(2226.99309386,4908.22205264)(2232.02854895,4902.93495178)
\curveto(2237.06388457,4897.64772239)(2239.58158225,4889.63305145)(2239.58164954,4878.89091492)
}
}
{
\newrgbcolor{curcolor}{0 0 0}
\pscustom[linestyle=none,fillstyle=solid,fillcolor=curcolor]
{
\newpath
\moveto(2275.2071073,4929.18197632)
\lineto(2275.2071073,4909.16625977)
\lineto(2299.06231665,4909.16625977)
\lineto(2299.06231665,4900.16548157)
\lineto(2275.2071073,4900.16548157)
\lineto(2275.2071073,4861.8964386)
\curveto(2275.2070837,4856.14767234)(2275.98337381,4852.45504908)(2277.53597998,4850.81855774)
\curveto(2279.13049591,4849.1820421)(2282.34056045,4848.36379036)(2287.16618323,4848.36380005)
\lineto(2299.06231665,4848.36380005)
\lineto(2299.06231665,4838.6706543)
\lineto(2287.16618323,4838.6706543)
\curveto(2278.22832091,4838.6706543)(2272.0599616,4840.3281386)(2268.66108679,4843.64311218)
\curveto(2265.26217787,4847.00003745)(2263.56273194,4853.0844735)(2263.5627439,4861.8964386)
\lineto(2263.5627439,4900.16548157)
\lineto(2255.06550574,4900.16548157)
\lineto(2255.06550574,4909.16625977)
\lineto(2263.5627439,4909.16625977)
\lineto(2263.5627439,4929.18197632)
\lineto(2275.2071073,4929.18197632)
}
}
{
\newrgbcolor{curcolor}{0 0 0}
\pscustom[linestyle=none,fillstyle=solid,fillcolor=curcolor]
{
\newpath
\moveto(2374.65626501,4876.81381226)
\lineto(2374.65626501,4871.14898682)
\lineto(2321.40690588,4871.14898682)
\curveto(2321.91042622,4863.17624503)(2324.30223902,4857.09180897)(2328.58235144,4852.8956604)
\curveto(2332.90437275,4848.74144501)(2338.90488555,4846.66434442)(2346.58390784,4846.66435242)
\curveto(2351.03179603,4846.66434442)(2355.33286289,4847.20984559)(2359.48712134,4848.30085754)
\curveto(2363.68322686,4849.39185024)(2367.83742802,4851.02835373)(2371.9497373,4853.21037292)
\lineto(2371.9497373,4842.25837708)
\curveto(2367.7954664,4840.49598511)(2363.53636116,4839.15321302)(2359.17240881,4838.23005676)
\curveto(2354.80834254,4837.30690139)(2350.38139079,4836.84532348)(2345.89154028,4836.84532166)
\curveto(2334.64578031,4836.84532348)(2325.72893437,4840.11833046)(2319.14097571,4846.66435242)
\curveto(2312.59494481,4853.21035839)(2309.32193783,4862.06426188)(2309.32194495,4873.22608948)
\curveto(2309.32193783,4884.76550261)(2312.4270983,4893.91313751)(2318.63743567,4900.66902161)
\curveto(2324.88970181,4907.46674334)(2333.30300821,4910.8656352)(2343.87738013,4910.8657074)
\curveto(2353.36066638,4910.8656352)(2360.85081697,4907.80243636)(2366.34785437,4901.67610168)
\curveto(2371.88672512,4895.59160262)(2374.65619257,4887.3041811)(2374.65626501,4876.81381226)
\moveto(2363.07484412,4880.21270752)
\curveto(2362.99085999,4886.5488718)(2361.20749081,4891.60524797)(2357.7247312,4895.3818512)
\curveto(2354.28382219,4899.158341)(2349.71000475,4901.04661426)(2344.00326514,4901.04667664)
\curveto(2337.54113264,4901.04661426)(2332.35887158,4899.22128344)(2328.45646643,4895.57067871)
\curveto(2324.59597041,4891.91996018)(2322.37200413,4886.77966075)(2321.78456091,4880.14976501)
\lineto(2363.07484412,4880.21270752)
}
}
{
\newrgbcolor{curcolor}{0 0 0}
\pscustom[linestyle=none,fillstyle=solid,fillcolor=curcolor]
{
\newpath
\moveto(2380.887573,4932.64381409)
\lineto(2460.38395667,4932.64381409)
\lineto(2460.38395667,4921.94358826)
\lineto(2427.02442908,4921.94358826)
\lineto(2427.02442908,4838.6706543)
\lineto(2414.24710059,4838.6706543)
\lineto(2414.24710059,4921.94358826)
\lineto(2380.887573,4921.94358826)
\lineto(2380.887573,4932.64381409)
}
}
{
\newrgbcolor{curcolor}{0 0 0}
\pscustom[linestyle=none,fillstyle=solid,fillcolor=curcolor]
{
\newpath
\moveto(2468.44059729,4909.16625977)
\lineto(2480.02201819,4909.16625977)
\lineto(2480.02201819,4838.6706543)
\lineto(2468.44059729,4838.6706543)
\lineto(2468.44059729,4909.16625977)
\moveto(2468.44059729,4936.60919189)
\lineto(2480.02201819,4936.60919189)
\lineto(2480.02201819,4921.94358826)
\lineto(2468.44059729,4921.94358826)
\lineto(2468.44059729,4936.60919189)
}
}
{
\newrgbcolor{curcolor}{0 0 0}
\pscustom[linestyle=none,fillstyle=solid,fillcolor=curcolor]
{
\newpath
\moveto(2559.07780432,4895.63362122)
\curveto(2561.97308962,4900.83680612)(2565.43492392,4904.67629508)(2569.46331763,4907.15209961)
\curveto(2573.49155649,4909.62776718)(2578.23322045,4910.8656352)(2583.68832373,4910.8657074)
\curveto(2591.03151698,4910.8656352)(2596.69633675,4908.28499508)(2600.68280005,4903.1237793)
\curveto(2604.66904607,4898.00439623)(2606.66222339,4890.70307297)(2606.66233801,4881.2197876)
\lineto(2606.66233801,4838.6706543)
\lineto(2595.01797461,4838.6706543)
\lineto(2595.01797461,4880.84213257)
\curveto(2595.01787164,4887.5979125)(2593.82196524,4892.61232704)(2591.43025183,4895.88539124)
\curveto(2589.03833965,4899.158341)(2585.38767802,4900.79484449)(2580.47825598,4900.79490662)
\curveto(2574.47765475,4900.79484449)(2569.73599079,4898.80166716)(2566.25324988,4894.81536865)
\curveto(2562.77036055,4890.82895785)(2561.02895299,4885.39492703)(2561.02902197,4878.51325989)
\lineto(2561.02902197,4838.6706543)
\lineto(2549.38465857,4838.6706543)
\lineto(2549.38465857,4880.84213257)
\curveto(2549.38460123,4887.63987413)(2548.18869483,4892.65428867)(2545.79693579,4895.88539124)
\curveto(2543.40506924,4899.158341)(2539.71244598,4900.79484449)(2534.71905493,4900.79490662)
\curveto(2528.80242271,4900.79484449)(2524.10272038,4898.78068635)(2520.61993384,4894.75242615)
\curveto(2517.13709014,4890.76601541)(2515.39568258,4885.3529654)(2515.39570593,4878.51325989)
\lineto(2515.39570593,4838.6706543)
\lineto(2503.75134253,4838.6706543)
\lineto(2503.75134253,4909.16625977)
\lineto(2515.39570593,4909.16625977)
\lineto(2515.39570593,4898.21426392)
\curveto(2518.03926514,4902.53625205)(2521.20736805,4905.72533578)(2524.90002417,4907.78152466)
\curveto(2528.59261457,4909.83757532)(2532.9776047,4910.8656352)(2538.05500769,4910.8657074)
\curveto(2543.17428029,4910.8656352)(2547.51730878,4909.56482474)(2551.0841062,4906.96327209)
\curveto(2554.69274717,4904.36158287)(2557.35731054,4900.58503635)(2559.07780432,4895.63362122)
}
}
{
\newrgbcolor{curcolor}{0 0 0}
\pscustom[linestyle=none,fillstyle=solid,fillcolor=curcolor]
{
\newpath
\moveto(2689.87232947,4876.81381226)
\lineto(2689.87232947,4871.14898682)
\lineto(2636.62297034,4871.14898682)
\curveto(2637.12649067,4863.17624503)(2639.51830347,4857.09180897)(2643.79841589,4852.8956604)
\curveto(2648.1204372,4848.74144501)(2654.12095,4846.66434442)(2661.79997229,4846.66435242)
\curveto(2666.24786048,4846.66434442)(2670.54892735,4847.20984559)(2674.70318579,4848.30085754)
\curveto(2678.89929131,4849.39185024)(2683.05349248,4851.02835373)(2687.16580176,4853.21037292)
\lineto(2687.16580176,4842.25837708)
\curveto(2683.01153085,4840.49598511)(2678.75242561,4839.15321302)(2674.38847327,4838.23005676)
\curveto(2670.024407,4837.30690139)(2665.59745525,4836.84532348)(2661.10760474,4836.84532166)
\curveto(2649.86184476,4836.84532348)(2640.94499882,4840.11833046)(2634.35704016,4846.66435242)
\curveto(2627.81100927,4853.21035839)(2624.53800229,4862.06426188)(2624.5380094,4873.22608948)
\curveto(2624.53800229,4884.76550261)(2627.64316276,4893.91313751)(2633.85350012,4900.66902161)
\curveto(2640.10576626,4907.46674334)(2648.51907267,4910.8656352)(2659.09344458,4910.8657074)
\curveto(2668.57673083,4910.8656352)(2676.06688142,4907.80243636)(2681.56391882,4901.67610168)
\curveto(2687.10278958,4895.59160262)(2689.87225702,4887.3041811)(2689.87232947,4876.81381226)
\moveto(2678.29090857,4880.21270752)
\curveto(2678.20692445,4886.5488718)(2676.42355526,4891.60524797)(2672.94079565,4895.3818512)
\curveto(2669.49988665,4899.158341)(2664.9260692,4901.04661426)(2659.21932959,4901.04667664)
\curveto(2652.75719709,4901.04661426)(2647.57493604,4899.22128344)(2643.67253088,4895.57067871)
\curveto(2639.81203486,4891.91996018)(2637.58806858,4886.77966075)(2637.00062537,4880.14976501)
\lineto(2678.29090857,4880.21270752)
}
}
{
\newrgbcolor{curcolor}{0.51764709 0.74509805 0.51764709}
\pscustom[linestyle=none,fillstyle=solid,fillcolor=curcolor]
{
\newpath
\moveto(1256.79159375,4718.59076074)
\curveto(1256.79159375,4699.17450001)(1241.05160449,4683.43451074)(1221.63534375,4683.43451074)
\curveto(1202.21908301,4683.43451074)(1186.47909375,4699.17450001)(1186.47909375,4718.59076074)
\curveto(1186.47909375,4738.00702148)(1202.21908301,4753.74701074)(1221.63534375,4753.74701074)
\curveto(1241.05160449,4753.74701074)(1256.79159375,4738.00702148)(1256.79159375,4718.59076074)
\closepath
}
}
{
\newrgbcolor{curcolor}{0.01176471 0.50196081 0.28235295}
\pscustom[linewidth=11.71875,linecolor=curcolor]
{
\newpath
\moveto(1256.79159375,4718.59076074)
\curveto(1256.79159375,4699.17450001)(1241.05160449,4683.43451074)(1221.63534375,4683.43451074)
\curveto(1202.21908301,4683.43451074)(1186.47909375,4699.17450001)(1186.47909375,4718.59076074)
\curveto(1186.47909375,4738.00702148)(1202.21908301,4753.74701074)(1221.63534375,4753.74701074)
\curveto(1241.05160449,4753.74701074)(1256.79159375,4738.00702148)(1256.79159375,4718.59076074)
\closepath
}
}
{
\newrgbcolor{curcolor}{0 0 0}
\pscustom[linestyle=none,fillstyle=solid,fillcolor=curcolor]
{
\newpath
\moveto(1385.64068579,4724.74464417)
\curveto(1385.64062725,4733.13693369)(1383.89921969,4739.64098602)(1380.41645789,4744.25682068)
\curveto(1376.97555108,4748.87254418)(1372.12898305,4751.18043371)(1365.87673926,4751.18049622)
\curveto(1359.66637955,4751.18043371)(1354.81981152,4748.87254418)(1351.33702063,4744.25682068)
\curveto(1347.8961429,4739.64098602)(1346.17571616,4733.13693369)(1346.17573523,4724.74464417)
\curveto(1346.17571616,4716.39424413)(1347.8961429,4709.91117262)(1351.33702063,4705.29541016)
\curveto(1354.81981152,4700.67961446)(1359.66637955,4698.37172493)(1365.87673926,4698.37173462)
\curveto(1372.12898305,4698.37172493)(1376.97555108,4700.67961446)(1380.41645789,4705.29541016)
\curveto(1383.89921969,4709.91117262)(1385.64062725,4716.39424413)(1385.64068579,4724.74464417)
\moveto(1397.22210669,4697.42759705)
\curveto(1397.22203657,4685.4265627)(1394.5574732,4676.50971676)(1389.22840857,4670.67703247)
\curveto(1383.89921969,4664.80242256)(1375.73768306,4661.8651086)(1364.74377417,4661.86508179)
\curveto(1360.67345862,4661.8651086)(1356.83396966,4662.17982081)(1353.22529578,4662.80921936)
\curveto(1349.61656965,4663.39670802)(1346.11277372,4664.31986384)(1342.71389746,4665.57868958)
\lineto(1342.71389746,4676.84539795)
\curveto(1346.11277372,4674.99909815)(1349.46970395,4673.63534524)(1352.78469824,4672.75413513)
\curveto(1356.09964117,4671.87295687)(1359.47755222,4671.43235978)(1362.91844153,4671.43234253)
\curveto(1370.51346037,4671.43235978)(1376.19926096,4673.4255371)(1379.97586035,4677.41188049)
\curveto(1383.752354,4681.35628479)(1385.64062725,4687.33581677)(1385.64068579,4695.35049438)
\lineto(1385.64068579,4701.07826233)
\curveto(1383.24881446,4696.92404876)(1380.18561562,4693.81888829)(1376.45108008,4691.76277161)
\curveto(1372.71644584,4689.70664875)(1368.24753246,4688.67858887)(1363.04432654,4688.67858887)
\curveto(1354.40019524,4688.67858887)(1347.434565,4691.97257666)(1342.14741492,4698.56056213)
\curveto(1336.86023475,4705.14852784)(1334.21665219,4713.87654646)(1334.2166593,4724.74464417)
\curveto(1334.21665219,4735.65463137)(1336.86023475,4744.4036308)(1342.14741492,4750.9916687)
\curveto(1347.434565,4757.57958198)(1354.40019524,4760.87356977)(1363.04432654,4760.87364197)
\curveto(1368.24753246,4760.87356977)(1372.71644584,4759.84550989)(1376.45108008,4757.78945923)
\curveto(1380.18561562,4755.73327035)(1383.24881446,4752.62810988)(1385.64068579,4748.47396851)
\lineto(1385.64068579,4759.17419434)
\lineto(1397.22210669,4759.17419434)
\lineto(1397.22210669,4697.42759705)
}
}
{
\newrgbcolor{curcolor}{0 0 0}
\pscustom[linestyle=none,fillstyle=solid,fillcolor=curcolor]
{
\newpath
\moveto(1481.37623572,4726.82174683)
\lineto(1481.37623572,4721.15692139)
\lineto(1428.12687659,4721.15692139)
\curveto(1428.63039692,4713.1841796)(1431.02220972,4707.09974354)(1435.30232214,4702.90359497)
\curveto(1439.62434345,4698.74937958)(1445.62485625,4696.67227899)(1453.30387854,4696.67228699)
\curveto(1457.75176673,4696.67227899)(1462.0528336,4697.21778016)(1466.20709204,4698.30879211)
\curveto(1470.40319756,4699.39978481)(1474.55739873,4701.0362883)(1478.66970801,4703.2183075)
\lineto(1478.66970801,4692.26631165)
\curveto(1474.5154371,4690.50391968)(1470.25633186,4689.16114759)(1465.89237952,4688.23799133)
\curveto(1461.52831325,4687.31483596)(1457.1013615,4686.85325805)(1452.61151099,4686.85325623)
\curveto(1441.36575101,4686.85325805)(1432.44890507,4690.12626503)(1425.86094641,4696.67228699)
\curveto(1419.31491552,4703.21829296)(1416.04190854,4712.07219646)(1416.04191565,4723.23402405)
\curveto(1416.04190854,4734.77343718)(1419.14706901,4743.92107208)(1425.35740637,4750.67695618)
\curveto(1431.60967251,4757.47467791)(1440.02297892,4760.87356977)(1450.59735083,4760.87364197)
\curveto(1460.08063708,4760.87356977)(1467.57078767,4757.81037093)(1473.06782507,4751.68403625)
\curveto(1478.60669583,4745.5995372)(1481.37616327,4737.31211567)(1481.37623572,4726.82174683)
\moveto(1469.79481482,4730.22064209)
\curveto(1469.7108307,4736.55680637)(1467.92746151,4741.61318254)(1464.4447019,4745.38978577)
\curveto(1461.0037929,4749.16627557)(1456.42997545,4751.05454883)(1450.72323584,4751.05461121)
\curveto(1444.26110334,4751.05454883)(1439.07884229,4749.22921801)(1435.17643713,4745.57861328)
\curveto(1431.31594111,4741.92789475)(1429.09197483,4736.78759532)(1428.50453162,4730.15769958)
\lineto(1469.79481482,4730.22064209)
}
}
{
\newrgbcolor{curcolor}{0 0 0}
\pscustom[linestyle=none,fillstyle=solid,fillcolor=curcolor]
{
\newpath
\moveto(1511.58863806,4779.18991089)
\lineto(1511.58863806,4759.17419434)
\lineto(1535.44384741,4759.17419434)
\lineto(1535.44384741,4750.17341614)
\lineto(1511.58863806,4750.17341614)
\lineto(1511.58863806,4711.90437317)
\curveto(1511.58861446,4706.15560691)(1512.36490458,4702.46298365)(1513.91751074,4700.82649231)
\curveto(1515.51202667,4699.18997667)(1518.72209121,4698.37172493)(1523.54771399,4698.37173462)
\lineto(1535.44384741,4698.37173462)
\lineto(1535.44384741,4688.67858887)
\lineto(1523.54771399,4688.67858887)
\curveto(1514.60985167,4688.67858887)(1508.44149236,4690.33607317)(1505.04261755,4693.65104675)
\curveto(1501.64370863,4697.00797202)(1499.9442627,4703.09240807)(1499.94427466,4711.90437317)
\lineto(1499.94427466,4750.17341614)
\lineto(1491.4470365,4750.17341614)
\lineto(1491.4470365,4759.17419434)
\lineto(1499.94427466,4759.17419434)
\lineto(1499.94427466,4779.18991089)
\lineto(1511.58863806,4779.18991089)
}
}
{
\newrgbcolor{curcolor}{0 0 0}
\pscustom[linestyle=none,fillstyle=solid,fillcolor=curcolor]
{
\newpath
\moveto(1563.95680212,4772.20329285)
\lineto(1563.95680212,4736.89254761)
\lineto(1579.94419836,4736.89254761)
\curveto(1585.86074655,4736.89249939)(1590.434564,4738.42409881)(1593.66566443,4741.48735046)
\curveto(1596.89665471,4744.5504965)(1598.51217738,4748.9145058)(1598.5122373,4754.57939148)
\curveto(1598.51217738,4760.20218373)(1596.89665471,4764.54521222)(1593.66566443,4767.60848999)
\curveto(1590.434564,4770.6716099)(1585.86074655,4772.20320932)(1579.94419836,4772.20329285)
\lineto(1563.95680212,4772.20329285)
\moveto(1551.24241614,4782.65174866)
\lineto(1579.94419836,4782.65174866)
\curveto(1590.47652563,4782.65165468)(1598.42825413,4780.25984189)(1603.79940771,4775.4763031)
\curveto(1609.21239251,4770.73455234)(1611.91891752,4763.7689221)(1611.91899084,4754.57939148)
\curveto(1611.91891752,4745.3058058)(1609.21239251,4738.29821393)(1603.79940771,4733.55659485)
\curveto(1598.42825413,4728.81488601)(1590.47652563,4726.44405403)(1579.94419836,4726.4440918)
\lineto(1563.95680212,4726.4440918)
\lineto(1563.95680212,4688.67858887)
\lineto(1551.24241614,4688.67858887)
\lineto(1551.24241614,4782.65174866)
}
}
{
\newrgbcolor{curcolor}{0 0 0}
\pscustom[linestyle=none,fillstyle=solid,fillcolor=curcolor]
{
\newpath
\moveto(1667.11956763,4748.3480835)
\curveto(1665.81870416,4749.10333313)(1664.39200881,4749.64883429)(1662.83947729,4749.98458862)
\curveto(1661.32880997,4750.36218197)(1659.65034485,4750.55100929)(1657.8040769,4750.55107117)
\curveto(1651.25801926,4750.55100929)(1646.22262391,4748.41096627)(1642.69787573,4744.13093567)
\curveto(1639.21503204,4739.89275579)(1637.47362448,4733.78733892)(1637.47364783,4725.81466675)
\lineto(1637.47364783,4688.67858887)
\lineto(1625.82928442,4688.67858887)
\lineto(1625.82928442,4759.17419434)
\lineto(1637.47364783,4759.17419434)
\lineto(1637.47364783,4748.22219849)
\curveto(1639.9073989,4752.50222499)(1643.07550181,4755.67032791)(1646.97796606,4757.72651672)
\curveto(1650.88036461,4759.82452907)(1655.62202857,4760.87356977)(1661.20297217,4760.87364197)
\curveto(1662.00019602,4760.87356977)(1662.88139021,4760.81062733)(1663.84655737,4760.68481445)
\curveto(1664.81162509,4760.60081919)(1665.8816466,4760.45395349)(1667.05662512,4760.24421692)
\lineto(1667.11956763,4748.3480835)
}
}
{
\newrgbcolor{curcolor}{0 0 0}
\pscustom[linestyle=none,fillstyle=solid,fillcolor=curcolor]
{
\newpath
\moveto(1736.92280554,4726.82174683)
\lineto(1736.92280554,4721.15692139)
\lineto(1683.67344641,4721.15692139)
\curveto(1684.17696675,4713.1841796)(1686.56877954,4707.09974354)(1690.84889197,4702.90359497)
\curveto(1695.17091327,4698.74937958)(1701.17142607,4696.67227899)(1708.85044836,4696.67228699)
\curveto(1713.29833655,4696.67227899)(1717.59940342,4697.21778016)(1721.75366187,4698.30879211)
\curveto(1725.94976738,4699.39978481)(1730.10396855,4701.0362883)(1734.21627783,4703.2183075)
\lineto(1734.21627783,4692.26631165)
\curveto(1730.06200692,4690.50391968)(1725.80290169,4689.16114759)(1721.43894934,4688.23799133)
\curveto(1717.07488307,4687.31483596)(1712.64793132,4686.85325805)(1708.15808081,4686.85325623)
\curveto(1696.91232084,4686.85325805)(1687.99547489,4690.12626503)(1681.40751624,4696.67228699)
\curveto(1674.86148534,4703.21829296)(1671.58847836,4712.07219646)(1671.58848547,4723.23402405)
\curveto(1671.58847836,4734.77343718)(1674.69363883,4743.92107208)(1680.9039762,4750.67695618)
\curveto(1687.15624233,4757.47467791)(1695.56954874,4760.87356977)(1706.14392065,4760.87364197)
\curveto(1715.62720691,4760.87356977)(1723.1173575,4757.81037093)(1728.6143949,4751.68403625)
\curveto(1734.15326565,4745.5995372)(1736.9227331,4737.31211567)(1736.92280554,4726.82174683)
\moveto(1725.34138464,4730.22064209)
\curveto(1725.25740052,4736.55680637)(1723.47403133,4741.61318254)(1719.99127173,4745.38978577)
\curveto(1716.55036272,4749.16627557)(1711.97654527,4751.05454883)(1706.26980566,4751.05461121)
\curveto(1699.80767316,4751.05454883)(1694.62541211,4749.22921801)(1690.72300696,4745.57861328)
\curveto(1686.86251094,4741.92789475)(1684.63854466,4736.78759532)(1684.05110144,4730.15769958)
\lineto(1725.34138464,4730.22064209)
}
}
{
\newrgbcolor{curcolor}{0 0 0}
\pscustom[linestyle=none,fillstyle=solid,fillcolor=curcolor]
{
\newpath
\moveto(1800.62062048,4757.09709167)
\lineto(1800.62062048,4746.14509583)
\curveto(1797.34755641,4747.82350348)(1793.94866455,4749.08235232)(1790.42393469,4749.92164612)
\curveto(1786.89911105,4750.76081743)(1783.24844942,4751.18043371)(1779.47193884,4751.18049622)
\curveto(1773.72315987,4751.18043371)(1769.40111219,4750.29923953)(1766.50578284,4748.53691101)
\curveto(1763.65236916,4746.77446278)(1762.22567381,4744.13088022)(1762.2256925,4740.6061554)
\curveto(1762.22567381,4737.92055928)(1763.2537337,4735.80149707)(1765.30987524,4734.2489624)
\curveto(1767.36597324,4732.73829823)(1771.49919359,4731.29062206)(1777.70954871,4729.90592957)
\lineto(1781.67492651,4729.0247345)
\curveto(1789.89936745,4727.26230578)(1795.73203374,4724.76558891)(1799.17294287,4721.53457642)
\curveto(1802.65570235,4718.34545984)(1804.39710991,4713.87654646)(1804.39717078,4708.12782288)
\curveto(1804.39710991,4701.58178946)(1801.79548898,4696.39952841)(1796.59230017,4692.58102417)
\curveto(1791.43096687,4688.76251212)(1784.31847093,4686.85325805)(1775.25479102,4686.85325623)
\curveto(1771.47821278,4686.85325805)(1767.53381975,4687.2309127)(1763.4216001,4687.98622131)
\curveto(1759.3513023,4688.69956968)(1755.05023543,4689.79057201)(1750.5183866,4691.25923157)
\lineto(1750.5183866,4703.2183075)
\curveto(1754.79846566,4700.99432667)(1759.01560927,4699.31586156)(1763.16983008,4698.1829071)
\curveto(1767.32401161,4697.09189527)(1771.43625115,4696.54639411)(1775.50656104,4696.54640198)
\curveto(1780.9615407,4696.54639411)(1785.15770349,4697.46954992)(1788.09506201,4699.31587219)
\curveto(1791.03233141,4701.20413481)(1792.50098838,4703.84771737)(1792.50103735,4707.24662781)
\curveto(1792.50098838,4710.39373134)(1791.43096687,4712.80652494)(1789.2909696,4714.48501587)
\curveto(1787.19284245,4716.16345518)(1782.55608256,4717.77897786)(1775.38067603,4719.33158875)
\lineto(1771.35235571,4720.27572632)
\curveto(1764.17688951,4721.78631333)(1758.99462846,4724.09420287)(1755.80555701,4727.19940186)
\curveto(1752.61646101,4730.34648543)(1751.02191915,4734.6475523)(1751.02192664,4740.10261536)
\curveto(1751.02191915,4746.73250115)(1753.37177031,4751.85181976)(1758.07148718,4755.46058655)
\curveto(1762.77117497,4759.06921977)(1769.44307382,4760.87356977)(1778.08720374,4760.87364197)
\curveto(1782.36725523,4760.87356977)(1786.39557152,4760.55885756)(1790.17216467,4759.92950439)
\curveto(1793.94866455,4759.30000872)(1797.43147967,4758.3558721)(1800.62062048,4757.09709167)
}
}
{
\newrgbcolor{curcolor}{0 0 0}
\pscustom[linestyle=none,fillstyle=solid,fillcolor=curcolor]
{
\newpath
\moveto(1883.20118689,4726.82174683)
\lineto(1883.20118689,4721.15692139)
\lineto(1829.95182776,4721.15692139)
\curveto(1830.4553481,4713.1841796)(1832.84716089,4707.09974354)(1837.12727332,4702.90359497)
\curveto(1841.44929462,4698.74937958)(1847.44980742,4696.67227899)(1855.12882971,4696.67228699)
\curveto(1859.5767179,4696.67227899)(1863.87778477,4697.21778016)(1868.03204321,4698.30879211)
\curveto(1872.22814873,4699.39978481)(1876.3823499,4701.0362883)(1880.49465918,4703.2183075)
\lineto(1880.49465918,4692.26631165)
\curveto(1876.34038827,4690.50391968)(1872.08128303,4689.16114759)(1867.71733069,4688.23799133)
\curveto(1863.35326442,4687.31483596)(1858.92631267,4686.85325805)(1854.43646216,4686.85325623)
\curveto(1843.19070218,4686.85325805)(1834.27385624,4690.12626503)(1827.68589758,4696.67228699)
\curveto(1821.13986669,4703.21829296)(1817.86685971,4712.07219646)(1817.86686682,4723.23402405)
\curveto(1817.86685971,4734.77343718)(1820.97202018,4743.92107208)(1827.18235754,4750.67695618)
\curveto(1833.43462368,4757.47467791)(1841.84793009,4760.87356977)(1852.422302,4760.87364197)
\curveto(1861.90558825,4760.87356977)(1869.39573884,4757.81037093)(1874.89277625,4751.68403625)
\curveto(1880.431647,4745.5995372)(1883.20111444,4737.31211567)(1883.20118689,4726.82174683)
\moveto(1871.61976599,4730.22064209)
\curveto(1871.53578187,4736.55680637)(1869.75241268,4741.61318254)(1866.26965308,4745.38978577)
\curveto(1862.82874407,4749.16627557)(1858.25492662,4751.05454883)(1852.54818701,4751.05461121)
\curveto(1846.08605451,4751.05454883)(1840.90379346,4749.22921801)(1837.00138831,4745.57861328)
\curveto(1833.14089229,4741.92789475)(1830.916926,4736.78759532)(1830.32948279,4730.15769958)
\lineto(1871.61976599,4730.22064209)
}
}
{
\newrgbcolor{curcolor}{0 0 0}
\pscustom[linestyle=none,fillstyle=solid,fillcolor=curcolor]
{
\newpath
\moveto(1960.55752539,4731.22772217)
\lineto(1960.55752539,4688.67858887)
\lineto(1948.97610449,4688.67858887)
\lineto(1948.97610449,4730.85006714)
\curveto(1948.97604533,4737.52192381)(1947.67523486,4742.51535754)(1945.07366919,4745.8303833)
\curveto(1942.47199299,4749.14529476)(1938.56956159,4750.80277906)(1933.36636328,4750.80284119)
\curveto(1927.11403716,4750.80277906)(1922.18354587,4748.80960173)(1918.57487463,4744.82330322)
\curveto(1914.96614586,4740.83689242)(1913.16179586,4735.4028616)(1913.16181921,4728.52119446)
\lineto(1913.16181921,4688.67858887)
\lineto(1901.51745581,4688.67858887)
\lineto(1901.51745581,4759.17419434)
\lineto(1913.16181921,4759.17419434)
\lineto(1913.16181921,4748.22219849)
\curveto(1915.93126331,4752.46026337)(1919.18328947,4755.62836628)(1922.91790747,4757.72651672)
\curveto(1926.69442088,4759.82452907)(1931.03744937,4760.87356977)(1935.94700598,4760.87364197)
\curveto(1944.04555404,4760.87356977)(1950.17195172,4758.3558721)(1954.32621741,4753.32054138)
\curveto(1958.48035406,4748.32704301)(1960.55745464,4740.96277731)(1960.55752539,4731.22772217)
}
}
{
\newrgbcolor{curcolor}{0 0 0}
\pscustom[linestyle=none,fillstyle=solid,fillcolor=curcolor]
{
\newpath
\moveto(2034.51496863,4756.46766663)
\lineto(2034.51496863,4745.64155579)
\curveto(2031.24189877,4747.44584883)(2027.94791097,4748.78862092)(2024.63299536,4749.6698761)
\curveto(2021.35993538,4750.59297092)(2018.04496677,4751.05454883)(2014.68807959,4751.05461121)
\curveto(2007.17690513,4751.05454883)(2001.34423885,4748.66273604)(1997.19006323,4743.87916565)
\curveto(1993.03583651,4739.13744649)(1990.95873593,4732.46554764)(1990.95875525,4723.8634491)
\curveto(1990.95873593,4715.26128018)(1993.03583651,4708.56840052)(1997.19006323,4703.78479004)
\curveto(2001.34423885,4699.04311097)(2007.17690513,4696.67227899)(2014.68807959,4696.67228699)
\curveto(2018.04496677,4696.67227899)(2021.35993538,4697.11287609)(2024.63299536,4697.99407959)
\curveto(2027.94791097,4698.91722609)(2031.24189877,4700.280979)(2034.51496863,4702.08534241)
\lineto(2034.51496863,4691.38511658)
\curveto(2031.2838604,4689.87449526)(2027.92693016,4688.74153131)(2024.44416785,4687.98622131)
\curveto(2021.00326155,4687.2309127)(2017.3316191,4686.85325805)(2013.42922949,4686.85325623)
\curveto(2002.81289582,4686.85325805)(1994.3786086,4690.18920747)(1988.12634253,4696.8611145)
\curveto(1981.87404347,4703.53300517)(1978.74790219,4712.53377436)(1978.7479093,4723.8634491)
\curveto(1978.74790219,4735.36089997)(1981.89502429,4744.4036308)(1988.18928503,4750.9916687)
\curveto(1994.5254743,4757.57958198)(2003.19055048,4760.87356977)(2014.18453955,4760.87364197)
\curveto(2017.75123538,4760.87356977)(2021.2340505,4760.49591512)(2024.63299536,4759.74067688)
\curveto(2028.03183423,4759.02725814)(2031.32582202,4757.93625582)(2034.51496863,4756.46766663)
}
}
{
\newrgbcolor{curcolor}{0 0 0}
\pscustom[linestyle=none,fillstyle=solid,fillcolor=curcolor]
{
\newpath
\moveto(2114.82960486,4726.82174683)
\lineto(2114.82960486,4721.15692139)
\lineto(2061.58024573,4721.15692139)
\curveto(2062.08376607,4713.1841796)(2064.47557886,4707.09974354)(2068.75569128,4702.90359497)
\curveto(2073.07771259,4698.74937958)(2079.07822539,4696.67227899)(2086.75724768,4696.67228699)
\curveto(2091.20513587,4696.67227899)(2095.50620274,4697.21778016)(2099.66046118,4698.30879211)
\curveto(2103.8565667,4699.39978481)(2108.01076787,4701.0362883)(2112.12307715,4703.2183075)
\lineto(2112.12307715,4692.26631165)
\curveto(2107.96880624,4690.50391968)(2103.709701,4689.16114759)(2099.34574866,4688.23799133)
\curveto(2094.98168239,4687.31483596)(2090.55473064,4686.85325805)(2086.06488013,4686.85325623)
\curveto(2074.81912015,4686.85325805)(2065.90227421,4690.12626503)(2059.31431555,4696.67228699)
\curveto(2052.76828466,4703.21829296)(2049.49527768,4712.07219646)(2049.49528479,4723.23402405)
\curveto(2049.49527768,4734.77343718)(2052.60043815,4743.92107208)(2058.81077551,4750.67695618)
\curveto(2065.06304165,4757.47467791)(2073.47634806,4760.87356977)(2084.05071997,4760.87364197)
\curveto(2093.53400622,4760.87356977)(2101.02415681,4757.81037093)(2106.52119421,4751.68403625)
\curveto(2112.06006497,4745.5995372)(2114.82953241,4737.31211567)(2114.82960486,4726.82174683)
\moveto(2103.24818396,4730.22064209)
\curveto(2103.16419984,4736.55680637)(2101.38083065,4741.61318254)(2097.89807104,4745.38978577)
\curveto(2094.45716204,4749.16627557)(2089.88334459,4751.05454883)(2084.17660498,4751.05461121)
\curveto(2077.71447248,4751.05454883)(2072.53221143,4749.22921801)(2068.62980627,4745.57861328)
\curveto(2064.76931026,4741.92789475)(2062.54534397,4736.78759532)(2061.95790076,4730.15769958)
\lineto(2103.24818396,4730.22064209)
}
}
{
\newrgbcolor{curcolor}{0 0 0}
\pscustom[linestyle=none,fillstyle=solid,fillcolor=curcolor]
{
\newpath
\moveto(2161.40705847,4786.49124146)
\curveto(2155.78416036,4776.83996921)(2151.60897837,4767.29369885)(2148.8815,4757.85240173)
\curveto(2146.15396674,4748.41096627)(2144.79021383,4738.84371509)(2144.79023718,4729.15061951)
\curveto(2144.79021383,4719.45744298)(2146.15396674,4709.84823017)(2148.8815,4700.32295227)
\curveto(2151.65094,4690.83961271)(2155.82612198,4681.29334235)(2161.40705847,4671.68411255)
\lineto(2151.33625769,4671.68411255)
\curveto(2145.0419836,4681.54511211)(2140.32130045,4691.23824817)(2137.17419409,4700.7635498)
\curveto(2134.06901789,4710.28882727)(2132.51643765,4719.75117437)(2132.51644873,4729.15061951)
\curveto(2132.51643765,4738.50802207)(2134.06901789,4747.92840755)(2137.17419409,4757.4118042)
\curveto(2140.27933883,4766.89506339)(2145.00002197,4776.58819944)(2151.33625769,4786.49124146)
\lineto(2161.40705847,4786.49124146)
}
}
{
\newrgbcolor{curcolor}{0 0 0}
\pscustom[linestyle=none,fillstyle=solid,fillcolor=curcolor]
{
\newpath
\moveto(2182.11514258,4786.49124146)
\lineto(2192.18594336,4786.49124146)
\curveto(2198.48016716,4776.58819944)(2203.17986949,4766.89506339)(2206.28506445,4757.4118042)
\curveto(2209.43215206,4747.92840755)(2211.00571311,4738.50802207)(2211.00575232,4729.15061951)
\curveto(2211.00571311,4719.75117437)(2209.43215206,4710.28882727)(2206.28506445,4700.7635498)
\curveto(2203.17986949,4691.23824817)(2198.48016716,4681.54511211)(2192.18594336,4671.68411255)
\lineto(2182.11514258,4671.68411255)
\curveto(2187.69602877,4681.29334235)(2191.85022994,4690.83961271)(2194.57775854,4700.32295227)
\curveto(2197.34720321,4709.84823017)(2198.73193693,4719.45744298)(2198.73196387,4729.15061951)
\curveto(2198.73193693,4738.84371509)(2197.34720321,4748.41096627)(2194.57775854,4757.85240173)
\curveto(2191.85022994,4767.29369885)(2187.69602877,4776.83996921)(2182.11514258,4786.49124146)
}
}
{
\newrgbcolor{curcolor}{0 0 0}
\pscustom[linestyle=none,fillstyle=solid,fillcolor=curcolor]
{
}
}
{
\newrgbcolor{curcolor}{0 0 0}
\pscustom[linestyle=none,fillstyle=solid,fillcolor=curcolor]
{
\newpath
\moveto(2278.29129004,4704.66598511)
\lineto(2291.57215857,4704.66598511)
\lineto(2291.57215857,4688.67858887)
\lineto(2278.29129004,4688.67858887)
\lineto(2278.29129004,4704.66598511)
\moveto(2278.29129004,4755.33470154)
\lineto(2291.57215857,4755.33470154)
\lineto(2291.57215857,4739.3473053)
\lineto(2278.29129004,4739.3473053)
\lineto(2278.29129004,4755.33470154)
}
}
{
\newrgbcolor{curcolor}{0 0 0}
\pscustom[linestyle=none,fillstyle=solid,fillcolor=curcolor]
{
}
}
{
\newrgbcolor{curcolor}{0 0 0}
\pscustom[linestyle=none,fillstyle=solid,fillcolor=curcolor]
{
\newpath
\moveto(2360.43125891,4782.65174866)
\lineto(2414.4359281,4782.65174866)
\lineto(2414.4359281,4771.95152283)
\lineto(2373.1456449,4771.95152283)
\lineto(2373.1456449,4744.25682068)
\lineto(2410.40760779,4744.25682068)
\lineto(2410.40760779,4733.55659485)
\lineto(2373.1456449,4733.55659485)
\lineto(2373.1456449,4688.67858887)
\lineto(2360.43125891,4688.67858887)
\lineto(2360.43125891,4782.65174866)
}
}
{
\newrgbcolor{curcolor}{0 0 0}
\pscustom[linestyle=none,fillstyle=solid,fillcolor=curcolor]
{
\newpath
\moveto(2424.63261389,4759.17419434)
\lineto(2436.21403479,4759.17419434)
\lineto(2436.21403479,4688.67858887)
\lineto(2424.63261389,4688.67858887)
\lineto(2424.63261389,4759.17419434)
\moveto(2424.63261389,4786.61712646)
\lineto(2436.21403479,4786.61712646)
\lineto(2436.21403479,4771.95152283)
\lineto(2424.63261389,4771.95152283)
\lineto(2424.63261389,4786.61712646)
}
}
{
\newrgbcolor{curcolor}{0 0 0}
\pscustom[linestyle=none,fillstyle=solid,fillcolor=curcolor]
{
\newpath
\moveto(2452.07554602,4759.17419434)
\lineto(2464.34933447,4759.17419434)
\lineto(2486.37921118,4700.00823975)
\lineto(2508.40908789,4759.17419434)
\lineto(2520.68287634,4759.17419434)
\lineto(2494.24702429,4688.67858887)
\lineto(2478.51139807,4688.67858887)
\lineto(2452.07554602,4759.17419434)
}
}
{
\newrgbcolor{curcolor}{0 0 0}
\pscustom[linestyle=none,fillstyle=solid,fillcolor=curcolor]
{
\newpath
\moveto(2596.96919226,4726.82174683)
\lineto(2596.96919226,4721.15692139)
\lineto(2543.71983313,4721.15692139)
\curveto(2544.22335347,4713.1841796)(2546.61516626,4707.09974354)(2550.89527869,4702.90359497)
\curveto(2555.21729999,4698.74937958)(2561.21781279,4696.67227899)(2568.89683508,4696.67228699)
\curveto(2573.34472327,4696.67227899)(2577.64579014,4697.21778016)(2581.80004858,4698.30879211)
\curveto(2585.9961541,4699.39978481)(2590.15035527,4701.0362883)(2594.26266455,4703.2183075)
\lineto(2594.26266455,4692.26631165)
\curveto(2590.10839364,4690.50391968)(2585.8492884,4689.16114759)(2581.48533606,4688.23799133)
\curveto(2577.12126979,4687.31483596)(2572.69431804,4686.85325805)(2568.20446753,4686.85325623)
\curveto(2556.95870755,4686.85325805)(2548.04186161,4690.12626503)(2541.45390295,4696.67228699)
\curveto(2534.90787206,4703.21829296)(2531.63486508,4712.07219646)(2531.63487219,4723.23402405)
\curveto(2531.63486508,4734.77343718)(2534.74002555,4743.92107208)(2540.95036292,4750.67695618)
\curveto(2547.20262905,4757.47467791)(2555.61593546,4760.87356977)(2566.19030737,4760.87364197)
\curveto(2575.67359362,4760.87356977)(2583.16374422,4757.81037093)(2588.66078162,4751.68403625)
\curveto(2594.19965237,4745.5995372)(2596.96911981,4737.31211567)(2596.96919226,4726.82174683)
\moveto(2585.38777136,4730.22064209)
\curveto(2585.30378724,4736.55680637)(2583.52041805,4741.61318254)(2580.03765845,4745.38978577)
\curveto(2576.59674944,4749.16627557)(2572.02293199,4751.05454883)(2566.31619238,4751.05461121)
\curveto(2559.85405988,4751.05454883)(2554.67179883,4749.22921801)(2550.76939368,4745.57861328)
\curveto(2546.90889766,4741.92789475)(2544.68493138,4736.78759532)(2544.09748816,4730.15769958)
\lineto(2585.38777136,4730.22064209)
}
}
{
\newrgbcolor{curcolor}{0 0 0}
\pscustom[linestyle=none,fillstyle=solid,fillcolor=curcolor]
{
\newpath
\moveto(2672.56314062,4779.56756592)
\lineto(2672.56314062,4767.16789246)
\curveto(2667.73748442,4769.4757035)(2663.18464779,4771.19613025)(2658.90461707,4772.32917786)
\curveto(2654.62447569,4773.46205816)(2650.49125533,4774.02854014)(2646.5049436,4774.02862549)
\curveto(2639.58123206,4774.02854014)(2634.2311245,4772.68576804)(2630.45460486,4770.00030518)
\curveto(2626.71999309,4767.31467966)(2624.85270065,4763.49617152)(2624.85272192,4758.54476929)
\curveto(2624.85270065,4754.39049825)(2626.09056867,4751.24337616)(2628.56632971,4749.10339355)
\curveto(2631.0840024,4747.00525173)(2635.82566636,4745.3058058)(2642.79133582,4744.00505066)
\lineto(2650.47032141,4742.43148804)
\curveto(2659.95360244,4740.62708428)(2666.94021349,4737.43800056)(2671.43017554,4732.86422729)
\curveto(2675.9619635,4728.33232729)(2678.22789141,4722.24789124)(2678.22796606,4714.61090088)
\curveto(2678.22789141,4705.50520168)(2675.16469257,4698.60251388)(2669.03836035,4693.90281677)
\curveto(2662.95385884,4689.20310922)(2654.01603208,4686.85325805)(2642.22485327,4686.85325623)
\curveto(2637.77688206,4686.85325805)(2633.0352181,4687.35679759)(2627.99984717,4688.36387634)
\curveto(2623.00638902,4689.37095573)(2617.82412797,4690.86059352)(2612.45304846,4692.83279419)
\lineto(2612.45304846,4705.92483521)
\curveto(2617.61431983,4703.02946563)(2622.670696,4700.84746098)(2627.62219214,4699.3788147)
\curveto(2632.57364019,4697.91014702)(2637.44118904,4697.17581853)(2642.22485327,4697.17582703)
\curveto(2649.48417626,4697.17581853)(2655.08605359,4698.60251388)(2659.03050208,4701.45591736)
\curveto(2662.97483965,4704.30929528)(2664.94703617,4708.37957319)(2664.94709753,4713.66676331)
\curveto(2664.94703617,4718.28251739)(2663.52034081,4721.8912174)(2660.6670072,4724.49287415)
\curveto(2657.85552104,4727.09445926)(2653.21876115,4729.04567496)(2646.75671362,4730.3465271)
\lineto(2639.01478552,4731.85714722)
\curveto(2629.53142217,4733.7453773)(2622.670696,4736.70367207)(2618.43258643,4740.73204041)
\curveto(2614.19444715,4744.76030464)(2612.07538494,4750.36218197)(2612.07539343,4757.53768921)
\curveto(2612.07538494,4765.84602269)(2614.99171808,4772.39203665)(2620.82440161,4777.17575073)
\curveto(2626.69901228,4781.95928782)(2634.77662566,4784.35110062)(2645.05726599,4784.35119629)
\curveto(2649.46319545,4784.35110062)(2653.95308964,4783.95246515)(2658.52696204,4783.1552887)
\curveto(2663.10072453,4782.35792329)(2667.77944605,4781.16201689)(2672.56314062,4779.56756592)
}
}
{
\newrgbcolor{curcolor}{0 0 0}
\pscustom[linestyle=none,fillstyle=solid,fillcolor=curcolor]
{
\newpath
\moveto(2709.25862097,4779.18991089)
\lineto(2709.25862097,4759.17419434)
\lineto(2733.11383032,4759.17419434)
\lineto(2733.11383032,4750.17341614)
\lineto(2709.25862097,4750.17341614)
\lineto(2709.25862097,4711.90437317)
\curveto(2709.25859737,4706.15560691)(2710.03488749,4702.46298365)(2711.58749365,4700.82649231)
\curveto(2713.18200958,4699.18997667)(2716.39207412,4698.37172493)(2721.2176969,4698.37173462)
\lineto(2733.11383032,4698.37173462)
\lineto(2733.11383032,4688.67858887)
\lineto(2721.2176969,4688.67858887)
\curveto(2712.27983458,4688.67858887)(2706.11147527,4690.33607317)(2702.71260046,4693.65104675)
\curveto(2699.31369154,4697.00797202)(2697.61424561,4703.09240807)(2697.61425757,4711.90437317)
\lineto(2697.61425757,4750.17341614)
\lineto(2689.11701941,4750.17341614)
\lineto(2689.11701941,4759.17419434)
\lineto(2697.61425757,4759.17419434)
\lineto(2697.61425757,4779.18991089)
\lineto(2709.25862097,4779.18991089)
}
}
{
\newrgbcolor{curcolor}{0 0 0}
\pscustom[linestyle=none,fillstyle=solid,fillcolor=curcolor]
{
\newpath
\moveto(2780.44659399,4724.11521912)
\curveto(2771.08910677,4724.11518368)(2764.60603525,4723.04516217)(2760.99735999,4720.90515137)
\curveto(2757.38863524,4718.76507611)(2755.58428524,4715.11441448)(2755.58430457,4709.95315552)
\curveto(2755.58428524,4705.8408947)(2756.92705734,4702.56788772)(2759.61262488,4700.13412476)
\curveto(2762.34010734,4697.74230051)(2766.0327306,4696.54639411)(2770.69050574,4696.54640198)
\curveto(2777.11060039,4696.54639411)(2782.25089981,4698.81232202)(2786.11141943,4703.3441925)
\curveto(2790.01380098,4707.91799529)(2791.96501668,4713.98145053)(2791.96507239,4721.53457642)
\lineto(2791.96507239,4724.11521912)
\lineto(2780.44659399,4724.11521912)
\moveto(2803.54649329,4728.89884949)
\lineto(2803.54649329,4688.67858887)
\lineto(2791.96507239,4688.67858887)
\lineto(2791.96507239,4699.3788147)
\curveto(2789.32143412,4695.09871795)(2786.02744633,4691.93061503)(2782.08309912,4689.87449646)
\curveto(2778.13866027,4687.86033712)(2773.31307306,4686.85325805)(2767.606323,4686.85325623)
\curveto(2760.38889164,4686.85325805)(2754.64014861,4688.86741619)(2750.36007666,4692.89573669)
\curveto(2746.12193814,4696.96601039)(2744.00287593,4702.40004121)(2744.00288367,4709.19784546)
\curveto(2744.00287593,4717.12857262)(2746.64645849,4723.10810461)(2751.93363928,4727.13645935)
\curveto(2757.26275036,4731.16473718)(2765.19349804,4733.17889532)(2775.72590613,4733.17893982)
\lineto(2791.96507239,4733.17893982)
\lineto(2791.96507239,4734.31190491)
\curveto(2791.96501668,4739.64098602)(2790.20262831,4743.75322556)(2786.67790198,4746.64863586)
\curveto(2783.19503644,4749.58589185)(2778.28552597,4751.05454883)(2771.94935583,4751.05461121)
\curveto(2767.92100386,4751.05454883)(2763.99759165,4750.57199011)(2760.17910742,4749.60693359)
\curveto(2756.36057536,4748.64175522)(2752.68893291,4747.19407906)(2749.16416907,4745.26390076)
\lineto(2749.16416907,4755.96412659)
\curveto(2753.40228059,4757.60056279)(2757.51452013,4758.81745)(2761.50090002,4759.61479187)
\curveto(2765.48722944,4760.45395349)(2769.36868003,4760.87356977)(2773.14526343,4760.87364197)
\curveto(2783.34190214,4760.87356977)(2790.95793761,4758.22998721)(2795.9933927,4752.94288635)
\curveto(2801.02872832,4747.65565697)(2803.546426,4739.64098602)(2803.54649329,4728.89884949)
}
}
{
\newrgbcolor{curcolor}{0 0 0}
\pscustom[linestyle=none,fillstyle=solid,fillcolor=curcolor]
{
\newpath
\moveto(2839.17195105,4779.18991089)
\lineto(2839.17195105,4759.17419434)
\lineto(2863.0271604,4759.17419434)
\lineto(2863.0271604,4750.17341614)
\lineto(2839.17195105,4750.17341614)
\lineto(2839.17195105,4711.90437317)
\curveto(2839.17192745,4706.15560691)(2839.94821756,4702.46298365)(2841.50082373,4700.82649231)
\curveto(2843.09533966,4699.18997667)(2846.3054042,4698.37172493)(2851.13102698,4698.37173462)
\lineto(2863.0271604,4698.37173462)
\lineto(2863.0271604,4688.67858887)
\lineto(2851.13102698,4688.67858887)
\curveto(2842.19316466,4688.67858887)(2836.02480535,4690.33607317)(2832.62593054,4693.65104675)
\curveto(2829.22702162,4697.00797202)(2827.52757569,4703.09240807)(2827.52758765,4711.90437317)
\lineto(2827.52758765,4750.17341614)
\lineto(2819.03034949,4750.17341614)
\lineto(2819.03034949,4759.17419434)
\lineto(2827.52758765,4759.17419434)
\lineto(2827.52758765,4779.18991089)
\lineto(2839.17195105,4779.18991089)
}
}
{
\newrgbcolor{curcolor}{0 0 0}
\pscustom[linestyle=none,fillstyle=solid,fillcolor=curcolor]
{
\newpath
\moveto(2938.62110876,4726.82174683)
\lineto(2938.62110876,4721.15692139)
\lineto(2885.37174963,4721.15692139)
\curveto(2885.87526997,4713.1841796)(2888.26708277,4707.09974354)(2892.54719519,4702.90359497)
\curveto(2896.8692165,4698.74937958)(2902.8697293,4696.67227899)(2910.54875159,4696.67228699)
\curveto(2914.99663978,4696.67227899)(2919.29770664,4697.21778016)(2923.45196509,4698.30879211)
\curveto(2927.64807061,4699.39978481)(2931.80227177,4701.0362883)(2935.91458105,4703.2183075)
\lineto(2935.91458105,4692.26631165)
\curveto(2931.76031015,4690.50391968)(2927.50120491,4689.16114759)(2923.13725256,4688.23799133)
\curveto(2918.77318629,4687.31483596)(2914.34623454,4686.85325805)(2909.85638403,4686.85325623)
\curveto(2898.61062406,4686.85325805)(2889.69377812,4690.12626503)(2883.10581946,4696.67228699)
\curveto(2876.55978856,4703.21829296)(2873.28678158,4712.07219646)(2873.2867887,4723.23402405)
\curveto(2873.28678158,4734.77343718)(2876.39194205,4743.92107208)(2882.60227942,4750.67695618)
\curveto(2888.85454556,4757.47467791)(2897.26785196,4760.87356977)(2907.84222388,4760.87364197)
\curveto(2917.32551013,4760.87356977)(2924.81566072,4757.81037093)(2930.31269812,4751.68403625)
\curveto(2935.85156887,4745.5995372)(2938.62103632,4737.31211567)(2938.62110876,4726.82174683)
\moveto(2927.03968787,4730.22064209)
\curveto(2926.95570374,4736.55680637)(2925.17233456,4741.61318254)(2921.68957495,4745.38978577)
\curveto(2918.24866594,4749.16627557)(2913.6748485,4751.05454883)(2907.96810889,4751.05461121)
\curveto(2901.50597639,4751.05454883)(2896.32371533,4749.22921801)(2892.42131018,4745.57861328)
\curveto(2888.56081416,4741.92789475)(2886.33684788,4736.78759532)(2885.74940466,4730.15769958)
\lineto(2927.03968787,4730.22064209)
}
}
{
\newrgbcolor{curcolor}{0.99607843 0.99607843 0.80784315}
\pscustom[linestyle=none,fillstyle=solid,fillcolor=curcolor]
{
\newpath
\moveto(2873.97915625,2516.6015625)
\lineto(5077.10415625,2516.6015625)
\lineto(5077.10415625,1597.87559509)
\lineto(2873.97915625,1597.87559509)
\lineto(2873.97915625,2516.6015625)
\closepath
}
}
{
\newrgbcolor{curcolor}{0.65882355 0 0.21176471}
\pscustom[linewidth=17.578125,linecolor=curcolor]
{
\newpath
\moveto(2873.97915625,2516.6015625)
\lineto(5077.10415625,2516.6015625)
\lineto(5077.10415625,1597.87559509)
\lineto(2873.97915625,1597.87559509)
\lineto(2873.97915625,2516.6015625)
\closepath
}
}
{
\newrgbcolor{curcolor}{0.70588237 0.65490198 0.89803922}
\pscustom[linestyle=none,fillstyle=solid,fillcolor=curcolor]
{
\newpath
\moveto(3685.50253125,2329.101625)
\curveto(3685.50253125,2257.90866897)(3627.78923728,2200.195375)(3556.59628125,2200.195375)
\curveto(3485.40332522,2200.195375)(3427.69003125,2257.90866897)(3427.69003125,2329.101625)
\curveto(3427.69003125,2400.29458103)(3485.40332522,2458.007875)(3556.59628125,2458.007875)
\curveto(3627.78923728,2458.007875)(3685.50253125,2400.29458103)(3685.50253125,2329.101625)
\closepath
}
}
{
\newrgbcolor{curcolor}{0.65882355 0 0.21176471}
\pscustom[linewidth=11.71875,linecolor=curcolor]
{
\newpath
\moveto(3685.50253125,2329.101625)
\curveto(3685.50253125,2257.90866897)(3627.78923728,2200.195375)(3556.59628125,2200.195375)
\curveto(3485.40332522,2200.195375)(3427.69003125,2257.90866897)(3427.69003125,2329.101625)
\curveto(3427.69003125,2400.29458103)(3485.40332522,2458.007875)(3556.59628125,2458.007875)
\curveto(3627.78923728,2458.007875)(3685.50253125,2400.29458103)(3685.50253125,2329.101625)
\closepath
}
}
{
\newrgbcolor{curcolor}{0 0 0}
\pscustom[linestyle=none,fillstyle=solid,fillcolor=curcolor]
{
\newpath
\moveto(3509.90446875,2258.789)
\lineto(3509.90446875,2278.01525)
\lineto(3543.59596875,2278.01525)
\lineto(3543.59596875,2383.667)
\lineto(3509.90446875,2383.667)
\lineto(3509.90446875,2402.709875)
\lineto(3606.03484375,2402.709875)
\lineto(3606.03484375,2383.667)
\lineto(3572.34346875,2383.667)
\lineto(3572.34346875,2278.01525)
\lineto(3606.03484375,2278.01525)
\lineto(3606.03484375,2258.789)
\lineto(3509.90446875,2258.789)
\closepath
}
}
{
\newrgbcolor{curcolor}{0 0 0}
\pscustom[linestyle=none,fillstyle=solid,fillcolor=curcolor]
{
\newpath
\moveto(3858.49148535,2291.6796875)
\lineto(3833.56625342,2335.4876709)
\lineto(3820.10800146,2323.88336182)
\lineto(3813.85952734,2291.6796875)
\lineto(3800.74459814,2291.6796875)
\lineto(3818.73471045,2384.3081665)
\lineto(3831.84963965,2384.3081665)
\lineto(3822.8545835,2338.09692383)
\lineto(3871.1257627,2384.3081665)
\lineto(3886.57528662,2384.3081665)
\lineto(3844.34658789,2345.03204346)
\lineto(3873.87234473,2291.6796875)
\lineto(3858.49148535,2291.6796875)
}
}
{
\newrgbcolor{curcolor}{0 0 0}
\pscustom[linestyle=none,fillstyle=solid,fillcolor=curcolor]
{
\newpath
\moveto(3933.54183936,2291.6796875)
\lineto(3942.81155371,2339.40155029)
\curveto(3943.31503159,2342.05652921)(3943.56680135,2344.29956896)(3943.56686377,2346.13067627)
\curveto(3943.56680135,2348.41943789)(3943.26925526,2350.29626707)(3942.67422461,2351.76116943)
\curveto(3942.07907091,2353.27172796)(3941.23220896,2354.46191232)(3940.13363623,2355.33172607)
\curveto(3939.03494553,2356.24718885)(3937.73032037,2356.88805735)(3936.21975684,2357.2543335)
\curveto(3934.70908316,2357.62047849)(3933.06113559,2357.80358378)(3931.27590918,2357.8036499)
\curveto(3928.57505609,2357.80358378)(3925.94291761,2357.27715608)(3923.37948584,2356.22436523)
\curveto(3920.81596961,2355.21722162)(3918.45848906,2353.72949117)(3916.30703711,2351.76116943)
\curveto(3914.20129117,2349.83850385)(3912.37023831,2347.43524698)(3910.81387305,2344.5513916)
\curveto(3909.25744846,2341.7132068)(3908.11304042,2338.44019982)(3907.38064551,2334.73236084)
\lineto(3899.00357031,2291.6796875)
\lineto(3886.64395117,2291.6796875)
\lineto(3897.97360205,2350.11322021)
\curveto(3898.29401948,2351.66955671)(3898.59156557,2353.29461612)(3898.86624121,2354.98840332)
\curveto(3899.18665775,2356.72784023)(3899.46131567,2358.35289964)(3899.69021582,2359.86358643)
\curveto(3899.96485521,2361.41991317)(3900.19373682,2362.74742649)(3900.37686133,2363.84613037)
\curveto(3900.55994739,2364.94468992)(3900.65150003,2365.6542229)(3900.65151953,2365.97473145)
\lineto(3912.32449316,2365.97473145)
\curveto(3912.32446199,2365.74577554)(3912.25579751,2365.1277952)(3912.11849951,2364.12078857)
\curveto(3911.98113958,2363.11363706)(3911.82092245,2361.94634087)(3911.63784766,2360.61889648)
\curveto(3911.45471188,2359.33709055)(3911.24871844,2358.00957722)(3911.0198667,2356.63635254)
\curveto(3910.79095522,2355.30877426)(3910.58496178,2354.18725439)(3910.40188574,2353.27178955)
\lineto(3910.60787939,2353.27178955)
\curveto(3913.72063979,2357.62047849)(3917.40563367,2361.0537026)(3921.66287207,2363.57147217)
\curveto(3925.92002945,2366.08909795)(3931.06986561,2367.34794679)(3937.112396,2367.34802246)
\curveto(3943.38369607,2367.34794679)(3948.19020982,2365.95176899)(3951.53195166,2363.15948486)
\curveto(3954.87355274,2360.4128341)(3956.54438847,2355.99541908)(3956.54446387,2349.90722656)
\curveto(3956.54438847,2348.57965501)(3956.45283583,2347.18347721)(3956.26980566,2345.71868896)
\curveto(3956.13240158,2344.29956896)(3955.90351997,2342.78895036)(3955.58316016,2341.18682861)
\lineto(3945.97012305,2291.6796875)
\lineto(3933.54183936,2291.6796875)
}
}
{
\newrgbcolor{curcolor}{0 0 0}
\pscustom[linestyle=none,fillstyle=solid,fillcolor=curcolor]
{
\newpath
\moveto(4040.79586768,2328.89587402)
\curveto(4039.55983167,2322.39559917)(4037.70589066,2316.71933531)(4035.23403906,2311.86706543)
\curveto(4032.76204794,2307.06053149)(4029.76369889,2303.05510337)(4026.23898291,2299.85076904)
\curveto(4022.75992171,2296.64641837)(4018.77738175,2294.2431615)(4014.29135107,2292.64099121)
\curveto(4009.85099908,2291.08459532)(4005.02159717,2290.30639786)(3999.80313086,2290.30639648)
\curveto(3995.6832276,2290.30639786)(3991.97534557,2290.83282555)(3988.67947363,2291.88568115)
\curveto(3985.38355528,2292.98431266)(3982.56831152,2294.63226023)(3980.23373389,2296.82952881)
\curveto(3977.94490305,2299.02678709)(3976.18251468,2301.79625453)(3974.94656348,2305.13793945)
\curveto(3973.75636964,2308.47959746)(3973.16127747,2312.39347294)(3973.16128516,2316.87957764)
\curveto(3973.16127747,2320.49588183)(3973.57326436,2324.50130995)(3974.39724707,2328.89587402)
\curveto(3976.91493582,2341.66743048)(3981.67567325,2351.25756982)(3988.67947363,2357.6663208)
\curveto(3995.72900392,2364.12071613)(4004.56383395,2367.34794679)(4015.18399023,2367.34802246)
\curveto(4024.29342848,2367.34794679)(4031.06832405,2365.24223601)(4035.50869727,2361.03088379)
\curveto(4039.9489304,2356.81939287)(4042.16908199,2350.45648419)(4042.16915869,2341.94213867)
\curveto(4042.16908199,2339.97370659)(4042.05464119,2337.91377213)(4041.82583594,2335.7623291)
\curveto(4041.6426543,2333.61079791)(4041.29933189,2331.32198184)(4040.79586768,2328.89587402)
\moveto(4027.81826758,2328.89587402)
\curveto(4028.27596844,2331.27620552)(4028.61929086,2333.47346895)(4028.84823584,2335.4876709)
\curveto(4029.07705407,2337.50178523)(4029.19149487,2339.37861441)(4029.19155859,2341.11816406)
\curveto(4029.19149487,2344.41400977)(4028.80239614,2347.13770089)(4028.02426123,2349.28924561)
\curveto(4027.29177753,2351.48645142)(4026.23892214,2353.2488398)(4024.86569189,2354.57641602)
\curveto(4023.49234286,2355.90386644)(4021.84439529,2356.84228103)(4019.92184424,2357.3916626)
\curveto(4018.04496061,2357.94091274)(4015.93924982,2358.21557067)(4013.60470557,2358.21563721)
\curveto(4010.62919654,2358.21557067)(4007.79106461,2357.7349193)(4005.09030127,2356.77368164)
\curveto(4002.38945869,2355.85809012)(3999.89464917,2354.25591887)(3997.60586523,2351.96716309)
\curveto(3995.36279335,2349.72406305)(3993.37152337,2346.725714)(3991.63204932,2342.97210693)
\curveto(3989.89252294,2339.21839729)(3988.47345698,2334.52632434)(3987.37484717,2328.89587402)
\curveto(3986.91706205,2326.65279706)(3986.5966278,2324.56997443)(3986.41354346,2322.6473999)
\curveto(3986.23041723,2320.72476344)(3986.13886459,2318.9394869)(3986.13888525,2317.29156494)
\curveto(3986.13886459,2313.85831522)(3986.52796332,2310.99729514)(3987.30618262,2308.70849609)
\curveto(3988.13013457,2306.41966299)(3989.22876628,2304.58861014)(3990.60208105,2303.21533203)
\curveto(3992.02112189,2301.84203085)(3993.6461813,2300.85783994)(3995.47726416,2300.26275635)
\curveto(3997.35406333,2299.71343191)(3999.34533331,2299.43877398)(4001.45108008,2299.43878174)
\curveto(4004.42650499,2299.43877398)(4007.26463692,2299.87364903)(4009.96548437,2300.7434082)
\curveto(4012.66624284,2301.65892557)(4015.1381642,2303.23820866)(4017.38125586,2305.48126221)
\curveto(4019.67002002,2307.72428815)(4021.70706632,2310.74552537)(4023.49240088,2314.54498291)
\curveto(4025.27761939,2318.34439472)(4026.71957352,2323.12802031)(4027.81826758,2328.89587402)
}
}
{
\newrgbcolor{curcolor}{0 0 0}
\pscustom[linestyle=none,fillstyle=solid,fillcolor=curcolor]
{
\newpath
\moveto(4112.48165869,2291.6796875)
\lineto(4098.47409033,2291.6796875)
\lineto(4094.97219824,2324.02069092)
\curveto(4094.8348241,2325.25661926)(4094.69749513,2326.65279706)(4094.56021094,2328.20922852)
\curveto(4094.4228372,2329.81136324)(4094.28550824,2331.34487)(4094.14822363,2332.80975342)
\curveto(4093.96507399,2334.50343618)(4093.82774503,2336.19716007)(4093.73623633,2337.89093018)
\curveto(4093.00377124,2336.19716007)(4092.22557378,2334.52632434)(4091.4016416,2332.87841797)
\curveto(4090.71495517,2331.45931081)(4089.98253403,2329.92580404)(4089.20437598,2328.27789307)
\curveto(4088.4261391,2326.67568522)(4087.6708298,2325.16506661)(4086.9384458,2323.74603271)
\lineto(4070.52761816,2291.6796875)
\lineto(4056.58871436,2291.6796875)
\lineto(4058.44265723,2365.97473145)
\lineto(4070.66494727,2365.97473145)
\lineto(4068.9483335,2324.29534912)
\curveto(4068.90253802,2323.19668479)(4068.81098537,2321.66317802)(4068.67367529,2319.69482422)
\curveto(4068.53632745,2317.7721907)(4068.37611032,2315.71225624)(4068.19302344,2313.51501465)
\curveto(4068.00989975,2311.31772939)(4067.8039063,2309.16624228)(4067.57504248,2307.06054688)
\curveto(4067.39191941,2304.95482071)(4067.1630378,2303.23820866)(4066.88839697,2301.91070557)
\curveto(4067.84968262,2303.9706298)(4068.74232089,2305.93901162)(4069.56631445,2307.81585693)
\curveto(4070.29871582,2309.41801205)(4071.05402512,2311.04307146)(4071.83224463,2312.69104004)
\curveto(4072.61042005,2314.38474292)(4073.25128855,2315.68936808)(4073.75485205,2316.60491943)
\lineto(4088.65505957,2345.71868896)
\lineto(4101.90731787,2345.71868896)
\lineto(4105.13455176,2316.53625488)
\curveto(4105.22604906,2315.6207036)(4105.3176017,2314.3389666)(4105.40920996,2312.69104004)
\curveto(4105.54648331,2311.08884778)(4105.66092411,2309.48667653)(4105.75253271,2307.88452148)
\curveto(4105.88980572,2306.0076761)(4106.00424652,2304.01640612)(4106.09585547,2301.91070557)
\curveto(4106.37045709,2303.28398498)(4106.80533215,2305.00059703)(4107.40048193,2307.06054688)
\curveto(4107.9955165,2309.12046596)(4108.636385,2311.22617674)(4109.32308936,2313.37768555)
\curveto(4110.00967464,2315.57492728)(4110.6734313,2317.6577499)(4111.31436133,2319.62615967)
\curveto(4111.9551683,2321.59451354)(4112.45870784,2323.15090847)(4112.82498145,2324.29534912)
\lineto(4127.86251807,2365.97473145)
\lineto(4139.947479,2365.97473145)
\lineto(4112.48165869,2291.6796875)
}
}
{
\newrgbcolor{curcolor}{0 0 0}
\pscustom[linestyle=none,fillstyle=solid,fillcolor=curcolor]
{
\newpath
\moveto(4179.36093115,2301.43005371)
\lineto(4205.45346045,2301.43005371)
\lineto(4203.53085303,2291.6796875)
\lineto(4136.10226416,2291.6796875)
\lineto(4138.02487158,2301.43005371)
\lineto(4167.00131201,2301.43005371)
\lineto(4183.00015234,2383.82751465)
\lineto(4163.15609717,2383.82751465)
\lineto(4165.01004004,2393.57788086)
\lineto(4197.21371436,2393.57788086)
\lineto(4179.36093115,2301.43005371)
}
}
{
\newrgbcolor{curcolor}{0 0 0}
\pscustom[linestyle=none,fillstyle=solid,fillcolor=curcolor]
{
\newpath
\moveto(4240.06039404,2326.21795654)
\curveto(4239.78571449,2324.75307972)(4239.55683288,2323.33401376)(4239.37374854,2321.96075439)
\curveto(4239.23639863,2320.58743447)(4239.16773415,2319.23703299)(4239.16775488,2317.9095459)
\curveto(4239.16773415,2315.16294038)(4239.4881684,2312.64524271)(4240.12905859,2310.35644531)
\curveto(4240.81568172,2308.11338689)(4241.84564895,2306.19078139)(4243.21896338,2304.58862305)
\curveto(4244.59222823,2302.98643889)(4246.28595213,2301.75047821)(4248.30014014,2300.8807373)
\curveto(4250.36004473,2300.010978)(4252.76330161,2299.57610294)(4255.50991797,2299.57611084)
\curveto(4258.02757857,2299.57610294)(4260.40794728,2299.85076087)(4262.65103125,2300.40008545)
\curveto(4264.89402678,2300.99516891)(4266.95396124,2301.79625453)(4268.83084082,2302.80334473)
\curveto(4270.75339592,2303.81041267)(4272.42423165,2305.00059703)(4273.84335303,2306.37390137)
\curveto(4275.26236358,2307.79295264)(4276.40677161,2309.3264594)(4277.27658057,2310.97442627)
\lineto(4287.50759863,2307.88452148)
\curveto(4286.36312152,2305.82457082)(4284.85250291,2303.74174819)(4282.97573828,2301.63604736)
\curveto(4281.09884456,2299.57610294)(4278.78714033,2297.69927377)(4276.04061865,2296.0055542)
\curveto(4273.29398176,2294.3576023)(4270.0667511,2292.98431266)(4266.35891699,2291.88568115)
\curveto(4262.69676335,2290.83282555)(4258.48534178,2290.30639786)(4253.72463965,2290.30639648)
\curveto(4249.37585382,2290.30639786)(4245.50775466,2290.85571371)(4242.12033057,2291.9543457)
\curveto(4238.73285909,2293.09875346)(4235.871839,2294.79247736)(4233.53726172,2297.03552246)
\curveto(4231.20265422,2299.27855685)(4229.41737769,2302.07091246)(4228.18142676,2305.41259766)
\curveto(4226.99123265,2308.75425539)(4226.39614047,2312.64524271)(4226.39614844,2317.08557129)
\curveto(4226.39614047,2320.88498056)(4226.83101553,2324.95907317)(4227.7007749,2329.30786133)
\curveto(4229.07405527,2336.44892984)(4231.13398974,2342.44562794)(4233.88058447,2347.29797363)
\curveto(4236.62714831,2352.15020808)(4239.80860265,2356.06408356)(4243.42495703,2359.03961182)
\curveto(4247.08703775,2362.01500535)(4251.04668955,2364.14360429)(4255.30392432,2365.42541504)
\curveto(4259.56108533,2366.70707829)(4263.88694771,2367.34794679)(4268.28152441,2367.34802246)
\curveto(4273.04221199,2367.34794679)(4277.09341644,2366.72996645)(4280.4351499,2365.49407959)
\curveto(4283.77675936,2364.30382142)(4286.50045049,2362.56432121)(4288.60623145,2360.27557373)
\curveto(4290.75764838,2358.03246539)(4292.3140433,2355.30877426)(4293.2754209,2352.10449219)
\curveto(4294.28242513,2348.90008926)(4294.78596466,2345.30664803)(4294.78604102,2341.32415771)
\curveto(4294.78596466,2337.15846282)(4294.30531329,2332.67238332)(4293.34408545,2327.86590576)
\lineto(4293.06942725,2326.21795654)
\lineto(4240.06039404,2326.21795654)
\moveto(4266.6335752,2358.21563721)
\curveto(4264.52781621,2358.21557067)(4262.2618883,2357.87224826)(4259.83578467,2357.18566895)
\curveto(4257.40959823,2356.54473494)(4255.05211768,2355.37743874)(4252.76333594,2353.68377686)
\curveto(4250.47448553,2351.98999096)(4248.36877475,2349.67828673)(4246.44619727,2346.74865723)
\curveto(4244.56934007,2343.86469391)(4243.10449779,2340.17970004)(4242.05166602,2335.69366455)
\lineto(4282.15176367,2335.69366455)
\curveto(4282.24325259,2336.60914697)(4282.31191708,2337.45600891)(4282.35775732,2338.23425293)
\curveto(4282.40346972,2339.05818016)(4282.42635788,2339.85926579)(4282.42642187,2340.63751221)
\curveto(4282.42635788,2343.97913471)(4282.01437099,2346.77149032)(4281.19045996,2349.0145874)
\curveto(4280.36642342,2351.25756982)(4279.24490354,2353.06573451)(4277.82589697,2354.43908691)
\curveto(4276.40677161,2355.8123138)(4274.73593588,2356.77361655)(4272.81338477,2357.32299805)
\curveto(4270.89072488,2357.91802458)(4268.83079042,2358.21557067)(4266.6335752,2358.21563721)
}
}
{
\newrgbcolor{curcolor}{0 0 0}
\pscustom[linestyle=none,fillstyle=solid,fillcolor=curcolor]
{
\newpath
\moveto(4357.06478857,2303.62731934)
\curveto(4353.90616808,2298.82079364)(4350.24406237,2295.31890505)(4346.07846045,2293.12164307)
\curveto(4341.95854819,2290.97015452)(4337.12914629,2289.89441096)(4331.59024023,2289.89440918)
\curveto(4324.35755261,2289.89441096)(4319.02461117,2291.7941283)(4315.5913999,2295.59356689)
\curveto(4312.15816296,2299.43877398)(4310.4415509,2305.2523668)(4310.44155859,2313.03436279)
\curveto(4310.4415509,2317.47464462)(4310.99086676,2322.53292813)(4312.08950781,2328.20922852)
\curveto(4317.17067015,2354.1185899)(4328.66052682,2367.07328886)(4346.5591123,2367.07336426)
\curveto(4352.09800339,2367.07328886)(4356.53830656,2366.08909795)(4359.88003516,2364.12078857)
\curveto(4363.26742581,2362.15233431)(4365.64779453,2358.92510365)(4367.02114844,2354.43908691)
\lineto(4367.15847754,2354.43908691)
\curveto(4367.20418945,2354.66790576)(4367.2957421,2355.26299794)(4367.43313574,2356.22436523)
\curveto(4367.61617635,2357.18560344)(4367.79928163,2358.23845883)(4367.98245215,2359.38293457)
\curveto(4368.21126852,2360.5272749)(4368.41726197,2361.60301846)(4368.60043311,2362.61016846)
\curveto(4368.78347254,2363.66295292)(4368.92080151,2364.39537406)(4369.01242041,2364.80743408)
\lineto(4374.57424902,2393.57788086)
\lineto(4386.93386816,2393.57788086)
\lineto(4370.11105322,2306.99188232)
\curveto(4369.79055161,2305.20659048)(4369.47011736,2303.48997842)(4369.14974951,2301.84204102)
\curveto(4368.87502518,2300.19408328)(4368.62325542,2298.68346468)(4368.39443945,2297.31018066)
\curveto(4368.1654922,2295.93688539)(4367.98238692,2294.74670103)(4367.84512305,2293.73962402)
\curveto(4367.70772899,2292.77831921)(4367.61617635,2292.09167439)(4367.57046484,2291.6796875)
\lineto(4355.76016211,2291.6796875)
\curveto(4355.71433278,2291.86279279)(4355.73722094,2292.36633232)(4355.82882666,2293.19030762)
\curveto(4355.92032623,2294.01427989)(4356.03476703,2294.9984708)(4356.17214941,2296.1428833)
\curveto(4356.30942496,2297.28728687)(4356.49253024,2298.52324755)(4356.72146582,2299.85076904)
\curveto(4356.95029346,2301.17827419)(4357.17917506,2302.43712303)(4357.40811133,2303.62731934)
\lineto(4357.06478857,2303.62731934)
\moveto(4324.99844336,2328.62121582)
\curveto(4324.49488158,2325.87459959)(4324.10578284,2323.35690192)(4323.831146,2321.06811523)
\curveto(4323.55646699,2318.8250461)(4323.41913802,2316.74222347)(4323.41915869,2314.81964111)
\curveto(4323.41913802,2309.46378837)(4324.47199342,2305.54991289)(4326.57772803,2303.07800293)
\curveto(4328.72919131,2300.6518465)(4332.09375093,2299.43877398)(4336.67141699,2299.43878174)
\curveto(4339.82994925,2299.43877398)(4342.73674566,2299.98808984)(4345.39181494,2301.08673096)
\curveto(4348.04679894,2302.18535326)(4350.42716766,2303.92485348)(4352.53292822,2306.30523682)
\curveto(4354.68436555,2308.73136723)(4356.56119473,2311.86704524)(4358.16342139,2315.71228027)
\curveto(4359.76553722,2319.55746724)(4361.11593871,2324.22665202)(4362.21462988,2329.71984863)
\curveto(4362.67233363,2332.05440298)(4363.01565604,2334.20589009)(4363.24459814,2336.17431641)
\curveto(4363.47341926,2338.14265373)(4363.58786006,2339.97370659)(4363.5879209,2341.66748047)
\curveto(4363.58786006,2344.64289137)(4363.24453765,2347.16058905)(4362.55795264,2349.22058105)
\curveto(4361.87124801,2351.3262343)(4360.86416894,2353.01995819)(4359.5367124,2354.30175781)
\curveto(4358.25491862,2355.58343219)(4356.65274737,2356.49895862)(4354.73019385,2357.04833984)
\curveto(4352.85331269,2357.64336665)(4350.67893742,2357.94091274)(4348.20706152,2357.940979)
\curveto(4345.2773315,2357.94091274)(4342.59941669,2357.41448505)(4340.17330908,2356.36169434)
\curveto(4337.79290295,2355.30877426)(4335.64141584,2353.61505037)(4333.71884131,2351.28051758)
\curveto(4331.84198116,2348.94586559)(4330.17114543,2345.90174021)(4328.7063291,2342.14813232)
\curveto(4327.28723718,2338.44019982)(4326.0512765,2333.93123216)(4324.99844336,2328.62121582)
}
}
{
\newrgbcolor{curcolor}{0 0 0}
\pscustom[linestyle=none,fillstyle=solid,fillcolor=curcolor]
{
\newpath
\moveto(4416.459625,2262.56591797)
\curveto(4412.33972668,2262.56594708)(4408.79206178,2263.00082214)(4405.81661963,2263.87054443)
\curveto(4402.84113999,2264.74032235)(4400.36921864,2265.99917119)(4398.40084814,2267.64709473)
\curveto(4396.47823132,2269.24929001)(4395.03627719,2271.17189551)(4394.07498145,2273.41491699)
\curveto(4393.11367169,2275.65797501)(4392.587244,2278.15278452)(4392.49569678,2280.89935303)
\lineto(4405.47329687,2282.6159668)
\curveto(4405.56483112,2279.18275175)(4406.73212731,2276.55061327)(4408.97518896,2274.71954346)
\curveto(4411.21820681,2272.84273124)(4414.44543747,2271.90431665)(4418.65689062,2271.90429688)
\curveto(4421.26610936,2271.90431665)(4423.69225439,2272.29341538)(4425.93533301,2273.07159424)
\curveto(4428.17833389,2273.84981031)(4430.19249203,2275.10865915)(4431.97781348,2276.84814453)
\curveto(4433.80882143,2278.54188326)(4435.38810452,2280.76203484)(4436.71566748,2283.50860596)
\curveto(4438.08890748,2286.25519341)(4439.14176287,2289.59686488)(4439.87423682,2293.53363037)
\curveto(4440.10306562,2294.72381287)(4440.35483539,2296.07421436)(4440.62954687,2297.58483887)
\curveto(4440.94992757,2299.09545157)(4441.24747366,2300.53740569)(4441.52218604,2301.91070557)
\curveto(4441.84256583,2303.51286658)(4442.16300008,2305.16081416)(4442.48348975,2306.85455322)
\lineto(4442.34616064,2306.85455322)
\curveto(4441.11014469,2305.02348519)(4439.64530241,2303.23820866)(4437.95162939,2301.49871826)
\curveto(4436.30363094,2299.75920823)(4434.38102544,2298.22570146)(4432.18380713,2296.89819336)
\curveto(4430.03227491,2295.57067482)(4427.60612988,2294.49493127)(4424.90536475,2293.67095947)
\curveto(4422.25030027,2292.89276002)(4419.27483938,2292.50366129)(4415.97897314,2292.50366211)
\curveto(4408.8378381,2292.50366129)(4403.5964493,2294.28893782)(4400.25479102,2297.85949707)
\curveto(4396.95888269,2301.47582028)(4395.31093512,2306.92320253)(4395.31094336,2314.20166016)
\curveto(4395.31093512,2316.39890106)(4395.44826408,2318.75638161)(4395.72293066,2321.27410889)
\curveto(4396.04335626,2323.79177697)(4396.47823132,2326.49257993)(4397.02755713,2329.37652588)
\curveto(4398.17195521,2335.37318629)(4399.68257381,2340.70612773)(4401.55941748,2345.37536621)
\curveto(4403.43623217,2350.0444973)(4405.77082456,2353.98126094)(4408.56320166,2357.18566895)
\curveto(4411.4013121,2360.43572226)(4414.7200954,2362.90764362)(4418.51956152,2364.60144043)
\curveto(4422.36474107,2366.2950914)(4426.80504425,2367.14195335)(4431.84048437,2367.14202881)
\curveto(4437.01316393,2367.14195335)(4441.27036182,2365.92888083)(4444.61209082,2363.50280762)
\curveto(4447.95370474,2361.12236708)(4450.21963265,2357.71203114)(4451.40988135,2353.27178955)
\lineto(4451.615875,2353.27178955)
\curveto(4451.84469206,2354.41613599)(4452.11934999,2355.69787299)(4452.43984961,2357.11700439)
\curveto(4452.80599481,2358.53600492)(4453.14931722,2359.86351824)(4453.46981787,2361.09954834)
\curveto(4453.79018572,2362.38121592)(4454.08773181,2363.47984763)(4454.36245703,2364.39544678)
\curveto(4454.63704767,2365.31090049)(4454.82015295,2365.83732819)(4454.91177344,2365.97473145)
\lineto(4466.65341162,2365.97473145)
\curveto(4466.5617794,2365.56267026)(4466.40156227,2364.85313728)(4466.17275977,2363.84613037)
\curveto(4465.94379906,2362.88475546)(4465.66914113,2361.71745926)(4465.34878516,2360.34423828)
\curveto(4465.02827263,2358.97087997)(4464.68495022,2357.43737321)(4464.31881689,2355.74371338)
\curveto(4463.9983054,2354.09570174)(4463.65498299,2352.37908969)(4463.28884863,2350.59387207)
\lineto(4452.23385596,2293.87695312)
\curveto(4450.21963265,2283.48572597)(4446.28286901,2275.68086317)(4440.42355322,2270.46234131)
\curveto(4434.56413073,2265.19808556)(4426.57616264,2262.56594708)(4416.459625,2262.56591797)
\moveto(4446.878021,2329.51385498)
\curveto(4447.61038233,2333.17592286)(4447.9765929,2336.51759432)(4447.97665381,2339.53887939)
\curveto(4447.9765929,2342.83472668)(4447.56460601,2345.64997044)(4446.74069189,2347.98461914)
\curveto(4445.96243476,2350.31915523)(4444.88669121,2352.21887257)(4443.51345801,2353.68377686)
\curveto(4442.18588824,2355.19433346)(4440.62949332,2356.27007701)(4438.84426855,2356.91101074)
\curveto(4437.05894025,2357.59759033)(4435.15922291,2357.94091274)(4433.14511084,2357.940979)
\curveto(4430.07805123,2357.94091274)(4427.30858379,2357.41448505)(4424.8367002,2356.36169434)
\curveto(4422.4105174,2355.35455058)(4420.23614213,2353.70660301)(4418.31356787,2351.41784668)
\curveto(4416.43670745,2349.17474719)(4414.78875988,2346.24506262)(4413.36972021,2342.62878418)
\curveto(4411.99640427,2339.05818016)(4410.80621992,2334.68654147)(4409.79916357,2329.51385498)
\curveto(4409.29560131,2326.81301418)(4408.90650258,2324.34109283)(4408.63186621,2322.0980835)
\curveto(4408.35718672,2319.90078965)(4408.21985776,2317.88663151)(4408.21987891,2316.05560303)
\curveto(4408.21985776,2313.35477569)(4408.49451569,2311.08884778)(4409.04385352,2309.2578125)
\curveto(4409.63892372,2307.42674207)(4410.50867383,2305.93901162)(4411.65310645,2304.7946167)
\curveto(4412.84326622,2303.65019555)(4414.28522034,2302.82622176)(4415.97897314,2302.32269287)
\curveto(4417.71844445,2301.86491901)(4419.70971443,2301.63603741)(4421.95278906,2301.63604736)
\curveto(4424.37889922,2301.63603741)(4426.85082057,2302.11668878)(4429.36856055,2303.07800293)
\curveto(4431.88621593,2304.0850706)(4434.26658464,2305.68724185)(4436.50967383,2307.88452148)
\curveto(4438.79844046,2310.08176871)(4440.83548676,2312.9427888)(4442.62081885,2316.46759033)
\curveto(4444.45181615,2320.03811861)(4445.87088212,2324.38686915)(4446.878021,2329.51385498)
}
}
{
\newrgbcolor{curcolor}{0 0 0}
\pscustom[linestyle=none,fillstyle=solid,fillcolor=curcolor]
{
\newpath
\moveto(4493.02059912,2326.21795654)
\curveto(4492.74591956,2324.75307972)(4492.51703796,2323.33401376)(4492.33395361,2321.96075439)
\curveto(4492.19660371,2320.58743447)(4492.12793922,2319.23703299)(4492.12795996,2317.9095459)
\curveto(4492.12793922,2315.16294038)(4492.44837347,2312.64524271)(4493.08926367,2310.35644531)
\curveto(4493.7758868,2308.11338689)(4494.80585403,2306.19078139)(4496.17916846,2304.58862305)
\curveto(4497.55243331,2302.98643889)(4499.2461572,2301.75047821)(4501.26034521,2300.8807373)
\curveto(4503.32024981,2300.010978)(4505.72350668,2299.57610294)(4508.47012305,2299.57611084)
\curveto(4510.98778365,2299.57610294)(4513.36815236,2299.85076087)(4515.61123633,2300.40008545)
\curveto(4517.85423186,2300.99516891)(4519.91416632,2301.79625453)(4521.7910459,2302.80334473)
\curveto(4523.713601,2303.81041267)(4525.38443673,2305.00059703)(4526.80355811,2306.37390137)
\curveto(4528.22256866,2307.79295264)(4529.36697669,2309.3264594)(4530.23678564,2310.97442627)
\lineto(4540.46780371,2307.88452148)
\curveto(4539.3233266,2305.82457082)(4537.81270799,2303.74174819)(4535.93594336,2301.63604736)
\curveto(4534.05904964,2299.57610294)(4531.74734541,2297.69927377)(4529.00082373,2296.0055542)
\curveto(4526.25418684,2294.3576023)(4523.02695618,2292.98431266)(4519.31912207,2291.88568115)
\curveto(4515.65696843,2290.83282555)(4511.44554686,2290.30639786)(4506.68484473,2290.30639648)
\curveto(4502.3360589,2290.30639786)(4498.46795974,2290.85571371)(4495.08053564,2291.9543457)
\curveto(4491.69306417,2293.09875346)(4488.83204408,2294.79247736)(4486.4974668,2297.03552246)
\curveto(4484.1628593,2299.27855685)(4482.37758276,2302.07091246)(4481.14163184,2305.41259766)
\curveto(4479.95143773,2308.75425539)(4479.35634555,2312.64524271)(4479.35635352,2317.08557129)
\curveto(4479.35634555,2320.88498056)(4479.7912206,2324.95907317)(4480.66097998,2329.30786133)
\curveto(4482.03426035,2336.44892984)(4484.09419482,2342.44562794)(4486.84078955,2347.29797363)
\curveto(4489.58735339,2352.15020808)(4492.76880772,2356.06408356)(4496.38516211,2359.03961182)
\curveto(4500.04724283,2362.01500535)(4504.00689463,2364.14360429)(4508.26412939,2365.42541504)
\curveto(4512.52129041,2366.70707829)(4516.84715279,2367.34794679)(4521.24172949,2367.34802246)
\curveto(4526.00241707,2367.34794679)(4530.05362151,2366.72996645)(4533.39535498,2365.49407959)
\curveto(4536.73696444,2364.30382142)(4539.46065556,2362.56432121)(4541.56643652,2360.27557373)
\curveto(4543.71785345,2358.03246539)(4545.27424838,2355.30877426)(4546.23562598,2352.10449219)
\curveto(4547.2426302,2348.90008926)(4547.74616974,2345.30664803)(4547.74624609,2341.32415771)
\curveto(4547.74616974,2337.15846282)(4547.26551836,2332.67238332)(4546.30429053,2327.86590576)
\lineto(4546.02963232,2326.21795654)
\lineto(4493.02059912,2326.21795654)
\moveto(4519.59378027,2358.21563721)
\curveto(4517.48802129,2358.21557067)(4515.22209338,2357.87224826)(4512.79598975,2357.18566895)
\curveto(4510.36980331,2356.54473494)(4508.01232275,2355.37743874)(4505.72354102,2353.68377686)
\curveto(4503.43469061,2351.98999096)(4501.32897983,2349.67828673)(4499.40640234,2346.74865723)
\curveto(4497.52954515,2343.86469391)(4496.06470287,2340.17970004)(4495.01187109,2335.69366455)
\lineto(4535.11196875,2335.69366455)
\curveto(4535.20345767,2336.60914697)(4535.27212215,2337.45600891)(4535.3179624,2338.23425293)
\curveto(4535.3636748,2339.05818016)(4535.38656296,2339.85926579)(4535.38662695,2340.63751221)
\curveto(4535.38656296,2343.97913471)(4534.97457607,2346.77149032)(4534.15066504,2349.0145874)
\curveto(4533.32662849,2351.25756982)(4532.20510862,2353.06573451)(4530.78610205,2354.43908691)
\curveto(4529.36697669,2355.8123138)(4527.69614096,2356.77361655)(4525.77358984,2357.32299805)
\curveto(4523.85092996,2357.91802458)(4521.7909955,2358.21557067)(4519.59378027,2358.21563721)
}
}
{
\newrgbcolor{curcolor}{0 0 0}
\pscustom[linestyle=none,fillstyle=solid,fillcolor=curcolor]
{
\newpath
\moveto(2885.69790625,2141.6015625)
\lineto(5065.38540625,2141.6015625)
}
}
{
\newrgbcolor{curcolor}{0.65882355 0 0.21176471}
\pscustom[linewidth=17.578125,linecolor=curcolor]
{
\newpath
\moveto(2885.69790625,2141.6015625)
\lineto(5065.38540625,2141.6015625)
}
}
{
\newrgbcolor{curcolor}{0.51764709 0.74509805 0.51764709}
\pscustom[linestyle=none,fillstyle=solid,fillcolor=curcolor]
{
\newpath
\moveto(3038.04159375,2002.15277246)
\curveto(3038.04159375,1982.73651172)(3022.30160449,1966.99652246)(3002.88534375,1966.99652246)
\curveto(2983.46908301,1966.99652246)(2967.72909375,1982.73651172)(2967.72909375,2002.15277246)
\curveto(2967.72909375,2021.5690332)(2983.46908301,2037.30902246)(3002.88534375,2037.30902246)
\curveto(3022.30160449,2037.30902246)(3038.04159375,2021.5690332)(3038.04159375,2002.15277246)
\closepath
}
}
{
\newrgbcolor{curcolor}{0.01176471 0.50196081 0.28235295}
\pscustom[linewidth=11.71875,linecolor=curcolor]
{
\newpath
\moveto(3038.04159375,2002.15277246)
\curveto(3038.04159375,1982.73651172)(3022.30160449,1966.99652246)(3002.88534375,1966.99652246)
\curveto(2983.46908301,1966.99652246)(2967.72909375,1982.73651172)(2967.72909375,2002.15277246)
\curveto(2967.72909375,2021.5690332)(2983.46908301,2037.30902246)(3002.88534375,2037.30902246)
\curveto(3022.30160449,2037.30902246)(3038.04159375,2021.5690332)(3038.04159375,2002.15277246)
\closepath
}
}
{
\newrgbcolor{curcolor}{0 0 0}
\pscustom[linestyle=none,fillstyle=solid,fillcolor=curcolor]
{
\newpath
\moveto(3166.89068579,2008.30665588)
\curveto(3166.89062725,2016.69894541)(3165.14921969,2023.20299774)(3161.66645789,2027.8188324)
\curveto(3158.22555108,2032.43455589)(3153.37898305,2034.74244543)(3147.12673926,2034.74250793)
\curveto(3140.91637955,2034.74244543)(3136.06981152,2032.43455589)(3132.58702063,2027.8188324)
\curveto(3129.1461429,2023.20299774)(3127.42571616,2016.69894541)(3127.42573523,2008.30665588)
\curveto(3127.42571616,1999.95625585)(3129.1461429,1993.47318433)(3132.58702063,1988.85742188)
\curveto(3136.06981152,1984.24162618)(3140.91637955,1981.93373664)(3147.12673926,1981.93374634)
\curveto(3153.37898305,1981.93373664)(3158.22555108,1984.24162618)(3161.66645789,1988.85742188)
\curveto(3165.14921969,1993.47318433)(3166.89062725,1999.95625585)(3166.89068579,2008.30665588)
\moveto(3178.47210669,1980.98960876)
\curveto(3178.47203657,1968.98857442)(3175.8074732,1960.07172848)(3170.47840857,1954.23904419)
\curveto(3165.14921969,1948.36443428)(3156.98768306,1945.42712032)(3145.99377417,1945.42709351)
\curveto(3141.92345862,1945.42712032)(3138.08396966,1945.74183253)(3134.47529578,1946.37123108)
\curveto(3130.86656965,1946.95871974)(3127.36277372,1947.88187556)(3123.96389746,1949.14070129)
\lineto(3123.96389746,1960.40740967)
\curveto(3127.36277372,1958.56110987)(3130.71970395,1957.19735696)(3134.03469824,1956.31614685)
\curveto(3137.34964117,1955.43496859)(3140.72755222,1954.99437149)(3144.16844153,1954.99435425)
\curveto(3151.76346037,1954.99437149)(3157.44926096,1956.98754882)(3161.22586035,1960.97389221)
\curveto(3165.002354,1964.91829651)(3166.89062725,1970.89782849)(3166.89068579,1978.9125061)
\lineto(3166.89068579,1984.64027405)
\curveto(3164.49881446,1980.48606048)(3161.43561562,1977.38090001)(3157.70108008,1975.32478333)
\curveto(3153.96644584,1973.26866047)(3149.49753246,1972.24060059)(3144.29432654,1972.24060059)
\curveto(3135.65019524,1972.24060059)(3128.684565,1975.53458838)(3123.39741492,1982.12257385)
\curveto(3118.11023475,1988.71053956)(3115.46665219,1997.43855818)(3115.4666593,2008.30665588)
\curveto(3115.46665219,2019.21664309)(3118.11023475,2027.96564252)(3123.39741492,2034.55368042)
\curveto(3128.684565,2041.1415937)(3135.65019524,2044.43558149)(3144.29432654,2044.43565369)
\curveto(3149.49753246,2044.43558149)(3153.96644584,2043.40752161)(3157.70108008,2041.35147095)
\curveto(3161.43561562,2039.29528207)(3164.49881446,2036.1901216)(3166.89068579,2032.03598022)
\lineto(3166.89068579,2042.73620605)
\lineto(3178.47210669,2042.73620605)
\lineto(3178.47210669,1980.98960876)
}
}
{
\newrgbcolor{curcolor}{0 0 0}
\pscustom[linestyle=none,fillstyle=solid,fillcolor=curcolor]
{
\newpath
\moveto(3262.62623572,2010.38375854)
\lineto(3262.62623572,2004.71893311)
\lineto(3209.37687659,2004.71893311)
\curveto(3209.88039692,1996.74619131)(3212.27220972,1990.66175526)(3216.55232214,1986.46560669)
\curveto(3220.87434345,1982.3113913)(3226.87485625,1980.23429071)(3234.55387854,1980.23429871)
\curveto(3239.00176673,1980.23429071)(3243.3028336,1980.77979188)(3247.45709204,1981.87080383)
\curveto(3251.65319756,1982.96179653)(3255.80739873,1984.59830002)(3259.91970801,1986.78031921)
\lineto(3259.91970801,1975.82832336)
\curveto(3255.7654371,1974.0659314)(3251.50633186,1972.72315931)(3247.14237952,1971.80000305)
\curveto(3242.77831325,1970.87684768)(3238.3513615,1970.41526977)(3233.86151099,1970.41526794)
\curveto(3222.61575101,1970.41526977)(3213.69890507,1973.68827675)(3207.11094641,1980.23429871)
\curveto(3200.56491552,1986.78030467)(3197.29190854,1995.63420817)(3197.29191565,2006.79603577)
\curveto(3197.29190854,2018.3354489)(3200.39706901,2027.4830838)(3206.60740637,2034.2389679)
\curveto(3212.85967251,2041.03668963)(3221.27297892,2044.43558149)(3231.84735083,2044.43565369)
\curveto(3241.33063708,2044.43558149)(3248.82078767,2041.37238265)(3254.31782507,2035.24604797)
\curveto(3259.85669583,2029.16154891)(3262.62616327,2020.87412739)(3262.62623572,2010.38375854)
\moveto(3251.04481482,2013.78265381)
\curveto(3250.9608307,2020.11881809)(3249.17746151,2025.17519426)(3245.6947019,2028.95179749)
\curveto(3242.2537929,2032.72828729)(3237.67997545,2034.61656055)(3231.97323584,2034.61662292)
\curveto(3225.51110334,2034.61656055)(3220.32884229,2032.79122973)(3216.42643713,2029.140625)
\curveto(3212.56594111,2025.48990647)(3210.34197483,2020.34960704)(3209.75453162,2013.7197113)
\lineto(3251.04481482,2013.78265381)
}
}
{
\newrgbcolor{curcolor}{0 0 0}
\pscustom[linestyle=none,fillstyle=solid,fillcolor=curcolor]
{
\newpath
\moveto(3292.83863806,2062.75192261)
\lineto(3292.83863806,2042.73620605)
\lineto(3316.69384741,2042.73620605)
\lineto(3316.69384741,2033.73542786)
\lineto(3292.83863806,2033.73542786)
\lineto(3292.83863806,1995.46638489)
\curveto(3292.83861446,1989.71761863)(3293.61490458,1986.02499537)(3295.16751074,1984.38850403)
\curveto(3296.76202667,1982.75198839)(3299.97209121,1981.93373664)(3304.79771399,1981.93374634)
\lineto(3316.69384741,1981.93374634)
\lineto(3316.69384741,1972.24060059)
\lineto(3304.79771399,1972.24060059)
\curveto(3295.85985167,1972.24060059)(3289.69149236,1973.89808489)(3286.29261755,1977.21305847)
\curveto(3282.89370863,1980.56998374)(3281.1942627,1986.65441979)(3281.19427466,1995.46638489)
\lineto(3281.19427466,2033.73542786)
\lineto(3272.6970365,2033.73542786)
\lineto(3272.6970365,2042.73620605)
\lineto(3281.19427466,2042.73620605)
\lineto(3281.19427466,2062.75192261)
\lineto(3292.83863806,2062.75192261)
}
}
{
\newrgbcolor{curcolor}{0 0 0}
\pscustom[linestyle=none,fillstyle=solid,fillcolor=curcolor]
{
\newpath
\moveto(3345.20680212,2055.76530457)
\lineto(3345.20680212,2020.45455933)
\lineto(3361.19419836,2020.45455933)
\curveto(3367.11074655,2020.45451111)(3371.684564,2021.98611053)(3374.91566443,2025.04936218)
\curveto(3378.14665471,2028.11250821)(3379.76217738,2032.47651752)(3379.7622373,2038.1414032)
\curveto(3379.76217738,2043.76419544)(3378.14665471,2048.10722394)(3374.91566443,2051.17050171)
\curveto(3371.684564,2054.23362162)(3367.11074655,2055.76522104)(3361.19419836,2055.76530457)
\lineto(3345.20680212,2055.76530457)
\moveto(3332.49241614,2066.21376038)
\lineto(3361.19419836,2066.21376038)
\curveto(3371.72652563,2066.2136664)(3379.67825413,2063.82185361)(3385.04940771,2059.03831482)
\curveto(3390.46239251,2054.29656406)(3393.16891752,2047.33093382)(3393.16899084,2038.1414032)
\curveto(3393.16891752,2028.86781752)(3390.46239251,2021.86022565)(3385.04940771,2017.11860657)
\curveto(3379.67825413,2012.37689773)(3371.72652563,2010.00606575)(3361.19419836,2010.00610352)
\lineto(3345.20680212,2010.00610352)
\lineto(3345.20680212,1972.24060059)
\lineto(3332.49241614,1972.24060059)
\lineto(3332.49241614,2066.21376038)
}
}
{
\newrgbcolor{curcolor}{0 0 0}
\pscustom[linestyle=none,fillstyle=solid,fillcolor=curcolor]
{
\newpath
\moveto(3436.03283667,2007.67723083)
\curveto(3426.67534945,2007.6771954)(3420.19227793,2006.60717389)(3416.58360266,2004.46716309)
\curveto(3412.97487792,2002.32708783)(3411.17052792,1998.6764262)(3411.17054724,1993.51516724)
\curveto(3411.17052792,1989.40290642)(3412.51330001,1986.12989944)(3415.19886755,1983.69613647)
\curveto(3417.92635002,1981.30431223)(3421.61897328,1980.10840583)(3426.27674841,1980.1084137)
\curveto(3432.69684306,1980.10840583)(3437.83714249,1982.37433374)(3441.69766211,1986.90620422)
\curveto(3445.60004366,1991.48000701)(3447.55125936,1997.54346225)(3447.55131506,2005.09658813)
\lineto(3447.55131506,2007.67723083)
\lineto(3436.03283667,2007.67723083)
\moveto(3459.13273596,2012.46086121)
\lineto(3459.13273596,1972.24060059)
\lineto(3447.55131506,1972.24060059)
\lineto(3447.55131506,1982.94082642)
\curveto(3444.9076768,1978.66072966)(3441.613689,1975.49262675)(3437.6693418,1973.43650818)
\curveto(3433.72490295,1971.42234884)(3428.89931573,1970.41526977)(3423.19256567,1970.41526794)
\curveto(3415.97513432,1970.41526977)(3410.22639129,1972.42942791)(3405.94631934,1976.45774841)
\curveto(3401.70818081,1980.52802211)(3399.5891186,1985.96205293)(3399.58912634,1992.75985718)
\curveto(3399.5891186,2000.69058434)(3402.23270116,2006.67011633)(3407.51988196,2010.69847107)
\curveto(3412.84899304,2014.7267489)(3420.77974072,2016.74090704)(3431.3121488,2016.74095154)
\lineto(3447.55131506,2016.74095154)
\lineto(3447.55131506,2017.87391663)
\curveto(3447.55125936,2023.20299774)(3445.78887099,2027.31523728)(3442.26414465,2030.21064758)
\curveto(3438.78127912,2033.14790357)(3433.87176864,2034.61656055)(3427.53559851,2034.61662292)
\curveto(3423.50724654,2034.61656055)(3419.58383432,2034.13400183)(3415.7653501,2033.16894531)
\curveto(3411.94681804,2032.20376694)(3408.27517559,2030.75609078)(3404.75041174,2028.82591248)
\lineto(3404.75041174,2039.52613831)
\curveto(3408.98852326,2041.16257451)(3413.1007628,2042.37946172)(3417.0871427,2043.17680359)
\curveto(3421.07347212,2044.01596521)(3424.9549227,2044.43558149)(3428.7315061,2044.43565369)
\curveto(3438.92814481,2044.43558149)(3446.54418029,2041.79199893)(3451.57963538,2036.50489807)
\curveto(3456.614971,2031.21766868)(3459.13266868,2023.20299774)(3459.13273596,2012.46086121)
}
}
{
\newrgbcolor{curcolor}{0 0 0}
\pscustom[linestyle=none,fillstyle=solid,fillcolor=curcolor]
{
\newpath
\moveto(3524.15234351,2031.91009521)
\curveto(3522.85148004,2032.66534485)(3521.42478469,2033.21084601)(3519.87225317,2033.54660034)
\curveto(3518.36158585,2033.92419369)(3516.68312073,2034.11302101)(3514.83685278,2034.11308289)
\curveto(3508.29079514,2034.11302101)(3503.25539978,2031.97297799)(3499.73065161,2027.69294739)
\curveto(3496.24780791,2023.45476751)(3494.50640035,2017.34935064)(3494.50642371,2009.37667847)
\lineto(3494.50642371,1972.24060059)
\lineto(3482.8620603,1972.24060059)
\lineto(3482.8620603,2042.73620605)
\lineto(3494.50642371,2042.73620605)
\lineto(3494.50642371,2031.78421021)
\curveto(3496.94017478,2036.06423671)(3500.10827769,2039.23233962)(3504.01074194,2041.28852844)
\curveto(3507.91314049,2043.38654079)(3512.65480445,2044.43558149)(3518.23574805,2044.43565369)
\curveto(3519.0329719,2044.43558149)(3519.91416608,2044.37263905)(3520.87933325,2044.24682617)
\curveto(3521.84440097,2044.16283091)(3522.91442248,2044.01596521)(3524.089401,2043.80622864)
\lineto(3524.15234351,2031.91009521)
}
}
{
\newrgbcolor{curcolor}{0 0 0}
\pscustom[linestyle=none,fillstyle=solid,fillcolor=curcolor]
{
\newpath
\moveto(3547.88166785,2062.75192261)
\lineto(3547.88166785,2042.73620605)
\lineto(3571.7368772,2042.73620605)
\lineto(3571.7368772,2033.73542786)
\lineto(3547.88166785,2033.73542786)
\lineto(3547.88166785,1995.46638489)
\curveto(3547.88164424,1989.71761863)(3548.65793436,1986.02499537)(3550.21054053,1984.38850403)
\curveto(3551.80505646,1982.75198839)(3555.015121,1981.93373664)(3559.84074377,1981.93374634)
\lineto(3571.7368772,1981.93374634)
\lineto(3571.7368772,1972.24060059)
\lineto(3559.84074377,1972.24060059)
\curveto(3550.90288146,1972.24060059)(3544.73452215,1973.89808489)(3541.33564734,1977.21305847)
\curveto(3537.93673842,1980.56998374)(3536.23729248,1986.65441979)(3536.23730444,1995.46638489)
\lineto(3536.23730444,2033.73542786)
\lineto(3527.74006628,2033.73542786)
\lineto(3527.74006628,2042.73620605)
\lineto(3536.23730444,2042.73620605)
\lineto(3536.23730444,2062.75192261)
\lineto(3547.88166785,2062.75192261)
}
}
{
\newrgbcolor{curcolor}{0 0 0}
\pscustom[linestyle=none,fillstyle=solid,fillcolor=curcolor]
{
\newpath
\moveto(3625.67860388,2057.59063721)
\curveto(3616.44699494,2057.59055186)(3609.10371004,2054.14969836)(3603.64872717,2047.26806641)
\curveto(3598.2356484,2040.38628439)(3595.5291234,2031.00786054)(3595.52914404,2019.13276672)
\curveto(3595.5291234,2007.29954075)(3598.2356484,1997.94209771)(3603.64872717,1991.06040955)
\curveto(3609.10371004,1984.17868374)(3616.44699494,1980.73783025)(3625.67860388,1980.73783875)
\curveto(3634.91011124,1980.73783025)(3642.2114345,1984.17868374)(3647.58259558,1991.06040955)
\curveto(3652.99557289,1997.94209771)(3655.70209789,2007.29954075)(3655.70217871,2019.13276672)
\curveto(3655.70209789,2031.00786054)(3652.99557289,2040.38628439)(3647.58259558,2047.26806641)
\curveto(3642.2114345,2054.14969836)(3634.91011124,2057.59055186)(3625.67860388,2057.59063721)
\moveto(3625.67860388,2067.91320801)
\curveto(3638.85450427,2067.91311234)(3649.38687288,2063.48616059)(3657.27574133,2054.63233948)
\curveto(3665.164445,2045.82031521)(3669.10883803,2033.98713613)(3669.10893225,2019.13276672)
\curveto(3669.10883803,2004.32026516)(3665.164445,1992.48708608)(3657.27574133,1983.63319397)
\curveto(3649.38687288,1974.82124071)(3638.85450427,1970.41526977)(3625.67860388,1970.41526794)
\curveto(3612.46064028,1970.41526977)(3601.88631003,1974.82124071)(3593.95558142,1983.63319397)
\curveto(3586.06677629,1992.44512445)(3582.12238326,2004.27830353)(3582.1223905,2019.13276672)
\curveto(3582.12238326,2033.98713613)(3586.06677629,2045.82031521)(3593.95558142,2054.63233948)
\curveto(3601.88631003,2063.48616059)(3612.46064028,2067.91311234)(3625.67860388,2067.91320801)
}
}
{
\newrgbcolor{curcolor}{0 0 0}
\pscustom[linestyle=none,fillstyle=solid,fillcolor=curcolor]
{
\newpath
\moveto(3724.18362402,2070.17913818)
\lineto(3724.18362402,2060.54893494)
\lineto(3713.10574316,2060.54893494)
\curveto(3708.95150524,2060.54884663)(3706.05615291,2059.70961407)(3704.41967749,2058.03123474)
\curveto(3702.82510756,2056.35268383)(3702.02783662,2053.33144662)(3702.0278623,2048.96751404)
\lineto(3702.0278623,2042.73620605)
\lineto(3721.09944128,2042.73620605)
\lineto(3721.09944128,2033.73542786)
\lineto(3702.0278623,2033.73542786)
\lineto(3702.0278623,1972.24060059)
\lineto(3690.3834989,1972.24060059)
\lineto(3690.3834989,2033.73542786)
\lineto(3679.30561804,2033.73542786)
\lineto(3679.30561804,2042.73620605)
\lineto(3690.3834989,2042.73620605)
\lineto(3690.3834989,2047.64572144)
\curveto(3690.38348487,2055.49247046)(3692.20881568,2061.19925186)(3695.85949683,2064.76608276)
\curveto(3699.51013895,2068.37469024)(3705.30084361,2070.17904025)(3713.23162817,2070.17913818)
\lineto(3724.18362402,2070.17913818)
}
}
{
\newrgbcolor{curcolor}{0 0 0}
\pscustom[linestyle=none,fillstyle=solid,fillcolor=curcolor]
{
\newpath
\moveto(3761.88618445,2070.05325317)
\curveto(3756.26328633,2060.40198093)(3752.08810435,2050.85571057)(3749.36062598,2041.41441345)
\curveto(3746.63309272,2031.97297799)(3745.26933981,2022.40572681)(3745.26936316,2012.71263123)
\curveto(3745.26933981,2003.01945469)(3746.63309272,1993.41024189)(3749.36062598,1983.88496399)
\curveto(3752.13006598,1974.40162443)(3756.30524796,1964.85535406)(3761.88618445,1955.24612427)
\lineto(3751.81538367,1955.24612427)
\curveto(3745.52110958,1965.10712383)(3740.80042643,1974.80025989)(3737.65332007,1984.32556152)
\curveto(3734.54814386,1993.85083899)(3732.99556363,2003.31318609)(3732.99557471,2012.71263123)
\curveto(3732.99556363,2022.07003379)(3734.54814386,2031.49041927)(3737.65332007,2040.97381592)
\curveto(3740.7584648,2050.4570751)(3745.47914795,2060.15021116)(3751.81538367,2070.05325317)
\lineto(3761.88618445,2070.05325317)
}
}
{
\newrgbcolor{curcolor}{0 0 0}
\pscustom[linestyle=none,fillstyle=solid,fillcolor=curcolor]
{
\newpath
\moveto(3782.59426855,2070.05325317)
\lineto(3792.66506934,2070.05325317)
\curveto(3798.95929314,2060.15021116)(3803.65899547,2050.4570751)(3806.76419043,2040.97381592)
\curveto(3809.91127803,2031.49041927)(3811.48483908,2022.07003379)(3811.4848783,2012.71263123)
\curveto(3811.48483908,2003.31318609)(3809.91127803,1993.85083899)(3806.76419043,1984.32556152)
\curveto(3803.65899547,1974.80025989)(3798.95929314,1965.10712383)(3792.66506934,1955.24612427)
\lineto(3782.59426855,1955.24612427)
\curveto(3788.17515475,1964.85535406)(3792.32935592,1974.40162443)(3795.05688452,1983.88496399)
\curveto(3797.82632918,1993.41024189)(3799.2110629,2003.01945469)(3799.21108984,2012.71263123)
\curveto(3799.2110629,2022.40572681)(3797.82632918,2031.97297799)(3795.05688452,2041.41441345)
\curveto(3792.32935592,2050.85571057)(3788.17515475,2060.40198093)(3782.59426855,2070.05325317)
}
}
{
\newrgbcolor{curcolor}{0 0 0}
\pscustom[linestyle=none,fillstyle=solid,fillcolor=curcolor]
{
}
}
{
\newrgbcolor{curcolor}{0 0 0}
\pscustom[linestyle=none,fillstyle=solid,fillcolor=curcolor]
{
\newpath
\moveto(3878.77041602,1988.22799683)
\lineto(3892.05128455,1988.22799683)
\lineto(3892.05128455,1972.24060059)
\lineto(3878.77041602,1972.24060059)
\lineto(3878.77041602,1988.22799683)
\moveto(3878.77041602,2038.89671326)
\lineto(3892.05128455,2038.89671326)
\lineto(3892.05128455,2022.90931702)
\lineto(3878.77041602,2022.90931702)
\lineto(3878.77041602,2038.89671326)
}
}
{
\newrgbcolor{curcolor}{0 0 0}
\pscustom[linestyle=none,fillstyle=solid,fillcolor=curcolor]
{
}
}
{
\newrgbcolor{curcolor}{0 0 0}
\pscustom[linestyle=none,fillstyle=solid,fillcolor=curcolor]
{
\newpath
\moveto(3960.91038489,2066.21376038)
\lineto(3973.62477087,2066.21376038)
\lineto(3973.62477087,2026.49703979)
\lineto(4015.79624915,2066.21376038)
\lineto(4032.16130042,2066.21376038)
\lineto(3985.5209043,2022.40577698)
\lineto(4035.49725317,1972.24060059)
\lineto(4018.75454687,1972.24060059)
\lineto(3973.62477087,2017.4962616)
\lineto(3973.62477087,1972.24060059)
\lineto(3960.91038489,1972.24060059)
\lineto(3960.91038489,2066.21376038)
}
}
{
\newrgbcolor{curcolor}{0 0 0}
\pscustom[linestyle=none,fillstyle=solid,fillcolor=curcolor]
{
\newpath
\moveto(4103.60104346,2014.78973389)
\lineto(4103.60104346,1972.24060059)
\lineto(4092.01962256,1972.24060059)
\lineto(4092.01962256,2014.41207886)
\curveto(4092.01956339,2021.08393553)(4090.71875293,2026.07736926)(4088.11718726,2029.39239502)
\curveto(4085.51551106,2032.70730648)(4081.61307966,2034.36479078)(4076.40988135,2034.36485291)
\curveto(4070.15755523,2034.36479078)(4065.22706394,2032.37161345)(4061.6183927,2028.38531494)
\curveto(4058.00966393,2024.39890414)(4056.20531393,2018.96487332)(4056.20533728,2012.08320618)
\lineto(4056.20533728,1972.24060059)
\lineto(4044.56097388,1972.24060059)
\lineto(4044.56097388,2042.73620605)
\lineto(4056.20533728,2042.73620605)
\lineto(4056.20533728,2031.78421021)
\curveto(4058.97478137,2036.02227509)(4062.22680754,2039.190378)(4065.96142554,2041.28852844)
\curveto(4069.73793895,2043.38654079)(4074.08096744,2044.43558149)(4078.99052405,2044.43565369)
\curveto(4087.08907211,2044.43558149)(4093.21546979,2041.91788381)(4097.36973547,2036.8825531)
\curveto(4101.52387213,2031.88905473)(4103.60097271,2024.52478902)(4103.60104346,2014.78973389)
}
}
{
\newrgbcolor{curcolor}{0 0 0}
\pscustom[linestyle=none,fillstyle=solid,fillcolor=curcolor]
{
\newpath
\moveto(4154.14387488,2034.61662292)
\curveto(4147.93351447,2034.61656055)(4143.024004,2032.18278613)(4139.41532874,2027.31529236)
\curveto(4135.80660399,2022.48965007)(4134.00225399,2015.85971285)(4134.00227332,2007.42546082)
\curveto(4134.00225399,1998.99113841)(4135.78562318,1992.34022038)(4139.35238623,1987.47268677)
\curveto(4142.96106156,1982.64708432)(4147.89155285,1980.23429071)(4154.14387488,1980.23429871)
\curveto(4160.31219472,1980.23429071)(4165.20072438,1982.66806513)(4168.80947852,1987.53562927)
\curveto(4172.41812439,1992.40316282)(4174.22247439,1999.03310004)(4174.22253394,2007.42546082)
\curveto(4174.22247439,2015.77578959)(4172.41812439,2022.384746)(4168.80947852,2027.25234985)
\curveto(4165.20072438,2032.16180531)(4160.31219472,2034.61656055)(4154.14387488,2034.61662292)
\moveto(4154.14387488,2044.43565369)
\curveto(4164.21462612,2044.43558149)(4172.12439299,2041.16257451)(4177.87319922,2034.61662292)
\curveto(4183.62187906,2028.07054659)(4186.49625057,2019.00683495)(4186.49632239,2007.42546082)
\curveto(4186.49625057,1995.88597794)(4183.62187906,1986.8222663)(4177.87319922,1980.23429871)
\curveto(4172.12439299,1973.68827675)(4164.21462612,1970.41526977)(4154.14387488,1970.41526794)
\curveto(4144.03108307,1970.41526977)(4136.10033539,1973.68827675)(4130.35160803,1980.23429871)
\curveto(4124.64481096,1986.8222663)(4121.79142026,1995.88597794)(4121.79142737,2007.42546082)
\curveto(4121.79142026,2019.00683495)(4124.64481096,2028.07054659)(4130.35160803,2034.61662292)
\curveto(4136.10033539,2041.16257451)(4144.03108307,2044.43558149)(4154.14387488,2044.43565369)
}
}
{
\newrgbcolor{curcolor}{0 0 0}
\pscustom[linestyle=none,fillstyle=solid,fillcolor=curcolor]
{
\newpath
\moveto(4198.89599585,2042.73620605)
\lineto(4210.47741675,2042.73620605)
\lineto(4224.95419287,1987.72445679)
\lineto(4239.36802649,2042.73620605)
\lineto(4253.02655005,2042.73620605)
\lineto(4267.50332617,1987.72445679)
\lineto(4281.91715979,2042.73620605)
\lineto(4293.49858069,2042.73620605)
\lineto(4275.05642676,1972.24060059)
\lineto(4261.3979032,1972.24060059)
\lineto(4246.22875952,2030.02182007)
\lineto(4230.99667334,1972.24060059)
\lineto(4217.33814978,1972.24060059)
\lineto(4198.89599585,2042.73620605)
}
}
{
\newrgbcolor{curcolor}{0 0 0}
\pscustom[linestyle=none,fillstyle=solid,fillcolor=curcolor]
{
\newpath
\moveto(4311.12248206,2070.17913818)
\lineto(4322.70390295,2070.17913818)
\lineto(4322.70390295,1972.24060059)
\lineto(4311.12248206,1972.24060059)
\lineto(4311.12248206,2070.17913818)
}
}
{
\newrgbcolor{curcolor}{0 0 0}
\pscustom[linestyle=none,fillstyle=solid,fillcolor=curcolor]
{
\newpath
\moveto(4407.17274451,2010.38375854)
\lineto(4407.17274451,2004.71893311)
\lineto(4353.92338538,2004.71893311)
\curveto(4354.42690571,1996.74619131)(4356.81871851,1990.66175526)(4361.09883093,1986.46560669)
\curveto(4365.42085224,1982.3113913)(4371.42136504,1980.23429071)(4379.10038733,1980.23429871)
\curveto(4383.54827552,1980.23429071)(4387.84934238,1980.77979188)(4392.00360083,1981.87080383)
\curveto(4396.19970635,1982.96179653)(4400.35390752,1984.59830002)(4404.4662168,1986.78031921)
\lineto(4404.4662168,1975.82832336)
\curveto(4400.31194589,1974.0659314)(4396.05284065,1972.72315931)(4391.68888831,1971.80000305)
\curveto(4387.32482203,1970.87684768)(4382.89787029,1970.41526977)(4378.40801978,1970.41526794)
\curveto(4367.1622598,1970.41526977)(4358.24541386,1973.68827675)(4351.6574552,1980.23429871)
\curveto(4345.11142431,1986.78030467)(4341.83841733,1995.63420817)(4341.83842444,2006.79603577)
\curveto(4341.83841733,2018.3354489)(4344.9435778,2027.4830838)(4351.15391516,2034.2389679)
\curveto(4357.4061813,2041.03668963)(4365.81948771,2044.43558149)(4376.39385962,2044.43565369)
\curveto(4385.87714587,2044.43558149)(4393.36729646,2041.37238265)(4398.86433386,2035.24604797)
\curveto(4404.40320461,2029.16154891)(4407.17267206,2020.87412739)(4407.17274451,2010.38375854)
\moveto(4395.59132361,2013.78265381)
\curveto(4395.50733949,2020.11881809)(4393.7239703,2025.17519426)(4390.24121069,2028.95179749)
\curveto(4386.80030169,2032.72828729)(4382.22648424,2034.61656055)(4376.51974463,2034.61662292)
\curveto(4370.05761213,2034.61656055)(4364.87535108,2032.79122973)(4360.97294592,2029.140625)
\curveto(4357.1124499,2025.48990647)(4354.88848362,2020.34960704)(4354.30104041,2013.7197113)
\lineto(4395.59132361,2013.78265381)
}
}
{
\newrgbcolor{curcolor}{0 0 0}
\pscustom[linestyle=none,fillstyle=solid,fillcolor=curcolor]
{
\newpath
\moveto(4472.31823706,2032.03598022)
\lineto(4472.31823706,2070.17913818)
\lineto(4483.89965796,2070.17913818)
\lineto(4483.89965796,1972.24060059)
\lineto(4472.31823706,1972.24060059)
\lineto(4472.31823706,1982.81494141)
\curveto(4469.8844041,1978.61876804)(4466.80022445,1975.49262675)(4463.06568884,1973.43650818)
\curveto(4459.3730163,1971.42234884)(4454.92508373,1970.41526977)(4449.72187781,1970.41526794)
\curveto(4441.20363139,1970.41526977)(4434.25898196,1973.81416163)(4428.88790869,1980.61195374)
\curveto(4423.55876683,1987.40972909)(4420.89420346,1996.34755585)(4420.89421057,2007.42546082)
\curveto(4420.89420346,2018.50329541)(4423.55876683,2027.44112217)(4428.88790869,2034.2389679)
\curveto(4434.25898196,2041.03668963)(4441.20363139,2044.43558149)(4449.72187781,2044.43565369)
\curveto(4454.92508373,2044.43558149)(4459.3730163,2043.40752161)(4463.06568884,2041.35147095)
\curveto(4466.80022445,2039.33724369)(4469.8844041,2036.23208323)(4472.31823706,2032.03598022)
\moveto(4432.8532865,2007.42546082)
\curveto(4432.85326743,1998.90721515)(4434.59467499,1992.21433549)(4438.0775144,1987.34680176)
\curveto(4441.60226686,1982.52119944)(4446.42785407,1980.10840583)(4452.55429053,1980.1084137)
\curveto(4458.68064944,1980.10840583)(4463.50623665,1982.52119944)(4467.03106665,1987.34680176)
\curveto(4470.55579015,1992.21433549)(4472.31817852,1998.90721515)(4472.31823706,2007.42546082)
\curveto(4472.31817852,2015.94363611)(4470.55579015,2022.61553495)(4467.03106665,2027.44117737)
\curveto(4463.50623665,2032.30867101)(4458.68064944,2034.74244543)(4452.55429053,2034.74250793)
\curveto(4446.42785407,2034.74244543)(4441.60226686,2032.30867101)(4438.0775144,2027.44117737)
\curveto(4434.59467499,2022.61553495)(4432.85326743,2015.94363611)(4432.8532865,2007.42546082)
}
}
{
\newrgbcolor{curcolor}{0 0 0}
\pscustom[linestyle=none,fillstyle=solid,fillcolor=curcolor]
{
\newpath
\moveto(4554.14349341,2008.30665588)
\curveto(4554.14343487,2016.69894541)(4552.40202731,2023.20299774)(4548.9192655,2027.8188324)
\curveto(4545.4783587,2032.43455589)(4540.63179067,2034.74244543)(4534.37954687,2034.74250793)
\curveto(4528.16918716,2034.74244543)(4523.32261913,2032.43455589)(4519.83982825,2027.8188324)
\curveto(4516.39895052,2023.20299774)(4514.67852378,2016.69894541)(4514.67854285,2008.30665588)
\curveto(4514.67852378,1999.95625585)(4516.39895052,1993.47318433)(4519.83982825,1988.85742188)
\curveto(4523.32261913,1984.24162618)(4528.16918716,1981.93373664)(4534.37954687,1981.93374634)
\curveto(4540.63179067,1981.93373664)(4545.4783587,1984.24162618)(4548.9192655,1988.85742188)
\curveto(4552.40202731,1993.47318433)(4554.14343487,1999.95625585)(4554.14349341,2008.30665588)
\moveto(4565.72491431,1980.98960876)
\curveto(4565.72484419,1968.98857442)(4563.06028081,1960.07172848)(4557.73121619,1954.23904419)
\curveto(4552.40202731,1948.36443428)(4544.24049067,1945.42712032)(4533.24658179,1945.42709351)
\curveto(4529.17626624,1945.42712032)(4525.33677728,1945.74183253)(4521.72810339,1946.37123108)
\curveto(4518.11937727,1946.95871974)(4514.61558133,1947.88187556)(4511.21670508,1949.14070129)
\lineto(4511.21670508,1960.40740967)
\curveto(4514.61558133,1958.56110987)(4517.97251157,1957.19735696)(4521.28750586,1956.31614685)
\curveto(4524.60244879,1955.43496859)(4527.98035984,1954.99437149)(4531.42124915,1954.99435425)
\curveto(4539.01626799,1954.99437149)(4544.70206858,1956.98754882)(4548.47866797,1960.97389221)
\curveto(4552.25516161,1964.91829651)(4554.14343487,1970.89782849)(4554.14349341,1978.9125061)
\lineto(4554.14349341,1984.64027405)
\curveto(4551.75162208,1980.48606048)(4548.68842324,1977.38090001)(4544.9538877,1975.32478333)
\curveto(4541.21925346,1973.26866047)(4536.75034008,1972.24060059)(4531.54713416,1972.24060059)
\curveto(4522.90300286,1972.24060059)(4515.93737261,1975.53458838)(4510.65022253,1982.12257385)
\curveto(4505.36304237,1988.71053956)(4502.71945981,1997.43855818)(4502.71946692,2008.30665588)
\curveto(4502.71945981,2019.21664309)(4505.36304237,2027.96564252)(4510.65022253,2034.55368042)
\curveto(4515.93737261,2041.1415937)(4522.90300286,2044.43558149)(4531.54713416,2044.43565369)
\curveto(4536.75034008,2044.43558149)(4541.21925346,2043.40752161)(4544.9538877,2041.35147095)
\curveto(4548.68842324,2039.29528207)(4551.75162208,2036.1901216)(4554.14349341,2032.03598022)
\lineto(4554.14349341,2042.73620605)
\lineto(4565.72491431,2042.73620605)
\lineto(4565.72491431,1980.98960876)
}
}
{
\newrgbcolor{curcolor}{0 0 0}
\pscustom[linestyle=none,fillstyle=solid,fillcolor=curcolor]
{
\newpath
\moveto(4649.87904333,2010.38375854)
\lineto(4649.87904333,2004.71893311)
\lineto(4596.6296842,2004.71893311)
\curveto(4597.13320454,1996.74619131)(4599.52501734,1990.66175526)(4603.80512976,1986.46560669)
\curveto(4608.12715107,1982.3113913)(4614.12766387,1980.23429071)(4621.80668616,1980.23429871)
\curveto(4626.25457435,1980.23429071)(4630.55564121,1980.77979188)(4634.70989966,1981.87080383)
\curveto(4638.90600518,1982.96179653)(4643.06020634,1984.59830002)(4647.17251562,1986.78031921)
\lineto(4647.17251562,1975.82832336)
\curveto(4643.01824472,1974.0659314)(4638.75913948,1972.72315931)(4634.39518713,1971.80000305)
\curveto(4630.03112086,1970.87684768)(4625.60416911,1970.41526977)(4621.1143186,1970.41526794)
\curveto(4609.86855863,1970.41526977)(4600.95171269,1973.68827675)(4594.36375403,1980.23429871)
\curveto(4587.81772313,1986.78030467)(4584.54471615,1995.63420817)(4584.54472327,2006.79603577)
\curveto(4584.54471615,2018.3354489)(4587.64987662,2027.4830838)(4593.86021399,2034.2389679)
\curveto(4600.11248013,2041.03668963)(4608.52578653,2044.43558149)(4619.10015845,2044.43565369)
\curveto(4628.5834447,2044.43558149)(4636.07359529,2041.37238265)(4641.57063269,2035.24604797)
\curveto(4647.10950344,2029.16154891)(4649.87897089,2020.87412739)(4649.87904333,2010.38375854)
\moveto(4638.29762244,2013.78265381)
\curveto(4638.21363832,2020.11881809)(4636.43026913,2025.17519426)(4632.94750952,2028.95179749)
\curveto(4629.50660051,2032.72828729)(4624.93278307,2034.61656055)(4619.22604346,2034.61662292)
\curveto(4612.76391096,2034.61656055)(4607.5816499,2032.79122973)(4603.67924475,2029.140625)
\curveto(4599.81874873,2025.48990647)(4597.59478245,2020.34960704)(4597.00733923,2013.7197113)
\lineto(4638.29762244,2013.78265381)
}
}
{
\newrgbcolor{curcolor}{0 0 0}
\pscustom[linestyle=none,fillstyle=solid,fillcolor=curcolor]
{
\newpath
\moveto(4667.56588721,2070.17913818)
\lineto(4694.25350928,2070.17913818)
\lineto(4694.25350928,2061.17835999)
\lineto(4679.14730811,2061.17835999)
\lineto(4679.14730811,1964.24690247)
\lineto(4694.25350928,1964.24690247)
\lineto(4694.25350928,1955.24612427)
\lineto(4667.56588721,1955.24612427)
\lineto(4667.56588721,2070.17913818)
}
}
{
\newrgbcolor{curcolor}{0 0 0}
\pscustom[linestyle=none,fillstyle=solid,fillcolor=curcolor]
{
\newpath
\moveto(4746.0551908,2070.17913818)
\lineto(4746.0551908,1955.24612427)
\lineto(4719.36756873,1955.24612427)
\lineto(4719.36756873,1964.24690247)
\lineto(4734.41082739,1964.24690247)
\lineto(4734.41082739,2061.17835999)
\lineto(4719.36756873,2061.17835999)
\lineto(4719.36756873,2070.17913818)
\lineto(4746.0551908,2070.17913818)
}
}
{
\newrgbcolor{curcolor}{0.51764709 0.74509805 0.51764709}
\pscustom[linestyle=none,fillstyle=solid,fillcolor=curcolor]
{
\newpath
\moveto(3038.04159375,1852.16070703)
\curveto(3038.04159375,1832.7444463)(3022.30160449,1817.00445703)(3002.88534375,1817.00445703)
\curveto(2983.46908301,1817.00445703)(2967.72909375,1832.7444463)(2967.72909375,1852.16070703)
\curveto(2967.72909375,1871.57696777)(2983.46908301,1887.31695703)(3002.88534375,1887.31695703)
\curveto(3022.30160449,1887.31695703)(3038.04159375,1871.57696777)(3038.04159375,1852.16070703)
\closepath
}
}
{
\newrgbcolor{curcolor}{0.01176471 0.50196081 0.28235295}
\pscustom[linewidth=11.71875,linecolor=curcolor]
{
\newpath
\moveto(3038.04159375,1852.16070703)
\curveto(3038.04159375,1832.7444463)(3022.30160449,1817.00445703)(3002.88534375,1817.00445703)
\curveto(2983.46908301,1817.00445703)(2967.72909375,1832.7444463)(2967.72909375,1852.16070703)
\curveto(2967.72909375,1871.57696777)(2983.46908301,1887.31695703)(3002.88534375,1887.31695703)
\curveto(3022.30160449,1887.31695703)(3038.04159375,1871.57696777)(3038.04159375,1852.16070703)
\closepath
}
}
{
\newrgbcolor{curcolor}{0 0 0}
\pscustom[linestyle=none,fillstyle=solid,fillcolor=curcolor]
{
\newpath
\moveto(3166.89068579,1858.31459045)
\curveto(3166.89062725,1866.70687998)(3165.14921969,1873.21093231)(3161.66645789,1877.82676697)
\curveto(3158.22555108,1882.44249047)(3153.37898305,1884.75038)(3147.12673926,1884.7504425)
\curveto(3140.91637955,1884.75038)(3136.06981152,1882.44249047)(3132.58702063,1877.82676697)
\curveto(3129.1461429,1873.21093231)(3127.42571616,1866.70687998)(3127.42573523,1858.31459045)
\curveto(3127.42571616,1849.96419042)(3129.1461429,1843.4811189)(3132.58702063,1838.86535645)
\curveto(3136.06981152,1834.24956075)(3140.91637955,1831.94167122)(3147.12673926,1831.94168091)
\curveto(3153.37898305,1831.94167122)(3158.22555108,1834.24956075)(3161.66645789,1838.86535645)
\curveto(3165.14921969,1843.4811189)(3166.89062725,1849.96419042)(3166.89068579,1858.31459045)
\moveto(3178.47210669,1830.99754333)
\curveto(3178.47203657,1818.99650899)(3175.8074732,1810.07966305)(3170.47840857,1804.24697876)
\curveto(3165.14921969,1798.37236885)(3156.98768306,1795.43505489)(3145.99377417,1795.43502808)
\curveto(3141.92345862,1795.43505489)(3138.08396966,1795.7497671)(3134.47529578,1796.37916565)
\curveto(3130.86656965,1796.96665431)(3127.36277372,1797.88981013)(3123.96389746,1799.14863586)
\lineto(3123.96389746,1810.41534424)
\curveto(3127.36277372,1808.56904444)(3130.71970395,1807.20529153)(3134.03469824,1806.32408142)
\curveto(3137.34964117,1805.44290316)(3140.72755222,1805.00230606)(3144.16844153,1805.00228882)
\curveto(3151.76346037,1805.00230606)(3157.44926096,1806.99548339)(3161.22586035,1810.98182678)
\curveto(3165.002354,1814.92623108)(3166.89062725,1820.90576306)(3166.89068579,1828.92044067)
\lineto(3166.89068579,1834.64820862)
\curveto(3164.49881446,1830.49399505)(3161.43561562,1827.38883458)(3157.70108008,1825.3327179)
\curveto(3153.96644584,1823.27659504)(3149.49753246,1822.24853516)(3144.29432654,1822.24853516)
\curveto(3135.65019524,1822.24853516)(3128.684565,1825.54252295)(3123.39741492,1832.13050842)
\curveto(3118.11023475,1838.71847413)(3115.46665219,1847.44649275)(3115.4666593,1858.31459045)
\curveto(3115.46665219,1869.22457766)(3118.11023475,1877.97357709)(3123.39741492,1884.56161499)
\curveto(3128.684565,1891.14952827)(3135.65019524,1894.44351606)(3144.29432654,1894.44358826)
\curveto(3149.49753246,1894.44351606)(3153.96644584,1893.41545618)(3157.70108008,1891.35940552)
\curveto(3161.43561562,1889.30321664)(3164.49881446,1886.19805617)(3166.89068579,1882.04391479)
\lineto(3166.89068579,1892.74414062)
\lineto(3178.47210669,1892.74414062)
\lineto(3178.47210669,1830.99754333)
}
}
{
\newrgbcolor{curcolor}{0 0 0}
\pscustom[linestyle=none,fillstyle=solid,fillcolor=curcolor]
{
\newpath
\moveto(3262.62623572,1860.39169312)
\lineto(3262.62623572,1854.72686768)
\lineto(3209.37687659,1854.72686768)
\curveto(3209.88039692,1846.75412589)(3212.27220972,1840.66968983)(3216.55232214,1836.47354126)
\curveto(3220.87434345,1832.31932587)(3226.87485625,1830.24222528)(3234.55387854,1830.24223328)
\curveto(3239.00176673,1830.24222528)(3243.3028336,1830.78772645)(3247.45709204,1831.8787384)
\curveto(3251.65319756,1832.9697311)(3255.80739873,1834.60623459)(3259.91970801,1836.78825378)
\lineto(3259.91970801,1825.83625793)
\curveto(3255.7654371,1824.07386597)(3251.50633186,1822.73109388)(3247.14237952,1821.80793762)
\curveto(3242.77831325,1820.88478225)(3238.3513615,1820.42320434)(3233.86151099,1820.42320251)
\curveto(3222.61575101,1820.42320434)(3213.69890507,1823.69621132)(3207.11094641,1830.24223328)
\curveto(3200.56491552,1836.78823924)(3197.29190854,1845.64214274)(3197.29191565,1856.80397034)
\curveto(3197.29190854,1868.34338347)(3200.39706901,1877.49101837)(3206.60740637,1884.24690247)
\curveto(3212.85967251,1891.0446242)(3221.27297892,1894.44351606)(3231.84735083,1894.44358826)
\curveto(3241.33063708,1894.44351606)(3248.82078767,1891.38031722)(3254.31782507,1885.25398254)
\curveto(3259.85669583,1879.16948348)(3262.62616327,1870.88206196)(3262.62623572,1860.39169312)
\moveto(3251.04481482,1863.79058838)
\curveto(3250.9608307,1870.12675266)(3249.17746151,1875.18312883)(3245.6947019,1878.95973206)
\curveto(3242.2537929,1882.73622186)(3237.67997545,1884.62449512)(3231.97323584,1884.6245575)
\curveto(3225.51110334,1884.62449512)(3220.32884229,1882.7991643)(3216.42643713,1879.14855957)
\curveto(3212.56594111,1875.49784104)(3210.34197483,1870.35754161)(3209.75453162,1863.72764587)
\lineto(3251.04481482,1863.79058838)
}
}
{
\newrgbcolor{curcolor}{0 0 0}
\pscustom[linestyle=none,fillstyle=solid,fillcolor=curcolor]
{
\newpath
\moveto(3292.83863806,1912.75985718)
\lineto(3292.83863806,1892.74414062)
\lineto(3316.69384741,1892.74414062)
\lineto(3316.69384741,1883.74336243)
\lineto(3292.83863806,1883.74336243)
\lineto(3292.83863806,1845.47431946)
\curveto(3292.83861446,1839.7255532)(3293.61490458,1836.03292994)(3295.16751074,1834.3964386)
\curveto(3296.76202667,1832.75992296)(3299.97209121,1831.94167122)(3304.79771399,1831.94168091)
\lineto(3316.69384741,1831.94168091)
\lineto(3316.69384741,1822.24853516)
\lineto(3304.79771399,1822.24853516)
\curveto(3295.85985167,1822.24853516)(3289.69149236,1823.90601946)(3286.29261755,1827.22099304)
\curveto(3282.89370863,1830.57791831)(3281.1942627,1836.66235436)(3281.19427466,1845.47431946)
\lineto(3281.19427466,1883.74336243)
\lineto(3272.6970365,1883.74336243)
\lineto(3272.6970365,1892.74414062)
\lineto(3281.19427466,1892.74414062)
\lineto(3281.19427466,1912.75985718)
\lineto(3292.83863806,1912.75985718)
}
}
{
\newrgbcolor{curcolor}{0 0 0}
\pscustom[linestyle=none,fillstyle=solid,fillcolor=curcolor]
{
\newpath
\moveto(3332.49241614,1916.22169495)
\lineto(3349.61277747,1916.22169495)
\lineto(3391.2807157,1837.60650635)
\lineto(3391.2807157,1916.22169495)
\lineto(3403.61744666,1916.22169495)
\lineto(3403.61744666,1822.24853516)
\lineto(3386.49708533,1822.24853516)
\lineto(3344.82914709,1900.86372375)
\lineto(3344.82914709,1822.24853516)
\lineto(3332.49241614,1822.24853516)
\lineto(3332.49241614,1916.22169495)
}
}
{
\newrgbcolor{curcolor}{0 0 0}
\pscustom[linestyle=none,fillstyle=solid,fillcolor=curcolor]
{
\newpath
\moveto(3455.7338407,1884.6245575)
\curveto(3449.5234803,1884.62449512)(3444.61396982,1882.1907207)(3441.00529456,1877.32322693)
\curveto(3437.39656981,1872.49758464)(3435.59221981,1865.86764742)(3435.59223914,1857.43339539)
\curveto(3435.59221981,1848.99907298)(3437.375589,1842.34815495)(3440.94235205,1837.48062134)
\curveto(3444.55102738,1832.65501889)(3449.48151867,1830.24222528)(3455.7338407,1830.24223328)
\curveto(3461.90216054,1830.24222528)(3466.7906902,1832.6759997)(3470.39944434,1837.54356384)
\curveto(3474.00809021,1842.41109739)(3475.81244021,1849.04103461)(3475.81249976,1857.43339539)
\curveto(3475.81244021,1865.78372416)(3474.00809021,1872.39268057)(3470.39944434,1877.26028442)
\curveto(3466.7906902,1882.16973988)(3461.90216054,1884.62449512)(3455.7338407,1884.6245575)
\moveto(3455.7338407,1894.44358826)
\curveto(3465.80459194,1894.44351606)(3473.71435881,1891.17050908)(3479.46316504,1884.6245575)
\curveto(3485.21184488,1878.07848116)(3488.08621639,1869.01476952)(3488.08628821,1857.43339539)
\curveto(3488.08621639,1845.89391251)(3485.21184488,1836.83020087)(3479.46316504,1830.24223328)
\curveto(3473.71435881,1823.69621132)(3465.80459194,1820.42320434)(3455.7338407,1820.42320251)
\curveto(3445.62104889,1820.42320434)(3437.69030121,1823.69621132)(3431.94157385,1830.24223328)
\curveto(3426.23477678,1836.83020087)(3423.38138608,1845.89391251)(3423.38139319,1857.43339539)
\curveto(3423.38138608,1869.01476952)(3426.23477678,1878.07848116)(3431.94157385,1884.6245575)
\curveto(3437.69030121,1891.17050908)(3445.62104889,1894.44351606)(3455.7338407,1894.44358826)
}
}
{
\newrgbcolor{curcolor}{0 0 0}
\pscustom[linestyle=none,fillstyle=solid,fillcolor=curcolor]
{
\newpath
\moveto(3553.60943579,1882.04391479)
\lineto(3553.60943579,1920.18707275)
\lineto(3565.19085669,1920.18707275)
\lineto(3565.19085669,1822.24853516)
\lineto(3553.60943579,1822.24853516)
\lineto(3553.60943579,1832.82287598)
\curveto(3551.17560283,1828.62670261)(3548.09142318,1825.50056132)(3544.35688757,1823.44444275)
\curveto(3540.66421503,1821.43028341)(3536.21628246,1820.42320434)(3531.01307654,1820.42320251)
\curveto(3522.49483012,1820.42320434)(3515.55018069,1823.8220962)(3510.17910742,1830.61988831)
\curveto(3504.84996556,1837.41766366)(3502.18540219,1846.35549042)(3502.1854093,1857.43339539)
\curveto(3502.18540219,1868.51122998)(3504.84996556,1877.44905674)(3510.17910742,1884.24690247)
\curveto(3515.55018069,1891.0446242)(3522.49483012,1894.44351606)(3531.01307654,1894.44358826)
\curveto(3536.21628246,1894.44351606)(3540.66421503,1893.41545618)(3544.35688757,1891.35940552)
\curveto(3548.09142318,1889.34517826)(3551.17560283,1886.2400178)(3553.60943579,1882.04391479)
\moveto(3514.14448523,1857.43339539)
\curveto(3514.14446616,1848.91514972)(3515.88587372,1842.22227007)(3519.36871313,1837.35473633)
\curveto(3522.89346559,1832.52913401)(3527.7190528,1830.1163404)(3533.84548926,1830.11634827)
\curveto(3539.97184817,1830.1163404)(3544.79743538,1832.52913401)(3548.32226538,1837.35473633)
\curveto(3551.84698888,1842.22227007)(3553.60937725,1848.91514972)(3553.60943579,1857.43339539)
\curveto(3553.60937725,1865.95157068)(3551.84698888,1872.62346952)(3548.32226538,1877.44911194)
\curveto(3544.79743538,1882.31660558)(3539.97184817,1884.75038)(3533.84548926,1884.7504425)
\curveto(3527.7190528,1884.75038)(3522.89346559,1882.31660558)(3519.36871313,1877.44911194)
\curveto(3515.88587372,1872.62346952)(3514.14446616,1865.95157068)(3514.14448523,1857.43339539)
}
}
{
\newrgbcolor{curcolor}{0 0 0}
\pscustom[linestyle=none,fillstyle=solid,fillcolor=curcolor]
{
\newpath
\moveto(3649.34498572,1860.39169312)
\lineto(3649.34498572,1854.72686768)
\lineto(3596.09562659,1854.72686768)
\curveto(3596.59914692,1846.75412589)(3598.99095972,1840.66968983)(3603.27107214,1836.47354126)
\curveto(3607.59309345,1832.31932587)(3613.59360625,1830.24222528)(3621.27262854,1830.24223328)
\curveto(3625.72051673,1830.24222528)(3630.0215836,1830.78772645)(3634.17584204,1831.8787384)
\curveto(3638.37194756,1832.9697311)(3642.52614873,1834.60623459)(3646.63845801,1836.78825378)
\lineto(3646.63845801,1825.83625793)
\curveto(3642.4841871,1824.07386597)(3638.22508186,1822.73109388)(3633.86112952,1821.80793762)
\curveto(3629.49706325,1820.88478225)(3625.0701115,1820.42320434)(3620.58026099,1820.42320251)
\curveto(3609.33450101,1820.42320434)(3600.41765507,1823.69621132)(3593.82969641,1830.24223328)
\curveto(3587.28366552,1836.78823924)(3584.01065854,1845.64214274)(3584.01066565,1856.80397034)
\curveto(3584.01065854,1868.34338347)(3587.11581901,1877.49101837)(3593.32615637,1884.24690247)
\curveto(3599.57842251,1891.0446242)(3607.99172892,1894.44351606)(3618.56610083,1894.44358826)
\curveto(3628.04938708,1894.44351606)(3635.53953767,1891.38031722)(3641.03657507,1885.25398254)
\curveto(3646.57544583,1879.16948348)(3649.34491327,1870.88206196)(3649.34498572,1860.39169312)
\moveto(3637.76356482,1863.79058838)
\curveto(3637.6795807,1870.12675266)(3635.89621151,1875.18312883)(3632.4134519,1878.95973206)
\curveto(3628.9725429,1882.73622186)(3624.39872545,1884.62449512)(3618.69198584,1884.6245575)
\curveto(3612.22985334,1884.62449512)(3607.04759229,1882.7991643)(3603.14518713,1879.14855957)
\curveto(3599.28469111,1875.49784104)(3597.06072483,1870.35754161)(3596.47328162,1863.72764587)
\lineto(3637.76356482,1863.79058838)
}
}
{
\newrgbcolor{curcolor}{0 0 0}
\pscustom[linestyle=none,fillstyle=solid,fillcolor=curcolor]
{
\newpath
\moveto(3655.5762937,1916.22169495)
\lineto(3735.07267737,1916.22169495)
\lineto(3735.07267737,1905.52146912)
\lineto(3701.71314978,1905.52146912)
\lineto(3701.71314978,1822.24853516)
\lineto(3688.93582129,1822.24853516)
\lineto(3688.93582129,1905.52146912)
\lineto(3655.5762937,1905.52146912)
\lineto(3655.5762937,1916.22169495)
}
}
{
\newrgbcolor{curcolor}{0 0 0}
\pscustom[linestyle=none,fillstyle=solid,fillcolor=curcolor]
{
\newpath
\moveto(3756.34724402,1815.70251465)
\curveto(3753.07419556,1807.3101956)(3749.88511183,1801.83420315)(3746.77998328,1799.27452087)
\curveto(3743.6747909,1796.71488454)(3739.52058973,1795.43505489)(3734.31736731,1795.43502808)
\lineto(3725.06481909,1795.43502808)
\lineto(3725.06481909,1805.12817383)
\lineto(3731.86260962,1805.12817383)
\curveto(3735.05167635,1805.12819095)(3737.5274124,1805.88350025)(3739.2898252,1807.394104)
\curveto(3741.05218915,1808.90473746)(3743.00340485,1812.47147584)(3745.14347815,1818.09432983)
\lineto(3747.22058081,1823.38150024)
\lineto(3718.7076261,1892.74414062)
\lineto(3730.98141455,1892.74414062)
\lineto(3753.01129126,1837.60650635)
\lineto(3775.04116797,1892.74414062)
\lineto(3787.31495642,1892.74414062)
\lineto(3756.34724402,1815.70251465)
}
}
{
\newrgbcolor{curcolor}{0 0 0}
\pscustom[linestyle=none,fillstyle=solid,fillcolor=curcolor]
{
\newpath
\moveto(3814.50611853,1832.82287598)
\lineto(3814.50611853,1795.43502808)
\lineto(3802.86175513,1795.43502808)
\lineto(3802.86175513,1892.74414062)
\lineto(3814.50611853,1892.74414062)
\lineto(3814.50611853,1882.04391479)
\curveto(3816.9398696,1886.2400178)(3820.00306844,1889.34517826)(3823.69572424,1891.35940552)
\curveto(3827.43027659,1893.41545618)(3831.87820915,1894.44351606)(3837.03953528,1894.44358826)
\curveto(3845.5996615,1894.44351606)(3852.54431092,1891.0446242)(3857.87350439,1884.24690247)
\curveto(3863.24452605,1877.44905674)(3865.93007024,1868.51122998)(3865.93014502,1857.43339539)
\curveto(3865.93007024,1846.35549042)(3863.24452605,1837.41766366)(3857.87350439,1830.61988831)
\curveto(3852.54431092,1823.8220962)(3845.5996615,1820.42320434)(3837.03953528,1820.42320251)
\curveto(3831.87820915,1820.42320434)(3827.43027659,1821.43028341)(3823.69572424,1823.44444275)
\curveto(3820.00306844,1825.50056132)(3816.9398696,1828.62670261)(3814.50611853,1832.82287598)
\moveto(3853.90812659,1857.43339539)
\curveto(3853.90806383,1865.95157068)(3852.14567546,1872.62346952)(3848.62095618,1877.44911194)
\curveto(3845.13808359,1882.31660558)(3840.33347719,1884.75038)(3834.20712256,1884.7504425)
\curveto(3828.08068182,1884.75038)(3823.25509461,1882.31660558)(3819.73034644,1877.44911194)
\curveto(3816.24750274,1872.62346952)(3814.50609518,1865.95157068)(3814.50611853,1857.43339539)
\curveto(3814.50609518,1848.91514972)(3816.24750274,1842.22227007)(3819.73034644,1837.35473633)
\curveto(3823.25509461,1832.52913401)(3828.08068182,1830.1163404)(3834.20712256,1830.11634827)
\curveto(3840.33347719,1830.1163404)(3845.13808359,1832.52913401)(3848.62095618,1837.35473633)
\curveto(3852.14567546,1842.22227007)(3853.90806383,1848.91514972)(3853.90812659,1857.43339539)
}
}
{
\newrgbcolor{curcolor}{0 0 0}
\pscustom[linestyle=none,fillstyle=solid,fillcolor=curcolor]
{
\newpath
\moveto(3945.42652869,1860.39169312)
\lineto(3945.42652869,1854.72686768)
\lineto(3892.17716956,1854.72686768)
\curveto(3892.68068989,1846.75412589)(3895.07250269,1840.66968983)(3899.35261511,1836.47354126)
\curveto(3903.67463642,1832.31932587)(3909.67514922,1830.24222528)(3917.35417151,1830.24223328)
\curveto(3921.8020597,1830.24222528)(3926.10312656,1830.78772645)(3930.25738501,1831.8787384)
\curveto(3934.45349053,1832.9697311)(3938.6076917,1834.60623459)(3942.72000098,1836.78825378)
\lineto(3942.72000098,1825.83625793)
\curveto(3938.56573007,1824.07386597)(3934.30662483,1822.73109388)(3929.94267249,1821.80793762)
\curveto(3925.57860621,1820.88478225)(3921.15165446,1820.42320434)(3916.66180396,1820.42320251)
\curveto(3905.41604398,1820.42320434)(3896.49919804,1823.69621132)(3889.91123938,1830.24223328)
\curveto(3883.36520849,1836.78823924)(3880.09220151,1845.64214274)(3880.09220862,1856.80397034)
\curveto(3880.09220151,1868.34338347)(3883.19736197,1877.49101837)(3889.40769934,1884.24690247)
\curveto(3895.65996548,1891.0446242)(3904.07327188,1894.44351606)(3914.6476438,1894.44358826)
\curveto(3924.13093005,1894.44351606)(3931.62108064,1891.38031722)(3937.11811804,1885.25398254)
\curveto(3942.65698879,1879.16948348)(3945.42645624,1870.88206196)(3945.42652869,1860.39169312)
\moveto(3933.84510779,1863.79058838)
\curveto(3933.76112367,1870.12675266)(3931.97775448,1875.18312883)(3928.49499487,1878.95973206)
\curveto(3925.05408587,1882.73622186)(3920.48026842,1884.62449512)(3914.77352881,1884.6245575)
\curveto(3908.31139631,1884.62449512)(3903.12913526,1882.7991643)(3899.2267301,1879.14855957)
\curveto(3895.36623408,1875.49784104)(3893.1422678,1870.35754161)(3892.55482458,1863.72764587)
\lineto(3933.84510779,1863.79058838)
}
}
{
\newrgbcolor{curcolor}{0 0 0}
\pscustom[linestyle=none,fillstyle=solid,fillcolor=curcolor]
{
\newpath
\moveto(3992.0039823,1920.06118774)
\curveto(3986.38108418,1910.4099155)(3982.2059022,1900.86364514)(3979.47842383,1891.42234802)
\curveto(3976.75089057,1881.98091256)(3975.38713766,1872.41366138)(3975.38716101,1862.7205658)
\curveto(3975.38713766,1853.02738927)(3976.75089057,1843.41817646)(3979.47842383,1833.89289856)
\curveto(3982.24786383,1824.409559)(3986.42304581,1814.86328864)(3992.0039823,1805.25405884)
\lineto(3981.93318152,1805.25405884)
\curveto(3975.63890743,1815.1150584)(3970.91822428,1824.80819446)(3967.77111792,1834.33349609)
\curveto(3964.66594172,1843.85877356)(3963.11336148,1853.32112066)(3963.11337256,1862.7205658)
\curveto(3963.11336148,1872.07796836)(3964.66594172,1881.49835384)(3967.77111792,1890.98175049)
\curveto(3970.87626265,1900.46500967)(3975.5969458,1910.15814573)(3981.93318152,1920.06118774)
\lineto(3992.0039823,1920.06118774)
}
}
{
\newrgbcolor{curcolor}{0 0 0}
\pscustom[linestyle=none,fillstyle=solid,fillcolor=curcolor]
{
\newpath
\moveto(4012.71206641,1920.06118774)
\lineto(4022.78286719,1920.06118774)
\curveto(4029.07709099,1910.15814573)(4033.77679332,1900.46500967)(4036.88198828,1890.98175049)
\curveto(4040.02907589,1881.49835384)(4041.60263693,1872.07796836)(4041.60267615,1862.7205658)
\curveto(4041.60263693,1853.32112066)(4040.02907589,1843.85877356)(4036.88198828,1834.33349609)
\curveto(4033.77679332,1824.80819446)(4029.07709099,1815.1150584)(4022.78286719,1805.25405884)
\lineto(4012.71206641,1805.25405884)
\curveto(4018.2929526,1814.86328864)(4022.44715377,1824.409559)(4025.17468237,1833.89289856)
\curveto(4027.94412703,1843.41817646)(4029.32886076,1853.02738927)(4029.3288877,1862.7205658)
\curveto(4029.32886076,1872.41366138)(4027.94412703,1881.98091256)(4025.17468237,1891.42234802)
\curveto(4022.44715377,1900.86364514)(4018.2929526,1910.4099155)(4012.71206641,1920.06118774)
}
}
{
\newrgbcolor{curcolor}{0 0 0}
\pscustom[linestyle=none,fillstyle=solid,fillcolor=curcolor]
{
}
}
{
\newrgbcolor{curcolor}{0 0 0}
\pscustom[linestyle=none,fillstyle=solid,fillcolor=curcolor]
{
\newpath
\moveto(4108.88821387,1838.2359314)
\lineto(4122.1690824,1838.2359314)
\lineto(4122.1690824,1822.24853516)
\lineto(4108.88821387,1822.24853516)
\lineto(4108.88821387,1838.2359314)
\moveto(4108.88821387,1888.90464783)
\lineto(4122.1690824,1888.90464783)
\lineto(4122.1690824,1872.91725159)
\lineto(4108.88821387,1872.91725159)
\lineto(4108.88821387,1888.90464783)
}
}
{
\newrgbcolor{curcolor}{0 0 0}
\pscustom[linestyle=none,fillstyle=solid,fillcolor=curcolor]
{
}
}
{
\newrgbcolor{curcolor}{0 0 0}
\pscustom[linestyle=none,fillstyle=solid,fillcolor=curcolor]
{
\newpath
\moveto(4191.02818274,1916.22169495)
\lineto(4208.14854407,1916.22169495)
\lineto(4249.8164823,1837.60650635)
\lineto(4249.8164823,1916.22169495)
\lineto(4262.15321326,1916.22169495)
\lineto(4262.15321326,1822.24853516)
\lineto(4245.03285193,1822.24853516)
\lineto(4203.3649137,1900.86372375)
\lineto(4203.3649137,1822.24853516)
\lineto(4191.02818274,1822.24853516)
\lineto(4191.02818274,1916.22169495)
}
}
{
\newrgbcolor{curcolor}{0 0 0}
\pscustom[linestyle=none,fillstyle=solid,fillcolor=curcolor]
{
\newpath
\moveto(4314.2696073,1884.6245575)
\curveto(4308.0592469,1884.62449512)(4303.14973643,1882.1907207)(4299.54106116,1877.32322693)
\curveto(4295.93233642,1872.49758464)(4294.12798641,1865.86764742)(4294.12800574,1857.43339539)
\curveto(4294.12798641,1848.99907298)(4295.9113556,1842.34815495)(4299.47811865,1837.48062134)
\curveto(4303.08679398,1832.65501889)(4308.01728527,1830.24222528)(4314.2696073,1830.24223328)
\curveto(4320.43792715,1830.24222528)(4325.3264568,1832.6759997)(4328.93521094,1837.54356384)
\curveto(4332.54385681,1842.41109739)(4334.34820681,1849.04103461)(4334.34826636,1857.43339539)
\curveto(4334.34820681,1865.78372416)(4332.54385681,1872.39268057)(4328.93521094,1877.26028442)
\curveto(4325.3264568,1882.16973988)(4320.43792715,1884.62449512)(4314.2696073,1884.6245575)
\moveto(4314.2696073,1894.44358826)
\curveto(4324.34035855,1894.44351606)(4332.25012542,1891.17050908)(4337.99893164,1884.6245575)
\curveto(4343.74761148,1878.07848116)(4346.62198299,1869.01476952)(4346.62205481,1857.43339539)
\curveto(4346.62198299,1845.89391251)(4343.74761148,1836.83020087)(4337.99893164,1830.24223328)
\curveto(4332.25012542,1823.69621132)(4324.34035855,1820.42320434)(4314.2696073,1820.42320251)
\curveto(4304.1568155,1820.42320434)(4296.22606781,1823.69621132)(4290.47734045,1830.24223328)
\curveto(4284.77054338,1836.83020087)(4281.91715268,1845.89391251)(4281.91715979,1857.43339539)
\curveto(4281.91715268,1869.01476952)(4284.77054338,1878.07848116)(4290.47734045,1884.6245575)
\curveto(4296.22606781,1891.17050908)(4304.1568155,1894.44351606)(4314.2696073,1894.44358826)
}
}
{
\newrgbcolor{curcolor}{0 0 0}
\pscustom[linestyle=none,fillstyle=solid,fillcolor=curcolor]
{
\newpath
\moveto(4412.14520239,1882.04391479)
\lineto(4412.14520239,1920.18707275)
\lineto(4423.72662329,1920.18707275)
\lineto(4423.72662329,1822.24853516)
\lineto(4412.14520239,1822.24853516)
\lineto(4412.14520239,1832.82287598)
\curveto(4409.71136943,1828.62670261)(4406.62718978,1825.50056132)(4402.89265417,1823.44444275)
\curveto(4399.19998163,1821.43028341)(4394.75204907,1820.42320434)(4389.54884314,1820.42320251)
\curveto(4381.03059672,1820.42320434)(4374.0859473,1823.8220962)(4368.71487402,1830.61988831)
\curveto(4363.38573217,1837.41766366)(4360.72116879,1846.35549042)(4360.7211759,1857.43339539)
\curveto(4360.72116879,1868.51122998)(4363.38573217,1877.44905674)(4368.71487402,1884.24690247)
\curveto(4374.0859473,1891.0446242)(4381.03059672,1894.44351606)(4389.54884314,1894.44358826)
\curveto(4394.75204907,1894.44351606)(4399.19998163,1893.41545618)(4402.89265417,1891.35940552)
\curveto(4406.62718978,1889.34517826)(4409.71136943,1886.2400178)(4412.14520239,1882.04391479)
\moveto(4372.68025183,1857.43339539)
\curveto(4372.68023276,1848.91514972)(4374.42164032,1842.22227007)(4377.90447974,1837.35473633)
\curveto(4381.42923219,1832.52913401)(4386.2548194,1830.1163404)(4392.38125586,1830.11634827)
\curveto(4398.50761477,1830.1163404)(4403.33320198,1832.52913401)(4406.85803198,1837.35473633)
\curveto(4410.38275548,1842.22227007)(4412.14514386,1848.91514972)(4412.14520239,1857.43339539)
\curveto(4412.14514386,1865.95157068)(4410.38275548,1872.62346952)(4406.85803198,1877.44911194)
\curveto(4403.33320198,1882.31660558)(4398.50761477,1884.75038)(4392.38125586,1884.7504425)
\curveto(4386.2548194,1884.75038)(4381.42923219,1882.31660558)(4377.90447974,1877.44911194)
\curveto(4374.42164032,1872.62346952)(4372.68023276,1865.95157068)(4372.68025183,1857.43339539)
}
}
{
\newrgbcolor{curcolor}{0 0 0}
\pscustom[linestyle=none,fillstyle=solid,fillcolor=curcolor]
{
\newpath
\moveto(4507.88075232,1860.39169312)
\lineto(4507.88075232,1854.72686768)
\lineto(4454.63139319,1854.72686768)
\curveto(4455.13491353,1846.75412589)(4457.52672632,1840.66968983)(4461.80683875,1836.47354126)
\curveto(4466.12886005,1832.31932587)(4472.12937285,1830.24222528)(4479.80839514,1830.24223328)
\curveto(4484.25628333,1830.24222528)(4488.5573502,1830.78772645)(4492.71160864,1831.8787384)
\curveto(4496.90771416,1832.9697311)(4501.06191533,1834.60623459)(4505.17422461,1836.78825378)
\lineto(4505.17422461,1825.83625793)
\curveto(4501.0199537,1824.07386597)(4496.76084846,1822.73109388)(4492.39689612,1821.80793762)
\curveto(4488.03282985,1820.88478225)(4483.6058781,1820.42320434)(4479.11602759,1820.42320251)
\curveto(4467.87026761,1820.42320434)(4458.95342167,1823.69621132)(4452.36546301,1830.24223328)
\curveto(4445.81943212,1836.78823924)(4442.54642514,1845.64214274)(4442.54643225,1856.80397034)
\curveto(4442.54642514,1868.34338347)(4445.65158561,1877.49101837)(4451.86192297,1884.24690247)
\curveto(4458.11418911,1891.0446242)(4466.52749552,1894.44351606)(4477.10186743,1894.44358826)
\curveto(4486.58515368,1894.44351606)(4494.07530427,1891.38031722)(4499.57234167,1885.25398254)
\curveto(4505.11121243,1879.16948348)(4507.88067987,1870.88206196)(4507.88075232,1860.39169312)
\moveto(4496.29933142,1863.79058838)
\curveto(4496.2153473,1870.12675266)(4494.43197811,1875.18312883)(4490.94921851,1878.95973206)
\curveto(4487.5083095,1882.73622186)(4482.93449205,1884.62449512)(4477.22775244,1884.6245575)
\curveto(4470.76561994,1884.62449512)(4465.58335889,1882.7991643)(4461.68095374,1879.14855957)
\curveto(4457.82045772,1875.49784104)(4455.59649143,1870.35754161)(4455.00904822,1863.72764587)
\lineto(4496.29933142,1863.79058838)
}
}
{
\newrgbcolor{curcolor}{0 0 0}
\pscustom[linestyle=none,fillstyle=solid,fillcolor=curcolor]
{
\newpath
\moveto(4514.1120603,1916.22169495)
\lineto(4593.60844397,1916.22169495)
\lineto(4593.60844397,1905.52146912)
\lineto(4560.24891638,1905.52146912)
\lineto(4560.24891638,1822.24853516)
\lineto(4547.47158789,1822.24853516)
\lineto(4547.47158789,1905.52146912)
\lineto(4514.1120603,1905.52146912)
\lineto(4514.1120603,1916.22169495)
}
}
{
\newrgbcolor{curcolor}{0 0 0}
\pscustom[linestyle=none,fillstyle=solid,fillcolor=curcolor]
{
\newpath
\moveto(4614.88301062,1815.70251465)
\curveto(4611.60996216,1807.3101956)(4608.42087844,1801.83420315)(4605.31574988,1799.27452087)
\curveto(4602.2105575,1796.71488454)(4598.05635633,1795.43505489)(4592.85313391,1795.43502808)
\lineto(4583.60058569,1795.43502808)
\lineto(4583.60058569,1805.12817383)
\lineto(4590.39837622,1805.12817383)
\curveto(4593.58744295,1805.12819095)(4596.063179,1805.88350025)(4597.8255918,1807.394104)
\curveto(4599.58795575,1808.90473746)(4601.53917145,1812.47147584)(4603.67924475,1818.09432983)
\lineto(4605.75634741,1823.38150024)
\lineto(4577.2433927,1892.74414062)
\lineto(4589.51718115,1892.74414062)
\lineto(4611.54705786,1837.60650635)
\lineto(4633.57693457,1892.74414062)
\lineto(4645.85072302,1892.74414062)
\lineto(4614.88301062,1815.70251465)
}
}
{
\newrgbcolor{curcolor}{0 0 0}
\pscustom[linestyle=none,fillstyle=solid,fillcolor=curcolor]
{
\newpath
\moveto(4673.04188513,1832.82287598)
\lineto(4673.04188513,1795.43502808)
\lineto(4661.39752173,1795.43502808)
\lineto(4661.39752173,1892.74414062)
\lineto(4673.04188513,1892.74414062)
\lineto(4673.04188513,1882.04391479)
\curveto(4675.4756362,1886.2400178)(4678.53883504,1889.34517826)(4682.23149084,1891.35940552)
\curveto(4685.96604319,1893.41545618)(4690.41397576,1894.44351606)(4695.57530188,1894.44358826)
\curveto(4704.1354281,1894.44351606)(4711.08007753,1891.0446242)(4716.409271,1884.24690247)
\curveto(4721.78029266,1877.44905674)(4724.46583685,1868.51122998)(4724.46591162,1857.43339539)
\curveto(4724.46583685,1846.35549042)(4721.78029266,1837.41766366)(4716.409271,1830.61988831)
\curveto(4711.08007753,1823.8220962)(4704.1354281,1820.42320434)(4695.57530188,1820.42320251)
\curveto(4690.41397576,1820.42320434)(4685.96604319,1821.43028341)(4682.23149084,1823.44444275)
\curveto(4678.53883504,1825.50056132)(4675.4756362,1828.62670261)(4673.04188513,1832.82287598)
\moveto(4712.44389319,1857.43339539)
\curveto(4712.44383043,1865.95157068)(4710.68144206,1872.62346952)(4707.15672278,1877.44911194)
\curveto(4703.67385019,1882.31660558)(4698.86924379,1884.75038)(4692.74288916,1884.7504425)
\curveto(4686.61644843,1884.75038)(4681.79086121,1882.31660558)(4678.26611304,1877.44911194)
\curveto(4674.78326934,1872.62346952)(4673.04186178,1865.95157068)(4673.04188513,1857.43339539)
\curveto(4673.04186178,1848.91514972)(4674.78326934,1842.22227007)(4678.26611304,1837.35473633)
\curveto(4681.79086121,1832.52913401)(4686.61644843,1830.1163404)(4692.74288916,1830.11634827)
\curveto(4698.86924379,1830.1163404)(4703.67385019,1832.52913401)(4707.15672278,1837.35473633)
\curveto(4710.68144206,1842.22227007)(4712.44383043,1848.91514972)(4712.44389319,1857.43339539)
}
}
{
\newrgbcolor{curcolor}{0 0 0}
\pscustom[linestyle=none,fillstyle=solid,fillcolor=curcolor]
{
\newpath
\moveto(4803.96229529,1860.39169312)
\lineto(4803.96229529,1854.72686768)
\lineto(4750.71293616,1854.72686768)
\curveto(4751.2164565,1846.75412589)(4753.60826929,1840.66968983)(4757.88838171,1836.47354126)
\curveto(4762.21040302,1832.31932587)(4768.21091582,1830.24222528)(4775.88993811,1830.24223328)
\curveto(4780.3378263,1830.24222528)(4784.63889317,1830.78772645)(4788.79315161,1831.8787384)
\curveto(4792.98925713,1832.9697311)(4797.1434583,1834.60623459)(4801.25576758,1836.78825378)
\lineto(4801.25576758,1825.83625793)
\curveto(4797.10149667,1824.07386597)(4792.84239143,1822.73109388)(4788.47843909,1821.80793762)
\curveto(4784.11437282,1820.88478225)(4779.68742107,1820.42320434)(4775.19757056,1820.42320251)
\curveto(4763.95181058,1820.42320434)(4755.03496464,1823.69621132)(4748.44700598,1830.24223328)
\curveto(4741.90097509,1836.78823924)(4738.62796811,1845.64214274)(4738.62797522,1856.80397034)
\curveto(4738.62796811,1868.34338347)(4741.73312858,1877.49101837)(4747.94346594,1884.24690247)
\curveto(4754.19573208,1891.0446242)(4762.60903849,1894.44351606)(4773.1834104,1894.44358826)
\curveto(4782.66669665,1894.44351606)(4790.15684724,1891.38031722)(4795.65388464,1885.25398254)
\curveto(4801.1927554,1879.16948348)(4803.96222284,1870.88206196)(4803.96229529,1860.39169312)
\moveto(4792.38087439,1863.79058838)
\curveto(4792.29689027,1870.12675266)(4790.51352108,1875.18312883)(4787.03076147,1878.95973206)
\curveto(4783.58985247,1882.73622186)(4779.01603502,1884.62449512)(4773.30929541,1884.6245575)
\curveto(4766.84716291,1884.62449512)(4761.66490186,1882.7991643)(4757.7624967,1879.14855957)
\curveto(4753.90200068,1875.49784104)(4751.6780344,1870.35754161)(4751.09059119,1863.72764587)
\lineto(4792.38087439,1863.79058838)
}
}
{
\newrgbcolor{curcolor}{0.51764709 0.74509805 0.51764709}
\pscustom[linestyle=none,fillstyle=solid,fillcolor=curcolor]
{
\newpath
\moveto(3038.04159375,1702.1686416)
\curveto(3038.04159375,1682.75238087)(3022.30160449,1667.0123916)(3002.88534375,1667.0123916)
\curveto(2983.46908301,1667.0123916)(2967.72909375,1682.75238087)(2967.72909375,1702.1686416)
\curveto(2967.72909375,1721.58490234)(2983.46908301,1737.3248916)(3002.88534375,1737.3248916)
\curveto(3022.30160449,1737.3248916)(3038.04159375,1721.58490234)(3038.04159375,1702.1686416)
\closepath
}
}
{
\newrgbcolor{curcolor}{0.01176471 0.50196081 0.28235295}
\pscustom[linewidth=11.71875,linecolor=curcolor]
{
\newpath
\moveto(3038.04159375,1702.1686416)
\curveto(3038.04159375,1682.75238087)(3022.30160449,1667.0123916)(3002.88534375,1667.0123916)
\curveto(2983.46908301,1667.0123916)(2967.72909375,1682.75238087)(2967.72909375,1702.1686416)
\curveto(2967.72909375,1721.58490234)(2983.46908301,1737.3248916)(3002.88534375,1737.3248916)
\curveto(3022.30160449,1737.3248916)(3038.04159375,1721.58490234)(3038.04159375,1702.1686416)
\closepath
}
}
{
\newrgbcolor{curcolor}{0 0 0}
\pscustom[linestyle=none,fillstyle=solid,fillcolor=curcolor]
{
\newpath
\moveto(3166.89068579,1708.32252502)
\curveto(3166.89062725,1716.71481455)(3165.14921969,1723.21886688)(3161.66645789,1727.83470154)
\curveto(3158.22555108,1732.45042504)(3153.37898305,1734.75831457)(3147.12673926,1734.75837708)
\curveto(3140.91637955,1734.75831457)(3136.06981152,1732.45042504)(3132.58702063,1727.83470154)
\curveto(3129.1461429,1723.21886688)(3127.42571616,1716.71481455)(3127.42573523,1708.32252502)
\curveto(3127.42571616,1699.97212499)(3129.1461429,1693.48905347)(3132.58702063,1688.87329102)
\curveto(3136.06981152,1684.25749532)(3140.91637955,1681.94960579)(3147.12673926,1681.94961548)
\curveto(3153.37898305,1681.94960579)(3158.22555108,1684.25749532)(3161.66645789,1688.87329102)
\curveto(3165.14921969,1693.48905347)(3166.89062725,1699.97212499)(3166.89068579,1708.32252502)
\moveto(3178.47210669,1681.00547791)
\curveto(3178.47203657,1669.00444356)(3175.8074732,1660.08759762)(3170.47840857,1654.25491333)
\curveto(3165.14921969,1648.38030342)(3156.98768306,1645.44298946)(3145.99377417,1645.44296265)
\curveto(3141.92345862,1645.44298946)(3138.08396966,1645.75770167)(3134.47529578,1646.38710022)
\curveto(3130.86656965,1646.97458888)(3127.36277372,1647.8977447)(3123.96389746,1649.15657043)
\lineto(3123.96389746,1660.42327881)
\curveto(3127.36277372,1658.57697901)(3130.71970395,1657.2132261)(3134.03469824,1656.33201599)
\curveto(3137.34964117,1655.45083773)(3140.72755222,1655.01024063)(3144.16844153,1655.01022339)
\curveto(3151.76346037,1655.01024063)(3157.44926096,1657.00341796)(3161.22586035,1660.98976135)
\curveto(3165.002354,1664.93416565)(3166.89062725,1670.91369763)(3166.89068579,1678.92837524)
\lineto(3166.89068579,1684.65614319)
\curveto(3164.49881446,1680.50192962)(3161.43561562,1677.39676915)(3157.70108008,1675.34065247)
\curveto(3153.96644584,1673.28452961)(3149.49753246,1672.25646973)(3144.29432654,1672.25646973)
\curveto(3135.65019524,1672.25646973)(3128.684565,1675.55045752)(3123.39741492,1682.13844299)
\curveto(3118.11023475,1688.7264087)(3115.46665219,1697.45442732)(3115.4666593,1708.32252502)
\curveto(3115.46665219,1719.23251223)(3118.11023475,1727.98151166)(3123.39741492,1734.56954956)
\curveto(3128.684565,1741.15746284)(3135.65019524,1744.45145063)(3144.29432654,1744.45152283)
\curveto(3149.49753246,1744.45145063)(3153.96644584,1743.42339075)(3157.70108008,1741.36734009)
\curveto(3161.43561562,1739.31115121)(3164.49881446,1736.20599074)(3166.89068579,1732.05184937)
\lineto(3166.89068579,1742.7520752)
\lineto(3178.47210669,1742.7520752)
\lineto(3178.47210669,1681.00547791)
}
}
{
\newrgbcolor{curcolor}{0 0 0}
\pscustom[linestyle=none,fillstyle=solid,fillcolor=curcolor]
{
\newpath
\moveto(3262.62623572,1710.39962769)
\lineto(3262.62623572,1704.73480225)
\lineto(3209.37687659,1704.73480225)
\curveto(3209.88039692,1696.76206046)(3212.27220972,1690.6776244)(3216.55232214,1686.48147583)
\curveto(3220.87434345,1682.32726044)(3226.87485625,1680.25015985)(3234.55387854,1680.25016785)
\curveto(3239.00176673,1680.25015985)(3243.3028336,1680.79566102)(3247.45709204,1681.88667297)
\curveto(3251.65319756,1682.97766567)(3255.80739873,1684.61416916)(3259.91970801,1686.79618835)
\lineto(3259.91970801,1675.8441925)
\curveto(3255.7654371,1674.08180054)(3251.50633186,1672.73902845)(3247.14237952,1671.81587219)
\curveto(3242.77831325,1670.89271682)(3238.3513615,1670.43113891)(3233.86151099,1670.43113708)
\curveto(3222.61575101,1670.43113891)(3213.69890507,1673.70414589)(3207.11094641,1680.25016785)
\curveto(3200.56491552,1686.79617381)(3197.29190854,1695.65007731)(3197.29191565,1706.81190491)
\curveto(3197.29190854,1718.35131804)(3200.39706901,1727.49895294)(3206.60740637,1734.25483704)
\curveto(3212.85967251,1741.05255877)(3221.27297892,1744.45145063)(3231.84735083,1744.45152283)
\curveto(3241.33063708,1744.45145063)(3248.82078767,1741.38825179)(3254.31782507,1735.26191711)
\curveto(3259.85669583,1729.17741805)(3262.62616327,1720.88999653)(3262.62623572,1710.39962769)
\moveto(3251.04481482,1713.79852295)
\curveto(3250.9608307,1720.13468723)(3249.17746151,1725.1910634)(3245.6947019,1728.96766663)
\curveto(3242.2537929,1732.74415643)(3237.67997545,1734.63242969)(3231.97323584,1734.63249207)
\curveto(3225.51110334,1734.63242969)(3220.32884229,1732.80709887)(3216.42643713,1729.15649414)
\curveto(3212.56594111,1725.50577561)(3210.34197483,1720.36547618)(3209.75453162,1713.73558044)
\lineto(3251.04481482,1713.79852295)
}
}
{
\newrgbcolor{curcolor}{0 0 0}
\pscustom[linestyle=none,fillstyle=solid,fillcolor=curcolor]
{
\newpath
\moveto(3292.83863806,1762.76779175)
\lineto(3292.83863806,1742.7520752)
\lineto(3316.69384741,1742.7520752)
\lineto(3316.69384741,1733.751297)
\lineto(3292.83863806,1733.751297)
\lineto(3292.83863806,1695.48225403)
\curveto(3292.83861446,1689.73348777)(3293.61490458,1686.04086451)(3295.16751074,1684.40437317)
\curveto(3296.76202667,1682.76785753)(3299.97209121,1681.94960579)(3304.79771399,1681.94961548)
\lineto(3316.69384741,1681.94961548)
\lineto(3316.69384741,1672.25646973)
\lineto(3304.79771399,1672.25646973)
\curveto(3295.85985167,1672.25646973)(3289.69149236,1673.91395403)(3286.29261755,1677.22892761)
\curveto(3282.89370863,1680.58585288)(3281.1942627,1686.67028893)(3281.19427466,1695.48225403)
\lineto(3281.19427466,1733.751297)
\lineto(3272.6970365,1733.751297)
\lineto(3272.6970365,1742.7520752)
\lineto(3281.19427466,1742.7520752)
\lineto(3281.19427466,1762.76779175)
\lineto(3292.83863806,1762.76779175)
}
}
{
\newrgbcolor{curcolor}{0 0 0}
\pscustom[linestyle=none,fillstyle=solid,fillcolor=curcolor]
{
\newpath
\moveto(3402.8621366,1758.99124146)
\lineto(3402.8621366,1745.58448792)
\curveto(3398.58196752,1749.57076924)(3394.00815008,1752.55004483)(3389.14067053,1754.52232361)
\curveto(3384.31501402,1756.49443786)(3379.17471459,1757.48053611)(3373.71975684,1757.48062134)
\curveto(3362.9775262,1757.48053611)(3354.75304712,1754.18654832)(3349.04629492,1747.59864807)
\curveto(3343.33948431,1741.05255877)(3340.48609361,1731.56923085)(3340.48611426,1719.14863586)
\curveto(3340.48609361,1706.76990872)(3343.33948431,1697.2865808)(3349.04629492,1690.69862366)
\curveto(3354.75304712,1684.15259125)(3362.9775262,1680.87958427)(3373.71975684,1680.8795929)
\curveto(3379.17471459,1680.87958427)(3384.31501402,1681.86568253)(3389.14067053,1683.83789062)
\curveto(3394.00815008,1685.81007556)(3398.58196752,1688.78935114)(3402.8621366,1692.77572632)
\lineto(3402.8621366,1679.49485779)
\curveto(3398.41412101,1676.47361334)(3393.69343787,1674.20768543)(3388.700073,1672.69706726)
\curveto(3383.74853204,1671.18644821)(3378.50332854,1670.43113891)(3372.96444678,1670.43113708)
\curveto(3358.73940178,1670.43113891)(3347.53564711,1674.7741674)(3339.35314917,1683.4602356)
\curveto(3331.17061221,1692.18824301)(3327.07935348,1704.08436453)(3327.07936072,1719.14863586)
\curveto(3327.07935348,1734.25477504)(3331.17061221,1746.15089656)(3339.35314917,1754.83703613)
\curveto(3347.53564711,1763.56497217)(3358.73940178,1767.92898148)(3372.96444678,1767.92907715)
\curveto(3378.5872518,1767.92898148)(3383.87441692,1767.17367217)(3388.82595801,1765.66314697)
\curveto(3393.81932275,1764.19439659)(3398.49804427,1761.97043031)(3402.8621366,1758.99124146)
}
}
{
\newrgbcolor{curcolor}{0 0 0}
\pscustom[linestyle=none,fillstyle=solid,fillcolor=curcolor]
{
\newpath
\moveto(3449.43959021,1734.63249207)
\curveto(3443.22922981,1734.63242969)(3438.31971934,1732.19865527)(3434.71104407,1727.3311615)
\curveto(3431.10231933,1722.50551921)(3429.29796932,1715.87558199)(3429.29798865,1707.44132996)
\curveto(3429.29796932,1699.00700755)(3431.08133851,1692.35608952)(3434.64810156,1687.48855591)
\curveto(3438.25677689,1682.66295346)(3443.18726818,1680.25015985)(3449.43959021,1680.25016785)
\curveto(3455.60791006,1680.25015985)(3460.49643971,1682.68393427)(3464.10519385,1687.55149841)
\curveto(3467.71383972,1692.41903196)(3469.51818972,1699.04896918)(3469.51824927,1707.44132996)
\curveto(3469.51818972,1715.79165874)(3467.71383972,1722.40061514)(3464.10519385,1727.26821899)
\curveto(3460.49643971,1732.17767445)(3455.60791006,1734.63242969)(3449.43959021,1734.63249207)
\moveto(3449.43959021,1744.45152283)
\curveto(3459.51034146,1744.45145063)(3467.42010833,1741.17844365)(3473.16891455,1734.63249207)
\curveto(3478.91759439,1728.08641573)(3481.7919659,1719.02270409)(3481.79203772,1707.44132996)
\curveto(3481.7919659,1695.90184708)(3478.91759439,1686.83813544)(3473.16891455,1680.25016785)
\curveto(3467.42010833,1673.70414589)(3459.51034146,1670.43113891)(3449.43959021,1670.43113708)
\curveto(3439.32679841,1670.43113891)(3431.39605072,1673.70414589)(3425.64732336,1680.25016785)
\curveto(3419.94052629,1686.83813544)(3417.08713559,1695.90184708)(3417.0871427,1707.44132996)
\curveto(3417.08713559,1719.02270409)(3419.94052629,1728.08641573)(3425.64732336,1734.63249207)
\curveto(3431.39605072,1741.17844365)(3439.32679841,1744.45145063)(3449.43959021,1744.45152283)
}
}
{
\newrgbcolor{curcolor}{0 0 0}
\pscustom[linestyle=none,fillstyle=solid,fillcolor=curcolor]
{
\newpath
\moveto(3559.52603125,1714.80560303)
\lineto(3559.52603125,1672.25646973)
\lineto(3547.94461035,1672.25646973)
\lineto(3547.94461035,1714.427948)
\curveto(3547.94455119,1721.09980467)(3546.64374072,1726.0932384)(3544.04217505,1729.40826416)
\curveto(3541.44049885,1732.72317562)(3537.53806745,1734.38065992)(3532.33486914,1734.38072205)
\curveto(3526.08254302,1734.38065992)(3521.15205173,1732.38748259)(3517.54338049,1728.40118408)
\curveto(3513.93465172,1724.41477328)(3512.13030172,1718.98074246)(3512.13032507,1712.09907532)
\lineto(3512.13032507,1672.25646973)
\lineto(3500.48596167,1672.25646973)
\lineto(3500.48596167,1742.7520752)
\lineto(3512.13032507,1742.7520752)
\lineto(3512.13032507,1731.80007935)
\curveto(3514.89976917,1736.03814423)(3518.15179533,1739.20624714)(3521.88641333,1741.30439758)
\curveto(3525.66292674,1743.40240993)(3530.00595523,1744.45145063)(3534.91551184,1744.45152283)
\curveto(3543.0140599,1744.45145063)(3549.14045758,1741.93375295)(3553.29472327,1736.89842224)
\curveto(3557.44885992,1731.90492387)(3559.5259605,1724.54065816)(3559.52603125,1714.80560303)
}
}
{
\newrgbcolor{curcolor}{0 0 0}
\pscustom[linestyle=none,fillstyle=solid,fillcolor=curcolor]
{
\newpath
\moveto(3633.48347449,1740.04554749)
\lineto(3633.48347449,1729.21943665)
\curveto(3630.21040463,1731.02372968)(3626.91641683,1732.36650178)(3623.60150122,1733.24775696)
\curveto(3620.32844124,1734.17085178)(3617.01347263,1734.63242969)(3613.65658545,1734.63249207)
\curveto(3606.14541099,1734.63242969)(3600.31274471,1732.2406169)(3596.15856909,1727.45704651)
\curveto(3592.00434237,1722.71532735)(3589.92724179,1716.0434285)(3589.92726111,1707.44132996)
\curveto(3589.92724179,1698.83916104)(3592.00434237,1692.14628138)(3596.15856909,1687.3626709)
\curveto(3600.31274471,1682.62099183)(3606.14541099,1680.25015985)(3613.65658545,1680.25016785)
\curveto(3617.01347263,1680.25015985)(3620.32844124,1680.69075695)(3623.60150122,1681.57196045)
\curveto(3626.91641683,1682.49510695)(3630.21040463,1683.85885986)(3633.48347449,1685.66322327)
\lineto(3633.48347449,1674.96299744)
\curveto(3630.25236625,1673.45237612)(3626.89543602,1672.31941217)(3623.41267371,1671.56410217)
\curveto(3619.9717674,1670.80879356)(3616.30012496,1670.43113891)(3612.39773535,1670.43113708)
\curveto(3601.78140168,1670.43113891)(3593.34711446,1673.76708833)(3587.09484839,1680.43899536)
\curveto(3580.84254933,1687.11088602)(3577.71640805,1696.11165522)(3577.71641516,1707.44132996)
\curveto(3577.71640805,1718.93878083)(3580.86353015,1727.98151166)(3587.15779089,1734.56954956)
\curveto(3593.49398016,1741.15746284)(3602.15905634,1744.45145063)(3613.15304541,1744.45152283)
\curveto(3616.71974124,1744.45145063)(3620.20255636,1744.07379598)(3623.60150122,1743.31855774)
\curveto(3627.00034009,1742.605139)(3630.29432788,1741.51413667)(3633.48347449,1740.04554749)
}
}
{
\newrgbcolor{curcolor}{0 0 0}
\pscustom[linestyle=none,fillstyle=solid,fillcolor=curcolor]
{
\newpath
\moveto(3653.49919104,1770.19500732)
\lineto(3665.08061194,1770.19500732)
\lineto(3665.08061194,1672.25646973)
\lineto(3653.49919104,1672.25646973)
\lineto(3653.49919104,1770.19500732)
}
}
{
\newrgbcolor{curcolor}{0 0 0}
\pscustom[linestyle=none,fillstyle=solid,fillcolor=curcolor]
{
\newpath
\moveto(3688.05462622,1700.07705688)
\lineto(3688.05462622,1742.7520752)
\lineto(3699.63604712,1742.7520752)
\lineto(3699.63604712,1700.51765442)
\curveto(3699.63602459,1693.84572731)(3700.93683505,1688.83131277)(3703.53848242,1685.47439575)
\curveto(3706.14007692,1682.15941393)(3710.04250832,1680.50192962)(3715.24578833,1680.50193787)
\curveto(3721.49803275,1680.50192962)(3726.42852404,1682.49510695)(3730.03727698,1686.48147583)
\curveto(3733.68788568,1690.46781626)(3735.51321649,1695.90184708)(3735.5132749,1702.78358459)
\lineto(3735.5132749,1742.7520752)
\lineto(3747.0946958,1742.7520752)
\lineto(3747.0946958,1672.25646973)
\lineto(3735.5132749,1672.25646973)
\lineto(3735.5132749,1683.08258057)
\curveto(3732.70178742,1678.80248369)(3729.42878044,1675.61339996)(3725.69424414,1673.51531982)
\curveto(3722.00157229,1671.4591988)(3717.70050542,1670.43113891)(3712.79103064,1670.43113708)
\curveto(3704.69240075,1670.43113891)(3698.54502226,1672.94883659)(3694.34887671,1677.98423767)
\curveto(3690.15269667,1683.0196273)(3688.05461527,1690.38389301)(3688.05462622,1700.07705688)
\moveto(3717.19700598,1744.45152283)
\lineto(3717.19700598,1744.45152283)
}
}
{
\newrgbcolor{curcolor}{0 0 0}
\pscustom[linestyle=none,fillstyle=solid,fillcolor=curcolor]
{
\newpath
\moveto(3816.01673865,1740.67497253)
\lineto(3816.01673865,1729.72297668)
\curveto(3812.74367458,1731.40138434)(3809.34478271,1732.66023318)(3805.82005286,1733.49952698)
\curveto(3802.29522922,1734.33869829)(3798.64456758,1734.75831457)(3794.86805701,1734.75837708)
\curveto(3789.11927804,1734.75831457)(3784.79723036,1733.87712039)(3781.901901,1732.11479187)
\curveto(3779.04848733,1730.35234364)(3777.62179198,1727.70876108)(3777.62181067,1724.18403625)
\curveto(3777.62179198,1721.49844014)(3778.64985186,1719.37937793)(3780.70599341,1717.82684326)
\curveto(3782.7620914,1716.31617908)(3786.89531175,1714.86850292)(3793.10566687,1713.48381042)
\lineto(3797.07104468,1712.60261536)
\curveto(3805.29548561,1710.84018664)(3811.1281519,1708.34346977)(3814.56906104,1705.11245728)
\curveto(3818.05182051,1701.92334069)(3819.79322808,1697.45442732)(3819.79328894,1691.70570374)
\curveto(3819.79322808,1685.15967032)(3817.19160714,1679.97740927)(3811.98841833,1676.15890503)
\curveto(3806.82708504,1672.34039298)(3799.7145891,1670.43113891)(3790.65090918,1670.43113708)
\curveto(3786.87433094,1670.43113891)(3782.92993791,1670.80879356)(3778.81771826,1671.56410217)
\curveto(3774.74742046,1672.27745054)(3770.44635359,1673.36845287)(3765.91450476,1674.83711243)
\lineto(3765.91450476,1686.79618835)
\curveto(3770.19458383,1684.57220753)(3774.41172744,1682.89374241)(3778.56594824,1681.76078796)
\curveto(3782.72012977,1680.66977613)(3786.83236931,1680.12427497)(3790.9026792,1680.12428284)
\curveto(3796.35765886,1680.12427497)(3800.55382166,1681.04743078)(3803.49118018,1682.89375305)
\curveto(3806.42844957,1684.78201567)(3807.89710655,1687.42559823)(3807.89715552,1690.82450867)
\curveto(3807.89710655,1693.9716122)(3806.82708504,1696.3844058)(3804.68708777,1698.06289673)
\curveto(3802.58896061,1699.74133604)(3797.95220072,1701.35685872)(3790.77679419,1702.9094696)
\lineto(3786.74847388,1703.85360718)
\curveto(3779.57300768,1705.36419419)(3774.39074662,1707.67208372)(3771.20167517,1710.77728271)
\curveto(3768.01257917,1713.92436629)(3766.41803731,1718.22543316)(3766.4180448,1723.68049622)
\curveto(3766.41803731,1730.31038201)(3768.76788848,1735.42970062)(3773.46760535,1739.03846741)
\curveto(3778.16729314,1742.64710063)(3784.83919198,1744.45145063)(3793.4833219,1744.45152283)
\curveto(3797.7633734,1744.45145063)(3801.79168968,1744.13673842)(3805.56828284,1743.50738525)
\curveto(3809.34478271,1742.87788958)(3812.82759783,1741.93375295)(3816.01673865,1740.67497253)
}
}
{
\newrgbcolor{curcolor}{0 0 0}
\pscustom[linestyle=none,fillstyle=solid,fillcolor=curcolor]
{
\newpath
\moveto(3838.29838538,1742.7520752)
\lineto(3849.87980627,1742.7520752)
\lineto(3849.87980627,1672.25646973)
\lineto(3838.29838538,1672.25646973)
\lineto(3838.29838538,1742.7520752)
\moveto(3838.29838538,1770.19500732)
\lineto(3849.87980627,1770.19500732)
\lineto(3849.87980627,1755.52940369)
\lineto(3838.29838538,1755.52940369)
\lineto(3838.29838538,1770.19500732)
}
}
{
\newrgbcolor{curcolor}{0 0 0}
\pscustom[linestyle=none,fillstyle=solid,fillcolor=curcolor]
{
\newpath
\moveto(3901.36677527,1734.63249207)
\curveto(3895.15641487,1734.63242969)(3890.24690439,1732.19865527)(3886.63822913,1727.3311615)
\curveto(3883.02950438,1722.50551921)(3881.22515438,1715.87558199)(3881.22517371,1707.44132996)
\curveto(3881.22515438,1699.00700755)(3883.00852357,1692.35608952)(3886.57528662,1687.48855591)
\curveto(3890.18396195,1682.66295346)(3895.11445324,1680.25015985)(3901.36677527,1680.25016785)
\curveto(3907.53509511,1680.25015985)(3912.42362477,1682.68393427)(3916.03237891,1687.55149841)
\curveto(3919.64102478,1692.41903196)(3921.44537478,1699.04896918)(3921.44543433,1707.44132996)
\curveto(3921.44537478,1715.79165874)(3919.64102478,1722.40061514)(3916.03237891,1727.26821899)
\curveto(3912.42362477,1732.17767445)(3907.53509511,1734.63242969)(3901.36677527,1734.63249207)
\moveto(3901.36677527,1744.45152283)
\curveto(3911.43752651,1744.45145063)(3919.34729338,1741.17844365)(3925.09609961,1734.63249207)
\curveto(3930.84477945,1728.08641573)(3933.71915096,1719.02270409)(3933.71922278,1707.44132996)
\curveto(3933.71915096,1695.90184708)(3930.84477945,1686.83813544)(3925.09609961,1680.25016785)
\curveto(3919.34729338,1673.70414589)(3911.43752651,1670.43113891)(3901.36677527,1670.43113708)
\curveto(3891.25398347,1670.43113891)(3883.32323578,1673.70414589)(3877.57450842,1680.25016785)
\curveto(3871.86771135,1686.83813544)(3869.01432065,1695.90184708)(3869.01432776,1707.44132996)
\curveto(3869.01432065,1719.02270409)(3871.86771135,1728.08641573)(3877.57450842,1734.63249207)
\curveto(3883.32323578,1741.17844365)(3891.25398347,1744.45145063)(3901.36677527,1744.45152283)
}
}
{
\newrgbcolor{curcolor}{0 0 0}
\pscustom[linestyle=none,fillstyle=solid,fillcolor=curcolor]
{
\newpath
\moveto(4011.45321631,1714.80560303)
\lineto(4011.45321631,1672.25646973)
\lineto(3999.87179541,1672.25646973)
\lineto(3999.87179541,1714.427948)
\curveto(3999.87173624,1721.09980467)(3998.57092578,1726.0932384)(3995.96936011,1729.40826416)
\curveto(3993.36768391,1732.72317562)(3989.46525251,1734.38065992)(3984.2620542,1734.38072205)
\curveto(3978.00972808,1734.38065992)(3973.07923679,1732.38748259)(3969.47056555,1728.40118408)
\curveto(3965.86183678,1724.41477328)(3964.05748678,1718.98074246)(3964.05751013,1712.09907532)
\lineto(3964.05751013,1672.25646973)
\lineto(3952.41314673,1672.25646973)
\lineto(3952.41314673,1742.7520752)
\lineto(3964.05751013,1742.7520752)
\lineto(3964.05751013,1731.80007935)
\curveto(3966.82695423,1736.03814423)(3970.07898039,1739.20624714)(3973.81359839,1741.30439758)
\curveto(3977.5901118,1743.40240993)(3981.93314029,1744.45145063)(3986.8426969,1744.45152283)
\curveto(3994.94124496,1744.45145063)(4001.06764264,1741.93375295)(4005.22190833,1736.89842224)
\curveto(4009.37604498,1731.90492387)(4011.45314556,1724.54065816)(4011.45321631,1714.80560303)
}
}
{
\newrgbcolor{curcolor}{0 0 0}
\pscustom[linestyle=none,fillstyle=solid,fillcolor=curcolor]
{
\newpath
\moveto(4062.49958777,1770.06912231)
\curveto(4056.87668965,1760.41785007)(4052.70150767,1750.87157971)(4049.9740293,1741.43028259)
\curveto(4047.24649604,1731.98884713)(4045.88274313,1722.42159595)(4045.88276648,1712.72850037)
\curveto(4045.88274313,1703.03532384)(4047.24649604,1693.42611103)(4049.9740293,1683.90083313)
\curveto(4052.7434693,1674.41749357)(4056.91865128,1664.87122321)(4062.49958777,1655.26199341)
\lineto(4052.42878699,1655.26199341)
\curveto(4046.1345129,1665.12299297)(4041.41382975,1674.81612903)(4038.26672339,1684.34143066)
\curveto(4035.16154718,1693.86670813)(4033.60896695,1703.32905523)(4033.60897803,1712.72850037)
\curveto(4033.60896695,1722.08590293)(4035.16154718,1731.50628841)(4038.26672339,1740.98968506)
\curveto(4041.37186812,1750.47294424)(4046.09255127,1760.1660803)(4052.42878699,1770.06912231)
\lineto(4062.49958777,1770.06912231)
}
}
{
\newrgbcolor{curcolor}{0 0 0}
\pscustom[linestyle=none,fillstyle=solid,fillcolor=curcolor]
{
\newpath
\moveto(4083.20767187,1770.06912231)
\lineto(4093.27847266,1770.06912231)
\curveto(4099.57269646,1760.1660803)(4104.27239879,1750.47294424)(4107.37759375,1740.98968506)
\curveto(4110.52468135,1731.50628841)(4112.0982424,1722.08590293)(4112.09828162,1712.72850037)
\curveto(4112.0982424,1703.32905523)(4110.52468135,1693.86670813)(4107.37759375,1684.34143066)
\curveto(4104.27239879,1674.81612903)(4099.57269646,1665.12299297)(4093.27847266,1655.26199341)
\lineto(4083.20767187,1655.26199341)
\curveto(4088.78855807,1664.87122321)(4092.94275924,1674.41749357)(4095.67028784,1683.90083313)
\curveto(4098.4397325,1693.42611103)(4099.82446622,1703.03532384)(4099.82449316,1712.72850037)
\curveto(4099.82446622,1722.42159595)(4098.4397325,1731.98884713)(4095.67028784,1741.43028259)
\curveto(4092.94275924,1750.87157971)(4088.78855807,1760.41785007)(4083.20767187,1770.06912231)
}
}
{
\newrgbcolor{curcolor}{0 0 0}
\pscustom[linestyle=none,fillstyle=solid,fillcolor=curcolor]
{
}
}
{
\newrgbcolor{curcolor}{0 0 0}
\pscustom[linestyle=none,fillstyle=solid,fillcolor=curcolor]
{
\newpath
\moveto(4179.38381934,1688.24386597)
\lineto(4192.66468787,1688.24386597)
\lineto(4192.66468787,1672.25646973)
\lineto(4179.38381934,1672.25646973)
\lineto(4179.38381934,1688.24386597)
\moveto(4179.38381934,1738.9125824)
\lineto(4192.66468787,1738.9125824)
\lineto(4192.66468787,1722.92518616)
\lineto(4179.38381934,1722.92518616)
\lineto(4179.38381934,1738.9125824)
}
}
{
\newrgbcolor{curcolor}{0 0 0}
\pscustom[linestyle=none,fillstyle=solid,fillcolor=curcolor]
{
}
}
{
\newrgbcolor{curcolor}{0 0 0}
\pscustom[linestyle=none,fillstyle=solid,fillcolor=curcolor]
{
\newpath
\moveto(4331.89350867,1758.99124146)
\lineto(4331.89350867,1745.58448792)
\curveto(4327.61333959,1749.57076924)(4323.03952215,1752.55004483)(4318.1720426,1754.52232361)
\curveto(4313.34638609,1756.49443786)(4308.20608666,1757.48053611)(4302.75112891,1757.48062134)
\curveto(4292.00889827,1757.48053611)(4283.78441919,1754.18654832)(4278.07766699,1747.59864807)
\curveto(4272.37085638,1741.05255877)(4269.51746568,1731.56923085)(4269.51748633,1719.14863586)
\curveto(4269.51746568,1706.76990872)(4272.37085638,1697.2865808)(4278.07766699,1690.69862366)
\curveto(4283.78441919,1684.15259125)(4292.00889827,1680.87958427)(4302.75112891,1680.8795929)
\curveto(4308.20608666,1680.87958427)(4313.34638609,1681.86568253)(4318.1720426,1683.83789062)
\curveto(4323.03952215,1685.81007556)(4327.61333959,1688.78935114)(4331.89350867,1692.77572632)
\lineto(4331.89350867,1679.49485779)
\curveto(4327.44549308,1676.47361334)(4322.72480994,1674.20768543)(4317.73144507,1672.69706726)
\curveto(4312.77990411,1671.18644821)(4307.53470061,1670.43113891)(4301.99581885,1670.43113708)
\curveto(4287.77077385,1670.43113891)(4276.56701918,1674.7741674)(4268.38452124,1683.4602356)
\curveto(4260.20198428,1692.18824301)(4256.11072555,1704.08436453)(4256.11073279,1719.14863586)
\curveto(4256.11072555,1734.25477504)(4260.20198428,1746.15089656)(4268.38452124,1754.83703613)
\curveto(4276.56701918,1763.56497217)(4287.77077385,1767.92898148)(4301.99581885,1767.92907715)
\curveto(4307.61862387,1767.92898148)(4312.90578899,1767.17367217)(4317.85733008,1765.66314697)
\curveto(4322.85069482,1764.19439659)(4327.52941634,1761.97043031)(4331.89350867,1758.99124146)
}
}
{
\newrgbcolor{curcolor}{0 0 0}
\pscustom[linestyle=none,fillstyle=solid,fillcolor=curcolor]
{
\newpath
\moveto(4378.47096228,1734.63249207)
\curveto(4372.26060188,1734.63242969)(4367.35109141,1732.19865527)(4363.74241614,1727.3311615)
\curveto(4360.1336914,1722.50551921)(4358.32934139,1715.87558199)(4358.32936072,1707.44132996)
\curveto(4358.32934139,1699.00700755)(4360.11271058,1692.35608952)(4363.67947363,1687.48855591)
\curveto(4367.28814896,1682.66295346)(4372.21864025,1680.25015985)(4378.47096228,1680.25016785)
\curveto(4384.63928213,1680.25015985)(4389.52781178,1682.68393427)(4393.13656592,1687.55149841)
\curveto(4396.74521179,1692.41903196)(4398.54956179,1699.04896918)(4398.54962134,1707.44132996)
\curveto(4398.54956179,1715.79165874)(4396.74521179,1722.40061514)(4393.13656592,1727.26821899)
\curveto(4389.52781178,1732.17767445)(4384.63928213,1734.63242969)(4378.47096228,1734.63249207)
\moveto(4378.47096228,1744.45152283)
\curveto(4388.54171353,1744.45145063)(4396.4514804,1741.17844365)(4402.20028662,1734.63249207)
\curveto(4407.94896646,1728.08641573)(4410.82333797,1719.02270409)(4410.82340979,1707.44132996)
\curveto(4410.82333797,1695.90184708)(4407.94896646,1686.83813544)(4402.20028662,1680.25016785)
\curveto(4396.4514804,1673.70414589)(4388.54171353,1670.43113891)(4378.47096228,1670.43113708)
\curveto(4368.35817048,1670.43113891)(4360.42742279,1673.70414589)(4354.67869543,1680.25016785)
\curveto(4348.97189836,1686.83813544)(4346.11850766,1695.90184708)(4346.11851477,1707.44132996)
\curveto(4346.11850766,1719.02270409)(4348.97189836,1728.08641573)(4354.67869543,1734.63249207)
\curveto(4360.42742279,1741.17844365)(4368.35817048,1744.45145063)(4378.47096228,1744.45152283)
}
}
{
\newrgbcolor{curcolor}{0 0 0}
\pscustom[linestyle=none,fillstyle=solid,fillcolor=curcolor]
{
\newpath
\moveto(4488.55740332,1714.80560303)
\lineto(4488.55740332,1672.25646973)
\lineto(4476.97598242,1672.25646973)
\lineto(4476.97598242,1714.427948)
\curveto(4476.97592326,1721.09980467)(4475.67511279,1726.0932384)(4473.07354712,1729.40826416)
\curveto(4470.47187092,1732.72317562)(4466.56943952,1734.38065992)(4461.36624121,1734.38072205)
\curveto(4455.11391509,1734.38065992)(4450.1834238,1732.38748259)(4446.57475256,1728.40118408)
\curveto(4442.96602379,1724.41477328)(4441.16167379,1718.98074246)(4441.16169714,1712.09907532)
\lineto(4441.16169714,1672.25646973)
\lineto(4429.51733374,1672.25646973)
\lineto(4429.51733374,1742.7520752)
\lineto(4441.16169714,1742.7520752)
\lineto(4441.16169714,1731.80007935)
\curveto(4443.93114124,1736.03814423)(4447.1831674,1739.20624714)(4450.9177854,1741.30439758)
\curveto(4454.69429881,1743.40240993)(4459.0373273,1744.45145063)(4463.94688391,1744.45152283)
\curveto(4472.04543197,1744.45145063)(4478.17182965,1741.93375295)(4482.32609534,1736.89842224)
\curveto(4486.48023199,1731.90492387)(4488.55733257,1724.54065816)(4488.55740332,1714.80560303)
}
}
{
\newrgbcolor{curcolor}{0 0 0}
\pscustom[linestyle=none,fillstyle=solid,fillcolor=curcolor]
{
\newpath
\moveto(4562.51484656,1740.04554749)
\lineto(4562.51484656,1729.21943665)
\curveto(4559.2417767,1731.02372968)(4555.9477889,1732.36650178)(4552.63287329,1733.24775696)
\curveto(4549.35981331,1734.17085178)(4546.0448447,1734.63242969)(4542.68795752,1734.63249207)
\curveto(4535.17678306,1734.63242969)(4529.34411678,1732.2406169)(4525.18994116,1727.45704651)
\curveto(4521.03571444,1722.71532735)(4518.95861386,1716.0434285)(4518.95863318,1707.44132996)
\curveto(4518.95861386,1698.83916104)(4521.03571444,1692.14628138)(4525.18994116,1687.3626709)
\curveto(4529.34411678,1682.62099183)(4535.17678306,1680.25015985)(4542.68795752,1680.25016785)
\curveto(4546.0448447,1680.25015985)(4549.35981331,1680.69075695)(4552.63287329,1681.57196045)
\curveto(4555.9477889,1682.49510695)(4559.2417767,1683.85885986)(4562.51484656,1685.66322327)
\lineto(4562.51484656,1674.96299744)
\curveto(4559.28373833,1673.45237612)(4555.92680809,1672.31941217)(4552.44404578,1671.56410217)
\curveto(4549.00313947,1670.80879356)(4545.33149703,1670.43113891)(4541.42910742,1670.43113708)
\curveto(4530.81277375,1670.43113891)(4522.37848653,1673.76708833)(4516.12622046,1680.43899536)
\curveto(4509.8739214,1687.11088602)(4506.74778012,1696.11165522)(4506.74778723,1707.44132996)
\curveto(4506.74778012,1718.93878083)(4509.89490222,1727.98151166)(4516.18916296,1734.56954956)
\curveto(4522.52535223,1741.15746284)(4531.19042841,1744.45145063)(4542.18441748,1744.45152283)
\curveto(4545.75111331,1744.45145063)(4549.23392843,1744.07379598)(4552.63287329,1743.31855774)
\curveto(4556.03171216,1742.605139)(4559.32569995,1741.51413667)(4562.51484656,1740.04554749)
}
}
{
\newrgbcolor{curcolor}{0 0 0}
\pscustom[linestyle=none,fillstyle=solid,fillcolor=curcolor]
{
\newpath
\moveto(4582.53056311,1770.19500732)
\lineto(4594.11198401,1770.19500732)
\lineto(4594.11198401,1672.25646973)
\lineto(4582.53056311,1672.25646973)
\lineto(4582.53056311,1770.19500732)
}
}
{
\newrgbcolor{curcolor}{0 0 0}
\pscustom[linestyle=none,fillstyle=solid,fillcolor=curcolor]
{
\newpath
\moveto(4617.08599829,1700.07705688)
\lineto(4617.08599829,1742.7520752)
\lineto(4628.66741919,1742.7520752)
\lineto(4628.66741919,1700.51765442)
\curveto(4628.66739666,1693.84572731)(4629.96820712,1688.83131277)(4632.56985449,1685.47439575)
\curveto(4635.17144899,1682.15941393)(4639.07388039,1680.50192962)(4644.2771604,1680.50193787)
\curveto(4650.52940482,1680.50192962)(4655.45989611,1682.49510695)(4659.06864905,1686.48147583)
\curveto(4662.71925775,1690.46781626)(4664.54458856,1695.90184708)(4664.54464697,1702.78358459)
\lineto(4664.54464697,1742.7520752)
\lineto(4676.12606787,1742.7520752)
\lineto(4676.12606787,1672.25646973)
\lineto(4664.54464697,1672.25646973)
\lineto(4664.54464697,1683.08258057)
\curveto(4661.73315949,1678.80248369)(4658.46015251,1675.61339996)(4654.72561621,1673.51531982)
\curveto(4651.03294436,1671.4591988)(4646.73187749,1670.43113891)(4641.82240271,1670.43113708)
\curveto(4633.72377283,1670.43113891)(4627.57639433,1672.94883659)(4623.38024878,1677.98423767)
\curveto(4619.18406874,1683.0196273)(4617.08598734,1690.38389301)(4617.08599829,1700.07705688)
\moveto(4646.22837805,1744.45152283)
\lineto(4646.22837805,1744.45152283)
}
}
{
\newrgbcolor{curcolor}{0 0 0}
\pscustom[linestyle=none,fillstyle=solid,fillcolor=curcolor]
{
\newpath
\moveto(4745.04811072,1740.67497253)
\lineto(4745.04811072,1729.72297668)
\curveto(4741.77504665,1731.40138434)(4738.37615478,1732.66023318)(4734.85142493,1733.49952698)
\curveto(4731.32660129,1734.33869829)(4727.67593965,1734.75831457)(4723.89942908,1734.75837708)
\curveto(4718.15065011,1734.75831457)(4713.82860243,1733.87712039)(4710.93327307,1732.11479187)
\curveto(4708.0798594,1730.35234364)(4706.65316405,1727.70876108)(4706.65318274,1724.18403625)
\curveto(4706.65316405,1721.49844014)(4707.68122393,1719.37937793)(4709.73736548,1717.82684326)
\curveto(4711.79346347,1716.31617908)(4715.92668382,1714.86850292)(4722.13703894,1713.48381042)
\lineto(4726.10241675,1712.60261536)
\curveto(4734.32685769,1710.84018664)(4740.15952397,1708.34346977)(4743.60043311,1705.11245728)
\curveto(4747.08319258,1701.92334069)(4748.82460015,1697.45442732)(4748.82466101,1691.70570374)
\curveto(4748.82460015,1685.15967032)(4746.22297921,1679.97740927)(4741.01979041,1676.15890503)
\curveto(4735.85845711,1672.34039298)(4728.74596117,1670.43113891)(4719.68228125,1670.43113708)
\curveto(4715.90570301,1670.43113891)(4711.96130998,1670.80879356)(4707.84909033,1671.56410217)
\curveto(4703.77879253,1672.27745054)(4699.47772566,1673.36845287)(4694.94587683,1674.83711243)
\lineto(4694.94587683,1686.79618835)
\curveto(4699.2259559,1684.57220753)(4703.44309951,1682.89374241)(4707.59732031,1681.76078796)
\curveto(4711.75150184,1680.66977613)(4715.86374138,1680.12427497)(4719.93405127,1680.12428284)
\curveto(4725.38903093,1680.12427497)(4729.58519373,1681.04743078)(4732.52255225,1682.89375305)
\curveto(4735.45982164,1684.78201567)(4736.92847862,1687.42559823)(4736.92852759,1690.82450867)
\curveto(4736.92847862,1693.9716122)(4735.85845711,1696.3844058)(4733.71845984,1698.06289673)
\curveto(4731.62033268,1699.74133604)(4726.98357279,1701.35685872)(4719.80816626,1702.9094696)
\lineto(4715.77984595,1703.85360718)
\curveto(4708.60437975,1705.36419419)(4703.42211869,1707.67208372)(4700.23304724,1710.77728271)
\curveto(4697.04395124,1713.92436629)(4695.44940938,1718.22543316)(4695.44941687,1723.68049622)
\curveto(4695.44940938,1730.31038201)(4697.79926055,1735.42970062)(4702.49897742,1739.03846741)
\curveto(4707.19866521,1742.64710063)(4713.87056405,1744.45145063)(4722.51469397,1744.45152283)
\curveto(4726.79474547,1744.45145063)(4730.82306175,1744.13673842)(4734.59965491,1743.50738525)
\curveto(4738.37615478,1742.87788958)(4741.8589699,1741.93375295)(4745.04811072,1740.67497253)
}
}
{
\newrgbcolor{curcolor}{0 0 0}
\pscustom[linestyle=none,fillstyle=solid,fillcolor=curcolor]
{
\newpath
\moveto(4767.32975745,1742.7520752)
\lineto(4778.91117834,1742.7520752)
\lineto(4778.91117834,1672.25646973)
\lineto(4767.32975745,1672.25646973)
\lineto(4767.32975745,1742.7520752)
\moveto(4767.32975745,1770.19500732)
\lineto(4778.91117834,1770.19500732)
\lineto(4778.91117834,1755.52940369)
\lineto(4767.32975745,1755.52940369)
\lineto(4767.32975745,1770.19500732)
}
}
{
\newrgbcolor{curcolor}{0 0 0}
\pscustom[linestyle=none,fillstyle=solid,fillcolor=curcolor]
{
\newpath
\moveto(4830.39814734,1734.63249207)
\curveto(4824.18778694,1734.63242969)(4819.27827646,1732.19865527)(4815.6696012,1727.3311615)
\curveto(4812.06087646,1722.50551921)(4810.25652645,1715.87558199)(4810.25654578,1707.44132996)
\curveto(4810.25652645,1699.00700755)(4812.03989564,1692.35608952)(4815.60665869,1687.48855591)
\curveto(4819.21533402,1682.66295346)(4824.14582531,1680.25015985)(4830.39814734,1680.25016785)
\curveto(4836.56646718,1680.25015985)(4841.45499684,1682.68393427)(4845.06375098,1687.55149841)
\curveto(4848.67239685,1692.41903196)(4850.47674685,1699.04896918)(4850.4768064,1707.44132996)
\curveto(4850.47674685,1715.79165874)(4848.67239685,1722.40061514)(4845.06375098,1727.26821899)
\curveto(4841.45499684,1732.17767445)(4836.56646718,1734.63242969)(4830.39814734,1734.63249207)
\moveto(4830.39814734,1744.45152283)
\curveto(4840.46889858,1744.45145063)(4848.37866545,1741.17844365)(4854.12747168,1734.63249207)
\curveto(4859.87615152,1728.08641573)(4862.75052303,1719.02270409)(4862.75059485,1707.44132996)
\curveto(4862.75052303,1695.90184708)(4859.87615152,1686.83813544)(4854.12747168,1680.25016785)
\curveto(4848.37866545,1673.70414589)(4840.46889858,1670.43113891)(4830.39814734,1670.43113708)
\curveto(4820.28535554,1670.43113891)(4812.35460785,1673.70414589)(4806.60588049,1680.25016785)
\curveto(4800.89908342,1686.83813544)(4798.04569272,1695.90184708)(4798.04569983,1707.44132996)
\curveto(4798.04569272,1719.02270409)(4800.89908342,1728.08641573)(4806.60588049,1734.63249207)
\curveto(4812.35460785,1741.17844365)(4820.28535554,1744.45145063)(4830.39814734,1744.45152283)
}
}
{
\newrgbcolor{curcolor}{0 0 0}
\pscustom[linestyle=none,fillstyle=solid,fillcolor=curcolor]
{
\newpath
\moveto(4940.48458838,1714.80560303)
\lineto(4940.48458838,1672.25646973)
\lineto(4928.90316748,1672.25646973)
\lineto(4928.90316748,1714.427948)
\curveto(4928.90310831,1721.09980467)(4927.60229785,1726.0932384)(4925.00073218,1729.40826416)
\curveto(4922.39905598,1732.72317562)(4918.49662458,1734.38065992)(4913.29342627,1734.38072205)
\curveto(4907.04110015,1734.38065992)(4902.11060886,1732.38748259)(4898.50193762,1728.40118408)
\curveto(4894.89320885,1724.41477328)(4893.08885885,1718.98074246)(4893.0888822,1712.09907532)
\lineto(4893.0888822,1672.25646973)
\lineto(4881.4445188,1672.25646973)
\lineto(4881.4445188,1742.7520752)
\lineto(4893.0888822,1742.7520752)
\lineto(4893.0888822,1731.80007935)
\curveto(4895.8583263,1736.03814423)(4899.11035246,1739.20624714)(4902.84497046,1741.30439758)
\curveto(4906.62148387,1743.40240993)(4910.96451236,1744.45145063)(4915.87406897,1744.45152283)
\curveto(4923.97261703,1744.45145063)(4930.09901471,1741.93375295)(4934.2532804,1736.89842224)
\curveto(4938.40741705,1731.90492387)(4940.48451763,1724.54065816)(4940.48458838,1714.80560303)
}
}
{
\newrgbcolor{curcolor}{0.99607843 0.99607843 0.80784315}
\pscustom[linestyle=none,fillstyle=solid,fillcolor=curcolor]
{
\newpath
\moveto(1538.04165625,477.5390625)
\lineto(2733.35415625,477.5390625)
\lineto(2733.35415625,8.7890625)
\lineto(1538.04165625,8.7890625)
\lineto(1538.04165625,477.5390625)
\closepath
}
}
{
\newrgbcolor{curcolor}{0.65882355 0 0.21176471}
\pscustom[linewidth=17.578125,linecolor=curcolor]
{
\newpath
\moveto(1538.04165625,477.5390625)
\lineto(2733.35415625,477.5390625)
\lineto(2733.35415625,8.7890625)
\lineto(1538.04165625,8.7890625)
\lineto(1538.04165625,477.5390625)
\closepath
}
}
{
\newrgbcolor{curcolor}{0.70588237 0.65490198 0.89803922}
\pscustom[linestyle=none,fillstyle=solid,fillcolor=curcolor]
{
\newpath
\moveto(1842.72909375,290.039125)
\curveto(1842.72909375,218.84616897)(1785.01579978,161.132875)(1713.82284375,161.132875)
\curveto(1642.62988772,161.132875)(1584.91659375,218.84616897)(1584.91659375,290.039125)
\curveto(1584.91659375,361.23208103)(1642.62988772,418.945375)(1713.82284375,418.945375)
\curveto(1785.01579978,418.945375)(1842.72909375,361.23208103)(1842.72909375,290.039125)
\closepath
}
}
{
\newrgbcolor{curcolor}{0.65882355 0 0.21176471}
\pscustom[linewidth=11.71875,linecolor=curcolor]
{
\newpath
\moveto(1842.72909375,290.039125)
\curveto(1842.72909375,218.84616897)(1785.01579978,161.132875)(1713.82284375,161.132875)
\curveto(1642.62988772,161.132875)(1584.91659375,218.84616897)(1584.91659375,290.039125)
\curveto(1584.91659375,361.23208103)(1642.62988772,418.945375)(1713.82284375,418.945375)
\curveto(1785.01579978,418.945375)(1842.72909375,361.23208103)(1842.72909375,290.039125)
\closepath
}
}
{
\newrgbcolor{curcolor}{0 0 0}
\pscustom[linestyle=none,fillstyle=solid,fillcolor=curcolor]
{
\newpath
\moveto(1667.13096875,219.726625)
\lineto(1667.13096875,238.95275)
\lineto(1700.82246875,238.95275)
\lineto(1700.82246875,344.6045)
\lineto(1667.13096875,344.6045)
\lineto(1667.13096875,363.6475)
\lineto(1763.26146875,363.6475)
\lineto(1763.26146875,344.6045)
\lineto(1729.56996875,344.6045)
\lineto(1729.56996875,238.95275)
\lineto(1763.26146875,238.95275)
\lineto(1763.26146875,219.726625)
\lineto(1667.13096875,219.726625)
\closepath
}
}
{
\newrgbcolor{curcolor}{0 0 0}
\pscustom[linestyle=none,fillstyle=solid,fillcolor=curcolor]
{
\newpath
\moveto(1960.14553809,317.50518799)
\curveto(1959.41303468,313.6599121)(1958.01685688,309.95203007)(1955.95700049,306.38153076)
\curveto(1953.89698795,302.85670025)(1951.21907315,299.72102223)(1947.92324805,296.9744873)
\curveto(1944.62728287,294.22786366)(1940.73629555,292.03060023)(1936.25027441,290.38269043)
\curveto(1931.76413655,288.73470509)(1926.70585303,287.91073131)(1921.07540869,287.9107666)
\lineto(1901.78066992,287.9107666)
\lineto(1894.91421484,252.6171875)
\lineto(1881.79928564,252.6171875)
\lineto(1899.78939795,345.2456665)
\lineto(1931.37509131,345.2456665)
\curveto(1936.36465685,345.24557388)(1940.69051923,344.74203434)(1944.35269141,343.73504639)
\curveto(1948.06050697,342.7278762)(1951.12752051,341.28592207)(1953.55374121,339.40917969)
\curveto(1955.97981058,337.53226372)(1957.78797527,335.26633581)(1958.97824072,332.61138916)
\curveto(1960.16834399,329.95628252)(1960.76343616,326.98082163)(1960.76351904,323.68499756)
\curveto(1960.76343616,321.71654467)(1960.55744272,319.65661021)(1960.14553809,317.50518799)
\moveto(1946.89327979,317.29919434)
\curveto(1947.16786871,318.76397194)(1947.30519767,320.13726158)(1947.30526709,321.41906738)
\curveto(1947.30519767,325.81352544)(1945.68013826,329.13230874)(1942.43008398,331.37542725)
\curveto(1939.17990062,333.61838824)(1934.28183423,334.73990811)(1927.73587012,334.73999023)
\lineto(1910.84439062,334.73999023)
\lineto(1903.77194189,298.27911377)
\lineto(1921.21273779,298.27911377)
\curveto(1924.96635282,298.27906811)(1928.30802428,298.75971948)(1931.23776221,299.72106934)
\curveto(1934.21316975,300.7281013)(1936.77664375,302.07850278)(1938.92819189,303.77227783)
\curveto(1941.07961796,305.46595057)(1942.81911817,307.48010871)(1944.14669775,309.8147583)
\curveto(1945.47414481,312.1492935)(1946.38967124,314.64410301)(1946.89327979,317.29919434)
}
}
{
\newrgbcolor{curcolor}{0 0 0}
\pscustom[linestyle=none,fillstyle=solid,fillcolor=curcolor]
{
\newpath
\moveto(2038.97244238,315.65124512)
\curveto(2036.91243115,316.01739265)(2034.71516772,316.31493874)(2032.38064551,316.54388428)
\curveto(2030.09175926,316.81847828)(2027.52828526,316.95580724)(2024.69021582,316.95587158)
\curveto(2020.84494234,316.95580724)(2017.22861295,316.17760978)(2013.8412168,314.62127686)
\curveto(2010.4994937,313.06481992)(2007.50114465,310.93622098)(2004.84616064,308.23547363)
\curveto(2002.23686768,305.58039137)(2000.01671609,302.46760152)(1998.18569922,298.89709473)
\curveto(1996.4003867,295.32649538)(1995.1186497,291.5270607)(1994.34048437,287.4987793)
\lineto(1987.54269385,252.6171875)
\lineto(1975.18307471,252.6171875)
\lineto(1984.52145361,300.7510376)
\curveto(1985.1622998,304.22998989)(1985.59717485,307.70899032)(1985.82608008,311.18804932)
\curveto(1986.10071439,314.71276749)(1986.23804335,317.84844551)(1986.23806738,320.59509277)
\curveto(1986.23804335,321.92253812)(1986.21515519,323.11272247)(1986.16940283,324.16564941)
\curveto(1986.12360255,325.26420958)(1986.07782623,326.17973601)(1986.03207373,326.91223145)
\lineto(1997.77371191,326.91223145)
\curveto(1997.91100531,325.76774911)(1998.00255795,324.41734763)(1998.04837012,322.86102295)
\curveto(1998.13988692,321.3503341)(1998.18566324,319.77105101)(1998.18569922,318.12316895)
\curveto(1998.18566324,316.65826115)(1998.13988692,315.21630703)(1998.04837012,313.79730225)
\curveto(1998.00255795,312.3781751)(1997.91100531,311.02777362)(1997.77371191,309.74609375)
\lineto(1998.11703467,309.74609375)
\curveto(1999.94805161,312.58416855)(2001.77910447,315.14764255)(2003.61019873,317.43652344)
\curveto(2005.4869865,319.72527469)(2007.54692097,321.67076835)(2009.7900083,323.27301025)
\curveto(2012.07877679,324.87511085)(2014.61936263,326.11107153)(2017.41177344,326.980896)
\curveto(2020.24985016,327.85057174)(2023.5457453,328.28544679)(2027.29946875,328.28552246)
\curveto(2029.86287766,328.28544679)(2032.35768717,328.17100599)(2034.78390479,327.94219971)
\curveto(2037.25575356,327.71324277)(2039.40724067,327.43858485)(2041.23837256,327.1182251)
\lineto(2038.97244238,315.65124512)
}
}
{
\newrgbcolor{curcolor}{0 0 0}
\pscustom[linestyle=none,fillstyle=solid,fillcolor=curcolor]
{
\newpath
\moveto(2068.15487646,287.15545654)
\curveto(2067.88019691,285.69057972)(2067.6513153,284.27151376)(2067.46823096,282.89825439)
\curveto(2067.33088105,281.52493447)(2067.26221657,280.17453299)(2067.2622373,278.8470459)
\curveto(2067.26221657,276.10044038)(2067.58265082,273.58274271)(2068.22354102,271.29394531)
\curveto(2068.91016414,269.05088689)(2069.94013137,267.12828139)(2071.3134458,265.52612305)
\curveto(2072.68671066,263.92393889)(2074.38043455,262.68797821)(2076.39462256,261.8182373)
\curveto(2078.45452715,260.948478)(2080.85778403,260.51360294)(2083.60440039,260.51361084)
\curveto(2086.12206099,260.51360294)(2088.5024297,260.78826087)(2090.74551367,261.33758545)
\curveto(2092.9885092,261.93266891)(2095.04844366,262.73375453)(2096.92532324,263.74084473)
\curveto(2098.84787834,264.74791267)(2100.51871407,265.93809703)(2101.93783545,267.31140137)
\curveto(2103.356846,268.73045264)(2104.50125404,270.2639594)(2105.37106299,271.91192627)
\lineto(2115.60208105,268.82202148)
\curveto(2114.45760394,266.76207082)(2112.94698534,264.67924819)(2111.0702207,262.57354736)
\curveto(2109.19332698,260.51360294)(2106.88162275,258.63677377)(2104.13510107,256.9430542)
\curveto(2101.38846418,255.2951023)(2098.16123352,253.92181266)(2094.45339941,252.82318115)
\curveto(2090.79124577,251.77032555)(2086.5798242,251.24389786)(2081.81912207,251.24389648)
\curveto(2077.47033624,251.24389786)(2073.60223708,251.79321371)(2070.21481299,252.8918457)
\curveto(2066.82734151,254.03625346)(2063.96632143,255.72997736)(2061.63174414,257.97302246)
\curveto(2059.29713664,260.21605685)(2057.51186011,263.00841246)(2056.27590918,266.35009766)
\curveto(2055.08571507,269.69175539)(2054.49062289,273.58274271)(2054.49063086,278.02307129)
\curveto(2054.49062289,281.82248056)(2054.92549795,285.89657317)(2055.79525732,290.24536133)
\curveto(2057.1685377,297.38642984)(2059.22847216,303.38312794)(2061.97506689,308.23547363)
\curveto(2064.72163073,313.08770808)(2067.90308507,317.00158356)(2071.51943945,319.97711182)
\curveto(2075.18152017,322.95250535)(2079.14117197,325.08110429)(2083.39840674,326.36291504)
\curveto(2087.65556776,327.64457829)(2091.98143013,328.28544679)(2096.37600684,328.28552246)
\curveto(2101.13669441,328.28544679)(2105.18789886,327.66746645)(2108.52963232,326.43157959)
\curveto(2111.87124178,325.24132142)(2114.59493291,323.50182121)(2116.70071387,321.21307373)
\curveto(2118.8521308,318.96996539)(2120.40852573,316.24627426)(2121.36990332,313.04199219)
\curveto(2122.37690755,309.83758926)(2122.88044708,306.24414803)(2122.88052344,302.26165771)
\curveto(2122.88044708,298.09596282)(2122.39979571,293.60988332)(2121.43856787,288.80340576)
\lineto(2121.16390967,287.15545654)
\lineto(2068.15487646,287.15545654)
\moveto(2094.72805762,319.15313721)
\curveto(2092.62229863,319.15307067)(2090.35637072,318.80974826)(2087.93026709,318.12316895)
\curveto(2085.50408065,317.48223494)(2083.1466001,316.31493874)(2080.85781836,314.62127686)
\curveto(2078.56896796,312.92749096)(2076.46325717,310.61578673)(2074.54067969,307.68615723)
\curveto(2072.66382249,304.80219391)(2071.19898021,301.11720004)(2070.14614844,296.63116455)
\lineto(2110.24624609,296.63116455)
\curveto(2110.33773502,297.54664697)(2110.4063995,298.39350891)(2110.45223975,299.17175293)
\curveto(2110.49795214,299.99568016)(2110.5208403,300.79676579)(2110.5209043,301.57501221)
\curveto(2110.5208403,304.91663471)(2110.10885341,307.70899032)(2109.28494238,309.9520874)
\curveto(2108.46090584,312.19506982)(2107.33938596,314.00323451)(2105.92037939,315.37658691)
\curveto(2104.50125404,316.7498138)(2102.8304183,317.71111655)(2100.90786719,318.26049805)
\curveto(2098.98520731,318.85552458)(2096.92527284,319.15307067)(2094.72805762,319.15313721)
}
}
{
\newrgbcolor{curcolor}{0 0 0}
\pscustom[linestyle=none,fillstyle=solid,fillcolor=curcolor]
{
\newpath
\moveto(2185.159271,264.56481934)
\curveto(2182.00065051,259.75829364)(2178.33854479,256.25640505)(2174.17294287,254.05914307)
\curveto(2170.05303062,251.90765452)(2165.22362871,250.83191096)(2159.68472266,250.83190918)
\curveto(2152.45203503,250.83191096)(2147.11909359,252.7316283)(2143.68588232,256.53106689)
\curveto(2140.25264538,260.37627398)(2138.53603333,266.1898668)(2138.53604102,273.97186279)
\curveto(2138.53603333,278.41214462)(2139.08534918,283.47042813)(2140.18399023,289.14672852)
\curveto(2145.26515257,315.0560899)(2156.75500925,328.01078886)(2174.65359473,328.01086426)
\curveto(2180.19248581,328.01078886)(2184.63278899,327.02659795)(2187.97451758,325.05828857)
\curveto(2191.36190823,323.08983431)(2193.74227695,319.86260365)(2195.11563086,315.37658691)
\lineto(2195.25295996,315.37658691)
\curveto(2195.29867187,315.60540576)(2195.39022452,316.20049794)(2195.52761816,317.16186523)
\curveto(2195.71065877,318.12310344)(2195.89376405,319.17595883)(2196.07693457,320.32043457)
\curveto(2196.30575095,321.4647749)(2196.51174439,322.54051846)(2196.69491553,323.54766846)
\curveto(2196.87795496,324.60045292)(2197.01528393,325.33287406)(2197.10690283,325.74493408)
\lineto(2202.66873145,354.51538086)
\lineto(2215.02835059,354.51538086)
\lineto(2198.20553564,267.92938232)
\curveto(2197.88503403,266.14409048)(2197.56459978,264.42747842)(2197.24423193,262.77954102)
\curveto(2196.96950761,261.13158328)(2196.71773784,259.62096468)(2196.48892187,258.24768066)
\curveto(2196.25997462,256.87438539)(2196.07686934,255.68420103)(2195.93960547,254.67712402)
\curveto(2195.80221141,253.71581921)(2195.71065877,253.02917439)(2195.66494727,252.6171875)
\lineto(2183.85464453,252.6171875)
\curveto(2183.8088152,252.80029279)(2183.83170336,253.30383232)(2183.92330908,254.12780762)
\curveto(2184.01480865,254.95177989)(2184.12924945,255.9359708)(2184.26663184,257.0803833)
\curveto(2184.40390738,258.22478687)(2184.58701266,259.46074755)(2184.81594824,260.78826904)
\curveto(2185.04477588,262.11577419)(2185.27365749,263.37462303)(2185.50259375,264.56481934)
\lineto(2185.159271,264.56481934)
\moveto(2153.09292578,289.55871582)
\curveto(2152.589364,286.81209959)(2152.20026527,284.29440192)(2151.92562842,282.00561523)
\curveto(2151.65094941,279.7625461)(2151.51362045,277.67972347)(2151.51364111,275.75714111)
\curveto(2151.51362045,270.40128837)(2152.56647584,266.48741289)(2154.67221045,264.01550293)
\curveto(2156.82367373,261.5893465)(2160.18823335,260.37627398)(2164.76589941,260.37628174)
\curveto(2167.92443167,260.37627398)(2170.83122808,260.92558984)(2173.48629736,262.02423096)
\curveto(2176.14128136,263.12285326)(2178.52165008,264.86235348)(2180.62741064,267.24273682)
\curveto(2182.77884797,269.66886723)(2184.65567715,272.80454524)(2186.25790381,276.64978027)
\curveto(2187.86001965,280.49496724)(2189.21042113,285.16415202)(2190.3091123,290.65734863)
\curveto(2190.76681606,292.99190298)(2191.11013847,295.14339009)(2191.33908057,297.11181641)
\curveto(2191.56790168,299.08015373)(2191.68234248,300.91120659)(2191.68240332,302.60498047)
\curveto(2191.68234248,305.58039137)(2191.33902007,308.09808905)(2190.65243506,310.15808105)
\curveto(2189.96573043,312.2637343)(2188.95865136,313.95745819)(2187.63119482,315.23925781)
\curveto(2186.34940104,316.52093219)(2184.74722979,317.43645862)(2182.82467627,317.98583984)
\curveto(2180.94779511,318.58086665)(2178.77341985,318.87841274)(2176.30154395,318.878479)
\curveto(2173.37181392,318.87841274)(2170.69389912,318.35198505)(2168.2677915,317.29919434)
\curveto(2165.88738537,316.24627426)(2163.73589826,314.55255037)(2161.81332373,312.21801758)
\curveto(2159.93646358,309.88336559)(2158.26562785,306.83924021)(2156.80081152,303.08563232)
\curveto(2155.3817196,299.37769982)(2154.14575893,294.86873216)(2153.09292578,289.55871582)
}
}
{
\newrgbcolor{curcolor}{0 0 0}
\pscustom[linestyle=none,fillstyle=solid,fillcolor=curcolor]
{
\newpath
\moveto(2261.03359961,262.36755371)
\lineto(2287.12612891,262.36755371)
\lineto(2285.20352148,252.6171875)
\lineto(2217.77493262,252.6171875)
\lineto(2219.69754004,262.36755371)
\lineto(2248.67398047,262.36755371)
\lineto(2259.31698584,317.16186523)
\lineto(2237.41299414,317.16186523)
\lineto(2239.26693701,326.91223145)
\lineto(2273.53054785,326.91223145)
\lineto(2261.03359961,262.36755371)
\moveto(2262.61288428,341.33178711)
\lineto(2265.15347266,354.51538086)
\lineto(2278.88638281,354.51538086)
\lineto(2276.34579443,341.33178711)
\lineto(2262.61288428,341.33178711)
}
}
{
\newrgbcolor{curcolor}{0 0 0}
\pscustom[linestyle=none,fillstyle=solid,fillcolor=curcolor]
{
\newpath
\moveto(2308.41213965,289.83337402)
\curveto(2309.87697301,297.2948772)(2312.07423643,303.49756875)(2315.00393652,308.44146729)
\curveto(2317.93360558,313.38525417)(2321.27527704,317.32201781)(2325.02896094,320.25177002)
\curveto(2328.82837007,323.18138696)(2332.90246268,325.24132142)(2337.25125098,326.43157959)
\curveto(2341.64574007,327.66746645)(2346.01737876,328.28544679)(2350.36618018,328.28552246)
\curveto(2354.85220879,328.28544679)(2358.74319611,327.71324277)(2362.03915381,326.56890869)
\curveto(2365.38076272,325.4244267)(2368.127342,323.84514362)(2370.2788999,321.83105469)
\curveto(2372.43031622,319.86260365)(2374.03248746,317.52801126)(2375.08541846,314.82727051)
\curveto(2376.13819825,312.12640533)(2376.66462595,309.21960892)(2376.66470312,306.10687256)
\lineto(2363.34378027,305.14556885)
\curveto(2363.43526906,309.26538525)(2362.2450847,312.53839223)(2359.77322363,314.96459961)
\curveto(2357.30124199,317.3906823)(2353.36447835,318.60375481)(2347.9629209,318.6038208)
\curveto(2344.02610878,318.60375481)(2340.59288467,317.96288631)(2337.66323828,316.68121338)
\curveto(2334.77929185,315.44518864)(2332.28448234,313.61413578)(2330.17880225,311.18804932)
\curveto(2328.07306077,308.80762203)(2326.31067239,305.83216114)(2324.89163184,302.26165771)
\curveto(2323.51831679,298.73683132)(2322.37390875,294.68562688)(2321.4584043,290.10803223)
\curveto(2320.58863222,285.7592442)(2320.15375716,281.77670424)(2320.15377783,278.16040039)
\curveto(2320.15375716,275.41379556)(2320.45130325,272.94187421)(2321.04641699,270.74462891)
\curveto(2321.64148761,268.59312367)(2322.60279036,266.78495898)(2323.93032812,265.32012939)
\curveto(2325.257817,263.85527441)(2326.95154089,262.73375453)(2329.01150488,261.95556641)
\curveto(2331.11718614,261.1773596)(2333.65777198,260.78826087)(2336.63327002,260.78826904)
\curveto(2341.57707558,260.78826087)(2345.90293796,262.00133339)(2349.61087012,264.42749023)
\curveto(2353.36447835,266.85362346)(2356.11105763,270.56150549)(2357.85061621,275.55114746)
\lineto(2370.75955176,274.72717285)
\curveto(2369.75240141,271.66013721)(2368.19600648,268.70756448)(2366.0903623,265.8694458)
\curveto(2363.98458492,263.03130062)(2361.42111092,260.5364911)(2358.39993262,258.38500977)
\curveto(2355.37863649,256.23351689)(2351.92252422,254.49401668)(2348.03158545,253.16650391)
\curveto(2344.14054958,251.88476636)(2339.90623985,251.24389786)(2335.32864355,251.24389648)
\curveto(2330.24743604,251.24389786)(2325.92157366,251.88476636)(2322.35104346,253.16650391)
\curveto(2318.78046752,254.49401668)(2315.87367111,256.32506953)(2313.63064551,258.65966797)
\curveto(2311.38759161,261.04003064)(2309.73964404,263.85527441)(2308.68679785,267.10540771)
\curveto(2307.67970958,270.40128837)(2307.17617004,274.01761776)(2307.17617773,277.95440674)
\curveto(2307.17617004,281.70803976)(2307.58815694,285.66769156)(2308.41213965,289.83337402)
}
}
{
\newrgbcolor{curcolor}{0 0 0}
\pscustom[linestyle=none,fillstyle=solid,fillcolor=curcolor]
{
\newpath
\moveto(2402.20791602,317.16186523)
\lineto(2404.06185889,326.91223145)
\lineto(2415.73483252,326.91223145)
\lineto(2423.49392676,346.27563477)
\lineto(2431.73367285,346.27563477)
\lineto(2427.95712256,326.91223145)
\lineto(2457.6202085,326.91223145)
\lineto(2455.76626562,317.16186523)
\lineto(2426.10317969,317.16186523)
\lineto(2417.38278174,272.39257812)
\curveto(2417.29119552,271.84324249)(2417.19964288,271.3168148)(2417.10812354,270.81329346)
\curveto(2417.06231391,270.30973573)(2417.03942575,269.82908435)(2417.03945898,269.37133789)
\curveto(2417.03942575,266.76207082)(2417.90917586,264.83946532)(2419.64871191,263.60351562)
\curveto(2421.38817628,262.36754396)(2424.11186741,261.74956362)(2427.81979346,261.74957275)
\curveto(2429.6508023,261.74956362)(2431.55051964,261.8182281)(2433.51895117,261.95556641)
\curveto(2435.48728328,262.13866235)(2437.38700062,262.3446558)(2439.21810889,262.57354736)
\curveto(2441.09488265,262.80241901)(2442.83438287,263.05418878)(2444.43661475,263.32885742)
\curveto(2446.08450169,263.60350464)(2447.48067949,263.85527441)(2448.62515234,264.08416748)
\lineto(2446.83987402,254.67712402)
\curveto(2445.7869556,254.40246404)(2444.41366595,254.08202979)(2442.72000098,253.71582031)
\curveto(2441.07199449,253.34960864)(2439.21805347,253.00628623)(2437.15817236,252.68585205)
\curveto(2435.09818455,252.36541773)(2432.87803296,252.0907598)(2430.49771094,251.86187744)
\curveto(2428.11729553,251.63299659)(2425.6911505,251.51855579)(2423.21926855,251.51855469)
\curveto(2416.85632047,251.51855579)(2412.09558304,252.64007566)(2408.93704199,254.88311768)
\curveto(2405.77845068,257.17193148)(2404.19916759,260.74248455)(2404.19918799,265.5947876)
\curveto(2404.19916759,266.41874841)(2404.22205576,267.28849851)(2404.26785254,268.20404053)
\curveto(2404.35938472,269.11955137)(2404.51960184,270.08085412)(2404.74850439,271.08795166)
\lineto(2413.74356055,317.16186523)
\lineto(2402.20791602,317.16186523)
}
}
{
\newrgbcolor{curcolor}{0 0 0}
\pscustom[linestyle=none,fillstyle=solid,fillcolor=curcolor]
{
\newpath
\moveto(2513.99380469,262.36755371)
\lineto(2540.08633398,262.36755371)
\lineto(2538.16372656,252.6171875)
\lineto(2470.7351377,252.6171875)
\lineto(2472.65774512,262.36755371)
\lineto(2501.63418555,262.36755371)
\lineto(2512.27719092,317.16186523)
\lineto(2490.37319922,317.16186523)
\lineto(2492.22714209,326.91223145)
\lineto(2526.49075293,326.91223145)
\lineto(2513.99380469,262.36755371)
\moveto(2515.57308936,341.33178711)
\lineto(2518.11367773,354.51538086)
\lineto(2531.84658789,354.51538086)
\lineto(2529.30599951,341.33178711)
\lineto(2515.57308936,341.33178711)
}
}
{
\newrgbcolor{curcolor}{0 0 0}
\pscustom[linestyle=none,fillstyle=solid,fillcolor=curcolor]
{
\newpath
\moveto(2627.77096533,289.83337402)
\curveto(2626.53492933,283.33309917)(2624.68098831,277.65683531)(2622.20913672,272.80456543)
\curveto(2619.7371456,267.99803149)(2616.73879655,263.99260337)(2613.21408057,260.78826904)
\curveto(2609.73501937,257.58391837)(2605.75247941,255.1806615)(2601.26644873,253.57849121)
\curveto(2596.82609673,252.02209532)(2591.99669482,251.24389786)(2586.77822852,251.24389648)
\curveto(2582.65832526,251.24389786)(2578.95044322,251.77032555)(2575.65457129,252.82318115)
\curveto(2572.35865294,253.92181266)(2569.54340917,255.56976023)(2567.20883154,257.76702881)
\curveto(2564.92000071,259.96428709)(2563.15761234,262.73375453)(2561.92166113,266.07543945)
\curveto(2560.7314673,269.41709746)(2560.13637512,273.33097294)(2560.13638281,277.81707764)
\curveto(2560.13637512,281.43338183)(2560.54836201,285.43880995)(2561.37234473,289.83337402)
\curveto(2563.89003348,302.60493048)(2568.6507709,312.19506982)(2575.65457129,318.6038208)
\curveto(2582.70410158,325.05821613)(2591.53893161,328.28544679)(2602.15908789,328.28552246)
\curveto(2611.26852614,328.28544679)(2618.04342171,326.17973601)(2622.48379492,321.96838379)
\curveto(2626.92402806,317.75689287)(2629.14417965,311.39398419)(2629.14425635,302.87963867)
\curveto(2629.14417965,300.91120659)(2629.02973885,298.85127213)(2628.80093359,296.6998291)
\curveto(2628.61775195,294.54829791)(2628.27442954,292.25948184)(2627.77096533,289.83337402)
\moveto(2614.79336523,289.83337402)
\curveto(2615.2510661,292.21370552)(2615.59438851,294.41096895)(2615.8233335,296.4251709)
\curveto(2616.05215173,298.43928523)(2616.16659253,300.31611441)(2616.16665625,302.05566406)
\curveto(2616.16659253,305.35150977)(2615.7774938,308.07520089)(2614.99935889,310.22674561)
\curveto(2614.26687519,312.42395142)(2613.2140198,314.1863398)(2611.84078955,315.51391602)
\curveto(2610.46744051,316.84136644)(2608.81949294,317.77978103)(2606.89694189,318.3291626)
\curveto(2605.02005827,318.87841274)(2602.91434748,319.15307067)(2600.57980322,319.15313721)
\curveto(2597.6042942,319.15307067)(2594.76616227,318.6724193)(2592.06539893,317.71118164)
\curveto(2589.36455634,316.79559012)(2586.86974683,315.19341887)(2584.58096289,312.90466309)
\curveto(2582.33789101,310.66156305)(2580.34662103,307.663214)(2578.60714697,303.90960693)
\curveto(2576.8676206,300.15589729)(2575.44855463,295.46382434)(2574.34994482,289.83337402)
\curveto(2573.89215971,287.59029706)(2573.57172546,285.50747443)(2573.38864111,283.5848999)
\curveto(2573.20551488,281.66226344)(2573.11396224,279.8769869)(2573.11398291,278.22906494)
\curveto(2573.11396224,274.79581522)(2573.50306097,271.93479514)(2574.28128027,269.64599609)
\curveto(2575.10523222,267.35716299)(2576.20386394,265.52611014)(2577.57717871,264.15283203)
\curveto(2578.99621954,262.77953085)(2580.62127895,261.79533994)(2582.45236182,261.20025635)
\curveto(2584.32916099,260.65093191)(2586.32043097,260.37627398)(2588.42617773,260.37628174)
\curveto(2591.40160265,260.37627398)(2594.23973457,260.81114903)(2596.94058203,261.6809082)
\curveto(2599.6413405,262.59642557)(2602.11326186,264.17570866)(2604.35635352,266.41876221)
\curveto(2606.64511768,268.66178815)(2608.68216398,271.68302537)(2610.46749854,275.48248291)
\curveto(2612.25271705,279.28189472)(2613.69467117,284.06552031)(2614.79336523,289.83337402)
}
}
{
\newrgbcolor{curcolor}{0 0 0}
\pscustom[linestyle=none,fillstyle=solid,fillcolor=curcolor]
{
\newpath
\moveto(2689.15707373,252.6171875)
\lineto(2698.42678809,300.33905029)
\curveto(2698.93026596,302.99402921)(2699.18203573,305.23706896)(2699.18209814,307.06817627)
\curveto(2699.18203573,309.35693789)(2698.88448964,311.23376707)(2698.28945898,312.69866943)
\curveto(2697.69430528,314.20922796)(2696.84744334,315.39941232)(2695.74887061,316.26922607)
\curveto(2694.65017991,317.18468885)(2693.34555475,317.82555735)(2691.83499121,318.1918335)
\curveto(2690.32431754,318.55797849)(2688.67636996,318.74108378)(2686.89114355,318.7411499)
\curveto(2684.19029047,318.74108378)(2681.55815199,318.21465608)(2678.99472021,317.16186523)
\curveto(2676.43120399,316.15472162)(2674.07372343,314.66699117)(2671.92227148,312.69866943)
\curveto(2669.81652554,310.77600385)(2667.98547269,308.37274698)(2666.42910742,305.4888916)
\curveto(2664.87268283,302.6507068)(2663.7282748,299.37769982)(2662.99587988,295.66986084)
\lineto(2654.61880469,252.6171875)
\lineto(2642.25918555,252.6171875)
\lineto(2653.58883643,311.05072021)
\curveto(2653.90925385,312.60705671)(2654.20679994,314.23211612)(2654.48147559,315.92590332)
\curveto(2654.80189212,317.66534023)(2655.07655005,319.29039964)(2655.3054502,320.80108643)
\curveto(2655.58008958,322.35741317)(2655.80897119,323.68492649)(2655.9920957,324.78363037)
\curveto(2656.17518176,325.88218992)(2656.26673441,326.5917229)(2656.26675391,326.91223145)
\lineto(2667.93972754,326.91223145)
\curveto(2667.93969637,326.68327554)(2667.87103188,326.0652952)(2667.73373389,325.05828857)
\curveto(2667.59637395,324.05113706)(2667.43615683,322.88384087)(2667.25308203,321.55639648)
\curveto(2667.06994626,320.27459055)(2666.86395281,318.94707722)(2666.63510107,317.57385254)
\curveto(2666.4061896,316.24627426)(2666.20019615,315.12475439)(2666.01712012,314.20928955)
\lineto(2666.22311377,314.20928955)
\curveto(2669.33587417,318.55797849)(2673.02086804,321.9912026)(2677.27810645,324.50897217)
\curveto(2681.53526382,327.02659795)(2686.68509998,328.28544679)(2692.72763037,328.28552246)
\curveto(2698.99893044,328.28544679)(2703.80544419,326.88926899)(2707.14718604,324.09698486)
\curveto(2710.48878712,321.3503341)(2712.15962285,316.93291908)(2712.15969824,310.84472656)
\curveto(2712.15962285,309.51715501)(2712.06807021,308.12097721)(2711.88504004,306.65618896)
\curveto(2711.74763596,305.23706896)(2711.51875435,303.72645036)(2711.19839453,302.12432861)
\lineto(2701.58535742,252.6171875)
\lineto(2689.15707373,252.6171875)
}
}
{
\newrgbcolor{curcolor}{0 0 0}
\pscustom[linestyle=none,fillstyle=solid,fillcolor=curcolor]
{
\newpath
\moveto(1549.76040625,102.5390625)
\lineto(2721.63540625,102.5390625)
}
}
{
\newrgbcolor{curcolor}{0.65882355 0 0.21176471}
\pscustom[linewidth=17.578125,linecolor=curcolor]
{
\newpath
\moveto(1549.76040625,102.5390625)
\lineto(2721.63540625,102.5390625)
}
}
{
\newrgbcolor{curcolor}{0.65882355 0 0.21176471}
\pscustom[linewidth=11.71875,linecolor=curcolor]
{
\newpath
\moveto(1090.10415625,1453.13675)
\curveto(943.56115625,1118.472625)(769.18615625,720.283625)(662.94403125,477.661)
}
}
{
\newrgbcolor{curcolor}{0.65882355 0 0.21176471}
\pscustom[linewidth=11.71875,linecolor=curcolor]
{
\newpath
\moveto(1185.95171875,1671.984375)
\lineto(1167.07284375,1424.390625)
\lineto(1016.79159375,1490.203125)
\lineto(1016.79159375,1490.203125)
\closepath
}
}
{
\newrgbcolor{curcolor}{0.65882355 0 0.21176471}
\pscustom[linewidth=11.71875,linecolor=curcolor]
{
\newpath
\moveto(1841.73296875,4380.73825)
\curveto(1708.57296875,3754.675875)(1532.22921875,2925.539)(1431.09646875,2450.05075)
}
}
{
\newrgbcolor{curcolor}{0.65882355 0 0.21176471}
\pscustom[linewidth=11.71875,linecolor=curcolor]
{
\newpath
\moveto(1891.30334375,4613.789125)
\lineto(1922.77984375,4367.484375)
\lineto(1762.30334375,4401.609375)
\lineto(1762.30334375,4401.609375)
\closepath
}
}
{
\newrgbcolor{curcolor}{0.65882355 0 0.21176471}
\pscustom[linewidth=11.71875,linecolor=curcolor]
{
\newpath
\moveto(2515.33846875,4414.617125)
\curveto(2889.49484375,3814.17175)(3388.84234375,3012.82025)(3695.97921875,2519.917875)
}
}
{
\newrgbcolor{curcolor}{0.65882355 0 0.21176471}
\pscustom[linewidth=11.71875,linecolor=curcolor]
{
\newpath
\moveto(2391.22521875,4613.789125)
\lineto(2584.79546875,4458.257875)
\lineto(2445.55334375,4371.49225)
\lineto(2445.55334375,4371.49225)
\closepath
}
}
{
\newrgbcolor{curcolor}{0.65882355 0 0.21176471}
\pscustom[linewidth=11.71875,linecolor=curcolor]
{
\newpath
\moveto(5229.04946875,2353.4765)
\curveto(5384.87371875,2301.351625)(5487.26034375,2203.382875)(5487.26034375,2059.57025)
\curveto(5487.26034375,1873.6405)(5316.09634375,1764.32825)(5079.58859375,1731.644625)
}
}
{
\newrgbcolor{curcolor}{0.65882355 0 0.21176471}
\pscustom[linestyle=none,fillstyle=solid,fillcolor=curcolor]
{
\newpath
\moveto(5079.58846875,2387.496125)
\lineto(5137.74971875,2326.1905)
\lineto(5216.70921875,2356.298625)
\lineto(5158.54796875,2417.604375)
\lineto(5079.58846875,2387.496125)
\closepath
}
}
{
\newrgbcolor{curcolor}{0.65882355 0 0.21176471}
\pscustom[linewidth=11.71875,linecolor=curcolor]
{
\newpath
\moveto(5079.58846875,2387.496125)
\lineto(5137.74971875,2326.1905)
\lineto(5216.70921875,2356.298625)
\lineto(5158.54796875,2417.604375)
\lineto(5079.58846875,2387.496125)
\closepath
}
}
{
\newrgbcolor{curcolor}{0 0 0}
\pscustom[linestyle=none,fillstyle=solid,fillcolor=curcolor]
{
\newpath
\moveto(5218.55900549,2550.33996582)
\curveto(5210.82275875,2550.33986465)(5204.9958145,2546.52135651)(5201.07815527,2538.88442993)
\curveto(5197.21002517,2531.29691495)(5195.27597559,2519.86618602)(5195.27600073,2504.59220886)
\curveto(5195.27597559,2489.36771188)(5197.21002517,2477.93698296)(5201.07815527,2470.29998779)
\curveto(5204.9958145,2462.71254139)(5210.82275875,2458.91882876)(5218.55900549,2458.9188385)
\curveto(5226.3447464,2458.91882876)(5232.17169065,2462.71254139)(5236.03985571,2470.29998779)
\curveto(5239.95747998,2477.93698296)(5241.91632507,2489.36771188)(5241.91639685,2504.59220886)
\curveto(5241.91632507,2519.86618602)(5239.95747998,2531.29691495)(5236.03985571,2538.88442993)
\curveto(5232.17169065,2546.52135651)(5226.3447464,2550.33986465)(5218.55900549,2550.33996582)
\moveto(5218.55900549,2562.24182129)
\curveto(5231.0063018,2562.24170822)(5240.50298114,2557.30740224)(5247.04907202,2547.43888855)
\curveto(5253.64460008,2537.61976934)(5256.94240257,2523.33755706)(5256.94248938,2504.59220886)
\curveto(5256.94240257,2485.89634084)(5253.64460008,2471.61412856)(5247.04907202,2461.74552917)
\curveto(5240.50298114,2451.92649566)(5231.0063018,2447.01698519)(5218.55900549,2447.01698303)
\curveto(5206.11161234,2447.01698519)(5196.59013748,2451.92649566)(5189.99455237,2461.74552917)
\curveto(5183.44851855,2471.61412856)(5180.17551156,2485.89634084)(5180.17552161,2504.59220886)
\curveto(5180.17551156,2523.33755706)(5183.44851855,2537.61976934)(5189.99455237,2547.43888855)
\curveto(5196.59013748,2557.30740224)(5206.11161234,2562.24170822)(5218.55900549,2562.24182129)
}
}
{
\newrgbcolor{curcolor}{0 0 0}
\pscustom[linestyle=none,fillstyle=solid,fillcolor=curcolor]
{
\newpath
\moveto(5283.4241178,2468.06838989)
\lineto(5299.1196897,2468.06838989)
\lineto(5299.1196897,2449.17419434)
\lineto(5283.4241178,2449.17419434)
\lineto(5283.4241178,2468.06838989)
}
}
{
\newrgbcolor{curcolor}{0 0 0}
\pscustom[linestyle=none,fillstyle=solid,fillcolor=curcolor]
{
\newpath
\moveto(5331.92417883,2468.06838989)
\lineto(5347.61975073,2468.06838989)
\lineto(5347.61975073,2449.17419434)
\lineto(5331.92417883,2449.17419434)
\lineto(5331.92417883,2468.06838989)
}
}
{
\newrgbcolor{curcolor}{0 0 0}
\pscustom[linestyle=none,fillstyle=solid,fillcolor=curcolor]
{
\newpath
\moveto(5435.7678678,2541.9342804)
\lineto(5409.06307959,2527.50328064)
\lineto(5435.7678678,2512.99789429)
\lineto(5431.45344519,2505.70800781)
\lineto(5406.45954871,2520.80848694)
\lineto(5406.45954871,2492.76473999)
\lineto(5397.97947668,2492.76473999)
\lineto(5397.97947668,2520.80848694)
\lineto(5372.9855802,2505.70800781)
\lineto(5368.67115759,2512.99789429)
\lineto(5395.3759458,2527.50328064)
\lineto(5368.67115759,2541.9342804)
\lineto(5372.9855802,2549.29855347)
\lineto(5397.97947668,2534.19807434)
\lineto(5397.97947668,2562.24182129)
\lineto(5406.45954871,2562.24182129)
\lineto(5406.45954871,2534.19807434)
\lineto(5431.45344519,2549.29855347)
\lineto(5435.7678678,2541.9342804)
}
}
{
\newrgbcolor{curcolor}{0 0 0}
\pscustom[linestyle=none,fillstyle=solid,fillcolor=curcolor]
{
\newpath
\moveto(5190.19696021,1756.31393433)
\lineto(5214.74453711,1756.31393433)
\lineto(5214.74453711,1841.04026794)
\lineto(5188.0397489,1835.68443298)
\lineto(5188.0397489,1849.37156677)
\lineto(5214.59576392,1854.72740173)
\lineto(5229.62185645,1854.72740173)
\lineto(5229.62185645,1756.31393433)
\lineto(5254.16943335,1756.31393433)
\lineto(5254.16943335,1743.66821289)
\lineto(5190.19696021,1743.66821289)
\lineto(5190.19696021,1756.31393433)
}
}
{
\newrgbcolor{curcolor}{0.65882355 0 0.21176471}
\pscustom[linewidth=11.71875,linecolor=curcolor]
{
\newpath
\moveto(1610.97921875,1453.13675)
\curveto(1757.52209375,1118.472625)(1931.89709375,720.283625)(2038.13921875,477.661)
}
}
{
\newrgbcolor{curcolor}{0.65882355 0 0.21176471}
\pscustom[linewidth=11.71875,linecolor=curcolor]
{
\newpath
\moveto(1515.13146875,1671.984375)
\lineto(1684.29159375,1490.203125)
\lineto(1534.01034375,1424.390625)
\lineto(1534.01034375,1424.390625)
\closepath
}
}
\end{pspicture}

    }
    \caption{Generic Fact model. Knowledge are used for a trusted information, Prediction for bioinformatics and bio-annotator observation and Assertion for an empiric observation. Predictions and Assertions results are facing to Knowledges.}
    \label{fig:fact}
\end{figure}

JBoss DROOLS is an open rule-engine written in java \cite{proctor2008drools}. Natively an oriented-object language it use a PHREAK algorithm \cite{phreak} to reasoning over structured data. It does a nice bridge between our java application and the logic.

\subsection{Reasoning over structured data}
Logic point of view a class is a structure. An instance of a class ( an object ) is a theory. An object can to contains others objects, that means a theory is a graph of smaller theories. Each theory are defined by a coherent set of values. A class C is a tuple \( \langle \mathcal{L}_{C}, \mathcal{A}_{C} , \mathcal{I}_{C} \rangle \). Where  \(\mathcal{L} \) is the vocabulary (class methods) ,  \(\mathcal{A} \)  a set of axioms (attributes) and \(\mathcal{I} \) a set of vocabulary to implement (\( \mathcal{I}_{C} \subset\mathcal{L}_{C} \)). The truth of a sub-theory can to be inferred to their upper theory following a Belnap logic (see Table~\ref{table:belnap}). In a such case a ``AND'' node have the following priority evaluation: FALSE > BOTH > NONE > TRUE. As example if one of sub-theory linked to a ``AND'' node is false the upper theory is false. Priority evaluation for ``OR'' node is: TRUE > NONE > BOTH > FALSE.
    
\begin{table}
    \caption{A four valued logic truth tables. T: true, F: false, B: both, N: none.\label{table:belnap}}
    \begin{minipage}{0.2\textwidth}
        \centering
        \begin{tabular}{|c||c|}\hline
            $F(\neg \alpha)$ & ~ \\ \hline \hline
            \textbf{T}       & F \\ \hline
            \textbf{B}       & B \\ \hline
            \textbf{N}       & N \\ \hline
            \textbf{F}       & T \\ \hline
        \end{tabular}
    \end{minipage}
    \begin{minipage}{0.35\textwidth}
        \centering
        \begin{tabular}{|c||c|c|c|c|c|c|}\hline
            $F(\land \alpha)$ & \textbf{T} & \textbf{B} & \textbf{N} & \textbf{F}  \\ \hline \hline
            \textbf{T}        & T          & B          & N          & F           \\ \hline
            \textbf{B}        & B          & B          & F          & F           \\ \hline
            \textbf{N}        & N          & F          & N          & F           \\ \hline
            \textbf{F}        & F          & F          & F          & F           \\ \hline
        \end{tabular}
    \end{minipage}
    \begin{minipage}{0.35\textwidth}
        \centering
        \begin{tabular}{|c||c|c|c|c|c|c|}\hline
            $F(\lor \alpha)$  & \textbf{T} & \textbf{B} & \textbf{N} & \textbf{F}  \\ \hline \hline
            \textbf{T}        & T          & T          & T          & T           \\ \hline
            \textbf{B}        & T          & B          & T          & B           \\ \hline
            \textbf{N}        & T          & T          & N          & N           \\ \hline
            \textbf{F}        & T          & B          & N          & F           \\ \hline
        \end{tabular}
    \end{minipage}
\end{table}


Point of view of a bio-annotator they are some cases not well cover by this logic. Indeed if a required pathway X has at least one reaction specific of the reaction block (not shared to another reactions blocks) then corresponding reaction block is present which infer the presence to the pathway X. That means we do not know yet what happen in this cases but the studied organisms has the required ability. By modifying ``AND'' node priority evaluation to TRUE > FALSE > BOTH > NONE in these cases we can to have an optimistic Belnap logic.

\subsection{Hierarchical knowledge representation}
Handling fact form a directed acyclic graph. A pathway contains one or more pathway variant. A pathway variant contains one or more reaction block. And reaction block contains one or more enzymatic reaction, these enzymatic reaction can to appear in one or more reaction block (as the reaction 3 in Figure~\ref{fig:metabolicNetworkReconstruction}). These knowledge need to be studied to know if is present or absent in the given organism.
Prediction fact will infer in a bottom-up way while assertion fact will do it in a top-down way. Once all facts are spreads over the DAG of knowledge a refined conclusion is put on each knowledge. This conclusion will to oppose an assertion fact with a prediction fact comparing a four logic value to an another following the conclusion table~\ref{table:conclusion}.


\begin{table}
    \caption{Conclusion table.}
    \begin{tabular}{|l||*{4}{l|}}\hline
        \backslashbox{\textbf{Assertion}}{\textbf{Prediction}} & \makebox[6em]{\textbf{TRUE}} & \makebox[6em]{\textbf{FALSE}} & \makebox[6em]{\textbf{BOTH}} &\makebox[6em]{\textbf{UNKNOWN}} \\ \hline \hline
        \textbf{REQUIRED}                                      & Confirmed P.                 & Unexpected A.                 & Contradictory A.             & Missing                        \\ \hline
        \textbf{AVOIDED}                                       & Unexpected A.                & Confirmed A.                  & Contradictory P.             & Confirmed A.                   \\ \hline
        \textbf{BOTH}                                          & Ambiguous P.                 & Ambiguous A.                  & Ambiguous C.                 & Ambiguous.                     \\ \hline
        \textbf{UNKNOWN}                                       & Unconfirmed P.               & Unconfirmed A.                & Unconfirmed C.               & Unknown                        \\ \hline
    \end{tabular}
    
    \begin{tabular}{|l|l|}
        \multicolumn{2}{c}{Legend}  \\ \hline
        A. & Absence       \\ \hline
        P. & Presence      \\ \hline
        C. & Contradictory \\ \hline
    \end{tabular}
    \label{table:conclusion}
\end{table}


\section{Conclusion and future work}


\bibliography{\jobname}


\end{document}

